\documentclass[11pt]{article}

\usepackage[a4paper,margin=28mm]{geometry}
\usepackage[T1]{fontenc}
\usepackage{amsmath,amssymb,amsthm,mathtools}
\usepackage{mathrsfs}
\usepackage{enumitem}
\usepackage{booktabs}
\usepackage{xcolor}
\usepackage{xurl}
\usepackage{hyperref}
\usepackage[nameinlink,noabbrev]{cleveref}

\hypersetup{
  colorlinks=true,
  linkcolor=blue!55!black,
  citecolor=green!40!black,
  urlcolor=blue!60!black,
  pdftitle={Selmer--Cartan and Motivic Moore--Reedy Towers}
}

\setlist[itemize]{leftmargin=2em,itemsep=0.2em,topsep=0.35em}
\setlist[enumerate]{leftmargin=2em,itemsep=0.2em,topsep=0.35em}
\emergencystretch=2em

\newtheorem{theorem}{Theorem}[section]
\newtheorem{proposition}[theorem]{Proposition}
\newtheorem{lemma}[theorem]{Lemma}
\newtheorem{corollary}[theorem]{Corollary}
\newtheorem{definition}[theorem]{Definition}
\newtheorem{construction}[theorem]{Construction}
\newtheorem{criterion}[theorem]{Criterion}
\newtheorem{principle}[theorem]{Principle}
\newtheorem{hypothesis}[theorem]{Hypothesis}
\newtheorem{problem}[theorem]{Problem}
\theoremstyle{remark}
\newtheorem{remark}[theorem]{Remark}
\newtheorem{warning}[theorem]{Warning}
\newtheorem{example}[theorem]{Example}
\newtheorem*{summary}{Summary}

\crefname{theorem}{theorem}{theorems}
\Crefname{theorem}{Theorem}{Theorems}
\crefname{proposition}{proposition}{propositions}
\Crefname{proposition}{Proposition}{Propositions}
\crefname{lemma}{lemma}{lemmas}
\Crefname{lemma}{Lemma}{Lemmas}
\crefname{corollary}{corollary}{corollaries}
\Crefname{corollary}{Corollary}{Corollaries}
\crefname{definition}{definition}{definitions}
\Crefname{definition}{Definition}{Definitions}
\crefname{construction}{construction}{constructions}
\Crefname{construction}{Construction}{Constructions}
\crefname{criterion}{criterion}{criteria}
\Crefname{criterion}{Criterion}{Criteria}
\crefname{hypothesis}{hypothesis}{hypotheses}
\Crefname{hypothesis}{Hypothesis}{Hypotheses}
\crefname{problem}{problem}{problems}
\Crefname{problem}{Problem}{Problems}
\crefname{remark}{remark}{remarks}
\Crefname{remark}{Remark}{Remarks}
\crefname{warning}{warning}{warnings}
\Crefname{warning}{Warning}{Warnings}

\newcommand{\Z}{\mathbf Z}
\newcommand{\F}{\mathbf F}
\newcommand{\N}{\mathbf N}
\newcommand{\Spec}{\operatorname{Spec}}
\newcommand{\REnd}{\mathbf R\!\operatorname{End}}
\newcommand{\cofib}{\operatorname{cofib}}
\newcommand{\Cobar}{\operatorname{Cobar}}
\newcommand{\DM}{\mathrm{DM}}
\newcommand{\bideg}{\operatorname{bideg}}
\newcommand{\Mot}{\mathrm{Mot}}
\newcommand{\rootc}{\mathrm{root}}
\newcommand{\MR}{\mathrm{MR}}
\newcommand{\Real}{\operatorname{Real}}
\newcommand{\Lin}{\operatorname{Lin}}
\newcommand{\ord}{\operatorname{ord}}
\newcommand{\supp}{\operatorname{supp}}


% ---- compatibility macros/packages ----
\usepackage{lmodern}
\usepackage{microtype}
\usepackage{amscd}
\usepackage{array}
\providecommand{\Q}{\mathbf Q}
\providecommand{\E}{\mathbf E}
\providecommand{\G}{\mathcal G}
\providecommand{\C}{\mathcal C}
\providecommand{\M}{\mathcal M}
\providecommand{\Hh}{\mathcal H}
\providecommand{\End}{\operatorname{End}}
\providecommand{\Hom}{\operatorname{Hom}}
\providecommand{\Perf}{\operatorname{Perf}}
\providecommand{\Fib}{\operatorname{Fib}}
\providecommand{\Tot}{\operatorname{Tot}}
\providecommand{\Gal}{\operatorname{Gal}}
\providecommand{\ad}{\operatorname{ad}}
\providecommand{\ob}{\operatorname{ob}}
\providecommand{\gr}{\operatorname{gr}}
\providecommand{\loc}{\operatorname{loc}}
\providecommand{\inv}{\operatorname{inv}}
\providecommand{\Inv}{\operatorname{Inv}}
\providecommand{\wt}{\operatorname{wt}}
\providecommand{\coker}{\operatorname{coker}}
\providecommand{\im}{\operatorname{im}}
\providecommand{\Alt}{\operatorname{Alt}}
\providecommand{\MC}{\operatorname{MC}}
\providecommand{\PD}{\mathrm{PD}}
\providecommand{\FS}{\mathrm{FS}}
\providecommand{\Sel}{\mathrm{Sel}}
\providecommand{\Lie}{\mathrm{Lie}}
\providecommand{\eps}{\varepsilon}
\providecommand{\zetaThree}{\zeta_3}
\newtheorem{recollection}[theorem]{Recollection}
\providecommand{\dR}{\mathrm{dR}}
\providecommand{\Fil}{\operatorname{Fil}}
\providecommand{\Nyg}{\mathcal N}
\providecommand{\DeltaPr}{\Delta}
\providecommand{\WCart}{\operatorname{WCart}}
\providecommand{\FWCart}{\operatorname{FWCart}}
\providecommand{\Spf}{\operatorname{Spf}}
\providecommand{\A}{\mathbf A}
\providecommand{\Conf}{\mathrm{Conf}}
\providecommand{\holim}{\operatorname*{holim}}
\providecommand{\hocolim}{\operatorname*{hocolim}}
\providecommand{\colim}{\operatorname*{colim}}
\providecommand{\Map}{\operatorname{Map}}
\providecommand{\Ord}{\operatorname{Ord}}
\providecommand{\Fin}{\operatorname{Fin}}
\providecommand{\AdMot}{\mathsf{AdMot}}
\providecommand{\Mfin}{\mathfrak M_{\Conf}^{\mathrm{fin}}}
\providecommand{\Ccat}{\mathcal C}
\providecommand{\FSCComp}{\mathsf{FSCComp}}
\providecommand{\Pic}{\operatorname{Pic}}
\providecommand{\fib}{\operatorname{fib}}
\providecommand{\id}{\mathrm{id}}
\providecommand{\RecOneMot}{\mathsf{Mot}^{\mathrm{rec},1}}
\providecommand{\FullMot}{\mathsf{Mot}^{\mathrm{full}}}
\providecommand{\tauone}{\tau_{1}}
\providecommand{\DepthFib}{\mathsf{DepthFib}}
\providecommand{\radical}{\operatorname{rad}}
\usepackage{longtable}
\usepackage{xr}
\newtheorem{firewall}{Firewall}[section]
\newtheorem*{observation}{Observation}
\newtheorem{maintheorem}{Main Theorem}
\providecommand{\themaintheorem}{\Alph{maintheorem}}

\newtheorem{proofgap}{Proof gap}

% ---- unified compatibility additions ----
\providecommand{\Aut}{\operatorname{Aut}}
\providecommand{\Cl}{\operatorname{Cl}}
\providecommand{\CH}{\operatorname{CH}}
\providecommand{\limop}{\operatorname*{lim}}
\renewcommand{\Ord}{\operatorname{ord}}
\providecommand{\R}{\mathcal R}
\providecommand{\SD}{\mathsf{SD}}
\providecommand{\Supp}{\mathsf{Supp}}
\providecommand{\tcofib}{\operatorname{tcofib}}
\providecommand{\val}{\operatorname{val}}
\providecommand{\Sha}{\operatorname{Sha}}
\providecommand{\rank}{\operatorname{rank}}
\crefname{recollection}{recollection}{recollections}
\Crefname{recollection}{Recollection}{Recollections}

\title{\textbf{Selmer--Cartan and Motivic Moore--Reedy Towers}\\[1.5mm]
\large Source geometry, full-motive realization, and a \(31\)-adic witness}
\author{}
\date{Unified core draft v0.1 --- 24 August 2026}

\begin{document}
\maketitle

\begin{abstract}
We construct a two-level Selmer--Cartan tower and its motivic realization.  At
the support level, finite ordered prime supports carry primitive transport,
Fox--Spencer and secondary classes, confluence, repeated-direction deformation,
and compatible higher obstruction data; the canonical \(5\!-
\!3\!-\!2\) descent is the retained finite transition calculation.  The
realization level promotes this source package to support-functorial motivic Moore--Reedy
objects, a role-separated objectification of support control, full-motive depth
thickenings, and derived stacks.  Pointed marked Morita equivalences preserve
the complete finite tower and its pseudo-perfect-module stack.
Finally, an explicit \(31\)-adic witness supplies an exact-
\(\lambda_{31}=2\) repeated branch, integral carry normalization, and a
Kummer--Bockstein--Massey local--global obstruction.  Its terminal order-\(31\)
source line admits an independent unipotent-gerbe provenance for the proper
\(31\)-fold operation law, together with the carrier-level motivic Moore
specialization, marked full-motive lifting presentation, classical shadow, and
stack globalization constructed in the structural parts.  The truncation
\(\tauone:\FullMot\to\RecOneMot\) separates valuation depth from support;
coefficient depth, confluence multiplicity, obstruction height, and geometric
realization remain independent coordinates throughout.
\end{abstract}

\tableofcontents

\section{Architecture and two-level convention}
\label{sec:unified-architecture}
\label{sec:new-architecture}
\label{sec:internal-provenance}
\label{P1C-arch:realization-boundary-final}
\label{P1C-int:two-level-tower-interface}
\label{sec:p2-self-contained-scart-input}
\label{P2R-int:two-level-tower-realization}
The organizing tower is
\[
 \boxed{\DepthFib\longrightarrow\FullMot
 \stackrel{\tauone}{\longrightarrow}\RecOneMot},
 \qquad
 (v_p(n))_p\longmapsto(\mathbf 1_{v_p(n)>0})_p .
\]
A new prime changes support, whereas a higher valuation of an existing prime
changes depth.  The source-side realization square is
\[
\begin{CD}
 \mathsf{SCart}^{\rm full} @>{\mathcal R_{\rm full}}>> \FullMot\\
 @V{\tauone^{\rm Cart}}VV @VV{\tauone}V\\
 \mathsf{SCart}^{\rm rec} @>{\mathcal R_{1}}>> \RecOneMot .
\end{CD}
\]
At every finite stage,
\(
 \mathsf{SCart}^{\rm full}
 :=\mathsf{SCart}^{\rm rec}\times_{\RecOneMot}\FullMot
\).
The finite support--coefficient-depth diagrams are algebraic controllers, not
substitutes for the valuation-depth fibre.  Likewise, support, confluence
multiplicity, coefficient exponent, obstruction height, and realization are
typed separately.  The adjective recursive records compatible finite-support
organization; it is not an intrinsic construction of motives.

Parts~I--II build the support-level source tower.  Parts~III--VI construct the
low motivic realization, the Moore--Reedy core, the full-motive thickening,
and the stack-level presentation.  Part~VII then gives a concrete
\(31\)-adic witness for the repeated and terminal source-obstruction branches,
together with its exact-order-\(31\) motivic specialization and terminal
unipotent-gerbe operation provenance.  The witness certifies non-vacuity of the
source obstruction; the general realization does not depend on the special
numerical choice \(31\).

\section{Internal dependency and firewall index}
\label{sec:internal-handle-index}
This is a minimal internal dependency index, not a cross-paper API.  The first
table lists only handles used across major parts of the present paper; labels
used solely inside one local proof are omitted.  The following list gives every
labeled warning/firewall, so that scope restrictions can be audited from one
place.

\subsection*{Stable internal handles}
\begin{longtable}{@{}p{0.46\linewidth}p{0.48\linewidth}@{}}
\toprule
\normalfont Handle & Role\\
\midrule
\endfirsthead
\toprule
\normalfont Handle & Role\\
\midrule
\endhead
\nolinkurl{P1C-def:filtered-confluence-lift}
  & filtered repeated-branch input, \cref{P1C-def:filtered-confluence-lift}\\
\nolinkurl{P1C-def:abstract-finite-pd-carrier}
  & finite integral Moore/PD carrier, \cref{P1C-def:abstract-finite-pd-carrier}\\
\nolinkurl{P1C-thm:universal-higher-obstruction-recursion}
  & universal finite obstruction recursion, \cref{P1C-thm:universal-higher-obstruction-recursion}\\
\nolinkurl{P1C-thm:formal-filtered-alignment}
  & pointed formal/filtered repeated-line comparison, \cref{P1C-thm:formal-filtered-alignment}\\
\nolinkurl{P1C-thm:finite-confluent-interface}
  & finite source package, \cref{P1C-thm:finite-confluent-interface}\\
\nolinkurl{P2L-thm:v48r4-motivic-seed}
  & primitive motivic Moore seed, \cref{P2L-thm:v48r4-motivic-seed}\\
\nolinkurl{P2R-def:v92-line-valued-lift}
  & line-valued root/Moore lifting object, \cref{P2R-def:v92-line-valued-lift}\\
\nolinkurl{P2M-thm:channel-complete-realization}
  & support-functorial motivic realization, \cref{P2M-thm:channel-complete-realization}\\
\nolinkurl{P2M-thm:role-separated-objectification}
  & independent controller/birth objects and their incidence maps,
    \cref{P2M-thm:role-separated-objectification}\\
\nolinkurl{P2M-thm:prime-power-all-support-comparison}
  & coefficient-faithful prime-power comparison, \cref{P2M-thm:prime-power-all-support-comparison}\\
\nolinkurl{P2M-thm:v38-finite-motivic-recursion-closure}
  & finite carrier-plus-ledger closure, \cref{P2M-thm:v38-finite-motivic-recursion-closure}\\
\nolinkurl{thm:genuine-successor-stage-motivic-functor}
  & marked full-motive lifting step, \cref{thm:genuine-successor-stage-motivic-functor}\\
\nolinkurl{thm:classical-low-sector-comparison}
  & classical \(0/1\)-motivic shadow, \cref{thm:classical-low-sector-comparison}\\
\nolinkurl{P2R-thm:v42-stack-globalization}
  & fixed-presentation derived stack, \cref{P2R-thm:v42-stack-globalization}\\
\nolinkurl{P2M-thm:marked-morita-presentation-independence}
  & marked Morita independence of the complete tower and stack,
    \cref{P2M-thm:marked-morita-presentation-independence}\\
\nolinkurl{thm:paper1-31adic-witness}
  & terminal arithmetic witness, \cref{thm:paper1-31adic-witness}\\
\nolinkurl{W31-thm:terminal-unipotent-gerbe-provenance}
  & independent proper-\(31\)-fold operation correspondence,
    \cref{W31-thm:terminal-unipotent-gerbe-provenance}\\
\nolinkurl{W31-thm:common-source-motivic-span}
  & carrier-level arithmetic--motivic specialization,
    \cref{W31-thm:common-source-motivic-span}\\
\bottomrule
\end{longtable}

\subsection*{Complete labeled firewall list}
\begingroup
\setlength{\parskip}{0.15em}
\noindent \nolinkurl{P1L-warn:pq-flat-scope}: \cref{P1L-warn:pq-flat-scope}.\par
\noindent \nolinkurl{P1C-warn:external-ray-alignment}: \cref{P1C-warn:external-ray-alignment}.\par
\noindent \nolinkurl{P1C-warn:pd-provenance-fixed-order-only}: \cref{P1C-warn:pd-provenance-fixed-order-only}.\par
\noindent \nolinkurl{P1C-warn:effective-weight-not-depth}: \cref{P1C-warn:effective-weight-not-depth}.\par
\noindent \nolinkurl{P1C-warn:cc-recursion-not-provenance}: \cref{P1C-warn:cc-recursion-not-provenance}.\par
\noindent \nolinkurl{P1C-warn:two-primary-integral-defect}: \cref{P1C-warn:two-primary-integral-defect}.\par
\noindent \nolinkurl{P2M-warn:tate-saturation-not-new-provenance}: \cref{P2M-warn:tate-saturation-not-new-provenance}.\par
\noindent \nolinkurl{P2M-warn:v33-squarefree-jet-scope}: \cref{P2M-warn:v33-squarefree-jet-scope}.\par
\noindent \nolinkurl{P2M-warn:v36-bar-jet-provenance-firewall}: \cref{P2M-warn:v36-bar-jet-provenance-firewall}.\par
\noindent \nolinkurl{P2M-warn:scalar-prime-three}: \cref{P2M-warn:scalar-prime-three}.\par
\noindent \nolinkurl{P2M-warn:v37-readout-not-new-provenance}: \cref{P2M-warn:v37-readout-not-new-provenance}.\par
\noindent \nolinkurl{P2M-warn:v38-finite-not-motivic-adelic}: \cref{P2M-warn:v38-finite-not-motivic-adelic}.\par
\noindent \nolinkurl{fw:cartan-layer-not-classical-motive}: \cref{fw:cartan-layer-not-classical-motive}.\par
\noindent \nolinkurl{warn:typed-31-coordinates}: \cref{warn:typed-31-coordinates}.\par
\noindent \nolinkurl{warn:v063-old-slot-marking}: \cref{warn:v063-old-slot-marking}.\par
\noindent \nolinkurl{warn:v074-two-bocksteins}: \cref{warn:v074-two-bocksteins}.\par
\noindent \nolinkurl{W31-warn:prime-power-motivic-firewall}: \cref{W31-warn:prime-power-motivic-firewall}.\par
\endgroup

\part{Support-level source tower: primitive Selmer--Cartan geometry}
\section{Arithmetic ternary class}

\subsection{The cancelled transverse candidate}

The exact-\(\lambda=3\) quartic calibration contains three marked arithmetic directions.  Let
\[
 x=\bar\chi,\qquad
 y=\bar\chi_{\mathrm{ac}},\qquad
 m=\alpha,
\]
where \(x\) is cyclotomic, \(y\) is anticyclotomic and \(m\) is the Kummer root class.  In the corresponding scalar dg algebra \(B\) and module \(M\), the calculation gives
\[
 H^2(B)=0,
 \qquad
 H^1(B)\smile H^1(M)=0.
\]
Consequently all mixed binary products vanish, every two-scalar/one-module triple product is defined, its ordinary indeterminacy is zero, and one obtains a trilinear map
\[
 T:H^1(B)\otimes H^1(B)\otimes H^1(M)\longrightarrow H^2(M).
\]
Define
\[
 \tau_R=T(x,y;m),\qquad
 \tau_L=T(y,x;m),\qquad
 \sigma=\langle x,m,y\rangle.
\]
The bicyclic--Heisenberg decomposition gives the swap relation
\begin{equation}\label{P1L-eq:swap}
 \tau_L=\sigma-\tau_R.
\end{equation}
A local universal-norm argument then proves \(\sigma=0\) for the explicit field \(\Q(\sqrt{-519})\) at \(5\) .

At first sight, \(\tau_R\) looks like the missing three-direction class.  The following observation explains why it is not.

\begin{lemma}[Cartan cancellation of the transverse Heisenberg triple]\label{P1L-prop:cartan-cancellation}
Let
\[
 D=E_{22},\qquad E=E_{21}\in \End_{\F_5}(\F_5^2)
\]
so that \(D^2=D\), \(DE=E\), \(ED=E^2=0\).  Put
\[
 X=xD,\qquad Y=yD,\qquad Z=mE.
\]
On the marked three-plane with formal parameters \(u,v,w\), the \(E\)-root component of the cubic formal-moduli coefficient of \(uvw\) is proportional, up to the global convention sign, to
\[
 \tau_R+\tau_L=\sigma.
\]
Hence in the exact-\(\lambda=3\) calibration it vanishes.
\end{lemma}

\begin{proof}
A transferred cubic operation is a sum of two binary-tree terms.  The coefficient matrix of each term is the ordered product of the input matrix labels.  Among the six permutations of \((X,Y,Z)\), a contribution can land back in the \(E\)-root line only if the two diagonal inputs precede the root input.  Thus only the orders \((X,Y,Z)\) and \((Y,X,Z)\) survive the root projection.  Their scalar cohomology values are \(\tau_R\) and \(\tau_L\).  Expanding the cubic Maurer--Cartan term on \(uX+vY+wZ\) therefore produces the sum \(\tau_R+\tau_L\).  Equation \eqref{P1L-eq:swap} identifies this sum with \(\sigma\), and the calibration gives \(\sigma=0\).
\end{proof}

This no-go suggests the correct design principle: the coefficient marking must remember the orientation of the three inputs so strongly that only one ordered cubic word can reach the chosen obstruction line.


\subsection{Restricted-ramification arithmetic calibration}

We now use the arithmetic triple linking theory of Amano--Mizusawa--Morishita \cite{AMM}.  Let
\[
 k=\Q(\zeta_3)=\Q(\sqrt{-3})
\]
and let \(\mathfrak p_1,\mathfrak p_2,\mathfrak p_3\) be distinct primes of \(k\) with norm congruent to \(1\) modulo \(9\).  Choose the normalized prime generators \(\pi_i\) used in cubic reciprocity.  Assume the pairwise cubic residue conditions
\begin{equation}\label{P1L-eq:pairwise-cubic}
 \left(\frac{\pi_i}{\pi_j}\right)_3=1
 \qquad (i\neq j).
\end{equation}
Let \(\chi_i\in H^1(G_S,\F_3)\) be the corresponding mod-3 Kummer characters on a restricted-ramification pro-3 Galois group containing the required auxiliary prime.

Amano--Mizusawa--Morishita show that the pairwise Milnor invariants vanish under \eqref{P1L-eq:pairwise-cubic}, while the ordered triple Milnor invariant \(\mu_3(123)\) may be nonzero.  They construct a Redei-type Heisenberg extension
\[
 K_{\{\mathfrak p_1,\mathfrak p_2\}}/k,
 \qquad
 \Gal(K_{\{\mathfrak p_1,\mathfrak p_2\}}/k)
 \simeq H_3(\F_3),
\]
of degree \(27\), ramified only at the first two marked primes with ramification index three.  The triple cubic residue symbol is
\[
 [\mathfrak p_1,\mathfrak p_2,\mathfrak p_3]_3
 =\zeta_3^{\mu_3(123)}
\]
and records the Frobenius class of \(\mathfrak p_3\) in the central Heisenberg direction.

The same paper gives a cohomological interpretation in terms of a triple Massey product.  Under the pairwise cubic residue hypotheses, the product
\[
 \langle\chi_1,\chi_2,\chi_3\rangle
 \in H^2(G_S,\F_3)
\]
is defined as a single class.  Evaluating it on the local relation at \(\mathfrak p_3\) recovers the exponent of the triple cubic residue symbol, up to the stated sign convention \cite[Theorem 7.5]{AMM}.

\begin{corollary}[Eisenstein calibration triple]\label{P1L-thm:amm-example}
Take the Eisenstein prime ideals
\[
 \mathfrak p_1=(17),\qquad
 \mathfrak p_2=(53),\qquad
 \mathfrak p_3=(71).
\]
Then the hypotheses above are satisfied and
\[
 [\mathfrak p_1,\mathfrak p_2,\mathfrak p_3]_3=\zeta_3^2.
\]
In particular
\[
 \langle\chi_1,\chi_2,\chi_3\rangle\neq0.
\]
\end{corollary}

\begin{proof}
The explicit computation is Example 6.4 of \cite{AMM}: for \((\pi_1,\pi_2)=(-17,-53)\), the listed third primes include \(\pi_3=-71\), and the resulting triple cubic residue symbol equals \(\zeta_3^2\).  Theorem 7.5 of the same paper identifies this nontrivial symbol with the evaluation of the triple Massey class, so the class cannot vanish.
\end{proof}

\begin{remark}
The three arithmetic primes in \cref{P1L-thm:amm-example} are genuinely distinct directions.  This is the feature missing from the repeated-direction exact-\(\lambda=2\) calibration.
\end{remark}


\subsection{Directed root-path realization}

The scalar triple class alone is not yet a marked deformation class.  We now embed it into an endomorphism-valued Galois cochain algebra so that the ordering of the three arithmetic inputs is visible at the coefficient level.

Let
\[
 V=\F_3^4
\]
with the trivial residual Galois action, and let
\[
 A=C^\bullet(G_S,\End_{\F_3}(V))
\]
with cup product followed by matrix multiplication.  Write \(E_{ij}\) for the standard matrix units.

\begin{construction}[Directed root path]\label{P1L-con:root-path}
Define the marked degree-one classes
\[
 X_1=\chi_1E_{12},\qquad
 X_2=\chi_2E_{23},\qquad
 X_3=\chi_3E_{34}.
\]
The coefficient path is
\[
 1\longrightarrow2\longrightarrow3\longrightarrow4.
\]
The key multiplication identities are
\[
 E_{12}E_{23}=E_{13},\qquad
 E_{23}E_{34}=E_{24},\qquad
 E_{12}E_{23}E_{34}=E_{14},
\]
whereas every permutation of the three matrix units different from the order \((12),(23),(34)\) has zero total product in the relevant binary-tree expansion.
\end{construction}

\subsection{Binary gates}

\begin{lemma}[Pairwise Cartan brackets vanish]\label{P1L-prop:pairwise-brackets}
For the marked classes of \cref{P1L-con:root-path}, every pairwise cup-commutator vanishes in cohomology:
\[
 [X_i,X_j]=0\in H^2(A)
 \qquad(i\neq j).
\]
Thus the quadratic Kuranishi tensor vanishes on the marked three-plane
\[
 P:=\langle X_1,X_2,X_3\rangle.
\]
\end{lemma}

\begin{proof}
For \(i<j\), each coefficient product is either zero or a matrix unit multiplying the scalar cup \(\chi_i\smile\chi_j\).  The pairwise cubic residue condition implies the corresponding cup class is zero.  Reversing the order either produces the reversed scalar cup or an incompatible matrix product, hence again zero in cohomology.  Therefore all pairwise commutators vanish.
\end{proof}

Choose cochains \(a_{12},a_{23}\in C^1(G_S,\F_3)\) with
\[
 da_{12}=\chi_1\smile\chi_2,
 \qquad
 da_{23}=\chi_2\smile\chi_3.
\]
The scalar Massey representative is
\begin{equation}\label{P1L-eq:massey-rep}
 m_{123}^{\rm AMM}
 :=
 [\chi_1\smile a_{23}+a_{12}\smile\chi_3]
 =\langle\chi_1,\chi_2,\chi_3\rangle.
\end{equation}

\subsection{The unique directed cubic coefficient}

\begin{proposition}[Directed cubic nonvanishing]\label{P1L-thm:111-jet}
Let \(m_3\) be a minimal \(A_\infty\) ternary operation obtained from a contraction of \(A\) preserving the matrix-unit weight decomposition.  Then the \(E_{14}\)-projection satisfies
\begin{equation}\label{P1L-eq:m3-root}
 \pi_{14}m_3(X_1,X_2,X_3)=m_{123}^{\rm AMM}E_{14}.
\end{equation}
For every nonidentity permutation \(\sigma\in S_3\),
\begin{equation}\label{P1L-eq:perm-zero}
 \pi_{14}m_3(X_{\sigma(1)},X_{\sigma(2)},X_{\sigma(3)})=0.
\end{equation}
Consequently, for
\[
 \xi=u_1X_1+u_2X_2+u_3X_3
\]
over the third-order Artinian test algebra
\[
 R_3=\F_3[u_1,u_2,u_3]/(u_1,u_2,u_3)^4,
\]
the quadratic obstruction on the marked three-plane is zero, while the \(E_{14}\)-component of the cubic formal-moduli obstruction is
\begin{equation}\label{P1L-eq:formal-111}
 \ob^{(3)}_{14}(\xi)
 =\pm u_1u_2u_3\,m_{123}^{\rm AMM}E_{14}
eq0
\end{equation}
for the arithmetic triple of \cref{P1L-thm:amm-example}.
\end{proposition}

\begin{proof}
The standard transferred ternary product is a signed sum of the two planar binary-tree terms
\[
 \mu(h\mu(-,-),-),
 \qquad
 \mu(-,h\mu(-,-)).
\]
For the ordered input \((X_1,X_2,X_3)\), the first binary products carry coefficient matrices \(E_{13}\) and \(E_{24}\), and the final multiplication lands in \(E_{14}\).  With the null-homotopies corresponding to \(a_{12}\) and \(a_{23}\), the scalar coefficient is exactly the defining-system expression \eqref{P1L-eq:massey-rep}.  This proves \eqref{P1L-eq:m3-root}.

For any other permutation, every planar binary tree contains at least one incompatible matrix multiplication.  For example \(E_{23}E_{12}=0\), \(E_{12}E_{34}=0\), \(E_{34}E_{23}=0\), and similarly after the first binary multiplication.  Hence the \(E_{14}\)-projection vanishes, proving \eqref{P1L-eq:perm-zero}.

Expanding the homotopy Maurer--Cartan equation on \(\xi\), the coefficient of \(u_1u_2u_3\) is the sum over the six ordered ternary evaluations.  Only the directed order survives the \(E_{14}\)-projection, so the coefficient is \(\pm m_{123}^{\rm AMM}E_{14}\).  Nonvanishing follows from \cref{P1L-thm:amm-example}.
\end{proof}

\section{The flat \texorpdfstring{\(pq\)}{pq} Cartan germ and its transverse Fox--Spencer obstruction}
\label{P1L-sec:pq-cartan-spencer}

The directed-root calculation has a Cartan-geometric reformulation which does
not use the Ore presentation.  The point is to regard the first two marked
arithmetic directions as a two-dimensional formal Cartan base and to ask
whether a third direction prolongs this flat base.  At the order relevant for
the present paper, the answer is controlled exactly by the ternary deformation
class \(\Theta_3\).

\subsection{The quadratic-flat \texorpdfstring{\(pq\)}{pq} slice}

Write
\[
 P_{12}:=\langle X_1,X_2\rangle\subset H^1(A).
\]
Let \(Q=Q_2+Q_3+\cdots\) be the minimal deformation structure fixed in
\cref{P1L-sec:hierarchy}.  Here and below ``flat'' refers to the Cartan curvature
seen by the binary Kuranishi tensor; it is not a claim that every higher
partition-Lie operation on the two-plane vanishes.

\begin{definition}[Marked \(pq\)-Cartan germ]\label{P1L-def:pq-cartan-germ}
The marked \(pq\)-Cartan germ is the formal two-direction slice
\[
 \mathfrak C_{12}
 :=\mathfrak M\big|_{P_{12}}
\]
of the residual formal-moduli problem, equipped with the ordered tangent frame
\((X_1,X_2)\).  Its Cartan curvature at the base point is
\[
 \mathcal R_{12}:=Q_2(X_1,X_2).
\]
We call the germ \emph{quadratic-flat} when \(\mathcal R_{12}=0\).
A choice of cochain \(a_{12}\) with
\[
 da_{12}=\chi_1\smile\chi_2
\]
is called a pointing of the flat binary face.
\end{definition}

\begin{lemma}[Flatness of the directed \(pq\)-face]
\label{P1L-prop:pq-cartan-flat}
For the directed-root marking of \cref{P1L-con:root-path}, the germ
\(\mathfrak C_{12}\) is quadratic-flat.  More precisely,
\[
 Q_2|_{P_{12}}=0
\]
in cohomology, and the chosen defining-system cochain \(a_{12}\) gives a
pointing of this vanishing.
\end{lemma}

\begin{proof}
This is the restriction of \cref{P1L-prop:pairwise-brackets} to the first two
marked directions.  The only potentially nonzero ordered coefficient product
is \(E_{12}E_{23}=E_{13}\), and its scalar coefficient is
\(\chi_1\smile\chi_2\), which is exact by the pairwise cubic-residue
hypothesis.  The reversed coefficient product is zero.  Hence the binary
Kuranishi class vanishes, while \(a_{12}\) records a chosen chain-level
trivialization.
\end{proof}

\subsection{Intrinsic Fox--Spencer prolongation on the flat face}

The preceding proposition allows one to restart the ordinary Cartan--Spencer
construction on the two-dimensional formal base itself.  Introduce formal
coordinates \(u_1,u_2\) dual to the ordered tangent frame and, for a finite
free test module \(E\), write
\[
 J^n_{12}(E)
 =E\otimes
 k[y_{a,b}:a+b\le n]
\]
for the formal \(n\)-jet chart, with the evident divided-power replacement in
positive characteristic when a PD model is required.  On the inverse jet
tower define the total derivatives
\[
 \mathcal T_1
 =\partial_{u_1}
  +\sum_{a,b\ge0}y_{a+1,b}\,\partial_{y_{a,b}},
 \qquad
 \mathcal T_2
 =\partial_{u_2}
  +\sum_{a,b\ge0}y_{a,b+1}\,\partial_{y_{a,b}}.
\]
Their contact generators are
\[
 \vartheta_{a,b}
 =dy_{a,b}-y_{a+1,b}\,du_1-y_{a,b+1}\,du_2.
\]

\begin{definition}[Intrinsic \(pq\)-Fox--Spencer complex]
\label{P1L-def:pq-fs}
The intrinsic Fox--Spencer complex of the flat marked face is the horizontal
Spencer complex generated by \(\mathcal T_1,\mathcal T_2\):
\[
 \operatorname{FSC}_{12}(E):
 J^\infty_{12}(E)
 \xrightarrow{\ \mathsf S_{12}\ }
 \Omega^1_{12}\otimes J^\infty_{12}(E)
 \xrightarrow{\ \mathsf S_{12}\ }
 \Omega^2_{12}\otimes J^\infty_{12}(E),
\]
where
\[
 \mathsf S_{12}
 =du_1\wedge\mathcal T_1+du_2\wedge\mathcal T_2.
\]
\end{definition}

\begin{lemma}[Spencer closure on the \(pq\)-face]
\label{P1L-prop:pq-spencer-closure}
On the formal flat jet chart one has
\[
 [\mathcal T_1,\mathcal T_2]=0,
 \qquad
 \mathsf S_{12}^2=0.
\]
Thus \(\mathfrak C_{12}\) supplies a genuine two-direction Cartan base on
which Fox--Spencer prolongation may be performed independently of the third
arithmetic direction.
\end{lemma}

\begin{proof}
Both identities are the standard coordinate identities for a two-dimensional
formal jet space.  The arithmetic input is not the coordinate calculation but
\cref{P1L-prop:pq-cartan-flat}: it identifies the marked binary deformation face
with a flat Cartan germ to the order at which the transverse third-jet problem
is posed.
\end{proof}

\begin{warning}[Scope of the flatness statement]\label{P1L-warn:pq-flat-scope}
The proposition does not identify \(\operatorname{FSC}_{12}\) with the finite
Fox gauge of the split global-descent paper, and it does not assert all-order
formality of the positive-characteristic two-direction formal-moduli problem.
It is the intrinsic flat prolongation model attached to the marked
quadratic-flat Cartan face.  A comparison with the global finite FSC carrier
is a separate theorem.
\end{warning}

\subsection{Spencer envelopes and finite global Fox gauges}

Two different objects here carry Fox--Spencer terminology.  The complex of
\cref{P1L-def:pq-fs} is the horizontal jet--Spencer complex of the formal Cartan
base.  By contrast, a global Fox--Spencer gauge is a bounded finite-projective chain
representative of the intrinsic continuous Galois cochain object.  A direct
quasi-isomorphism between these two objects is therefore ill-typed.  The
correct comparison is obtained after applying the same formal Cartan
prolongation to the coefficient carrier on both sides.

\begin{definition}[$pq$-Spencer envelope]\label{P1L-def:pq-spencer-envelope}
Let $C^\bullet$ be a cochain complex over the residual coefficient field, or a
bounded finite-projective complex before reduction.  Its \emph{$pq$-Spencer
envelope} is the total complex
\[
 \mathbb{Sp}_{12}(C^\bullet)
 :=
 \operatorname{Tot}\!\left(
   \Omega^\bullet_{12}\widehat\otimes
   J^\infty_{12}(C^\bullet),
   d_C+\mathsf S_{12}
 \right),
\]
where $J^\infty_{12}$ is formed termwise and completed in the formal jet
filtration.  In positive characteristic one uses the corresponding completed
PD jet algebra.  A marked envelope additionally records cocycle lifts of
$X_1,X_2$ and a chosen binary pointing representing $a_{12}$.
\end{definition}

\begin{lemma}[Homotopy invariance of the Spencer envelope]
\label{P1L-lem:pq-spencer-envelope-invariance}
Let $f:C^\bullet\to D^\bullet$ be a quasi-isomorphism between bounded-below
cochain complexes for which the completed jet filtration is degreewise
complete and the associated graded terms are flat over the coefficient ring.
Then
\[
 \mathbb{Sp}_{12}(f):
 \mathbb{Sp}_{12}(C^\bullet)
 \xrightarrow{\sim}
 \mathbb{Sp}_{12}(D^\bullet)
\]
is a filtered quasi-isomorphism.  The same conclusion holds for bounded
finite-projective gauges and after derived coefficient specialization.
\end{lemma}

\begin{proof}
Filter both envelopes by total jet order.  On the associated graded, the
Spencer direction is the ordinary two-variable Koszul--Spencer differential
and the map is obtained from $f$ by tensoring with finite free homogeneous jet
pieces.  Hence every associated-graded map is a quasi-isomorphism.  Completeness
and the boundedness hypotheses give the filtered comparison.  For a bounded
finite-projective gauge, ordinary and derived completed tensor products agree
on each finite jet stage, and passage to the inverse limit gives the stated
completed result.
\end{proof}

\begin{definition}[Marked $pq$-compatible Fox gauge]
\label{P1L-def:pq-compatible-fox-gauge}
A global finite Fox gauge $(C^\bullet_{\mathrm{FS}},\alpha)$ is
\emph{marked $pq$-compatible} if the homotopy-algebra transfer along
\[
 \alpha:C^\bullet_{\mathrm{FS}}
 \xrightarrow{\sim}
 C^\bullet(G_{F,S},\operatorname{End}V)
\]
is chosen compatibly with the directed matrix-weight decomposition, carries
marked cocycles $\widetilde X_1,\widetilde X_2$ to $X_1,X_2$, and carries a
chosen binary null-homotopy $\widetilde a_{12}$ to the pointing $a_{12}$ up to
chain homotopy.
\end{definition}

\begin{proposition}[Finite-gauge realization of the flat $pq$ Cartan face]
\label{P1L-thm:pq-finite-gauge-comparison}
Let $A=C^\bullet(G_{F,S},\operatorname{End}V)$ be the intrinsic Galois cochain
algebra of the directed-root construction, and let
$(C^\bullet_{\mathrm{FS}},\alpha)$ be any marked $pq$-compatible finite Fox
gauge.  Then there is an equivalence of marked filtered Spencer envelopes
\[
 \boxed{
 \mathbb{Sp}_{12}(C^\bullet_{\mathrm{FS}})
 \simeq
 \mathbb{Sp}_{12}(A).}
\]
Under this equivalence the transferred binary Kuranishi tensor, its chosen
null-homotopy on $P_{12}$, and the marked ternary deformation class agree up to
the standard homotopy-transfer equivalence.  In particular, the quadratic
flatness of the $pq$ face and the class
$\Theta_3(Q)(X_1,X_2,X_3)$ are independent of the compatible finite Fox gauge.
\end{proposition}

\begin{proof}
By definition, a global Fox--Spencer gauge is a finite chain representative
of the intrinsic continuous cochain object, so
$\alpha$ is a quasi-isomorphism.  Its multiplicative-compatibility theorem
transfers the cup product and the full chosen $A_\infty/L_\infty$ package to
finite gauges and states that different coherent transfer choices give
equivalent homotopy-algebra diagrams.  Apply
\cref{P1L-lem:pq-spencer-envelope-invariance} to the underlying complexes and use
the marked compatibility to identify the ordered tangent frame and the binary
pointing.  Homotopy transfer then identifies the minimal tensors $Q_2$ and
$Q_3$ modulo the usual coordinate changes.  Therefore their invariant
restriction to the marked face, and in particular the class
$[Q_3]=\Theta_3(Q)$, is preserved.
\end{proof}


\subsection{The third direction as a transverse Spencer lifting problem}

Set
\[
 R_{123}
 :=\F_3[u_1,u_2,u_3]/(u_1,u_2,u_3)^4
\]
and remove only the fully mixed cubic monomial:
\[
 R_{\partial}
 :=R_{123}/(u_1u_2u_3).
\]
Then
\[
 0\longrightarrow
 I_{123}:=\F_3\,u_1u_2u_3
 \longrightarrow R_{123}\longrightarrow R_{\partial}
 \longrightarrow0
\]
is a square-zero small extension.  The pairwise vanishing data and the chosen
cochains \(a_{12},a_{23}\) give a pointed deformation over the boundary ring
\(R_{\partial}\); the complementary binary branch is the marked zero branch
in the directed incidence model.

\begin{definition}[Transverse Spencer obstruction]
\label{P1L-def:transverse-spencer}
The \emph{transverse Spencer obstruction of \(X_3\) along the flat
\(12\)-face} is the obstruction to lifting this pointed boundary deformation
across
\(R_{123}\twoheadrightarrow R_{\partial}\), projected to the directed top
root line:
\[
 \operatorname{Sp}^{\perp}_{12\mid3}(X_3)
 \in
 I_{123}\otimes H^2(A)_{E_{14}}.
\]
After dividing by the distinguished generator \(u_1u_2u_3\), we also write
\(\operatorname{sp}^{\perp}_{12\mid3}(X_3)\) for the resulting class in
\(H^2(A)_{E_{14}}\).
\end{definition}

\begin{proposition}[Transverse Spencer equals the directed third-jet class]
\label{P1L-thm:transverse-spencer-massey}
For the pointed defining system of \cref{P1L-eq:massey-rep},
\[
 \boxed{
 \operatorname{Sp}^{\perp}_{12\mid3}(X_3)
 =\pm u_1u_2u_3\,m_{123}^{\rm AMM}E_{14}.}
\]
Equivalently,
\[
 \operatorname{sp}^{\perp}_{12\mid3}(X_3)
 =\pm m_{123}^{\rm AMM}E_{14}.
\]
For the Eisenstein triple of \cref{P1L-thm:amm-example} this obstruction is
nonzero.
\end{proposition}

\begin{proof}
The obstruction to a small-extension lift is the coefficient of the missing
kernel monomial in the Maurer--Cartan equation.  Here that monomial is
\(u_1u_2u_3\).  By \cref{P1L-thm:111-jet}, its directed \(E_{14}\)-coefficient is
exactly
\(\pm m_{123}^{\rm AMM}E_{14}\), with
\[
 m_{123}^{\rm AMM}
 =[\chi_1\smile a_{23}+a_{12}\smile\chi_3].
\]
The same theorem shows that every nonidentity ordering has zero
\(E_{14}\)-projection, so no additional ordered term enters the marked top
root.  Nonvanishing follows from \cref{P1L-thm:amm-example}.
\end{proof}

\begin{proposition}[Cartan--Spencer interpretation of \(\Theta_3\)]
\label{P1L-thm:theta3-spencer}
Fix the marked quadratic tensor \(Q_2\) and the ordered flat face
\(P_{12}\).  The class of the transverse obstruction modulo changes of
quadratic pointing is the restriction of the invariant ternary deformation
class:
\[
 \boxed{
 [\operatorname{sp}^{\perp}_{12\mid3}(X_3)]
 =\Theta_3(Q)(X_1,X_2,X_3)
 \quad\text{on the marked }E_{14}\text{-line}.}
\]
Thus the nonzero \((1,1,1)\) third jet is the first obstruction to extending
the pointed flat \(pq\)-Cartan germ coherently in the transverse third
direction.
\end{proposition}

\begin{proof}
An identity-linear change of minimal coordinates with quadratic Taylor term
\(F_2\) changes the raw ternary tensor by
\[
 Q_3\longmapsto Q_3+[Q_2,F_2].
\]
Changing the pairwise pointing changes the small-extension representative by
the corresponding \([Q_2,-]\)-boundary.  Therefore the invariant class of
the transverse obstruction lies in the deformation quotient
\[
 H^3_{\mathrm{def}}(H,Q_2)
 =\frac{\ker([Q_2,-]:C^{(3)}\to C^{(4)})}
 {[Q_2,C^{(2)}]}
\]
and is precisely the evaluation of \([Q_3]=\Theta_3(Q)\) on the ordered
triple.  The representative on the directed top-root line is identified in
\cref{P1L-thm:transverse-spencer-massey}.
\end{proof}

\section{Secondary class and Toda normal form}

In characteristic zero one may pass from an associative \(A_\infty\) model to a homotopy-Lie model by antisymmetrization.  In characteristic \(3\), however, the correct formal-moduli object is a partition Lie algebra, and one must distinguish primary operations on tangent homotopy from secondary coherence data \cite{BrantnerMathew,BCN}.  The same distinction applies to multiweight \((1,1,1)\).

Let \(\mathfrak g^\pi\) denote the partition-Lie object corresponding to the residual derived representation formal-moduli problem of the trivial rank-four representation together with the fixed ordered root-path marking.

\begin{proposition}[Primary \((1,1,1)\) operations vanish]\label{P1L-prop:primary-111}
Every natural primary partition-Lie operation on tangent homotopy that is homogeneous of marked multiweight \((1,1,1)\) vanishes on \((X_1,X_2,X_3)\).
\end{proposition}

\begin{proof}
Use the corresponding classification of primary operations.  A positive-characteristic unary or divided-power decoration contributes a common factor to all input multiplicities.  Since
\[
 \gcd(1,1,1)=1,
\]
no nontrivial such factor is possible.  Thus only ordinary length-three Lyndon Lie words remain.  Every length-three Lie word contains an inner binary bracket of two marked tangent classes, and all such brackets vanish by \cref{P1L-prop:pairwise-brackets}.  Hence every primary operation of the required multiweight vanishes.
\end{proof}

\begin{theorem}[Three-direction secondary Cartan class]\label{P1L-thm:secondary-111}
For the arithmetic triple of \cref{P1L-thm:amm-example}, the nonzero formal-moduli coefficient \eqref{P1L-eq:formal-111} determines a canonical nonzero secondary coherence class
\[
 \kappa^\pi_{1,1,1}(X_1,X_2,X_3)\neq0
\]
in the partition-Lie object \(\mathfrak g^\pi\).  It is not represented by a primary operation on tangent homotopy.
\end{theorem}

\begin{proof}
The derived representation formal-moduli functor retains its complete system of Artinian infinitesimal tests.  By \cref{P1L-thm:111-jet}, its restriction to the marked third-order test algebra has a nonzero coefficient on the \(u_1u_2u_3E_{14}\)-line.  Transport this marked third-jet coefficient across the formal-moduli/partition-Lie equivalence and denote the resulting coherence class by \(\kappa^\pi_{1,1,1}\).  If the class were zero, the partition-Lie object would recover a marked third jet with vanishing \((1,1,1)\) obstruction coefficient, contradicting \eqref{P1L-eq:formal-111}.  Finally \cref{P1L-prop:primary-111} shows that no primary operation can carry this information, so it is intrinsically secondary.
\end{proof}

\subsection{A one-cell normal form}

After choosing a point-set rectification of the partition-Lie object, the marked quotient of the third jet has an immediate one-cell model.

\begin{corollary}[Directed PD-cell normal form]\label{P1L-cor:pd-cell}
Normalize the one-dimensional obstruction line generated by \(m_{123}^{\rm AMM}E_{14}\).  After a compatible Brantner--Campos--Nuiten type rectification, the marked third-jet quotient admits a cofibrant normalized-chain presentation
\[
 T^{\PD}_{3,123}
 =
 \F_3[u_1,u_2,u_3]\langle\eta_{123}\rangle_{\PD}/(u_1,u_2,u_3)^4,
 \qquad
 |\eta_{123}|=-1,
\]
with
\begin{equation}\label{P1L-eq:pd-cell}
 d\eta_{123}=u_1u_2u_3.
\end{equation}
The continuous dual \(\eta_{123}^\vee\) represents the canonical secondary class \(\kappa^\pi_{1,1,1}\) in the localized homotopy category.
\end{corollary}

\begin{remark}
The literal chain \(\eta_{123}\) depends on the chosen rectification; the marked third-jet obstruction line and the resulting secondary class do not.
\end{remark}


\subsection{Universal Toda skeleton and controller dictionary}

We now compare the preceding arithmetic class with the formal controller
\[
 B_{pqr}=[D_p,H_{qr}]+[D_q,H_{rp}]+[D_r,H_{pq}]
\]
from the universal controller above.  The comparison is structural: the arithmetic example has residual coefficient order \(3\), while the Ore controller is built over \(\Z\) with exact order \(pqr\).

\begin{lemma}[Universal ternary Toda cycle]\label{P1L-lem:toda-cycle}
Let \((L,d,[-,-])\) be a cohomologically graded differential graded Lie algebra.  Let \(x_1,x_2,x_3\) be closed elements of degree one and suppose pairwise brackets are null-homologous, with chosen degree-one homotopies
\[
 dh_{ij}=[x_i,x_j].
\]
Then the cyclic expression
\begin{equation}\label{P1L-eq:toda}
 \beta(x_1,x_2,x_3)
 =
 [x_1,h_{23}]+[x_2,h_{31}]+[x_3,h_{12}]
\end{equation}
is closed.  Its class changes by the usual secondary indeterminacy when the \(h_{ij}\) are changed.
\end{lemma}

\begin{proof}
Apply \(d\) to \eqref{P1L-eq:toda}, use \(dx_i=0\) and substitute \(dh_{ij}=[x_i,x_j]\).  The resulting sum is the graded Jacobi identity.  Changing a pairwise homotopy by a closed term changes \(\beta\) by the standard sum of brackets with those closed corrections, which is precisely the secondary indeterminacy.
\end{proof}

\begin{proposition}[Homotopy-transfer interpretation]\label{P1L-prop:hpt-controller}
For a contraction of a dg Lie algebra onto cohomology, the transferred ternary bracket \(\ell_3\) is represented, up to convention signs and projection/inclusion maps, by the cyclic Toda expression \eqref{P1L-eq:toda}.  Therefore the Ore-controller pattern
\[
 [D_p,H_{qr}]+[D_q,H_{rp}]+[D_r,H_{pq}]
\]
is the universal chain skeleton of a transferred ternary Cartan class after pairwise brackets have been null-homotoped.
\end{proposition}

\begin{proof}
The tree formula for homotopy transfer of a dg Lie algebra in arity three is the sum over the three rooted binary trees obtained by first bracketing a pair, applying the contracting homotopy, and then bracketing with the remaining input.  Identifying \(D_i\) with the included closed inputs and \(H_{ij}\) with the chosen homotopies gives the displayed controller expression.
\end{proof}

\subsubsection{Controller dictionary}

The comparison can be summarized as follows:
\[
\boxed{
\begin{array}{ccc}
\text{pairwise-zero arithmetic linking}
&\longleftrightarrow&
 dh_{ij}=[x_i,x_j]
\\[1mm]
\text{triple Milnor/Massey class}
&\longleftrightarrow&
 [x_1,h_{23}]+[x_2,h_{31}]+[x_3,h_{12}]
\\[1mm]
\text{marked }(1,1,1)\text{ third jet}
&\longleftrightarrow&
 B_{pqr}\text{ secondary skeleton}
\\[1mm]
\text{PD top cell }d\eta=u_1u_2u_3
&\longleftrightarrow&
\text{pointed top-cell rectification.}
\end{array}}
\]
The first three rows are realized by the restricted-ramification Eisenstein calibration; exact composite order requires the separate ray-principal construction below.


\section{Exact-order arithmetic and Moore comparison}

The preceding sections used \(m_{123}^{\rm AMM}\) for the residual
restricted-ramification mod-\(3\) calibration.  From this section onward,
\(m_{123}\) denotes the pointed \(A_N\)-valued cocycle produced by the
ray-class construction.  The matrix-path formula is the same, but the two
arithmetic branches are logically independent.

The integral controller carries more information than the residual
Eisenstein calibration.  For a fixed odd coefficient order \(N\), its top
relation is
\[
 dC_{pqr}=N B_{pqr}.
\]
The arithmetic task is therefore to construct three marked Kummer
directions with pairwise linking trivializations whose pointed secondary
class has exact order \(N\), and then to compare that class with an integral
Moore antecedent.  The ray-principal construction below solves the
exact-order arithmetic existence problem; the Moore comparison of
\cref{P1L-sec:moore-comparison} supplies the external integral antecedent.

\begin{remark}[Relation with the Kim--Morishita triple symbol]
\label{P1L-rem:kim-morishita-relation}
Kim--Morishita define a composite-modulus triple symbol for tame characters
with disjoint ramification and suitable local pairwise-linking conditions
\cite{KimMorishita}.  Their global/local cochain construction is compatible
with the Heisenberg picture motivating the present pointed class.  It is not,
however, used in the ray-principal existence proof below: no one-place
factorization and no local unramified-trivialization proposition is an input
to \cref{P1L-thm:ray-class-primitive}.  We retain the symbol only as a comparison
interface with the exact-order pointed seed.
\end{remark}

\subsection{Unconditional ray-class realization of a primitive pointed class}

The finite ray-class endpoint can in fact be closed for the pointed problem.  The construction uses only Kummer theory, ray class fields, local duality and Chebotarev; we use the standard class-field conventions of \cite{NeukirchANT}.  We state it for arbitrary odd $N$ because no step uses squarefreeness until the final CRT readout.

\begin{lemma}[A deep ray modulus killing the $N$-adic places]\label{P1L-lem:deep-ray-modulus}
Let $F$ be a number field containing $\mu_N$, with $N$ odd.  There is an integral modulus $\mathfrak m_N$, supported only at the finite places $w\mid N$, such that
\[
 1+\mathfrak m_{N,w}\subset (F_w^\times)^N
 \qquad (w\mid N).
\]
Consequently, if $\alpha\in F^\times$ satisfies $\alpha\equiv1\pmod{\mathfrak m_N}$, the Kummer character of $F(\sqrt[N]{\alpha})/F$ is trivial at every place above $N$.
\end{lemma}

\begin{proof}
For every finite extension $F_w/\Q_\ell$, the quotient $F_w^\times/(F_w^\times)^N$ is finite.  Hence $(F_w^\times)^N$ is open and contains a sufficiently deep principal-unit subgroup $1+\mathfrak p_w^{a_w}$.  Take $\mathfrak m_N=\prod_{w\mid N}\mathfrak p_w^{a_w}$.  The last assertion is immediate from local Kummer theory.
\end{proof}

\begin{lemma}[Two ray-principal Kummer directions with zero cup]\label{P1L-lem:first-two-rays}
Let $F\supset\mu_N$ and fix $\mathfrak m_N$ as in \cref{P1L-lem:deep-ray-modulus}.  There exist distinct finite primes $v_1,v_2\nmid N$ and generators
\[
 v_i=(\pi_i)
\]
such that, for the order-$N$ Kummer characters $\chi_i$ attached to $\pi_i$,
\[
 S(\chi_i)=\{v_i\},\qquad
 \chi_1\smile\chi_2=0\in H^2(F,A_N).
\]
Moreover $\chi_2$ is locally trivial at $v_1$ and $\chi_1$ is locally trivial at $v_2$.
\end{lemma}

\begin{proof}
Let $H_{\mathfrak m_N}/F$ be the ray class field for $\mathfrak m_N$.  Choose $v_1\nmid N$ splitting completely in $H_{\mathfrak m_N}$.  Its ray class is trivial, hence $v_1=(\pi_1)$ with $\pi_1\equiv1\pmod{\mathfrak m_N}$.  The valuation of $\pi_1$ at $v_1$ is one, so its Kummer character has exact order $N$ and is tamely ramified at $v_1$; by \cref{P1L-lem:deep-ray-modulus} it is trivial above $N$, and at every other finite place it is unramified.

Put $K_1=F(\sqrt[N]{\pi_1})$ and enlarge the ray modulus to
\[
 \mathfrak m_1:=\mathfrak m_N v_1.
\]
Choose $v_2$ splitting completely in the compositum $K_1H_{\mathfrak m_1}$.  Then $v_2=(\pi_2)$ with $\pi_2\equiv1\pmod{\mathfrak m_1}$, so the Kummer character $\chi_2$ is trivial at $v_1$ and at all places above $N$.  Splitting of $v_2$ in $K_1$ gives $\loc_{v_2}\chi_1=0$.

At $v_1$ and $v_2$ the local cup vanishes because one factor is zero.  At a place not dividing $Nv_1v_2$ both characters are unramified, hence their cup inflates from the procyclic unramified quotient and vanishes in degree two.  The places above $N$ contribute zero by construction.  Global reciprocity gives injectivity of the $N$-torsion Brauer localization map, so $\chi_1\smile\chi_2=0$ globally.
\end{proof}

\begin{lemma}[Canonical Heisenberg lift of the first two directions]\label{P1L-prop:ray-heisenberg-lift}
For $v_1,v_2,\chi_1,\chi_2$ as in \cref{P1L-lem:first-two-rays}, the map
\[
 (\chi_1,\chi_2):G_F\longrightarrow A_N^2\simeq \overline D_N
\]
admits a lift
\[
 \rho_{12}:G_F\longrightarrow D_N.
\]
Every such lift is surjective.  If $L/F$ is the finite extension cut out by $\rho_{12}$ and $z\in D_N$ is the central element with central coordinate one, then
\[
 \Gal(L/F)\simeq D_N,\qquad [D_N,D_N]=\langle z\rangle\simeq A_N.
\]
Fix one lift once and for all; equivalently, fix its central-coordinate cochain $c_{12}$ as part of the pointed two-skeleton.
\end{lemma}

\begin{proof}
For the central extension
\[
 1\longrightarrow A_N\longrightarrow D_N\longrightarrow A_N^2\longrightarrow1,
\]
the obstruction to lifting \((\chi_1,\chi_2)\) is the cup class
\(\chi_1\smile\chi_2\), which vanishes by \cref{P1L-lem:first-two-rays}.
Thus a lift exists.  The two characters are jointly surjective: tame inertia at $v_1$ maps to $(1,0)$ after choosing a generator, while tame inertia at $v_2$ maps to $(0,1)$.  In $D_N$ the commutator of lifts of these two elements is the central generator $z$.  Hence the image contains the two abelian generators and the full center, so it is all of $D_N$.
\end{proof}

\begin{theorem}[Primitive ray-class three-prime construction]\label{P1L-thm:ray-class-primitive}
Let $N>1$ be odd and let $F$ be a number field containing $\mu_N$.  Then there exist three distinct finite primes
\[
 v_1,v_2,v_3\nmid N
\]
and order-$N$ tame Kummer characters
\[
 \chi_1,\chi_2,\chi_3:G_F\longrightarrow A_N
\]
with pairwise disjoint ramification supports $S(\chi_i)=\{v_i\}$ such that:
\begin{enumerate}
\item all three pairwise cup classes vanish in $H^2(F,A_N)$;
\item the first two characters admit a fixed surjective Heisenberg lift $\rho_{12}:G_F\twoheadrightarrow D_N$;
\item with $c_{12}$ the central-coordinate nullhomotopy fixed by this lift and with any $c_{23}$ satisfying $dc_{23}=\chi_2\smile\chi_3$, the pointed triple-Massey cocycle
\[
 m_{123}:=c_{12}\smile\chi_3+\chi_1\smile c_{23}
\]
(up to the global sign convention) has exact additive order $N$ in $H^2(F,A_N)$.
\end{enumerate}
More precisely, the third prime may be chosen so that
\[
 \inv_{v_3}(\loc_{v_3}m_{123})\in A_N^\times,
\]
and after normalizing the central generator and local reciprocity orientation this invariant is $1$.
\end{theorem}

\begin{proof}
Start with \cref{P1L-lem:first-two-rays,P1L-prop:ray-heisenberg-lift} and let $L/F$ be the resulting $D_N$-extension.  Let $R_L$ be its finite ramification set.  Choose a ray modulus $\mathfrak m_2$ which contains $\mathfrak m_N$, $v_1$, $v_2$, and every place of $R_L$, taking exponent one away from $N$ and keeping the deep exponents of \cref{P1L-lem:deep-ray-modulus} above $N$.  Let $H_{\mathfrak m_2}/F$ be the corresponding ray class field.

The intersection $L\cap H_{\mathfrak m_2}$ is abelian over $F$.  Since $z\in[D_N,D_N]$, the central generator acts trivially on every abelian quotient of $D_N$ and therefore fixes $L\cap H_{\mathfrak m_2}$.  Hence the pair
\[
 (z,1)\in \Gal(L/F)\times\Gal(H_{\mathfrak m_2}/F)
\]
defines an element of the fiber-product Galois group of $LH_{\mathfrak m_2}/F$.  Chebotarev supplies infinitely many primes $v_3$ unramified in this compositum with Frobenius $(z,1)$.

Because $v_3$ splits completely in the ray class field, its ray class is trivial.  Thus
\[
 v_3=(\pi_3),\qquad \pi_3\equiv1\pmod{\mathfrak m_2}.
\]
Let $\chi_3$ be the Kummer character of $\pi_3$.  It is tame of exact order $N$ at $v_3$, trivial at every place in $R_L$ and at every place above $N$, and unramified elsewhere.  Since the Frobenius of $v_3$ in $L/F$ is central, the two abelian coordinates vanish on the decomposition group at $v_3$:
\[
 \loc_{v_3}\chi_1=\loc_{v_3}\chi_2=0.
\]
The local audit of \cref{P1L-lem:first-two-rays} now gives
\[
 \chi_1\smile\chi_3=\chi_2\smile\chi_3=0.
\]
Choose $c_{23}$ with $dc_{23}=\chi_2\smile\chi_3$.

On the decomposition group at $v_3$, the fixed cochain $c_{12}$ is an unramified $A_N$-character whose value on Frobenius is the central coordinate of $z$, hence a generator of $A_N$.  The character $\chi_3$ is the primitive tame Kummer character of a local uniformizer.  Local Tate duality therefore gives
\[
 \inv_{v_3}(\loc c_{12}\smile\loc\chi_3)\in A_N^\times;
\]
a choice of the central generator and reciprocity orientation makes it equal to $1$.  Since $\loc_{v_3}\chi_1=0$, the second term $\chi_1\smile c_{23}$ has zero localization at $v_3$.  Therefore
\[
 \inv_{v_3}(\loc_{v_3}m_{123})=1.
\]
The coefficient group $A_N$ is killed by $N$, while the localization of $[m_{123}]$ has order exactly $N$.  Hence $[m_{123}]$ itself has exact order $N$.
\end{proof}

\begin{proposition}[Primewise Cartan readout of the exact-order seed]
\label{P1L-prop:ray-class-cartan-readout}
Keep the pointed ray-class datum of \cref{P1L-thm:ray-class-primitive}.  For every
prime \(\ell\mid N\), reduce the three characters and the fixed defining
system modulo \(\ell\), put \(k=\F_\ell\), and form
\[
 X_{1,\ell}=\bar\chi_1E_{12},\qquad
 X_{2,\ell}=\bar\chi_2E_{23},\qquad
 X_{3,\ell}=\bar\chi_3E_{34}
\]
in \(C^\bullet(G_F,\End_k(k^4))\).  Then all pairwise Cartan brackets vanish,
and the directed \(E_{14}\)-component of the cubic obstruction is
\[
 \boxed{
 \operatorname{ob}^{(3)}_{14}
 =\pm\bar m_{123,\ell}E_{14}
e0.}
\]
Its invariant class is the primewise secondary Cartan class
\[
 \Theta_{3,\ell}(X_{1,\ell},X_{2,\ell},X_{3,\ell})\ne0.
\]
Collectively, before primary reduction, the \(A_N\)-valued directed top-root
class \(m_{123}E_{14}\) has exact additive order \(N\).
\end{proposition}

\begin{proof}
Pairwise cup vanishing is part of \cref{P1L-thm:ray-class-primitive}, hence it
survives every primary reduction.  The matrix-unit calculation of
\cref{P1L-thm:111-jet,P1L-thm:theta3-spencer} is coefficient-independent: among the
six orderings only
\[
 E_{12}E_{23}E_{34}=E_{14}
\]
survives on the top-root line, and its scalar coefficient is precisely the
reduced defining-system cocycle \(\bar m_{123,\ell}\).  The normalized local
invariant of \(m_{123}\) is a unit of \(A_N\), so its reduction modulo every
\(\ell\mid N\) is nonzero.  The Cartan--Spencer quotient therefore gives the
displayed nonzero \(\Theta_{3,\ell}\).  The final exact-order statement is
\cref{P1L-thm:ray-class-primitive}.
\end{proof}

\begin{proposition}[Absolute-Galois Massey firewall]\label{P1L-thm:absolute-massey-firewall}
Assume that $N$ is squarefree and that $F$ is a number field containing
$\mu_N$.  For every prime $\ell\mid N$, the reduction of a defined
three-fold Massey product of classes in
$H^1(G_F,\F_\ell)$ contains $0$.  Consequently, by the Chinese remainder
decomposition of $A_N$, the unpointed triple product determined by
$(\chi_1,\chi_2,\chi_3)$ in the situation of
\cref{P1L-thm:ray-class-primitive} contains $0$.

The exact-order class $[m_{123}]$ of
\cref{P1L-thm:ray-class-primitive} is therefore not an essential unpointed
absolute-Galois triple Massey product.  It is the value selected by the fixed
defining system $(c_{12},c_{23})$, with $c_{12}$ fixed by the chosen
surjective Heisenberg lift.  Its nonzero local invariant and exact order are
statements in the pointed/marked strong-descent fibre.
\end{proposition}

\begin{proof}
Min\'a\v{c}--T\^an prove that the absolute Galois group of every field has
the vanishing triple Massey product property for every prime coefficient
\cite{MinacTan}.  Thus each defined \(\F_\ell\)-valued triple product contains
\(0\).  For higher-length Massey products over number fields, see
Harpaz--Wittenberg \cite{HarpazWittenberg}.  Since the pairwise cup products in
\cref{P1L-thm:ray-class-primitive} vanish, every prime reduction is defined and
contains $0$.  For squarefree $N$ the cochain algebra with coefficients
$A_N$ is the product of its $\F_\ell$-cochain algebras.  Componentwise
defining systems giving zero therefore assemble, by CRT, to an $A_N$ defining
system giving zero.

On the other hand, the theorem fixes the central-coordinate cochain
$c_{12}$ coming from $\rho_{12}$ and then fixes a compatible $c_{23}$.
The resulting cocycle $m_{123}$ is one selected value of the multi-valued
Massey product.  Changing the defining system is precisely the ordinary
Massey indeterminacy, so the absolute vanishing theorem does not force this
chosen pointed value to be zero.
\end{proof}

\begin{corollary}[Pointed Massey value lies in the ordinary indeterminacy]
\label{P1L-cor:pointed-massey-decomposable}
Assume \(N\) is squarefree and keep the pointed defining system
\((c_{12},c_{23})\) of \cref{P1L-thm:ray-class-primitive}.  There exist
\(\mu,\nu\in H^1(G_F,A_N)\) such that
\[
 \boxed{
 [m_{123}]
 =\mu\smile\chi_3+\chi_1\smile\nu.}
\]
Hence the exact-order pointed value is cohomologically decomposable; its
nontrivial information is the chosen point of the defining-system torsor.
At the distinguished place \(v_3\), where
\(\loc_{v_3}\chi_1=0\), one has
\[
 \inv_{v_3}(\loc_{v_3}m_{123})
 =\inv_{v_3}(\loc_{v_3}\mu\smile\loc_{v_3}\chi_3),
\]
so the nonzero local invariant is carried by the first indeterminacy
summand.
\end{corollary}

\begin{proof}
By \cref{P1L-thm:absolute-massey-firewall}, choose an \(A_N\)-defining system
\((c_{12}^0,c_{23}^0)\) whose triple-Massey value represents zero in
\(H^2(G_F,A_N)\).  Then
\(c_{12}-c_{12}^0\) and \(c_{23}-c_{23}^0\) are cocycles.  Put
\(\mu=[c_{12}-c_{12}^0]\) and \(\nu=[c_{23}-c_{23}^0]\).  Subtracting the
two defining-system values gives the displayed decomposition in cohomology.
The local statement follows from \(\loc_{v_3}\chi_1=0\).
\end{proof}


\begin{remark}[Restricted-ramification firewall]
\label{P1L-rem:amm-restricted-firewall}
The AMM class lives on the restricted-ramification pro-\(3\) quotient
\(\mathfrak G_{k,S}(3)\) and is detected by its degree-\(27\) Heisenberg
extension \cite{AMM}; it is therefore compatible with
\cref{P1L-thm:absolute-massey-firewall}, which concerns the absolute Galois
triple product.
\end{remark}

\begin{corollary}[Infinite congruence-eight admissible family]
\label{P1L-prop:amm-eight-mod-nine-family}
Let
\[
 F_0=\mathbf Q(\zeta_3)
\]
and let $p_1,p_2,p_3$ be distinct rational primes with
\[
 p_i\equiv8\pmod 9.
\]
Then each $\mathfrak p_i=(p_i)$ is a prime of $F_0$ with
$\operatorname N\mathfrak p_i=p_i^2\equiv1\pmod9$.  After replacing the
prime generator by $\pi_i=-p_i$, one has the standard primary normalization
\[
 \pi_i\equiv1\pmod{(3\sqrt{-3})}.
\]
Moreover, for every $i\ne j$,
\[
 \boxed{
 \left(\frac{\pi_i}{\pi_j}\right)_3=1.}
\]
Hence every such triple satisfies the pairwise hypotheses for the AMM triple
cubic residue symbol and restricted-ramification triple-Massey construction.
There are infinitely many admissible triples of this form.
\end{corollary}

The congruence condition gives an infinite admissible family, not automatic triple-symbol nonvanishing; the explicit \((17,53,71)\) triple remains the certified nonzero calibration.

\begin{proof}
Since $p_j\equiv2\pmod3$, it is inert in $\mathbf Q(\zeta_3)$ and
$\mathcal O_{F_0}/(p_j)\simeq\mathbf F_{p_j^2}$.  The rational residue class
of $\pi_i=-p_i$ lies in the subfield $\mathbf F_{p_j}^{\times}$.  Therefore
\[
 \pi_i^{(p_j^2-1)/3}
 =\left(\pi_i^{p_j-1}\right)^{(p_j+1)/3}
 =1,
\]
because $3\mid p_j+1$.  This is exactly the triviality of the cubic residue
symbol.  The norm congruence follows from $p_j^2\equiv1\pmod9$.  Finally,
$(3\sqrt{-3})$ is, up to a unit, the third power of the prime above $3$, and
its intersection with $\mathbf Z$ is $9\mathbf Z$; hence
$-p_i\equiv1\pmod9$ gives the displayed primary normalization.  Dirichlet's
theorem supplies infinitely many rational primes in the class $8\pmod9$.
\end{proof}

\paragraph{Why the ray-class theorem is pointed.}
The mechanism of \cref{P1L-thm:ray-class-primitive} is exactly the distinction in
\cref{P1L-thm:absolute-massey-firewall}.  The chosen Heisenberg lift fixes a
point of the defining-system torsor before the top class is evaluated.
Forgetting that point returns to the ordinary multi-valued Massey product,
which contains $0$.  Hence ``pointed'' is not a disclaimer attached to the
theorem; it is the structure that carries the nonzero selected value.


\begin{corollary}[The \(N=105\) calibration]
\label{P1L-cor:N105}
Take
\[
 \boxed{N=105=3\cdot5\cdot7,\qquad F=\Q(\zeta_{105}).}
\]
Then \cref{P1L-thm:ray-class-primitive} produces three distinct arithmetic
primes \(v_1,v_2,v_3\) and a pointed class
\[
 [m_{123}]\in H^2(F,\Z/105\Z)
\]
of exact order \(105\).  Under the directed root path
\[
 X_1=\chi_1E_{12},\qquad X_2=\chi_2E_{23},\qquad X_3=\chi_3E_{34},
\]
the marked top-root third jet has exact order \(105\).  Its reductions
modulo \(3,5,7\) are nonzero and lie in secondary coherence relative to the
fixed two-skeleton.
\end{corollary}

\begin{proof}
The field \(\Q(\zeta_{105})\) contains \(\mu_{105}\), so
\cref{P1L-thm:ray-class-primitive} applies.  CRT converts exact order \(105\)
into nonvanishing on all three primary components.  The directed-root and primewise secondary statements are
\cref{P1L-prop:ray-class-cartan-readout}.
\end{proof}


\subsection{Integral Moore and residual Galois realization}\label{P1L-sec:moore-comparison}

The ray-class theorem and \cref{P1L-prop:ray-class-cartan-readout} leave one
distinguished exact-order cyclic arithmetic line,
\[
 \langle[m_{123}]\rangle\simeq A_N.
\]
The next task is to identify that same line with the universal Moore
generator before introducing any motivic enhancement.  The exact-order
arithmetic class and the integral theta/root-stack class should not be
compared by choosing an arbitrary isomorphism between two cyclic groups of
order $N$.  Both are already realizations of the same pointed Moore sector of the universal controller.  This gives a chain-level comparison span and separates the part of the carry problem which is now closed from the genuinely stronger internal Galois lift.

\subsubsection{The constant cyclic Moore sector}

Write
\begin{equation}\label{P1L-eq:cyclic-moore}
 \mathcal K_N^{\mathrm{cyc}}
 :=[\Z C\xrightarrow{\,N\,}\Z B],
\end{equation}
with $C$ in cohomological degree $1$ and $B$ in degree $2$.  It is the constant-coefficient restriction of the top-support controller complex
\[
 [\Hh C_{pqr}\xrightarrow{\,N\,}\Hh B_{pqr}]
\]
from the universal Ore model.  Thus
\[
 H^2(\mathcal K_N^{\mathrm{cyc}})\simeq \Z/N\Z,
 \qquad [B]\longleftrightarrow1.
\]
The point of retaining the degree-one cell $C$ is that it remembers the chosen nullhomotopy of $NB$ before reduction modulo $N$.

\begin{lemma}[Pointed Galois transgression of Moore sector]\label{P1L-prop:galois-moore-map}
Let $(\chi_1,\chi_2,\chi_3;c_{12},c_{23})$ be the pointed ray-class datum supplied by \cref{P1L-thm:ray-class-primitive}.  Let
\[
 m_{123}=c_{12}\smile\chi_3+\chi_1\smile c_{23}
 \in Z^2(G_{F,S},A_N)
\]
be the chosen cocycle representative.  There is a chain map
\begin{equation}\label{P1L-eq:galois-moore-map}
 \Theta_N^{\mathrm{Gal}}:
 \mathcal K_N^{\mathrm{cyc}}
 \longrightarrow C^\bullet(G_{F,S},A_N)
\end{equation}
defined by
\[
 C\longmapsto0,
 \qquad
 B\longmapsto m_{123}.
\]
It induces an isomorphism
\[
 (\Z/N)[B]\xrightarrow{\sim}\langle[m_{123}]\rangle
 \subset H^2(G_{F,S},A_N).
\]
\end{lemma}

\begin{proof}
The defining-system identity gives $dm_{123}=0$.  The only nonzero differential in \eqref{P1L-eq:cyclic-moore} is $dC=NB$, whose image under \eqref{P1L-eq:galois-moore-map} is
\[
 N m_{123}=0
\]
because the target coefficients are $A_N=\Z/N\Z$.  Hence the formulas define a chain map.  By \cref{P1L-thm:ray-class-primitive}, the cohomology class of $m_{123}$ has exact additive order $N$, so the induced map from the cyclic group generated by $[B]$ is injective and has image exactly $\langle[m_{123}]\rangle$.
\end{proof}

\subsubsection{The abstract Moore carrier and realization boundary}

The cyclic Moore sector has two geometric realizations.  First, after fixing the theta-twisted pairwise norm comparisons, the pointed Picard rectification complex has a canonical exact-order class
\[
 \omega_\Theta\in H^2\C^\bullet_{\mathrm{Pic,pt}}(S,N),
 \qquad
 H^2\C^\bullet_{\mathrm{Pic,pt}}(S,N)
 \simeq (\Hh/N)\omega_\Theta,
\]
whose multiplicative form is the tautological theta root $\xi$.  Restriction to constant coefficients gives a chain map
\begin{equation}\label{P1L-eq:pic-moore-map}
 \Theta_N^{\Theta}:\mathcal K_N^{\mathrm{cyc}}
 \longrightarrow\C^\bullet_{\mathrm{Pic,pt}}(S,N),
 \qquad C\mapsto0,\quad B\mapsto\omega_\Theta,
\end{equation}
and an isomorphism on the constant cyclic $H^2$ line.

Second, the ranked root-stack localization sector contains actual integral classes $\Gamma,\Omega$ with
\begin{equation}\label{P1L-eq:root-carry-v04}
 \boxed{d\Gamma=N\Omega.}
\end{equation}
Let
\[
 \mathcal K_N^{\mathrm{root}}
 :=[\Z\Gamma\xrightarrow{\,N\,}\Z\Omega].
\]
The assignment
\begin{equation}\label{P1L-eq:root-moore-iso}
 \Theta_N^{\mathrm{root}}:
 \mathcal K_N^{\mathrm{cyc}}\xrightarrow{\sim}
 \mathcal K_N^{\mathrm{root}},
 \qquad C\mapsto\Gamma,\quad B\mapsto\Omega,
\end{equation}
is an isomorphism of two-term complexes.  In the full prime-sensitive root-stack incidence algebra the controller map may differ from this normalized cyclic map by the factor
\[
 g_S=\gcd(p^2-1,q^2-1,r^2-1),
\]
but $\gcd(g_S,N)=1$, so it induces the same oriented cyclic line in $H^2$ up to the unique unit forced by the full incidence normalization .

\begin{corollary}[Integral geometric carry for the arithmetic cyclic class]\label{P1L-cor:external-carry-closed}
The pointed arithmetic class of \cref{P1L-thm:ray-class-primitive} admits a natural \emph{external arithmetic-geometric} integral carry realization.  Under the normalized cyclic comparison, the residual class $[m_{123}]$ corresponds to the root-stack Moore class $[\Omega]$, and the actual geometric class $\Gamma$ satisfies
\[
 \boxed{d\Gamma=N\Omega.}
\]
For $N=105$ this gives
\[
 d\Gamma=105\,\Omega,
 \qquad
 [m_{123}]\longleftrightarrow[\Omega],
 \qquad
 \operatorname{ord}[m_{123}]=\operatorname{ord}[\Omega]=105.
\]
Thus the integral carry existence problem is closed at the common cyclic-controller level.  What is not produced is an integral Galois cochain whose image under a natural comparison functor is $\Gamma$.
\end{corollary}

\subsubsection{Compatibility of the Heisenberg centre}

The comparison above uses the universal Moore generator rather than an arbitrary choice of a generator in each target.  The central normalizations in the two arithmetic constructions are compatible as well.  Fix the primitive root used to identify
\[
 A_N\xrightarrow{\sim}\mu_N,
 \qquad 1\longmapsto\zeta_N.
\]
The ray-class construction fixes the standard Heisenberg group
\[
 D_N=U_3(A_N)
\]
and its central generator $z$ with central coordinate $1$.  The theta construction fixes a symplectic level-$N$ pair $(P,Q)$ and has commutator $e_N(P,Q)=\zeta_N$; on the finite-flat root base this becomes the tautological root $\xi$ and specializes to $\zeta_N$ on the cyclotomic generic point .

\begin{lemma}[Central-generator compatibility]\label{P1L-prop:center-compatibility}
Under $1\mapsto\zeta_N$, the normalization used in \cref{P1L-thm:ray-class-primitive}
\[
 \rho_{12}(\operatorname{Frob}_{v_3})=z,
 \qquad
 \inv_{v_3}(\loc m_{123})=1
\]
and the theta/root-stack normalization
\[
 \Delta_{pqr}=\xi,
 \qquad
 \omega_\Theta\longleftrightarrow1
\]
select the same oriented $A_N$-central line.  In particular, the cyclic-line comparison contains no additional undetermined scalar unit after these normalizations have been fixed.
\end{lemma}

\begin{proof}
In the ray-class extension, the value of the central-coordinate cochain on $\operatorname{Frob}_{v_3}$ is $1\in A_N$ by construction, and local Tate duality pairs it with the primitive tame third character to give invariant $1$.  In the theta extension, the chosen symplectic basis has commutator $\zeta_N$, so its additive character is again $1\in A_N$.  The pointed Tate--norm defect is the tautological root $\xi$, whose cyclotomic specialization is $\zeta_N$.  Thus both conventions are the image of the same chosen generator of $A_N$.
\end{proof}


\subsubsection{CRT ownership is impossible, but support labelling is harmless}

The arithmetic direction triad and the CRT coefficient triad have different
roles.  The following multilinearity argument makes the separation
structural rather than conventional.  For a three-prime squarefree
coefficient order write
\[
 N=\ell_1\ell_2\ell_3,\qquad
 A_N\simeq A_{\ell_1}\times A_{\ell_2}\times A_{\ell_3},
\]
with orthogonal CRT idempotents \(e_1,e_2,e_3\).  The symbols \(p,q,r\)
remain the ordered controller-direction labels.

\begin{lemma}[Orthogonal-CRT directions kill multilinear three-jets]\label{P1L-lem:crt-diagonal-no-go}
Let $T:V_1\otimes_{A_N}V_2\otimes_{A_N}V_3\to W$ be any $A_N$-trilinear operation.  If
\[
 x_1=e_1x_1,\qquad x_2=e_2x_2,\qquad x_3=e_3x_3,
\]
then $T(x_1,x_2,x_3)=0$.  The same conclusion holds for every $A_N$-multilinear Taylor coefficient of a strong-Cartan third jet.
\end{lemma}

\begin{proof}
By $A_N$-multilinearity,
\[
 T(x_1,x_2,x_3)=e_1e_2e_3T(x_1,x_2,x_3)=0.
\]
\end{proof}

\begin{corollary}[No CRT ownership of the three arithmetic directions]
\label{P1L-cor:two-triads-no-ownership}
For a nonzero three-direction Cartan jet, the three arithmetic directions
cannot be assigned to three mutually orthogonal CRT coefficient factors.
In particular, the arithmetic triad \((v_1,v_2,v_3)\) and the coefficient
triad \((\ell_1,\ell_2,\ell_3)\mid N\) admit no canonical
identification by primary ownership.  Ordered support labels are nevertheless harmless provided every
primary coefficient component sees all three arithmetic directions.
\end{corollary}

\begin{proof}
A primary-ownership assignment would put the three inputs in orthogonal CRT
idempotent summands, forcing every trilinear third-jet coefficient to vanish
by \cref{P1L-lem:crt-diagonal-no-go}.  This contradicts the nonzero primewise directed class of
\cref{P1L-prop:ray-class-cartan-readout}.
\end{proof}


The distinction matters because the universal Ore controller itself contains no relation asserting that left whiskering by $D_\ell$ multiplies an edge by $\ell$ or $\ell^2$.  Its defining data are the squarefree support rule, Ore semilinearity, the equations
\[
 dH_{\ell m}=[D_\ell,D_m],
\qquad
 dC_{pqr}=NB_{pqr},
\]
and the cyclic definition of $B_{pqr}$.  The factors $\ell$ in the theta incidence model and $\ell^2$ in the ramification-corrected root-stack model are \emph{target-specific realization laws} .  After specialization $q\mapsto1$, the Habiro substitutions act trivially on constant coefficients.  Therefore a residual Galois realization need not manufacture those numerical weights internally.


\subsubsection{Full residual controller realization}\label{P1L-sec:full-galois-controller}

The preceding observation allows the individual controller generators to be realized directly.  The only apparent obstacle is grading: the raw controller has
\[
 |D_\ell|=0,\qquad |H_{\ell m}|=-1,\qquad |B_{pqr}|=-1,
\]
whereas the arithmetic directions and pairwise defining-system cochains have ordinary cohomological degree one.  The support grading supplies the required desuspension.

Put
\[
 A_N=\Z/N\Z,
 \qquad
 \mathscr A_N=C^\bullet\!\left(G_{F,S},\End_{A_N}(A_N^4)\right).
\]
Let
\[
 \Lambda^-_S:=\Lambda_{A_N}(\epsilon_p,\epsilon_q,\epsilon_r),
 \qquad |\epsilon_\ell|=-1,
\]
with zero differential, and define the support-desuspended arithmetic algebra
\[
 \boxed{\mathscr G^{\sharp}_{S,N}:=\mathscr A_N\otimes_{A_N}\Lambda^-_S}
\]
with the standard graded tensor-product multiplication.  Thus an ordinary $n$-cochain carrying support $T$ has total degree $n-|T|$, and repeated support vanishes automatically.

Let $(\chi_1,\chi_2,\chi_3;c_{12},c_{23})$ be the pointed ray-class datum of \cref{P1L-thm:ray-class-primitive}.  Set
\[
 X_1=\chi_1E_{12},\qquad
 X_2=\chi_2E_{23},\qquad
 X_3=\chi_3E_{34},
\]
\[
 h_{12}=c_{12}E_{13},\qquad
 h_{23}=c_{23}E_{24}.
\]
Then
\[
 dh_{12}=X_1X_2,\qquad dh_{23}=X_2X_3,
\]
while both ordered products between $X_3$ and $X_1$ vanish literally by matrix-unit multiplication.

Write
\[
 \mathscr U^{\mathrm{const}}_{S,N}
 :=\mathscr U_S\otimes_{\mathcal H,\,\operatorname{ev}_1}A_N,
\]
so the Ore substitutions become the identity on coefficients.

\begin{proposition}[Residual controller realization]\label{P1L-thm:full-galois-controller}
There is a dg-algebra homomorphism
\[
 \boxed{
 \rho_{\mathrm{Gal}}:
 \mathscr U^{\mathrm{const}}_{S,N}
 \longrightarrow
 \mathscr G^{\sharp}_{S,N}}
\]
defined on generators by
\[
 \begin{aligned}
 D_p&\longmapsto X_1\epsilon_p,
 &D_q&\longmapsto X_2\epsilon_q,
 &D_r&\longmapsto X_3\epsilon_r,\\
 H_{pq}&\longmapsto-h_{12}\epsilon_p\epsilon_q,
 &H_{qr}&\longmapsto-h_{23}\epsilon_q\epsilon_r,
 &H_{rp}&\longmapsto0,\\
 C_{pqr}&\longmapsto0.
 \end{aligned}
\]
It satisfies
\[
 \boxed{
 \rho_{\mathrm{Gal}}(B_{pqr})
 =m_{123}E_{14}\epsilon_p\epsilon_q\epsilon_r,}
\]
where
\[
 m_{123}=c_{12}\smile\chi_3+\chi_1\smile c_{23}.
\]
For the primitive ray-class construction this top class has exact additive order $N$.
\end{proposition}

\begin{proof}
The degrees match:
\[
 \deg(X_i\epsilon_i)=1-1=0,
 \qquad
 \deg(h_{ij}\epsilon_i\epsilon_j)=1-2=-1.
\]
Since $\epsilon_i^2=0$, every repeated-support relation in the controller maps to zero.  Constant coefficients are central and the specialized Ore substitutions are the identity.

For the pair $(p,q)$, the graded tensor-product sign gives
\[
 (X_1\epsilon_p)(X_2\epsilon_q)
 =-X_1X_2\epsilon_p\epsilon_q,
\]
whereas $X_2X_1=0$ because $E_{23}E_{12}=0$.  Hence
\[
 [\rho(D_p),\rho(D_q)]
 =-X_1X_2\epsilon_p\epsilon_q
 =d(-h_{12}\epsilon_p\epsilon_q).
\]
The $(q,r)$ calculation is identical.  For $(r,p)$ both matrix products vanish, so $H_{rp}\mapsto0$ respects its differential.

For the cyclic class, the only nonzero complementary products are
\[
 [\rho(D_p),\rho(H_{qr})]
 =\chi_1\smile c_{23}\,E_{14}\epsilon_p\epsilon_q\epsilon_r
\]
and
\[
 [\rho(D_r),\rho(H_{pq})]
 =c_{12}\smile\chi_3\,E_{14}\epsilon_p\epsilon_q\epsilon_r.
\]
The middle term is zero.  Their sum is the displayed $m_{123}$-class.  Finally,
\[
 d\rho(C_{pqr})=0
 =N\,m_{123}E_{14}\epsilon_p\epsilon_q\epsilon_r
 =\rho(dC_{pqr})
\]
because the target coefficients are $A_N$.  Exact order is \cref{P1L-thm:ray-class-primitive}.
\end{proof}

\begin{remark}[What the exterior support factor does]\label{P1L-rem:support-desuspension}
The symbols $\epsilon_T$ do not add arithmetic information.  They implement, multiplicatively, the same support-degree shift that places $D_\ell,H_{\ell m},C_{pqr}$ in total obstruction degree one and $B_{pqr}$ in degree two in the controller's additive support totalization.  In particular \cref{P1L-thm:full-galois-controller} resolves the degree mismatch without identifying a Toda degree with an ordinary Galois cochain degree by fiat.
\end{remark}

\begin{remark}[The residual incidence problem is closed]\label{P1L-rem:residual-incidence-closed}
At constant coefficients the pointed ray-class source realizes all of the controller generators $D_\ell,H_{\ell m},C_{pqr}$ and the cyclic class $B_{pqr}$ in one dg algebra.  Thus the full residual incidence comparison is not an additional existence problem.  What is absent from the Galois target is only a nonzero integral image of the carry cell $C_{pqr}$.
\end{remark}


\subsubsection{Common-source arithmetic--geometric incidence}\label{P1L-sec:common-source-incidence}

The root-stack construction of  provides a second full controller realization.  Put
\[
 \mathscr U^{\mathrm{const}}_{S,\Z}
 :=\mathscr U_S\otimes_{\mathcal H,\,\operatorname{ev}_1}\Z
\]
and define the constant-coefficient root target
\[
 \boxed{
 \mathscr R^{\mathrm{const}}_{S,N}
 :=
 \mathscr T^{\mathrm{root}}_{S,N}
 \otimes_{\mathcal H,\,\operatorname{ev}_1}\Z.}
\]
The prime-sensitive root-stack map base-changes to
\[
 \rho_{\mathrm{root}}:
 \mathscr U^{\mathrm{const}}_{S,\Z}
 \longrightarrow
 \mathscr R^{\mathrm{const}}_{S,N}.
\]  It has
\[
 \rho_{\mathrm{root}}(B_{pqr})=g_S\Omega,
 \qquad
 \rho_{\mathrm{root}}(C_{pqr})=g_S\Gamma,
 \qquad
 d\Gamma=N\Omega,
\]
where
\[
 g_S=\gcd(p^2-1,q^2-1,r^2-1),\qquad \gcd(g_S,N)=1.
\]
The Galois map of \cref{P1L-thm:full-galois-controller} is viewed as a $\Z$-algebra map by reduction $\Z\to A_N$.

\begin{proposition}[Common-source incidence span]\label{P1L-thm:common-source-full-incidence}
There is a diagram of dg algebras
\[
 \boxed{
 \mathscr G^{\sharp}_{S,N}
 \ \xleftarrow{\ \rho_{\mathrm{Gal}}\ }\ 
 \mathscr U^{\mathrm{const}}_{S,\Z}
 \ \xrightarrow{\ \rho_{\mathrm{root}}\ }\ 
 \mathscr R^{\mathrm{const}}_{S,N}.}
\]
Equivalently, the product target carries the diagonal realization
\[
 \rho_{\mathrm{mix}}
 :=(\rho_{\mathrm{Gal}},\rho_{\mathrm{root}}):
 \mathscr U^{\mathrm{const}}_{S,\Z}
 \longrightarrow
 \mathscr G^{\sharp}_{S,N}\times\mathscr R^{\mathrm{const}}_{S,N}.
\]
On the top pair it satisfies
\[
 \boxed{
 B_{pqr}\longmapsto
 \bigl(m_{123}E_{14}\epsilon_p\epsilon_q\epsilon_r,\,g_S\Omega\bigr),}
\]
\[
 \boxed{
 C_{pqr}\longmapsto(0,g_S\Gamma),
 \qquad
 d(0,g_S\Gamma)
 =N\bigl(m_{123}E_{14}\epsilon_{pqr},g_S\Omega\bigr).}
\]
The resulting top cohomology class has exact order $N$.
\end{proposition}

\begin{proof}
Both arrows are dg-algebra maps from the same constant-coefficient controller, so their product is one.  The displayed differential follows from $Nm_{123}=0$ in the Galois factor and $d\Gamma=N\Omega$ in the root-stack factor.  The Galois component already has exact order $N$, while multiplication by $g_S$ is a unit modulo $N$ on the root-stack cyclic line.  Hence the product class also has exact order $N$.
\end{proof}

\begin{remark}[Why a span is the correct comparison notion]\label{P1L-rem:span-not-direct-map}
The Galois and root-stack targets encode different lower-support geometry.  There is no reason for a canonical target-to-target map to send a Kummer defining-system cochain to a root-stack unit loop or a prime-power correspondence.  What is canonical after the markings are fixed is that both are realizations of the same universal pointed controller.  The common-source span records exactly this statement and avoids manufacturing an artificial direct map between unrelated carriers.
\end{remark}

\paragraph{No ordinary integral Galois carry on the Kummer lattice.}

The remaining distinction between the two factors is not merely technical.  The root-stack carry cannot be replaced by an ordinary integral Kummer-lattice carry inside Galois cochains.

The residual Galois Moore map realizes the exact-order class but sends the
integral Moore antecedent to zero; exact order in
\(H^2(G_{F,S},A_N)\) alone does not produce an integral Galois lift.

\begin{lemma}[Internal Kummer carry obstruction]\label{P1L-thm:internal-galois-carry-no-go}
Let $N$ be odd and squarefree, and let $m_{123}$ be the primitive pointed class of \cref{P1L-thm:ray-class-primitive}.  Fix $\ell\mid N$ and identify its $\ell$-primary coefficient line with $\mu_\ell$ using the chosen roots of unity in $F$.  At the distinguished arithmetic place $v_3$, the localization of $m_{123}$ has nonzero invariant in
\[
 H^2(F_{v_3},\mu_\ell)\simeq\Z/\ell\Z.
\]
There do not exist a cocycle
\[
 \widetilde b\in Z^2(G_{F,S},\Z_\ell(1))
\]
and a cochain
\[
 \widetilde C\in C^1(G_{F,S},\Z_\ell(1))
\]
such that the class $[\widetilde b]$ reduces to the $\ell$-primary component of $m_{123}$ and
\[
 d\widetilde C=N\widetilde b.
\]
In particular, the natural $\Z_\ell(1)$ Galois lattice cannot itself carry the controller's exact-order integral Moore relation.
\end{lemma}

\begin{proof}
Local Tate duality and the local invariant map give
\[
 H^2(F_{v_3},\Z_\ell(1))\simeq\Z_\ell,
\]
compatibly with reduction to
\[
 H^2(F_{v_3},\mu_\ell)\simeq\Z/\ell\Z
\]
\cite{MilneCFT}.  Since the residual localization is primitive, every integral lift has local invariant congruent to a nonzero element modulo $\ell$; hence its invariant is an $\ell$-adic unit.  In particular the localized lift is not $\ell$-power torsion.

If $d\widetilde C=N\widetilde b$, then $N[\widetilde b]=0$ in global cohomology and therefore after localization.  Since $N/\ell$ is a unit in $\Z_\ell$, this would make the localized class $\ell$-torsion.  But $H^2(F_{v_3},\Z_\ell(1))\simeq\Z_\ell$ is torsion-free and the class is nonzero, a contradiction.
\end{proof}

\begin{corollary}[The integral target must be relative or Moore-like]\label{P1L-cor:moore-target-necessary}
For the primitive three-direction arithmetic class, an integral realization of $dC=NB$ cannot live in ordinary Kummer-lattice Galois cohomology.  Any natural integral enhancement must enlarge the target by a relative/cofiber/stacky direction which creates an exact-order Moore class.  The ranked root-stack localization is therefore not merely one convenient external repair: it has the correct unavoidable homological shape.
\end{corollary}



\section{FSC transition compression and the canonical
\texorpdfstring{\(5\!-\!3\!-\!2\)}{5-3-2} descent}
\label{P1L-sec:ordered-fsc-compression}

Fix a CRT factor \(k=\F_\ell\), with \(\ell\in\{3,5,7\}\).  The
finite-controller defining-system tangent fits into
\begin{equation}\label{P1L-eq:v31-D-F-K}
 0\longrightarrow \mathsf F_\ell
 \longrightarrow \mathsf D_\ell
 \longrightarrow K_{\chi,\ell}
 \longrightarrow0,
 \qquad
 (\dim\mathsf D_\ell,\dim\mathsf F_\ell,\dim K_{\chi,\ell})=(13,6,7).
\end{equation}
The balanced theta--Moore residue is the canonical reversal-fixed line
\begin{equation}\label{P1L-eq:v31-Ltheta}
 L_{\Theta,\ell}
 =\mathsf P_\ell^{\sigma=+1}
 =\ker(\mathcal E_H)
 \subset\mathsf F_\ell\subset\mathsf D_\ell.
\end{equation}
Quotienting by this retained secondary state gives
\begin{equation}\label{P1L-eq:v31-Qcomp}
 \mathsf Q^{\rm comp}_{3,\ell}
 :=\mathsf D_\ell/L_{\Theta,\ell},
 \qquad
 \dim\mathsf Q^{\rm comp}_{3,\ell}=12,
\end{equation}
with the canonical filtration
\begin{equation}\label{P1L-eq:v31-Q-filtration}
 0\longrightarrow
 \mathsf W_{5,\ell}:=\mathsf F_\ell/L_{\Theta,\ell}
 \longrightarrow\mathsf Q^{\rm comp}_{3,\ell}
 \longrightarrow K_{\chi,\ell}
 \longrightarrow0.
\end{equation}

\begin{remark}[Canonical \(5\!-\!3\!-\!2\) descent]
\label{P1L-rem:canonical-5-3-2-descent}
On \(\mathsf F_\ell\simeq H^1(\Gamma_\ell,k)^{\oplus2}\), strict marked
path reversal exchanges the two summands through endpoint reversal.  Its
eigenspace dimensions are \((3_+,3_-)\), and the fixed line
\(L_{\Theta,\ell}\) lies in the positive eigenspace.  Therefore the induced
involution on \(\mathsf W_{5,\ell}\) has type \((2_+,3_-)\).  Since
\(2\in k^\times\), the idempotents
\[
 e_\pm=\frac{1\pm\sigma}{2}
\]
are canonical relative to the marked reversal.  Setting
\[
 \mathsf W_{2,\ell}:=e_+\mathsf W_{5,\ell},
 \qquad
 \mathsf W_{3,\ell}:=e_-\mathsf W_{5,\ell}
\]
gives the canonical split exact sequence
\[
 0\longrightarrow\mathsf W_{3,\ell}
 \longrightarrow\mathsf W_{5,\ell}
 \longrightarrow\mathsf W_{2,\ell}
 \longrightarrow0,
 \qquad
 (\dim\mathsf W_{3,\ell},\dim\mathsf W_{5,\ell},
 \dim\mathsf W_{2,\ell})=(3,5,2).
\]
This is the canonical \(5\!-\!3\!-\!2\) descent used below.  Its
canonicity is relative to the marked reversal and orientation; changing the
orientation interchanges the signs.  No particular finite carrier
realization is part of this statement.
\end{remark}

The finite FSC calculation stops at this transition quotient and its marked-reversal
splitting.  Later arguments use only these data.

\part{Support-level source tower: confluence and repeated directions}
\section{Minimal confluence and the PD direction hull}

\subsection{Split slots and the strict diagonal}
Let $D_p,D_0,D_1$ be degree-zero marked directions in a dg-Lie/Toda carrier.  Let $H_{ab}$ be degree-$-1$ binary nullhomotopies with
\[
 H_{ba}=-H_{ab},\qquad dH_{ab}=[D_a,D_b].
\]
The associated alternating ternary cycle is
\begin{equation}\label{P1C-eq:confluent-split-B}
 B_{p01}:=[D_p,H_{01}]+[D_0,H_{1p}]+[D_1,H_{p0}].
\end{equation}
By the graded Jacobi identity, $dB_{p01}=0$.

\begin{definition}[Strict collision map]
The strict collision identifies the two marked slots by
\[
 D_0,D_1\longmapsto D_q,
 \qquad
 H_{p0},H_{p1}\longmapsto H_{pq},
 \qquad
 H_{01}\longmapsto H_{qq}:=0.
\]
It forgets the distinction between the two $q$-slots rather than recording their common origin.
\end{definition}

\begin{theorem}[Strict alternating collision vanishes]\label{P1C-thm:strict-collision-vanishes}
Under the strict collision map,
\[
 B_{p01}\longmapsto B_{pqq}=0.
\]
Hence repeated-prime depth cannot be represented by literal identification inside the ordinary alternating ternary Toda--Cartan skeleton.
\end{theorem}

\begin{proof}
Using $H_{qq}=0$ and $H_{qp}=-H_{pq}$,
\[
 B_{pqq}
 =[D_p,H_{qq}]+[D_q,H_{qp}]+[D_q,H_{pq}]
 =0-[D_q,H_{pq}]+[D_q,H_{pq}]=0.
\]
\end{proof}

The vanishing theorem is a no-go for a naive quotient, not a no-go for repeated-input deformation theory.  The missing information reappears one coherence level higher.

\subsection{The divided-power confluence cell}
Let $k$ be a commutative ring in which $2$ is a unit, let $O=k\kappa$ be a rank-one marked obstruction line, and truncate all parameter algebras modulo fourth powers of their maximal ideals.

\begin{definition}[Split one-cell third jet]
The split $(1,1,1)$ third-jet cell is
\[
 T_{111}(\kappa)
 :=k[u_0,u_1,v]\langle\eta_s\rangle_{\rm PD}/(u_0,u_1,v)^4,
 \qquad
 d\eta_s=u_0u_1v\,\kappa.
\]
\end{definition}

\begin{definition}[Normalized repeated-input cell]
The normalized $(2,1)$ confluence cell is
\[
 T_{21}(\kappa)
 :=k[u,v]\langle\eta_c\rangle_{\rm PD}/(u,v)^4,
 \qquad
 d\eta_c=\gamma_2(u)v\,\kappa.
\]
Since $2$ is invertible, this is equivalent after rescaling the obstruction cell to the polynomial normal form $d\eta=u^2v\,\kappa$.
\end{definition}

The divided-power normalization is preferable because polarization has coefficient one.

\begin{theorem}[Confluence--polarization identity]\label{P1C-thm:confluence-polarization}
The divided-power identity
\begin{equation}\label{P1C-eq:pd-polarization}
 \gamma_2(u_0+u_1)-\gamma_2(u_0)-\gamma_2(u_1)=u_0u_1
\end{equation}
canonically identifies the split cross term with the polarization of the repeated-input relation.  Equivalently,
\[
 \operatorname{Pol}_{01}\bigl(d\eta_c\bigr)
 =u_0u_1v\,\kappa
 =d\eta_s.
\]
On the diagonal $u_0=u_1=u$, the split cubic monomial becomes $u^2v=2\gamma_2(u)v$.
\end{theorem}

\begin{proof}
The standard divided-power addition formula gives
\[
 \gamma_2(x+y)=\gamma_2(x)+xy+\gamma_2(y).
\]
Subtracting the pure terms and multiplying by $v\kappa$ gives the first assertion.  The diagonal statement follows from $u^2=2\gamma_2(u)$ when $2$ is a unit.
\end{proof}

\begin{corollary}[Diagonal base change]\label{P1C-cor:diagonal-base-change}
Let
\[
 \Delta^*:k[u_0,u_1,v]/(u_0,u_1,v)^4
 \longrightarrow k[u,v]/(u,v)^4
\]
send $u_0,u_1$ to $u$ and $v$ to $v$.  Then the base change of $T_{111}(\kappa)$ along $\Delta^*$ is isomorphic to $T_{21}(\kappa)$ after the unit rescaling
\[
 \eta_s\longmapsto 2\eta_c.
\]
Hence the normalized repeated-input cell is the genuine diagonal confluence of the split one-cell third jet.
\end{corollary}

\begin{proof}
After base change, $d\eta_s=u^2v\kappa=2\gamma_2(u)v\kappa=d(2\eta_c)$.  Since $2$ is a unit, the cell rescaling is an isomorphism.
\end{proof}

\begin{corollary}[The confluence defect is secondary]\label{P1C-cor:confluence-secondary}
The strict alternating shadow of the repeated direction is zero by \cref{P1C-thm:strict-collision-vanishes}, while the normalized repeated-input PD cell of \cref{P1C-thm:confluence-polarization} may be nonzero.  Therefore a nonzero confluence class cannot be represented by the ordinary primary alternating ternary Lie operation; it belongs to secondary/Postnikov/divided-power coherence.
\end{corollary}

This is the first structural distinction of the theory: collision lowers the number of reduced support directions without lowering the third infinitesimal order.

\subsection{Primitive filtered confluence and the first genuine $q^2$ line}

The diagonal PD relation has two different integral enhancements, and they must be separated from the outset.  The first lives on the obstruction line itself.

\begin{definition}[Primitive filtered confluence lift]\label{P1C-def:filtered-confluence-lift}
Let
\[
 W_q=\Z_q\widetilde\kappa
\]
be a free rank-one marked obstruction module, with $\widetilde\kappa$ primitive, and let
\[
 \bar\kappa:=\widetilde\kappa\bmod q\ne0.
\]
A \emph{filtered confluence third jet} on this line is a one-cell marked quotient of the form
\begin{equation}\label{P1C-eq:filtered-confluence-normal-form}
 \widetilde T_{21}
 =
 \Z_q[u,v]\langle\widetilde\eta\rangle_{\rm PD}/(u,v)^4,
 \qquad
 d\widetilde\eta=(a\,uv+u^2v)\widetilde\kappa,
 \qquad a\in q\Z_q.
\end{equation}
The coefficient $a$ records a positive-filtration binary carry; the cubic coefficient has been normalized to the primitive basis $\widetilde\kappa$.
\end{definition}

\begin{theorem}[Primitive filtered depth-two line]\label{P1C-thm:filtered-q2-line}
For a primitive filtered confluence lift, put
\[
 L^{\rm conf}_{q^2}:=W_q/q^2W_q.
\]
Then
\[
 \boxed{
 L^{\rm conf}_{q^2}\simeq\Z/q^2,
 \qquad
 \Ord(\widetilde\kappa\bmod q^2)=q^2,
 \qquad
 (\widetilde\kappa\bmod q^2)\bmod q=\bar\kappa\ne0.}
\]
Moreover reduction of \eqref{P1C-eq:filtered-confluence-normal-form} modulo $q$ is
\[
 d\bar\eta=u^2v\,\bar\kappa,
\]
so the depth-two line reduces to the normalized repeated-input secondary cell.
\end{theorem}

\begin{proof}
Primitivity identifies $W_q$ with $\Z_q$ by $\widetilde\kappa\mapsto1$.  Hence
$W_q/q^2W_q\simeq\Z/q^2$, and the class of $\widetilde\kappa$ has exact order
$q^2$.  Since $a\in q\Z_q$, the binary term vanishes modulo $q$, leaving the
displayed repeated-input cubic relation.
\end{proof}

\begin{corollary}[Filtered depth does not require a minimal cocycle lift]\label{P1C-cor:filtered-vs-minimal}
Existence of $L^{\rm conf}_{q^2}$ requires no vanishing of a ternary Hochschild--Bockstein class.  It is a statement about the primitive filtered obstruction line.  Whether the same residual ternary class is represented by a cocycle on a prescribed minimal $q$-adic cohomology model is a strictly stronger lifting problem.
\end{corollary}

\begin{remark}[Realized seed]
The marked repeated arithmetic branch input (I3) of \cref{P1C-def:imported-hypotheses} is precisely the data isolated in \cref{P1C-def:filtered-confluence-lift}: a free primitive $q$-adic line, a nonzero residual repeated point and a distinguished lower carry in positive $q$-filtration.  No further source theorem is used in the formal arguments below.
\end{remark}

\subsection{The stronger minimal-model $q^2$ lifting gate}

Now fix a finite free graded $R=\Z_q$-module $H_R$ with an associative product $m_2$, and put $\bar H=H_R/qH_R$.  This subsection asks the stronger question of lifting a chosen residual ternary \emph{Hochschild cocycle} on a fixed cohomology model.  On ternary cochains write
\[
 C_R^{(3),j}:=\Hom_R^j(H_R^{\otimes3},H_R),
 \qquad
 \partial_{m_2}:=[m_2,-].
\]

\begin{definition}[First confluence--Bockstein class]\label{P1C-def:first-confluence-bockstein}
Let
\[
 \bar\tau\in\ker\bigl(
 \partial_{\bar m_2}:C_{\F_q}^{(3),j}\to C_{\F_q}^{(4),j+1}
 \bigr)
\]
represent the residual repeated-input ternary class.  Choose a lift $\widetilde\tau\in C_R^{(3),j}$.  Since $\bar\tau$ is closed, $\partial_{m_2}\widetilde\tau$ is divisible by $q$.  Define
\[
 \boxed{
 \beta_q^{\rm conf}(\bar\tau)
 :=\left[q^{-1}\partial_{m_2}\widetilde\tau\bmod q\right]
 \in H\bigl(C_{\F_q}^{(\ge4)},\partial_{\bar m_2}\bigr).}
\]
Only the arity-four component is needed at the first step.
\end{definition}

\begin{theorem}[Depth-two lifting criterion]\label{P1C-thm:depth-two-lifting-criterion}
The class $\beta_q^{\rm conf}(\bar\tau)$ is independent of the chosen lift.  Moreover,
\[
 \boxed{
 \beta_q^{\rm conf}(\bar\tau)=0
 \iff
 \bar\tau\text{ admits a ternary Hochschild-cocycle lift modulo }q^2.}
\]
Thus the first \emph{minimal-model ternary-cocycle} lifting step is governed by a canonical degree-one obstruction.  This does not contradict \cref{P1C-thm:filtered-q2-line}: filtered obstruction-line depth and minimal-cohomology cocycle depth are different notions.
\end{theorem}

\begin{proof}
If $\widetilde\tau'=\widetilde\tau+qF$, then
\[
 q^{-1}\partial_{m_2}\widetilde\tau'
 -q^{-1}\partial_{m_2}\widetilde\tau
 =\partial_{m_2}F,
\]
so the residual cohomology class is independent of the lift.  Vanishing means that after replacing $\widetilde\tau$ by $\widetilde\tau-qF$, its Hochschild differential is divisible by $q^2$, which is exactly the cocycle-lifting condition modulo $q^2$.
\end{proof}

\begin{remark}[Binary carry is a separate coordinate]
A $q$-divisible binary carry and the ternary Bockstein class are not related by the formal $A_\infty$ identities at the truncated three-jet level: associativity of $m_2$ appears before the relation $\partial_{m_2}m_3=0$.  Any arithmetic identification between those two coordinates is extra structure.
\end{remark}

\begin{principle}[Two depth notions]\label{P1C-prin:two-depth-notions}
At the first repeated-prime stage there are two legitimate but inequivalent depth questions:
\[
 \boxed{
 \begin{array}{ll}
 \text{filtered/derived depth:}
   & \widetilde\kappa\bmod q^2\text{ on a primitive obstruction line},\\[1mm]
 \text{minimal-model depth:}
   & \bar\tau\text{ lifted as a Hochschild cocycle modulo }q^2.
 \end{array}}
\]
The first is already present once a primitive filtered confluence lift is fixed.  The second is controlled by $\beta_q^{\rm conf}$.
\end{principle}

\subsection{Finite $p\times q^2$ assembly and the universal Moore line}

Let $p\ne q$ be odd primes.  The distinguished support label $p$ and the repeated arithmetic place $q$ do not by themselves prescribe a coefficient decomposition, but once a primitive $p$-primary coefficient line is chosen it combines canonically with the depth-two $q$-line by finite CRT.

\begin{definition}[Depth-two confluence coefficient line]\label{P1C-def:pq2-confluence-line}
Let
\[
 L_p=(\Z/p)b_p
\]
with $b_p$ a chosen generator, and let $L^{\rm conf}_{q^2}$ be as in
\cref{P1C-thm:filtered-q2-line}.  Define
\[
 \boxed{
 L^{\rm conf}_{p;q^2}
 :=
 L_p\times L^{\rm conf}_{q^2},
 \qquad
 b^{\rm conf}_{p;q^2}
 :=
 (b_p,\widetilde\kappa\bmod q^2).}
\]
\end{definition}

\begin{theorem}[Finite CRT confluence order]\label{P1C-thm:pq2-crt-confluence}
The pointed line $\bigl(L^{\rm conf}_{p;q^2},b^{\rm conf}_{p;q^2}\bigr)$ is cyclic of exact additive order $pq^2$.  Equivalently, there is a unique pointed isomorphism
\[
 \boxed{
 \Z/(pq^2)\xrightarrow{\ \sim\ }L^{\rm conf}_{p;q^2},
 \qquad
 1\longmapsto b^{\rm conf}_{p;q^2}.}
\]
Its $q$-primary reduction modulo $q$ is the repeated-input secondary generator $\bar\kappa$.
\end{theorem}

\begin{proof}
The two components have coprime exact orders $p$ and $q^2$.  Their product generator therefore has order $\operatorname{lcm}(p,q^2)=pq^2$.  The Chinese remainder theorem gives the displayed pointed cyclic identification, and the final assertion is \cref{P1C-thm:filtered-q2-line}.
\end{proof}

We now compare this realized finite-depth line with the universal Moore presentation.

\begin{definition}[Primitive confluence line]
A primitive confluence line is a free cyclic degree-$d$ summand $\Z B_{2,1}$ in a semi-free controller, with $dB_{2,1}=0$, whose image in the relevant third-jet realization is the normalized repeated-input secondary class.
\end{definition}

\begin{definition}[Depth-two Moore attachment]
For distinct odd primes $p,q$, the universal $(p;q^2)$ Moore extension adjoins a generator $C_{2,1}$ one degree lower with
\begin{equation}\label{P1C-eq:pq2-moore-attachment}
 dC_{2,1}=pq^2 B_{2,1}.
\end{equation}
\end{definition}

\begin{theorem}[Formal exact order $pq^2$]\label{P1C-thm:formal-pq2-order}
If $\Z B_{2,1}$ is a primitive free direct summand before the attachment and no other differential hits that line, then after \eqref{P1C-eq:pq2-moore-attachment}
\[
 H^d(\Z B_{2,1}\oplus\Z C_{2,1})\simeq\Z/(pq^2),
\]
and $[B_{2,1}]$ has exact additive order $pq^2$.
\end{theorem}

\begin{proof}
On the two-cell summand the only new boundary in degree $d$ is $pq^2B_{2,1}$.  Primitivity excludes a smaller positive multiple as a boundary, so the quotient of the free cyclic line is exactly $\Z/(pq^2)$.
\end{proof}

\begin{theorem}[Formal/filtered alignment at depth two]\label{P1C-thm:formal-filtered-alignment}
Let $L^{\rm Moore}_{p;q^2}$ be the cohomology line generated by $[B_{2,1}]$ in
\cref{P1C-thm:formal-pq2-order}.  Once the primitive generators $b_p$ and
$\widetilde\kappa$ are fixed, there is a unique pointed isomorphism
\[
 \boxed{
 \Phi^{\rm conf}_{p;q^2}:
 L^{\rm Moore}_{p;q^2}
 \xrightarrow{\ \sim\ }
 L^{\rm conf}_{p;q^2},
 \qquad
 [B_{2,1}]\longmapsto
 (b_p,\widetilde\kappa\bmod q^2).}
\]
Under $q$-primary reduction, $\Phi^{\rm conf}_{p;q^2}$ sends the universal Moore
generator to the residual repeated-input secondary generator $\bar\kappa$.
Thus the universal exact-order Moore line and the primitive filtered confluence line
are aligned inside the finite theory.
\end{theorem}

\begin{proof}
Both pointed groups are cyclic of exact order $pq^2$ by
\cref{P1C-thm:formal-pq2-order,P1C-thm:pq2-crt-confluence}.  A homomorphism between
cyclic groups is determined by the image of a generator, and sending one chosen
generator to the other gives an isomorphism.  The residual statement follows from
\cref{P1C-thm:filtered-q2-line}.
\end{proof}

\begin{warning}[Comparison boundary]\label{P1C-warn:external-ray-alignment}
The formal/filtered alignment does not identify every independently constructed
exact-order \(pq^2\) carrier.  Such a comparison must preserve the marked root
projector, coefficient assembly, and secondary marking.  On the Qi branch,
\cref{P1C-thm:qi-carry-comparison-package} removes the principal-unit ambiguity;
no stronger carrier comparison is used here.
\end{warning}

\begin{corollary}[Minimal $pqq$ model is internally closed]\label{P1C-cor:minimal-pqq-closed}
Within the finite formal/filtered theory, the minimal repeated-prime model now has all four pieces:
\[
 \boxed{
 (1,1,1)\xrightarrow{\rm confluence}(2,1)
 \xrightarrow{\rm primitive\ filtered\ lift}q^2
 \xrightarrow{\rm CRT\ with\ }p pq^2
 \xleftrightarrow{\rm pointed\ alignment}
 \text{Moore}(pq^2).}
\]
The Hochschild--Bockstein class remains only the stronger minimal-cohomology lifting
gate described in \cref{P1C-prin:two-depth-notions}.
\end{corollary}

\subsection{PD direction hull and provenance extension}\label{P1C-sec:pd-direction-hull}

The minimal $q=r$ confluence calculation suggests a broader but still finite closure principle. The correct object is not obtained by adjoining a new primitive direction symbol for every sum or product such as $q+r$ or $qr$. Instead one keeps the free marked direction lattice, applies the divided-power functor, and extends any already-defined additive direction provenance functorially. The resulting theorem separates three issues which must not be conflated: existence of an order-$m$ obstruction line, PD closure of its direction dependence, and preservation of exact coefficient order.

\begin{definition}[Marked direction lattice and PD hull]\label{P1C-def:pd-direction-hull}
For a finite marked support set $S$ put
\[
D_S:=\bigoplus_{i\in S}\Z e_i,
\qquad
\Gamma^\bullet(D_S):=\bigoplus_{m\ge0}\Gamma^m(D_S).
\]
The degree-$m$ piece $\Gamma^m(D_S)$ is called the \emph{degree-$m$ PD direction hull}. The symbols $e_i$ remember marked direction slots only; they are not coefficient prime factors and do not identify arithmetic places by themselves.
\end{definition}

\begin{lemma}[Integral PD monomial basis]\label{P1C-thm:pd-monomial-basis}
For every $m\ge0$, the abelian group $\Gamma^m(D_S)$ is finite free with basis
\[
\boxed{
\gamma_\alpha(e):=\prod_{i\in S}\gamma_{\alpha_i}(e_i),
\qquad
\alpha_i\ge0,
\quad
\sum_{i\in S}\alpha_i=m.}
\]
Thus squarefree sectors, repeated-input sectors and mixed multiplicity profiles are basis vectors of one integral degree-$m$ object.
\end{lemma}

\begin{proof}
This is the standard basis theorem for the divided-power algebra of a free abelian group. The divided-power relations reorder products into the displayed monomials, and these monomials form the usual integral basis of the degree-$m$ divided-power module.
\end{proof}

\begin{proposition}[Integral linear-combination formula]\label{P1C-thm:pd-integral-linear-combination}
Let
\[
x=\sum_{i\in S}a_i e_i\in D_S,
\qquad a_i\in\Z.
\]
Then for every $m\ge0$,
\[
\boxed{
\gamma_m(x)=\sum_{|\alpha|=m}\left(\prod_{i\in S}a_i^{\alpha_i}\right)\gamma_\alpha(e).}
\]
In particular, for two marked directions,
\[
\boxed{
\gamma_m(ae_q+be_r)=\sum_{i+j=m}a^i b^j\,\gamma_i(e_q)\gamma_j(e_r).}
\]
No factorial denominator is introduced.
\end{proposition}

\begin{proof}
Iterate the divided-power addition law
\[
\gamma_m(x+y)=\sum_{a+b=m}\gamma_a(x)\gamma_b(y)
\]
and use homogeneity $\gamma_j(ne)=n^j\gamma_j(e)$ for $n\in\Z$. The coefficient of the basis monomial $\gamma_\alpha(e)$ is exactly $\prod_i a_i^{\alpha_i}$.
\end{proof}

\begin{corollary}[The cubic $q+r$ identity]\label{P1C-cor:pd-cubic-q-plus-r}
For distinct marked slots $p,q,r$,
\[
\boxed{
e_p\gamma_2(e_q+e_r)=e_p\gamma_2(e_q)+e_pe_qe_r+e_p\gamma_2(e_r).}
\]
Hence the squarefree $pqr$ cell, the two repeated cells $pqq$ and $prr$, and the polarized $p(q+r)^2$ cell all lie in the same degree-three PD direction hull.
\end{corollary}

\begin{corollary}[No new generator for an integral linear combination]\label{P1C-cor:no-new-linear-direction-generator}
Let $e_*$ generate a free rank-one module and let
\[
\iota_x:\Z e_*\longrightarrow D_S,
\qquad e_*\longmapsto x=\sum_i a_i e_i.
\]
Then the degree-$m$ direction attached to the formal linear combination $x$ is simply
\[
\Gamma^m(\iota_x)\bigl(\gamma_m(e_*)\bigr)=\gamma_m(x)\in\Gamma^m(D_S).
\]
Thus the formal/algebraic theory requires no additional direction generator for $q+r$, or for any integral linear combination of the marked directions.  Product directions are handled by the PD-extension theorem below, rather than by adjoining separate primitive symbols.
\end{corollary}

\subsubsection{Naturality: addition, diagonal confluence, deletion and permutation}

\begin{proposition}[PD direction functoriality]\label{P1C-thm:pd-direction-functoriality}
Every homomorphism of marked direction lattices
\[
\varphi:D_S\longrightarrow D_T
\]
induces functorially a graded divided-power map
\[
\Gamma^\bullet(\varphi):\Gamma^\bullet(D_S)\longrightarrow\Gamma^\bullet(D_T),
\qquad
\Gamma^m(\varphi)(\gamma_m(x))=\gamma_m(\varphi(x)).
\]
Consequently PD polarization commutes strictly with composition of direction-linear maps.
\end{proposition}

\begin{proof}
The divided-power algebra is functorial in the underlying module. The displayed identity is part of the universal divided-power structure and composition follows from functoriality.
\end{proof}

\begin{corollary}[Addition--confluence compatibility]\label{P1C-cor:pd-addition-confluence-compatibility}
Let $\delta:D_{\{q,r\}}\to\Z e_q$ be the diagonal map
\[
\delta(e_q)=e_q,
\qquad
\delta(e_r)=e_q.
\]
Then
\[
\boxed{
\Gamma^m(\delta)\bigl(\gamma_m(e_q+e_r)\bigr)=\gamma_m(2e_q)=2^m\gamma_m(e_q).}
\]
Applying $\Gamma^m(\delta)$ termwise to the expansion of \cref{P1C-thm:pd-integral-linear-combination} gives the same result. Equivalently,
\[
\sum_{a+b=m}\gamma_a(e_q)\gamma_b(e_q)=\sum_{a=0}^m\binom ma\gamma_m(e_q)=2^m\gamma_m(e_q).
\]
Thus direction addition and diagonal confluence commute strictly.
\end{corollary}

\begin{proof}
The first equality is \cref{P1C-thm:pd-direction-functoriality}. For the termwise expansion use
\[
\gamma_a(e_q)\gamma_b(e_q)=\binom{a+b}{a}\gamma_{a+b}(e_q)
\]
from the block-confluence law, and then the binomial theorem.
\end{proof}

\begin{corollary}[Strict support deletion and permutation]\label{P1C-cor:pd-support-deletion-permutation}
For $T\subseteq S$ define
\[
\pi_T:D_S\to D_T,
\qquad
e_i\longmapsto
\begin{cases}
e_i,&i\in T,\\
0,&i\notin T.
\end{cases}
\]
Then $\Gamma^m(\pi_T)$ gives strict support deletion on the PD hull. In particular
\[
\gamma_2(e_q+e_r)\longmapsto\gamma_2(e_q)
\]
when $r$ is deleted. Every permutation $\sigma:S\to S$ similarly induces $e_i\mapsto e_{\sigma(i)}$ and hence a strict action on every $\Gamma^m(D_S)$. Support deletion, permutation, integral linear combination and diagonal confluence therefore belong to one functorial package.
\end{corollary}

\subsubsection{Primitive vectors and exact Moore order}

The formal closure under all integral linear combinations does not mean that every such combination preserves an exact coefficient order. The relevant invariant is the integral content of the PD vector.

\begin{definition}[Content]\label{P1C-def:pd-content}
For a finite free abelian group $F$ and $0\ne v\in F$, define
\[
\operatorname{cont}(v):=\max\{c\ge1:v\in cF\}.
\]
Equivalently, it is the greatest common divisor of the coordinates of $v$ in any integral basis. The vector is primitive precisely when $\operatorname{cont}(v)=1$.
\end{definition}

\begin{lemma}[Content of an integral PD direction]\label{P1C-thm:pd-direction-content}
Let $x=\sum_i a_i e_i\ne0$ and put $d=\gcd_i(a_i)$. Then
\[
\boxed{\operatorname{cont}\bigl(\gamma_m(x)\bigr)=d^m.}
\]
Hence $\gamma_m(x)$ is primitive if and only if the direction vector $x$ is primitive in $D_S$. In particular $\gamma_m(e_q+e_r)$ is primitive for every $m$.
\end{lemma}

\begin{proof}
In the basis of \cref{P1C-thm:pd-monomial-basis}, every coefficient of $\gamma_m(x)$ is a monomial $\prod_i a_i^{\alpha_i}$ of total degree $m$, hence is divisible by $d^m$. Conversely the list of coefficients contains $a_i^m$ for every $i$, and
\[
\gcd_i(a_i^m)=d^m.
\]
Therefore the common content is exactly $d^m$.
\end{proof}

\begin{proposition}[Primitive-vector Moore-order theorem]\label{P1C-thm:primitive-vector-moore-order}
Let $F$ be finite free, let $v\in F$ have content $c$, and let $[b]$ be an element of exact additive order $N$ in an abelian group $H$. Then the element
\[
[b]\otimes v\in H\otimes_\Z F
\]
has exact order
\[
\boxed{\frac{N}{\gcd(N,c)}.}
\]
In particular a primitive $v$ preserves the exact order $N$.
\end{proposition}

\begin{proof}
Write $v=cv_0$ with $v_0$ primitive. Since $F$ is free and $v_0$ is primitive there exists a homomorphism $\lambda:F\to\Z$ with $\lambda(v_0)=1$. Thus
\[
\id_H\otimes\lambda:H\otimes F\to H
\]
sends $[b]\otimes v_0$ to $[b]$, so $[b]\otimes v_0$ has exact order $N$. The class $[b]\otimes v=c([b]\otimes v_0)$ then has order $N/\gcd(N,c)$ by the scalar-order lemma.
\end{proof}

\begin{corollary}[Moore order of a linear-combination direction]\label{P1C-cor:pd-linear-combination-moore-order}
With $x=\sum_i a_i e_i$ and $d=\gcd_i(a_i)$ as above, if $[b]$ has exact order $N$, then
\[
\boxed{\operatorname{ord}\bigl([b]\otimes\gamma_m(x)\bigr)=\frac{N}{\gcd(N,d^m)}.}
\]
Thus the $q+r$ direction preserves exact Moore order because $e_q+e_r$ is primitive, whereas a nonprimitive integral combination may lose coefficient depth. Formal linear-combination closure and exact-order preservation are therefore distinct assertions.
\end{corollary}


\subsubsection{PD extension of genuine additive provenance}

\begin{definition}[Additive direction-provenance datum]\label{P1C-def:additive-direction-provenance}
Let $S$ be finite.  An \emph{additive direction-provenance datum} consists of a finite free abelian direction module $M_S$ and a homomorphism
\[
 \rho_S:D_S\longrightarrow M_S
\]
which is natural for the marked direction maps under consideration.  Let $W$ be a $\Z$-linear dg or stable root carrier with a Moore pair
\[
 db=0,\qquad dc=Nb,\qquad \Ord[b]=N.
\]
The root factor $W$ and the direction factor $M_S$ are kept separate.  If $W,b,c$ and $\rho_S$ come from an independently constructed arithmetic, geometric or motivic source, we call the datum \emph{genuine}.  This terminology deliberately does not identify a direction vector with the root Moore class itself.
\end{definition}

\begin{theorem}[PD extension of genuine additive provenance]\label{P1C-thm:pd-extension-genuine-additive-provenance}
For every additive direction-provenance datum of \cref{P1C-def:additive-direction-provenance} there is a unique natural graded PD extension
\begin{equation}\label{P1C-eq:pd-provenance-extension}
 \boxed{
 \rho_S^{\rm PD}:=\Gamma^\bullet(\rho_S):
 \Gamma^\bullet(D_S)\longrightarrow\Gamma^\bullet(M_S)}
\end{equation}
satisfying, for every $x\in D_S$ and $m\ge0$,
\begin{equation}\label{P1C-eq:pd-provenance-gamma}
 \boxed{
 \rho_{S,m}^{\rm PD}\bigl(\gamma_m(x)\bigr)
 =\gamma_m\bigl(\rho_S(x)\bigr).}
\end{equation}
Consequently every homogeneous PD polynomial $P\in\Gamma^m(D_S)$ has a canonical root-tagged provenance class
\begin{equation}\label{P1C-eq:pd-provenance-class}
 \boxed{
 \operatorname{Prov}^{\rm PD}_{S,m}(P)
 :=[b]\otimes\rho_{S,m}^{\rm PD}(P)
 \in H^\bullet(W)\otimes_\Z\Gamma^m(M_S).}
\end{equation}
If $v=\rho_{S,m}^{\rm PD}(P)$ has integral content $c(v)$, then
\begin{equation}\label{P1C-eq:pd-provenance-order}
 \boxed{
 \Ord\bigl(\operatorname{Prov}^{\rm PD}_{S,m}(P)\bigr)
 =\frac{N}{\gcd(N,c(v))}.}
\end{equation}
In particular primitive PD directions preserve exact order $N$.

For two direction vectors $y,z\in M_S$, multiplication is the multiplication in the PD direction factor and
\begin{equation}\label{P1C-eq:pd-product-from-addition}
 \boxed{
 yz=\gamma_2(y+z)-\gamma_2(y)-\gamma_2(z).}
\end{equation}
Hence the mixed product direction requires no new primitive provenance generator once additive provenance and PD compatibility are available.  More generally, for a multi-index $\alpha$,
\begin{equation}\label{P1C-eq:ordinary-monomial-pd-normalization}
 \boxed{
 \prod_i y_i^{\alpha_i}
 =\left(\prod_i\alpha_i!\right)
  \prod_i\gamma_{\alpha_i}(y_i).}
\end{equation}
Thus all finite polynomial mixed directions at a fixed infinitesimal order are obtained from the same PD extension; the factorial scalar in \eqref{P1C-eq:ordinary-monomial-pd-normalization}, rather than a new provenance generator, is what can reduce exact Moore order.
\end{theorem}

\begin{proof}
Existence and uniqueness of \eqref{P1C-eq:pd-provenance-extension} are the functoriality and universal property of the divided-power algebra, giving \eqref{P1C-eq:pd-provenance-gamma}; see \cref{P1C-thm:pd-direction-functoriality}.  Tensoring the fixed root class with the direction tag gives \eqref{P1C-eq:pd-provenance-class}.  The order formula is \cref{P1C-thm:primitive-vector-moore-order} applied to $v$.  Identity \eqref{P1C-eq:pd-product-from-addition} is the degree-two divided-power addition law.  Finally $y^a=a!\gamma_a(y)$ and multiplicativity give \eqref{P1C-eq:ordinary-monomial-pd-normalization}.
\end{proof}

\begin{corollary}[No separate product-direction provenance upgrade]\label{P1C-cor:no-separate-product-provenance-upgrade}
Assume genuine additive provenance has been fixed on marked directions $q,r$ and that the relevant order-two obstruction/root line already exists.  Then
\[
 \boxed{
 \rho_S(e_q)\rho_S(e_r)
 =\gamma_2\bigl(\rho_S(e_q)+\rho_S(e_r)\bigr)
  -\gamma_2\bigl(\rho_S(e_q)\bigr)
  -\gamma_2\bigl(\rho_S(e_r)\bigr)}
\]
in the PD direction factor.  In particular the mixed direction $e_qe_r$ (the formal $qr$ direction) is not a new primitive provenance generator.  The same statement applies at arbitrary finite support: $e_pe_qe_r$, the normalized repeated direction $e_p\gamma_2(e_q)$, and every fixed-order homogeneous PD polynomial are values of the single map $\rho_S^{\rm PD}$.
\end{corollary}

\begin{remark}[Ordinary repeated monomials carry factorial content]\label{P1C-rem:ordinary-monomial-factorial-content}
The phrase ``all polynomial mixed directions'' concerns provenance existence, not automatic preservation of exact coefficient order.  For example
\[
 q^2r=2\,\gamma_2(q)r,
 \qquad
 q^mr=m!\,\gamma_m(q)r.
\]
The normalized PD basis vector is primitive, whereas the ordinary monomial carries the displayed factorial.  Resonance between this factorial and $N$ is measured by \eqref{P1C-eq:pd-provenance-order} and is the same arithmetic content later encoded by the confluence-carry theory.
\end{remark}

\begin{warning}[Fixed-order closure, not higher-obstruction creation]\label{P1C-warn:pd-provenance-fixed-order-only}
\Cref{P1C-thm:pd-extension-genuine-additive-provenance} is a fixed-order extension theorem.  Once an order-\(m\) obstruction line and its additive direction provenance exist, every element of \(\Gamma^m(D_S)\) acquires canonical mixed-direction provenance.  It neither creates an order-\((m+1)\) obstruction nor identifies coefficient depth \(N\mapsto Nq\) with direction multiplicity.
\end{warning}


\section{Higher confluent obstruction calculus and the $q^3/(3,1)$ firewall}\label{P1C-sec:iterated-confluence}

The first confluence identifies two marked slots and produces the repeated-input shape $(2,1)$.  Beyond that seed there are three different operations: producing the next formal-moduli obstruction, polarizing a fixed obstruction order by divided powers, and deepening the coefficient order on a fixed primitive line.  Only the first is genuine higher recursion.  We construct it before returning to the multiplicity--depth firewall and the factorial confluence calculus.

\subsection{Compatible higher-obstruction transport}\label{P1C-subsec:higher-obstruction-recursion}

Let $k$ be a commutative ring and let $D_{S,k}:=D_S\otimes_\Z k$.  For a finite ceiling $M$ put
\[
 \mathscr P_{S,\le M}
 :=\bigoplus_{j=0}^{M}\Gamma_k^j(D_{S,k}),
\]
with the divided-power product truncated in total degree $>M$, and write
\[
 \mathscr J_{S,\le M}:=\bigoplus_{j=1}^{M}\Gamma_k^j(D_{S,k})
\]
for its positive-degree ideal.  Everything below is finite; no completion is used.

\begin{definition}[Strict pointed formal-moduli controller]\label{P1C-def:strict-formal-moduli-controller}
A \emph{strict pointed formal-moduli controller} is a dg associative $k$-algebra $(\mathscr C,d)$ together with a chosen Maurer--Cartan background $\xi_0\in\mathscr C^1$,
\[
 d\xi_0+\xi_0^2=0.
\]
Twisting by $\xi_0$ replaces the differential by
\[
 d_{\xi_0}(a)
 :=da+\xi_0a-(-1)^{|a|}a\xi_0.
\]
We work in the twisted controller and again denote its differential by $d$, so that the background point is zero.  When $2$ is invertible, the commutator dg-Lie formulation is equivalent through $dX+\frac12[X,X]=dX+X^2$; the associative form is used here because it remains integral without dividing by $2$.
\end{definition}

\begin{definition}[Marked finite jet and curvature]\label{P1C-def:marked-finite-mc-jet}
Fix $1\le n<M$.  A \emph{marked $n$-jet} is an element
\[
 X_{\le n}=X_1+\cdots+X_n,
 \qquad
 X_j\in\mathscr C^1\otimes_k\Gamma_k^j(D_{S,k}),
\]
such that the curvature
\[
 \operatorname{Curv}(X_{\le n})
 :=dX_{\le n}+X_{\le n}^2
\]
has zero homogeneous components in PD degrees $1,\ldots,n$.  Write
\[
 \operatorname{pr}_j:
 \mathscr C\otimes_k\mathscr P_{S,\le M}
 \longrightarrow
 \mathscr C\otimes_k\Gamma_k^j(D_{S,k})
\]
for the homogeneous projection.
\end{definition}

\begin{definition}[Next recursive obstruction]\label{P1C-def:next-recursive-obstruction}
For a marked $n$-jet define
\begin{equation}\label{P1C-eq:next-recursive-obstruction}
 \boxed{
 \Omega_{n+1}(X_{\le n})
 :=\operatorname{pr}_{n+1}\operatorname{Curv}(X_{\le n})
 \in
 \mathscr C^2\otimes_k\Gamma_k^{n+1}(D_{S,k}).}
\end{equation}
Its cohomology class, when defined by the theorem below, is denoted
\[
 \mathfrak o_{n+1}(X_{\le n})
 :=[\Omega_{n+1}(X_{\le n})].
\]
\end{definition}

\begin{theorem}[Universal higher-obstruction recursion]\label{P1C-thm:universal-higher-obstruction-recursion}
Let $X_{\le n}$ be a marked $n$-jet in a strict pointed formal-moduli controller.  Then:
\begin{enumerate}[label=\textup{(\roman*)}]
\item \textbf{Cocycle.}
\[
 \boxed{d\Omega_{n+1}(X_{\le n})=0.}
\]
Hence
\[
 \mathfrak o_{n+1}(X_{\le n})
 \in
 H^2(\mathscr C)\otimes_k\Gamma_k^{n+1}(D_{S,k})
\]
is canonically defined.

\item \textbf{Lifting criterion.}  An extension
\[
 X_{\le n+1}=X_{\le n}+X_{n+1},
 \qquad
 X_{n+1}\in\mathscr C^1\otimes_k\Gamma_k^{n+1}(D_{S,k}),
\]
satisfying the Maurer--Cartan equation through PD degree $n+1$ exists if and only if
\[
 \boxed{\mathfrak o_{n+1}(X_{\le n})=0.}
\]
When it exists, the set of fillers modulo gauge transformations which are the identity through degree $n$ is an affine torsor under
\[
 H^1(\mathscr C)\otimes_k\Gamma_k^{n+1}(D_{S,k}).
\]

\item \textbf{Gauge invariance.}  If two marked $n$-jets are related by a gauge transformation congruent to the identity modulo $\mathscr J_{S,\le M}$, then their degree-$(n+1)$ obstruction cocycles agree under the induced identification.  In particular the class $\mathfrak o_{n+1}$ depends only on the pointed gauge class of the lower jet.

\item \textbf{Controller and parameter naturality.}  For every dg-algebra map $F:\mathscr C\to\mathscr C'$ and every filtration-preserving algebra map
\[
 \Phi:\mathscr P_{S,\le M}\longrightarrow\mathscr P_{T,\le M},
\]
one has
\begin{equation}\label{P1C-eq:higher-obstruction-naturality}
 \boxed{
 \Omega_{n+1}\bigl((F\otimes\Phi)X_{\le n}\bigr)
 =
 (F\otimes\operatorname{gr}_{n+1}\Phi)
 \Omega_{n+1}(X_{\le n}).}
\end{equation}
In particular a direction-linear map $\varphi:D_S\to D_T$ gives $\Phi=\Gamma^{\le M}(\varphi)$, so addition, diagonal confluence, deletion, permutation and every integral direction-linear map commute with formation of the next recursive obstruction.  A tangent-identity filtered parameter automorphism acts trivially on the associated graded and therefore fixes $\mathfrak o_{n+1}$.
\end{enumerate}
\end{theorem}

\begin{proof}
Put $\mathcal F(X):=dX+X^2$.  The standard dg-algebra Bianchi identity is
\[
 d\mathcal F(X)+X\mathcal F(X)-\mathcal F(X)X=0.
\]
Because $X_{\le n}$ has strictly positive PD degree and its curvature has no component of degree at most $n$, the two product terms have PD degree at least $n+2$.  Taking the degree-$(n+1)$ component gives $d\Omega_{n+1}=0$, proving (i).

For a candidate $X_{n+1}$, all products involving $X_{n+1}$ have degree at least $n+2$, so
\[
 \operatorname{pr}_{n+1}\mathcal F(X_{\le n}+X_{n+1})
 =\Omega_{n+1}(X_{\le n})+dX_{n+1}.
\]
This proves the lifting criterion.  Two fillers differ by a degree-$(n+1)$ $1$-cocycle, and an order-$(n+1)$ gauge changes a filler by an exact $1$-cochain, giving the stated $H^1$ torsor.

For (iii), if $g=1+$ positive PD degree is a truncated gauge unit, then curvature transforms by conjugation,
\[
 \mathcal F(X^g)=g\mathcal F(X)g^{-1}.
\]
Since the first possible curvature term has degree $n+1$, conjugation by $g$ cannot change that homogeneous component.  Finally curvature is functorial under dg-algebra maps and under base change of the truncated parameter algebra; taking the first nonzero homogeneous term gives \eqref{P1C-eq:higher-obstruction-naturality}.  The direction-linear and tangent-identity cases follow from the induced maps on the associated graded.
\end{proof}

\begin{theorem}[Universal recursive cell attachment]\label{P1C-thm:universal-recursive-cell-attachment}
Let $X_{\le n}$ be a marked $n$-jet.  If $\mathfrak o_{n+1}(X_{\le n})$ vanishes, no new formal cell is required: any solution of
\[
 dX_{n+1}=-\Omega_{n+1}
\]
gives the next jet.  If the marked obstruction is nonzero, form the semi-free PD-weighted dg extension by one degree-$1$, weight-$(n+1)$ generator $h_{n+1}$ with
\begin{equation}\label{P1C-eq:universal-recursive-cell}
 \boxed{dh_{n+1}=-\Omega_{n+1}(X_{\le n}).}
\end{equation}
Then $d^2h_{n+1}=0$ and
\[
 X_{\le n}^{\rm cell}:=X_{\le n}+h_{n+1}
\]
satisfies the Maurer--Cartan equation through degree $n+1$ in the enlarged controller.

If $\Omega'_{n+1}=\Omega_{n+1}+dY$, the two semi-free attachments are filtered-dg isomorphic by
\[
 h'_{n+1}\longmapsto h_{n+1}-Y.
\]
Hence the marked cell attachment depends only on the recursive obstruction class and the chosen lower branch, up to filtered dg isomorphism.  Iterating to any finite ceiling $M$ produces a finite branchwise semi-free recursive controller; no inverse limit is involved.
\end{theorem}

\begin{proof}
The cocycle statement in \cref{P1C-thm:universal-higher-obstruction-recursion} gives $d^2h_{n+1}=0$.  In PD degree $n+1$, the curvature of $X_{\le n}+h_{n+1}$ is
\[
 \Omega_{n+1}+dh_{n+1}=0,
\]
because every product containing $h_{n+1}$ has degree at least $n+2$.  The change-of-representative formula is immediate from the differential.  At a finite ceiling only finitely many cells are adjoined along a fixed marked branch.
\end{proof}

\begin{definition}[Recursive repeated-sector coefficient]\label{P1C-def:recursive-repeated-sector}
Assume $S$ contains marked directions $q$ and $v$.  For $m\ge2$ the PD monomial
\[
 \vartheta_{m,1}:=\gamma_m(e_q)e_v
 \in\Gamma_k^{m+1}(D_{S,k})
\]
is a basis vector.  Let
\[
 \operatorname{coeff}_{m,1}:\Gamma_k^{m+1}(D_{S,k})\to k\,\vartheta_{m,1}
\]
be the corresponding basis projection.  For a marked $m$-jet define the recursive repeated-sector class
\begin{equation}\label{P1C-eq:recursive-repeated-sector}
 \boxed{
 \mathfrak o^{\rm rec}_{m,1}(X_{\le m})
 :=
 (\id\otimes\operatorname{coeff}_{m,1})
 \mathfrak o_{m+1}(X_{\le m}).}
\end{equation}
\end{definition}

\begin{theorem}[Higher repeated-sector recursion criterion]\label{P1C-thm:higher-repeated-sector-recursion}
For every $m\ge2$, the class $\mathfrak o^{\rm rec}_{m,1}$ is canonical relative to the pointed lower jet and gauge invariant.  Under a confluence or direction map, the full class $\mathfrak o_{m+1}$ transforms by \eqref{P1C-eq:higher-obstruction-naturality}; whenever the marked source line $\vartheta_{m,1}$ is carried to a scalar multiple of a marked target basis line, the corresponding repeated-sector coefficient is transported by that same PD scalar.  It gives the precise next repeated-input gate:
\begin{enumerate}[label=\textup{(\roman*)}]
\item if $\mathfrak o^{\rm rec}_{m,1}=0$, there is no intrinsic new $(m,1)$ obstruction line forced at that branch;
\item if it is nonzero and spans a primitive rank-one marked summand of the relevant integral or $q$-adic obstruction lattice, that summand is the next primitive repeated obstruction line, denoted $k\kappa_{m,1}$ (or $\Z_q\kappa_{m,1}$);
\item once such a line exists, \cref{P1C-thm:pd-extension-genuine-additive-provenance} and the confluence calculus below supply all of its fixed-order mixed-direction polarizations, but they do not decide whether the next class $\mathfrak o^{\rm rec}_{m+1,1}$ is nonzero.
\end{enumerate}
For $m=2$, in a strict controller realization in which the normalized one-cell $(2,1)$ seed is the degree-three Maurer--Cartan defect, this coefficient recovers that cubic repeated sector.  For $m=3$ it defines a canonical quartic $(3,1)$ obstruction gate without asserting that the gate is nonzero.
\end{theorem}

\begin{proof}
The PD monomial basis makes the coefficient projection canonical.  Gauge invariance and naturality are inherited from \cref{P1C-thm:universal-higher-obstruction-recursion}.  The first two assertions are the distinction between vanishing, nonzero torsion/content and a primitive rank-one obstruction summand.  The final assertion is exactly the fixed-order scope of the PD-extension theorem.  At $m=2$, the statement is the defining identification in any strict controller model chosen to realize the normalized one-cell third-jet seed.
\end{proof}



\subsection{Carry-compatible filtered transport at all finite arities}\label{P1C-subsec:carry-compatible-recursion}

The homogeneous recursion of \cref{P1C-thm:universal-higher-obstruction-recursion} is the correct
residual theory, but an integral mixed-characteristic branch need not be homogeneous in parameter
degree.  Lower parameter-degree terms may survive with positive $q$-filtration.  The cubic Qi
relation is the minimal example: its binary carry has parameter degree two and positive
$q$-filtration, while its repeated cubic coefficient is primitive.  The correct integral recursion is
therefore filtered by the sum of parameter degree and coefficient filtration.  This subsection builds
that recursion uniformly at every finite arity.

\begin{definition}[Effective-weight filtration]\label{P1C-def:effective-weight-relation-module}
Fix an odd prime $q$, put $\mathcal O=\Z_q$, and let $\widetilde{\mathscr C}$ be a
$q$-torsion-free $\mathcal O$-linear dg algebra.  For a finite marked set $S$ and a finite ceiling
$M$ put
\[
 \mathscr J^{+}_{S,\le M}
 :=\bigoplus_{j=1}^{M}\Gamma^j_{\mathcal O}(D_{S,\mathcal O}),
 \qquad
 \mathscr A_{S,\le M}:=
 \widetilde{\mathscr C}\otimes_{\mathcal O}\mathscr J^{+}_{S,\le M},
\]
with products of parameter degree $>M$ truncated to zero.  For $N\ge1$ define the decreasing
\emph{effective-weight filtration}
\begin{equation}\label{P1C-eq:effective-weight-filtration}
 \boxed{
 \mathsf F_{\rm eff}^{N}\mathscr A^i_{S,\le M}
 :=\bigoplus_{j=1}^{M}
 q^{\max(N-j,0)}
 \bigl(\widetilde{\mathscr C}^{i}\otimes_{\mathcal O}
       \Gamma^j_{\mathcal O}(D_{S,\mathcal O})\bigr).}
\end{equation}
Thus a term of parameter degree $j<N$ can contribute at effective weight $N$ only after acquiring at
least $N-j$ units of $q$-filtration.  For finite coefficient depth $r$ we write
\[
 \mathscr A_{S,\le M;r}
 :=\mathscr A_{S,\le M}\otimes_{\mathcal O}\mathcal O/q^r
\]
and define the finite-depth effective filtration by its image in the ambient
coefficient reduction:
\begin{equation}\label{P1C-eq:finite-depth-effective-image-filtration}
 \boxed{
 \mathsf F_{{\rm eff},r}^{N}\mathscr A_{S,\le M;r}
 :=\operatorname{im}\!\left(
 \mathsf F_{\rm eff}^{N}\mathscr A_{S,\le M}
 \longrightarrow \mathscr A_{S,\le M;r}
 \right).}
\end{equation}
This is an image filtration in the ambient reduction, not the intrinsic tensor
$\mathsf F_{\rm eff}^{N}\mathscr A\otimes_{\mathcal O}\mathcal O/q^r$.
No inverse limit in $r$ is part of this definition.
\end{definition}

\begin{proposition}[Multiplicativity and the effective associated graded]
\label{P1C-prop:effective-weight-multiplicative}
The effective filtration satisfies
\[
 \boxed{
 \mathsf F_{\rm eff}^{A}\mathscr A\,
 \mathsf F_{\rm eff}^{B}\mathscr A
 \subseteq
 \mathsf F_{\rm eff}^{A+B}\mathscr A,
 \qquad
 d\mathsf F_{\rm eff}^{N}\mathscr A^i
 \subseteq
 \mathsf F_{\rm eff}^{N}\mathscr A^{i+1}.}
\]
For $1\le N\le M$, parameter degree gives a canonical decomposition
\begin{equation}\label{P1C-eq:effective-associated-graded}
 \operatorname{gr}^{N}_{\rm eff}\mathscr A^i
 \cong
 \bigoplus_{j=1}^{N}
 \frac{q^{N-j}
 (\widetilde{\mathscr C}^{i}\otimes\Gamma^j(D_S))}
 {q^{N+1-j}
 (\widetilde{\mathscr C}^{i}\otimes\Gamma^j(D_S))}.
\end{equation}
In particular there is a canonical top-parameter projection
\begin{equation}\label{P1C-eq:effective-top-projection}
 \boxed{
 \operatorname{top}_N:
 \operatorname{gr}^{N}_{\rm eff}\mathscr A^i
 \longrightarrow
 (\widetilde{\mathscr C}^{i}/q)
 \otimes_{\F_q}\Gamma^N_{\F_q}(D_{S,\F_q}).}
\end{equation}
At finite depth $r$, every summand with $N-j\ge r$ vanishes automatically.
\end{proposition}

\begin{proof}
For homogeneous terms of parameter degrees $j_1,j_2$ one has
\[
 \max(A-j_1,0)+\max(B-j_2,0)
 \ge \max(A+B-j_1-j_2,0),
\]
so multiplication has the required filtration.  The differential is $\mathcal O$-linear and preserves
parameter degree.  Since the parameter-degree decomposition is direct and
$\widetilde{\mathscr C}$ is $q$-torsion-free, the successive quotient in degree $j\le N$ is exactly the
quotient displayed in \eqref{P1C-eq:effective-associated-graded}; degrees $j>N$ occur with filtration zero
in both consecutive stages and hence disappear.  The top projection is the $j=N$ summand.  The
finite-depth statement follows from
\eqref{P1C-eq:finite-depth-effective-image-filtration}: in parameter degree
$j$, the image of $q^{N-j}V_j$ in $V_j/q^rV_j$ is zero whenever
$N-j\ge r$.  This is vanishing in the ambient quotient; the intrinsic tensor
$q^{N-j}V_j\otimes_{\mathcal O}\mathcal O/q^r$ need not vanish.
\end{proof}

\begin{definition}[Carry-compatible filtered branch and obstruction]\label{P1C-def:cc-filtered-branch}
Let
\[
 X\in \mathsf F_{\rm eff}^{1}\mathscr A^1_{S,\le M}
\]
be a pointed degree-one element and put
\[
 \operatorname{Curv}(X):=dX+X^2.
\]
We say that $X$ is solved through effective weight $N-1$ if
\[
 \operatorname{Curv}(X)\in\mathsf F_{\rm eff}^{N}\mathscr A^2.
\]
Its next carry-compatible curvature symbol is
\[
 \Omega_N^{\rm cc}(X)
 :=[\operatorname{Curv}(X)]
 \in\operatorname{gr}^{N}_{\rm eff}\mathscr A^2.
\]
When this symbol is a cocycle, its cohomology class is denoted
\[
 \mathfrak o_N^{\rm cc}(X)
 :=[\Omega_N^{\rm cc}(X)]
 \in H^2(\operatorname{gr}^{N}_{\rm eff}\mathscr A).
\]
\end{definition}

\begin{theorem}[Effective-weight higher-obstruction recursion]
\label{P1C-thm:effective-weight-higher-recursion}
Let $X$ be solved through effective weight $N-1$, with $1\le N\le M$.  Then:
\begin{enumerate}[label=\textup{(\roman*)}]
\item \textbf{Cocycle.}  The Bianchi identity implies
\[
 \boxed{d\Omega_N^{\rm cc}(X)=0}
\]
in $\operatorname{gr}^{N}_{\rm eff}\mathscr A$.  Hence
$\mathfrak o_N^{\rm cc}(X)$ is canonically defined.
\item \textbf{Lifting criterion.}  There exists
$Y_N\in\mathsf F_{\rm eff}^{N}\mathscr A^1$ such that
\[
 \operatorname{Curv}(X+Y_N)\in\mathsf F_{\rm eff}^{N+1}\mathscr A^2
\]
if and only if
\[
 \boxed{\mathfrak o_N^{\rm cc}(X)=0.}
\]
When nonempty, the set of such corrections modulo effective-weight-$N$ gauge is an affine torsor
under $H^1(\operatorname{gr}^{N}_{\rm eff}\mathscr A)$.
\item \textbf{Gauge invariance.}  If $g\in1+\mathsf F_{\rm eff}^{1}\mathscr A^0$ is a pointed filtered
gauge unit, then
\[
 \mathfrak o_N^{\rm cc}(X^g)=\mathfrak o_N^{\rm cc}(X).
\]
\item \textbf{Naturality.}  Every filtered dg map of controllers together with a linear map of marked
direction lattices carries $\mathfrak o_N^{\rm cc}$ to the corresponding target class.  In particular
support deletion, permutation, integral direction maps and diagonal confluence commute with the
carry-compatible recursion.
\item \textbf{Residual top face.}  Reduction modulo $q$ followed by
\eqref{P1C-eq:effective-top-projection} gives precisely the homogeneous order-$N$ Maurer--Cartan
obstruction of \cref{P1C-thm:universal-higher-obstruction-recursion} for the residual branch:
\begin{equation}\label{P1C-eq:cc-to-residual-obstruction}
 \boxed{
 \operatorname{top}_N(\mathfrak o_N^{\rm cc}(X))
 =\mathfrak o_N(\bar X).}
\end{equation}
Consequently the repeated coefficient of the right-hand side is
$\mathfrak o^{\rm rec}_{N-1,1}$.
\end{enumerate}
\end{theorem}

\begin{proof}
The Bianchi identity reads
\[
 d\operatorname{Curv}(X)+[X,\operatorname{Curv}(X)]=0.
\]
Since $X\in\mathsf F_{\rm eff}^{1}$ and
$\operatorname{Curv}(X)\in\mathsf F_{\rm eff}^{N}$, multiplicativity puts the bracket in
$\mathsf F_{\rm eff}^{N+1}$; hence the curvature symbol is $d$-closed on the associated graded.
For $Y_N\in\mathsf F_{\rm eff}^{N}$,
\[
 \operatorname{Curv}(X+Y_N)
 =\operatorname{Curv}(X)+dY_N+XY_N+Y_NX+Y_N^2.
\]
The last three terms lie in $\mathsf F_{\rm eff}^{N+1}$, so the weight-$N$ symbol changes by the
coboundary $d[Y_N]$.  This proves the lifting criterion and the usual $H^1$ torsor statement.
Curvature transforms by conjugation under a filtered gauge unit; conjugation by an element congruent
to $1$ modulo $\mathsf F_{\rm eff}^{1}$ changes a weight-$N$ class only in weight $N+1$, proving
(iii).  Functoriality of curvature and of the effective filtration gives (iv).  Finally every lower
parameter-degree component of effective weight $N$ is $q$-divisible, so modulo $q$ only the
parameter-degree-$N$ component remains; this is exactly the homogeneous recursion of the residual
controller.
\end{proof}

\begin{definition}[Effective closed-lift Bockstein]\label{P1C-def:effective-closed-lift-bockstein}
For every $N$ the short exact sequence of complexes
\[
 0\longrightarrow
 \mathsf F_{\rm eff}^{N+1}\mathscr A
 \longrightarrow
 \mathsf F_{\rm eff}^{N}\mathscr A
 \longrightarrow
 \operatorname{gr}^{N}_{\rm eff}\mathscr A
 \longrightarrow0
\]
has a connecting map
\begin{equation}\label{P1C-eq:effective-closed-lift-bockstein}
 \boxed{
 \beta_N^{\rm eff}:
 H^2(\operatorname{gr}^{N}_{\rm eff}\mathscr A)
 \longrightarrow
 H^3(\mathsf F_{\rm eff}^{N+1}\mathscr A).}
\end{equation}
We call $\beta_N^{\rm eff}(\mathfrak o_N^{\rm cc})$ the \emph{effective closed-lift
Bockstein}.  It is a characteristic-lift gate for the full carry-compatible relation, not for the
existence of its top residual obstruction line.
\end{definition}

\begin{proposition}[Closed-lift criterion for an effective obstruction]
\label{P1C-prop:effective-closed-lift-criterion}
The carry-compatible obstruction class $\mathfrak o_N^{\rm cc}$ admits a closed lift
\[
 \widetilde\Omega_N\in
 Z^2(\mathsf F_{\rm eff}^{N}\mathscr A)
\]
representing its full associated-graded symbol if and only if
\[
 \boxed{\beta_N^{\rm eff}(\mathfrak o_N^{\rm cc})=0.}
\]
When the gate vanishes, the set of cohomology classes of closed lifts is a torsor under the image of
\[
 H^2(\mathsf F_{\rm eff}^{N+1}\mathscr A)
 \longrightarrow
 H^2(\mathsf F_{\rm eff}^{N}\mathscr A).
\]
In particular the gate vanishes automatically whenever
$H^3(\mathsf F_{\rm eff}^{N+1}\mathscr A)=0$.
\end{proposition}

\begin{proof}
This is the exactness of the cohomology sequence associated with
\eqref{P1C-eq:effective-closed-lift-bockstein}.  The torsor statement is the usual description of the
fibre of the map on $H^2$.
\end{proof}


\begin{theorem}[Cohomological-dimension-two vanishing of effective Bocksteins]
\label{P1C-thm:cd2-effective-bockstein-vanishing}
Let $G$ be a profinite group with $\operatorname{cd}_q(G)\le2$, let $T$ be a finite free
$\Z_q$-module with continuous $G$-action, and put
\[
 \mathscr M_T:=C^\bullet_{\rm cont}(G,T).
\]
For a finite marked direction set $S$ and finite parameter ceiling $M$, equip
\[
 \mathscr A_T:=\mathscr M_T\otimes_{\Z_q}\mathscr J^+_{S,\le M}
\]
with the effective filtration of \cref{P1C-def:effective-weight-relation-module}.  Then:
\begin{enumerate}[label=\textup{(\roman*)}]
\item for every finite depth $r\ge1$, after replacing $T$ by $T/q^rT$ one has
\[
 H^3\!\left(\mathsf F_{\rm eff}^{N+1}\mathscr A_{T/q^rT}\right)=0
 \qquad(N\ge1),
\]
and hence every finite-depth effective closed-lift Bockstein vanishes;
\item if the inverse system $\{H^2(G,T/q^rT)\}_r$ is Mittag--Leffler---for example, if every
$H^2(G,T/q^rT)$ is finite---then for the $q$-adic coefficient line itself
\begin{equation}\label{P1C-eq:cd2-effective-h3-vanishing}
 \boxed{
 H^3\!\left(\mathsf F_{\rm eff}^{N+1}\mathscr A_T\right)=0,}
\end{equation}
and consequently
\begin{equation}\label{P1C-eq:cd2-effective-beta-zero}
 \boxed{\beta_N^{\rm eff}=0.}
\end{equation}
\end{enumerate}
\end{theorem}

\begin{proof}
At finite depth the coefficient modules are discrete finite $q$-primary modules, so
$\operatorname{cd}_q(G)\le2$ gives $H^3(G,T/q^rT)=0$ directly.  Parameter degree is a finite
direct sum and every divided-power factor is finite free, so the same vanishing holds for each
finite-depth effective-filtration term.

For the $q$-adic coefficient line, the standard Milnor exact sequence for continuous cohomology gives
\[
 0\longrightarrow
 \varprojlim\nolimits_r^{1}H^2(G,T/q^rT)
 \longrightarrow H^3_{\rm cont}(G,T)
 \longrightarrow
 \varprojlim_r H^3(G,T/q^rT)
 \longrightarrow0.
\]
The right term is zero by cohomological dimension, while the Mittag--Leffler hypothesis kills the
$\varprojlim^1$ term.  Thus $H^3_{\rm cont}(G,T)=0$.  Now
\[
 \mathsf F_{\rm eff}^{N+1}\mathscr A_T
 =\bigoplus_{j=1}^{M}
 q^{\max(N+1-j,0)}\mathscr M_T\otimes_{\Z_q}\Gamma^j(D_S)
\]
is a finite direct sum of copies of the same continuous cochain complex up to finite-free tensor
factors and multiplication by powers of $q$.  Hence its degree-three cohomology vanishes, and
\cref{P1C-prop:effective-closed-lift-criterion} gives the Bockstein statement.
\end{proof}

\begin{corollary}[Root-projected effective Bocksteins vanish]
\label{P1C-cor:root-projected-effective-bockstein-zero}
Suppose a marked arithmetic controller admits a dg coefficient projection
\[
 \pi_T:\widetilde{\mathscr C}\longrightarrow C^\bullet_{\rm cont}(G,T)
\]
to a finite-free $q$-adic coefficient line with $\operatorname{cd}_q(G)\le2$, and assume the finite-level system $\{H^2(G,T/q^rT)\}_r$ is Mittag--Leffler.  Then for every
carry-compatible obstruction,
\[
 \boxed{
 \pi_{T,*}\!\left(
 \beta_N^{\rm eff}(\mathfrak o_N^{\rm cc})
 \right)=0.}
\]
Equivalently, the projected effective obstruction always has a closed filtered lift.  This is a
statement about the chosen coefficient face only: complementary coefficient corners can still carry
nonzero effective Bocksteins.
\end{corollary}

\begin{proof}
Naturality of the short exact sequence defining \cref{P1C-def:effective-closed-lift-bockstein} makes the
connecting homomorphism commute with $\pi_T$.  The target connecting map is zero by
\cref{P1C-thm:cd2-effective-bockstein-vanishing}.
\end{proof}

\begin{definition}[Carry-compatible recursive relation and relative cell]
\label{P1C-def:carry-compatible-recursive-cell}
Let a fixed integral lower branch be represented by $X$ as above.  A \emph{strict carry-compatible
order-$N$ relation} is a closed cocycle
\[
 \widetilde\Omega_N\in
 Z^2(\mathsf F_{\rm eff}^{N}\mathscr A)
\]
whose class in $\operatorname{gr}^{N}_{\rm eff}$ represents the desired obstruction symbol.  A
\emph{relative recursive cell} is a new degree-one generator $h_N$ assigned effective weight $N$ and
satisfying
\[
 \boxed{dh_N=-\widetilde\Omega_N.}
\]
The effective filtration is extended multiplicatively by the rule
$h_N\in\mathsf F_{\rm eff}^{N}$, and all lower cells and carry markings are left unchanged.
\end{definition}

\begin{theorem}[Carry-compatible filtered recursive attachment]
\label{P1C-thm:carry-compatible-filtered-recursion}
Let $X$ be solved through effective weight $N-1$.  Assume:
\begin{enumerate}[label=\textup{(A\arabic*)}]
\item the effective closed-lift gate vanishes,
\[
 \beta_N^{\rm eff}(\mathfrak o_N^{\rm cc}(X))=0,
\]
and choose a closed lift
\[
 \widetilde\Omega_N\in
 Z^2(\mathsf F_{\rm eff}^{N}\mathscr A)
\]
of the \emph{full} effective obstruction symbol;
\item the top residual face of this class contains a nonzero marked repeated line
\[
 \F_q\bar o_N\otimes\F_q\vartheta,
 \qquad
 \vartheta\in\Gamma^N(D_S)\text{ primitive},
\]
and a marked free rank-one $\mathcal O$-lattice
\[
 W_N^{\rm mark}
 \subseteq
 H^2(\widetilde{\mathscr C})\otimes_{\mathcal O}\mathcal O\vartheta
\]
reduces onto that top line.
\end{enumerate}
Then:
\begin{enumerate}[label=\textup{(\roman*)}]
\item the top repeated coefficient of every lift of the nonzero residual generator to
$W_N^{\rm mark}$ is primitive;
\item adjoining a degree-one generator $h_N$ of effective weight $N$ with
\[
 \boxed{dh_N=-\widetilde\Omega_N}
\]
produces a relative integral recursive branch solved through effective weight $N$;
\item every closed lift can be written, after separating parameter degree, in the form
\begin{equation}\label{P1C-eq:general-carry-compatible-relation}
 \boxed{
 \widetilde\Omega_N
 =\omega_N+
 \sum_{j<N}q^{N-j}\omega_j
 +\Omega_{>N},}
\end{equation}
where $\omega_N$ has top parameter degree $N$ and reduces to the residual obstruction,
$\Omega_{>N}\in\mathsf F_{\rm eff}^{N+1}$, and the lower parameter-degree summands are arithmetic
carry corrections of the same effective weight;
\item if two closed lifts satisfy
\[
 \widetilde\Omega'_N
 =u\widetilde\Omega_N+dY,
 \qquad u\in1+q\mathcal O,
\]
then their relative attachments are filtered-dg isomorphic by
$h'_N\mapsto u h_N-Y$.  In general, different closed lifts can differ by the higher-filtration
$H^2$-torsor of \cref{P1C-prop:effective-closed-lift-criterion}; therefore the \emph{top primitive repeated
line} is canonical up to principal-unit pointing, while the complete carry-corrected relation may
retain higher-filtration moduli until additional arithmetic provenance fixes them.
\item Reducing the construction modulo any finite $q^r$ gives the corresponding finite-depth relative
cell.  These reductions are compatible with $\mathcal O/q^{r+1}\to\mathcal O/q^r$, but no inverse
limit is formed here.
\end{enumerate}
\end{theorem}

\begin{proof}
The first assertion is the primitive-lift criterion in a free rank-one $\mathcal O$-module.  Since
$\widetilde\Omega_N$ is an actual closed lift of the full weight-$N$ symbol, the semi-free differential
satisfies $d^2h_N=0$.  Assigning $h_N$ effective weight $N$ gives
\[
 \operatorname{Curv}(X+h_N)
 =\operatorname{Curv}(X)-\widetilde\Omega_N+Xh_N+h_NX+h_N^2,
\]
and the final three terms have effective weight at least $N+1$, so the branch advances by one
weight.  Formula \eqref{P1C-eq:general-carry-compatible-relation} is simply the parameter-degree
splitting of a lift in $\mathsf F_{\rm eff}^{N}$, with the top residual coefficient primitive by (A2).
The stated change of lift gives
$d(u h_N-Y)=-u\widetilde\Omega_N-dY=-\widetilde\Omega'_N$.  The remaining freedom is exactly the
closed-lift torsor from \cref{P1C-prop:effective-closed-lift-criterion}.  Finite coefficient reduction is
functorial.
\end{proof}

\begin{example}[The cubic Qi cell is effective-weight three]
\label{P1C-ex:qi-effective-weight-three}
For the canonical finite Qi relation of \cref{P1C-thm:finite-qi-bifiltered-relation},
\[
 d\eta_{c,r}
 =\bigl(\alpha_{q^r}e_pe_q+e_p\gamma_2(e_q)\bigr)\kappa_{\rm Qi}^{\rm can},
 \qquad \alpha_{q^r}\in qR_r
\]
whenever the binary carry is visible.  The first term has parameter degree two and at least one unit
of coefficient filtration, while the second has parameter degree three and filtration zero.  Hence the
whole relation has effective weight at least three, and its weight-three top residual face is precisely
the repeated cubic obstruction.  Thus the Qi third jet is the first nontrivial carry-compatible
recursive cell, not an exception to the recursion.
\end{example}

\begin{theorem}[All finite arity recursive Confluent mechanism]
\label{P1C-thm:all-finite-arity-confluent-recursion}
Fix a finite ceiling $M$.  Starting from any pointed integral lower branch in a $q$-torsion-free
controller, the effective filtration defines, for every $2\le N\le M$, a canonical carry-compatible
obstruction
\[
 \mathfrak o_N^{\rm cc}
 \in H^2(\operatorname{gr}^{N}_{\rm eff}\mathscr A)
\]
and its effective closed-lift Bockstein
\[
 \beta_N^{\rm eff}(\mathfrak o_N^{\rm cc})
 \in H^3(\mathsf F_{\rm eff}^{N+1}\mathscr A).
\]
At each stage the recursion separates four logically distinct outcomes.
\begin{enumerate}[label=\textup{(\alph*)}]
\item If $\mathfrak o_N^{\rm cc}=0$, the branch extends inside the current controller; the extensions
form the $H^1(\operatorname{gr}^{N}_{\rm eff}\mathscr A)$ torsor of
\cref{P1C-thm:effective-weight-higher-recursion}.
\item If $\mathfrak o_N^{\rm cc}
e0$ but its marked residual top face fails the arithmetic/provenance
detector, the ledger records a genuine new obstruction with no selected intrinsic source line.
\item If the residual top line survives but
$\beta_N^{\rm eff}(\mathfrak o_N^{\rm cc})\ne0$, the primitive residual line exists while the full
mixed-characteristic relation has a characteristic closed-lift obstruction.
\item If the residual top line has a marked free rank-one integral lift and the effective Bockstein
vanishes, \cref{P1C-thm:carry-compatible-filtered-recursion} attaches the relative integral recursive cell
and advances the branch by one effective weight.
\end{enumerate}
After each nonzero repeated line is produced, the fixed-order PD-extension theorem supplies all mixed
direction polarizations at that arity.  Hence for every finite ceiling the Confluent theory gives a
complete \emph{recursive gate mechanism}: every failure is localized either to residual arithmetic
nonvanishing/provenance, to the marked integral obstruction lattice, or to the effective closed-lift
Bockstein.  No theorem asserts that these gates succeed at every arity.

The mechanism is natural under support deletion, permutation and confluence, is compatible with every
finite coefficient reduction $\mathcal O/q^r$, and never identifies jet order, confluence multiplicity or
coefficient depth.  Moreover, whenever a detected repeated sector is projected to a finite Galois coefficient line
$T/q^rT$ with $\operatorname{cd}_q(G)\le2$, outcome~(c) disappears on that finite-depth face.  For the
$q$-adic coefficient line $T$ itself, the same conclusion holds when the finite-level $H^2$ system is
Mittag--Leffler.  Under these hypotheses \cref{P1C-thm:cd2-effective-bockstein-vanishing} forces the projected
effective Bockstein to vanish at every finite arity.  Thus the arithmetic rank-one Galois root faces used
below have only the residual source/provenance and integral-lattice gates; any surviving characteristic
$H^3$ gate must be transverse to that coefficient line, fail the Mittag--Leffler hypothesis, or belong to a
controller of larger cohomological dimension.
\end{theorem}

\begin{proof}
Apply \cref{P1C-thm:effective-weight-higher-recursion} successively in effective weights
$2,\ldots,M$.  The residual top-face identity isolates the arithmetic repeated sector.  Exactness of
\eqref{P1C-eq:effective-closed-lift-bockstein} gives the characteristic closed-lift gate.  Once the residual
line, its marked integral lattice and a closed full effective lift are available,
\cref{P1C-thm:carry-compatible-filtered-recursion} provides the relative cell.  Naturality and finite-depth
compatibility are inherited from these constructions.  Fixed-order mixed directions are supplied only
afterwards by \cref{P1C-thm:pd-extension-genuine-additive-provenance}.
\end{proof}

\begin{corollary}[The quartic repeated/root effective lift is automatic]
\label{P1C-cor:lambda3-carry-compatible-quartic-line}
Keep the exact-$\lambda=3$ calibration of \cref{P1C-thm:lambda3-heisenberg-central-gate}.  If the equivalent
Artin/class-group conditions of \cref{P1C-cor:lambda3-quartic-handoff} hold, then the nonzero residual
quartic coefficient $q_4E$ has a primitive lift in
\[
 WE=H^2(G_{K,S},\Z_5(1))E,
\]
and therefore determines the canonical unpointed primitive quartic repeated line
\[
 \boxed{
 \mathcal L^{\rm cc}_{3,1}
 :=WE\otimes_{\Z_5}\Z_5\gamma_3(e_q)e_v.}
\]
Let
\[
 \mathscr M_E:=C^\bullet_{\rm cont}(G_{K,S},\Z_5(1))E
\]
be the root coefficient face of the marked endomorphism controller.  Since $\operatorname{cd}_5(G_{K,S})=2$ and the finite-level groups $H^2(G_{K,S},\Z/5^r(1))$ are finite, the inverse system is Mittag--Leffler; hence \cref{P1C-thm:cd2-effective-bockstein-vanishing} gives
\begin{equation}\label{P1C-eq:lambda3-effective-bockstein-gate}
 \boxed{
 \beta_{4,E}^{\rm eff}
 \bigl(\pi_E\mathfrak o_4^{\rm cc}\bigr)=0.}
\end{equation}
Hence the repeated/root effective obstruction admits a closed effective-weight-four lift
$\widetilde\Omega_{4,E}$ with primitive top coefficient $q_4E$, and one may adjoin the root-facing
relative cell
\[
 \boxed{dh_{4,E}=-\widetilde\Omega_{4,E}.}
\]
Thus there is no additional characteristic-$H^3$ obstruction on the quartic repeated/root face.
Different root-facing closed lifts can still differ by the higher-filtration $H^2$ torsor of
\cref{P1C-prop:effective-closed-lift-criterion}; fixing a preferred arithmetic pointing is provenance data.

Consequently the single intrinsic arithmetic existence test for the exact-$\lambda=3$ repeated/root
quartic line is the degree-$25$ Artin/class-group certificate.  A full End-valued carry-corrected
quartic relation can still have transverse effective Bocksteins in the complementary matrix corners;
the present statement does not identify or kill those transverse classes and does not require an
ordinary integral four-fold Massey defining system in the unextended Galois cochain algebra.
\end{corollary}

\begin{proof}
The primitive line and nonzero residual root projection are
\cref{P1C-cor:lambda3-quartic-handoff}.  The coefficient projection
\[
 C^\bullet(G_{K,S},\operatorname{End}_{\Z_5}(\Z_5\oplus\Z_5(1)))
 \longrightarrow \mathscr M_E
\]
is a dg projection because the $E_{21}$ coefficient line is the $\Z_5(1)$ weight summand.  The restricted global Galois group has $5$-cohomological dimension two and the finite-level $H^2$ groups on this coefficient line are finite, hence Mittag--Leffler.  Therefore \cref{P1C-cor:root-projected-effective-bockstein-zero} proves
\eqref{P1C-eq:lambda3-effective-bockstein-gate}.  The closed-lift criterion then supplies
$\widetilde\Omega_{4,E}$, and \cref{P1C-thm:carry-compatible-filtered-recursion} attaches the relative
root-facing cell.  The final firewall is exactly the scope statement of
\cref{P1C-cor:root-projected-effective-bockstein-zero}.
\end{proof}

\begin{definition}[Recursive stage status]\label{P1C-def:recursive-stage-status}
For each finite recursive stage $n$, define
\[
 \operatorname{stat}_n
 :=(\operatorname{res}_n,\operatorname{lat}_n,
    \operatorname{bock}_n,\operatorname{cell}_n)
\]
where the four entries record respectively: survival of the marked residual detector, existence of a marked free integral source lattice, vanishing/nonvanishing of the effective closed-lift Bockstein, and whether the branch extends internally or by a relative recursive cell.  This status tuple is bookkeeping only; it does not identify any of the support, multiplicity, depth, jet or PD-order coordinates.
\end{definition}

\begin{theorem}[Finite branchwise higher-recursion tower]
\label{P1C-thm:finite-branchwise-higher-recursion}
Fix a finite ceiling $M\ge3$ and a pointed lower branch.  The homogeneous obstruction classes
$\mathfrak o_n$ and the effective-weight classes $\mathfrak o_n^{\rm cc}$ assemble into the finite
recursive ledger
\[
 \boxed{
 \mathfrak R^{\rm Conf}_{S,\le M}
 :=\{\mathfrak o_n,\;\mathfrak o_n^{\rm cc},\;
       \mathfrak o^{\rm rec}_{n-1,1},\;
       \beta_n^{\rm eff},\;\operatorname{stat}_n\}_{3\le n\le M}.}
\]
The top residual face of $\mathfrak o_n^{\rm cc}$ is $\mathfrak o_n$, so the filtered ledger refines
rather than replaces the homogeneous recursion.  The ledger is natural under direction-linear maps
and finite coefficient reductions.  Together with
\cref{P1C-thm:finite-confluent-interface} it gives the finite higher-support package.
\end{theorem}

\begin{warning}[Effective weight is not coefficient depth]\label{P1C-warn:effective-weight-not-depth}
The equality ``parameter-degree loss $+$ $q$-filtration gain $=$ effective weight'' is a filtration rule
inside the recursive controller, not an identification of jet order with coefficient depth.  A class may
have effective weight $N$ at many finite coefficient depths $q^r$, and the same coefficient line may
deepen in $r$ without producing any new jet-order obstruction.  The support/multiplicity/depth
firewall therefore remains unchanged.
\end{warning}

\begin{warning}[Input boundary for carry-compatible attachment]
\label{P1C-warn:cc-recursion-not-provenance}
The carry-compatible attachment theorem starts from a marked integral obstruction
lattice; it does not derive that lattice or its arithmetic pointing.  Its conclusion
is the characteristic-lift mechanism forced by lower nonzero carries.
\end{warning}

\subsection{Coefficient depth continues without a new arity}

\begin{proposition}[The same primitive line has arbitrary finite coefficient depth]\label{P1C-prop:primitive-line-arbitrary-depth}
Let $W_q=\Z_q\widetilde\kappa$ be the primitive filtered confluence line of \cref{P1C-def:filtered-confluence-lift}.  For every finite $r\ge1$ put
\[
 L^{\rm coeff}_{q^r}:=W_q/q^rW_q.
\]
Then
\[
 \boxed{
 L^{\rm coeff}_{q^r}\simeq\Z/q^r,
 \qquad
 \Ord(\widetilde\kappa\bmod q^r)=q^r.}
\]
In particular exact coefficient depth $q^3$ exists on the same primitive line even if the formal-moduli germ carries no independent quartic $(3,1)$ invariant.
\end{proposition}

\begin{proof}
Primitivity identifies $W_q$ with $\Z_q$ and $\widetilde\kappa$ with a $q$-adic unit generator.  Reduction modulo $q^r$ gives the displayed cyclic group and exact order.
\end{proof}

\begin{principle}[Multiplicity--depth firewall]\label{P1C-prin:multiplicity-depth-firewall}
The exponent $r$ in a coefficient line of order $q^r$ and the repeated-input multiplicity $m$ in a formal-moduli multiweight $(m,1)$ are independent coordinates.  The minimal $q^2/(2,1)$ model supplies a natural pointed alignment at $r=m=2$; no rule $r=m$ is propagated beyond that seed without an additional theorem.
\end{principle}

\subsection{The divided-power confluence calculus}

Fix a marked obstruction line $k\kappa_m$ and work in the free divided-power parameter algebra.  For a composition $\lambda=(m_1,\ldots,m_s)$ of $m$ define the normalized confluence profile by the one-cell relation
\begin{equation}\label{P1C-eq:partition-profile-cell}
 d\eta_{\lambda;1}
 =\gamma_{m_1}(u_1)\cdots\gamma_{m_s}(u_s)v\,\kappa_m.
\end{equation}
The fully split profile is $(1,\ldots,1)$, while the fully confluent profile is $(m)$.

\begin{theorem}[Binary block-confluence coefficient]\label{P1C-thm:block-confluence-coefficient}
Suppose two blocks of multiplicities $a$ and $b$ are identified by $u_a,u_b\mapsto u$.  Then the base change of \eqref{P1C-eq:partition-profile-cell} contains the factor
\[
 \boxed{
 \gamma_a(u)\gamma_b(u)
 =\binom{a+b}{a}\gamma_{a+b}(u).}
\]
Thus the raw block collision differs from the normalized $(a+b)$-block cell by the scalar $\binom{a+b}{a}$.
\end{theorem}

\begin{proof}
This is the multiplication axiom in a divided-power algebra:
$\gamma_a(u)\gamma_b(u)=\frac{(a+b)!}{a!b!}\gamma_{a+b}(u)$.
\end{proof}

\begin{corollary}[The second confluence has coefficient three]\label{P1C-cor:second-confluence-three}
Let the partially confluent quartic cell have relation
\[
 d\eta_{211}=\gamma_2(u)wv\,\kappa_3.
\]
On the diagonal $w=u$,
\[
 \boxed{
 d\eta_{211}\longmapsto
 3\gamma_3(u)v\,\kappa_3.}
\]
Hence comparison with the normalized $(3,1)$ cell
\[
 d\eta_{31}=\gamma_3(u)v\,\kappa_3
\]
is a unit rescaling over $\Z_q$ exactly when $q\ne3$.  At $q=3$ the map on the one-cell obstruction summand is multiplication by $3$ and has cokernel $\F_3$.
\end{corollary}

\begin{proof}
Apply \cref{P1C-thm:block-confluence-coefficient} with $(a,b)=(2,1)$.  Multiplication by $3$ is a unit of $\Z_q$ precisely when $q\ne3$; over $\Z_3$ its cokernel is $\Z_3/3\Z_3$.
\end{proof}

\begin{theorem}[Merge-tree coherence]\label{P1C-thm:merge-tree-coherence}
Start from a profile $\lambda=(m_1,\ldots,m_s)$ of total multiplicity $m$ and merge blocks pairwise until a single block of size $m$ remains.  For every binary merge tree, the product of the block-confluence coefficients is
\[
 \boxed{
 \frac{m!}{m_1!\cdots m_s!}.}
\]
In particular, starting from $m$ singleton slots, every merge tree has total scalar $m!$.  Hence the formal divided-power confluence calculus has no scalar associator ambiguity: different parenthesizations give the same normalized comparison coefficient.
\end{theorem}

\begin{proof}
At an internal vertex merging blocks of sizes $a,b$, attach the factor
$\frac{(a+b)!}{a!b!}$.  Multiplying over the tree cancels every factorial attached to an internal child against the corresponding denominator at its parent.  Only the root factorial $m!$ and the leaf factorials $m_i!$ remain.  The result is therefore independent of the tree.
\end{proof}

\subsection{Prime-primary confluence carries}

\begin{definition}[Confluence-carry index]\label{P1C-def:confluence-carry-index}
For an odd prime $q$ and a block merge $(a,b)$ define
\[
 \boxed{
 \chi_q(a,b):=v_q\binom{a+b}{a}.}
\]
For complete confluence of $m$ singleton slots define the factorial defect
\[
 \boxed{
 \delta_q(m):=v_q(m!).}
\]
\end{definition}

\begin{theorem}[Prime-primary order loss under raw confluence]\label{P1C-thm:confluence-order-loss}
Let $x$ be a normalized target generator of exact order $q^r$.  Under a raw block collision of sizes $a,b$, the corresponding scalar image has exact order
\[
 \boxed{
 q^{\max(r-\chi_q(a,b),0)}.}
\]
For complete confluence from $m$ singleton slots the raw diagonal image has order
\[
 \boxed{
 q^{\max(r-\delta_q(m),0)}.}
\]
Thus a positive confluence-carry index measures exactly how many $q$-power layers are lost by using the unrenormalized diagonal collision.
\end{theorem}

\begin{proof}
The first statement combines \cref{P1C-thm:block-confluence-coefficient} with the elementary scalar-order lemma in the appendix.  The second uses \cref{P1C-thm:merge-tree-coherence} and the same lemma.
\end{proof}

\begin{proposition}[First factorial resonance]\label{P1C-prop:first-factorial-resonance}
For an odd prime $q$,
\[
 \delta_q(m)=v_q(m!)
 =\sum_{j\ge1}\left\lfloor\frac{m}{q^j}\right\rfloor.
\]
Consequently
\[
 \boxed{
 \delta_q(m)=0\quad(1\le m<q),
 \qquad
 \delta_q(q)=1.}
\]
In particular the multiplicity-three collision $2+1\to3$ is resonant exactly at the prime $q=3$.
\end{proposition}

\begin{proof}
The displayed valuation is Legendre's formula.  If $m<q$ all floors vanish.  At $m=q$ only the first floor contributes and equals one.  The final statement is also immediate from $\binom31=3$.
\end{proof}

\begin{definition}[Base-$q$ digit sum]\label{P1C-def:base-q-digit-sum}
For $n\ge0$ write
\[
 n=\sum_{j\ge0}n_jq^j,
 \qquad 0\le n_j<q,
\]
and put
\[
 \boxed{s_q(n):=\sum_{j\ge0}n_j.}
\]
\end{definition}

\begin{theorem}[Kummer carry formula for a block merge]\label{P1C-thm:kummer-block-carry}
For every $a,b\ge0$,
\[
 \boxed{
 \chi_q(a,b)
 =\frac{s_q(a)+s_q(b)-s_q(a+b)}{q-1}.}
\]
Equivalently, $\chi_q(a,b)$ is the number of carries occurring in the ordinary base-$q$ addition of $a$ and $b$.
\end{theorem}

\begin{proof}
Legendre's formula in digit-sum form is
\[
 v_q(n!)=\frac{n-s_q(n)}{q-1}.
\]
Subtracting the formulas for $a!$ and $b!$ from that for $(a+b)!$ gives the displayed identity.  To identify the right-hand side with carries, write
\[
 a_j+b_j+c_{j-1}=d_j+qc_j,
 \qquad c_{-1}=0,
\]
where $d_j$ is the $j$-th digit of $a+b$.  For a binary sum each $c_j$ is $0$ or $1$.  Summing over all digit positions after adjoining enough terminal zero digits gives
\[
 s_q(a)+s_q(b)-s_q(a+b)=(q-1)\sum_jc_j.
\]
Thus $\chi_q(a,b)=\sum_jc_j$.
\end{proof}

\begin{corollary}[Carry-free block criterion]\label{P1C-cor:carry-free-block}
For a block merge $(a,b)$ the following are equivalent:
\begin{enumerate}[label=\textup{(\roman*)}]
\item $\chi_q(a,b)=0$;
\item $\binom{a+b}{a}$ is a unit in $\Z_q$;
\item the base-$q$ addition of $a$ and $b$ is carry-free;
\item if $a=\sum a_jq^j$ and $b=\sum b_jq^j$, then $a_j+b_j<q$ for every $j$.
\end{enumerate}
Hence a raw block confluence preserves full $q$-power depth exactly in the carry-free case.
\end{corollary}

\begin{definition}[Profile carry mass]\label{P1C-def:profile-carry-mass}
Let $\lambda=(m_1,\ldots,m_s)$ be a multiplicity profile of total mass
$m=\sum_i m_i$.  Define its multinomial scalar and $q$-carry mass by
\[
 M(\lambda):=\frac{m!}{m_1!\cdots m_s!},
 \qquad
 \boxed{C_q(\lambda):=v_q(M(\lambda)).}
\]
\end{definition}

\begin{theorem}[Multinomial digit formula and carry conservation]\label{P1C-thm:multinomial-carry}
For every multiplicity profile $\lambda=(m_1,\ldots,m_s)$ of total mass $m$,
\[
 \boxed{
 C_q(\lambda)
 =\frac{\sum_i s_q(m_i)-s_q(m)}{q-1}.}
\]
If the $m_i$ are added simultaneously in base $q$ and $c_j\ge0$ denotes the carry emitted from the $j$-th digit column, then
\[
 \boxed{C_q(\lambda)=\sum_{j\ge0}c_j.}
\]
Moreover, for every binary merge tree from $\lambda$ to the one-block profile $(m)$,
\[
 \boxed{
 C_q(\lambda)
 =\sum_{v\in\operatorname{Int}(T)}\chi_q(a_v,b_v),}
\]
where $(a_v,b_v)$ are the two block sizes merged at the internal vertex $v$.  Thus total prime-primary carry is conserved under reparenthesization.
\end{theorem}

\begin{proof}
The digit formula follows by applying Legendre's formula to the multinomial coefficient.  For the simultaneous addition write
\[
 \sum_i m_{i,j}+c_{j-1}=d_j+qc_j,
 \qquad c_{-1}=0,
\]
where $d_j$ is the $j$-th digit of $m$.  Summing over digit positions gives
\[
 \sum_i s_q(m_i)-s_q(m)=(q-1)\sum_jc_j.
\]
Finally, \cref{P1C-thm:merge-tree-coherence} says that the product of the binary merge scalars is $M(\lambda)$; taking $q$-adic valuations gives the tree formula.
\end{proof}

\begin{corollary}[Carry-free profile classification]\label{P1C-cor:carry-free-profile}
For a multiplicity profile $\lambda=(m_1,\ldots,m_s)$ the following are equivalent:
\begin{enumerate}[label=\textup{(\roman*)}]
\item $C_q(\lambda)=0$;
\item $M(\lambda)$ is a $q$-adic unit;
\item the simultaneous base-$q$ addition $m_1+\cdots+m_s=m$ is carry-free;
\item at every digit position $j$, the sum of the $j$-th digits of the $m_i$ is strictly less than $q$.
\end{enumerate}
For such a profile every merge tree has total $q$-valuation zero, although individual integral merge coefficients may differ by units.
\end{corollary}

\begin{definition}[Confluence-carry complex]\label{P1C-def:confluence-carry-complex}
For a profile $\lambda$ define the two-term $q$-adic complex
\[
 \boxed{
 \mathsf K_q(\lambda)
 :=[\Z_q\xrightarrow{\ M(\lambda)\ }\Z_q],}
\]
with the source in degree $-1$ and the target in degree $0$.
\end{definition}

\begin{theorem}[Finite Moore defect of a confluence profile]\label{P1C-thm:confluence-carry-complex}
For every profile $\lambda$,
\[
 H^{-1}(\mathsf K_q(\lambda))=0,
 \qquad
 \boxed{H^0(\mathsf K_q(\lambda))\simeq\Z/q^{C_q(\lambda)}.}
\]
In particular $\mathsf K_q(\lambda)$ is acyclic exactly for carry-free profiles.  If a normalized target line has exact order $q^r$, raw confluence along $\lambda$ leaves exact order
\[
 \boxed{q^{\max(r-C_q(\lambda),0)}.}
\]
Thus $C_q(\lambda)$ is simultaneously the carry mass, the $q$-adic length of the confluence defect complex, and the number of coefficient-depth layers lost by raw confluence.
\end{theorem}

\begin{proof}
Write $M(\lambda)=q^{C_q(\lambda)}u$ with $u\in\Z_q^\times$.  Multiplication by $M(\lambda)$ on the domain $\Z_q$ is injective and has cokernel isomorphic to $\Z/q^{C_q(\lambda)}$.  The order statement is the scalar-order lemma applied to $M(\lambda)$.
\end{proof}

\begin{definition}[$q$-typical multiplicity profile]\label{P1C-def:q-typical-profile}
Write the base-$q$ expansion
\[
 m=\sum_{j\ge0}m_jq^j,
 \qquad 0\le m_j<q.
\]
The $q$-typical profile of $m$ is the multiset
\[
 \boxed{
 \lambda_q^{\rm typ}(m)
 :=\{\underbrace{q^j,\ldots,q^j}_{m_j\text{ copies}}:j\ge0\}.}
\]
It has $s_q(m)$ blocks.
\end{definition}

\begin{theorem}[$q$-typical normal form]\label{P1C-thm:q-typical-normal-form}
The profile $\lambda_q^{\rm typ}(m)$ is carry-free:
\[
 \boxed{C_q(\lambda_q^{\rm typ}(m))=0.}
\]
It is the unique profile of $m$ by powers of $q$ in which every block size $q^j$ occurs fewer than $q$ times.  Consequently the confluence map from the $q$-typical split to the normalized multiplicity-$m$ line is a $q$-adic unit equivalence.

Starting instead from $m$ singleton blocks, repeatedly replace $q$ blocks of size $q^j$ by one block of size $q^{j+1}$.  This process terminates at $\lambda_q^{\rm typ}(m)$ and uses exactly
\[
 \boxed{
 \delta_q(m)=v_q(m!)=\frac{m-s_q(m)}{q-1}}
\]
$q$-ary carry moves.  Each such elementary $q$-ary carry has $q$-valuation exactly one.  Hence all nonunit $q$-primary content of the full singleton confluence scalar $m!$ is concentrated in these canonical carry moves; the final $q$-typical-to-one-block comparison is a unit.
\end{theorem}

\begin{proof}
The digit blocks in $\lambda_q^{\rm typ}(m)$ add without carry, so the first assertion is \cref{P1C-cor:carry-free-profile}.  Uniqueness is uniqueness of the base-$q$ expansion.  The replacement rule is exactly the usual carrying algorithm and therefore terminates at that expansion.

Each replacement of $q$ blocks of size $q^j$ by one block of size $q^{j+1}$ reduces the number of blocks by $q-1$.  Starting with $m$ blocks and ending with $s_q(m)$ blocks, the number of replacements is $(m-s_q(m))/(q-1)=\delta_q(m)$.  Its divided-power scalar is
\[
 \frac{(q^{j+1})!}{(q^j!)^q},
\]
and Legendre's formula gives valuation one.  Summing over the $\delta_q(m)$ moves accounts for the full valuation of $m!$; the remaining comparison is carry-free and hence a unit.
\end{proof}

\begin{definition}[Finite confluence category]\label{P1C-def:finite-confluence-category}
Fix $m$.  Let $\mathsf{Conf}(m)$ be the category whose objects are finite ordered multiplicity profiles of total mass $m$ and whose generating arrows merge two blocks.  For a coarsening $f:\lambda\to\mu$, let $M(f)$ be the product of the corresponding block-confluence coefficients.
\end{definition}

\begin{theorem}[Carry valuation is a functorial length]\label{P1C-thm:carry-functorial-length}
The scalar $M(f)$ depends only on the coarsening $f$, not on its factorization into elementary merges, and
\[
 M(g\circ f)=M(g)M(f).
\]
Therefore
\[
 \boxed{\ell_q(f):=v_q(M(f))}
\]
is an additive length:
\[
 \ell_q(g\circ f)=\ell_q(g)+\ell_q(f).
\]
The zero-length arrows form a wide subcategory on which the induced maps of rank-one $\Z_q$ confluence lines are isomorphisms.  Positive-length arrows form a composition ideal.  Thus the finite confluence category carries a canonical prime-primary resonance filtration by the integer $\ell_q$.
\end{theorem}

\begin{proof}
For a coarsening, each target block is the union of a specified collection of source blocks.  Applying \cref{P1C-thm:merge-tree-coherence} inside each target block gives
\[
 M(f)=\prod_{B\in\mu}
 \frac{(\sum_{i\in B}m_i)!}{\prod_{i\in B}m_i!},
\]
which is independent of the merge factorization.  The formula is multiplicative under refinement of block partitions, proving the composition law.  Taking valuations gives additivity.  The final assertions follow because a scalar in $\Z_q$ is invertible exactly when its valuation is zero.
\end{proof}

\subsection{Integral resonance divisors and serial Moore composition}

The prime-primary lengths above come from one integral scalar. Keeping that scalar before
localization exposes a useful distinction between carries at different primes and repeated carries
at the same prime.

\begin{definition}[Integral confluence defect]\label{P1C-def:integral-confluence-defect}
For a confluence morphism $f:\lambda\to\mu$ with positive integral scalar $M(f)$, define
\[
 \boxed{
 \mathsf K_{\Z}(f):=[\Z\xrightarrow{\ M(f)\ }\Z]}
\]
with the source in degree $-1$ and the target in degree $0$. Its defect group is
\[
 \mathsf D(f):=H^0(\mathsf K_{\Z}(f)).
\]
Define the resonance divisor
\[
 \boxed{
 \operatorname{Res}(f):=\sum_{\ell\mid M(f)}v_\ell(M(f))[\ell],}
\]
where the sum runs over the finitely many prime divisors of $M(f)$.
\end{definition}

\begin{theorem}[Integral Moore defect and canonical primary decomposition]
\label{P1C-thm:integral-primary-decomposition}
For every confluence morphism $f$,
\[
 H^{-1}(\mathsf K_{\Z}(f))=0,
 \qquad
 \boxed{\mathsf D(f)\simeq\Z/M(f)\Z.}
\]
Moreover the primary torsion subgroups give a canonical decomposition
\[
 \boxed{
 \mathsf D(f)
 \simeq
 \bigoplus_{\ell\mid M(f)}\Z/\ell^{c_\ell(f)}\Z,
 \qquad
 c_\ell(f):=v_\ell(M(f)).}
\]
After extension of scalars to $\Z_\ell$, the complex $\mathsf K_{\Z}(f)$ is isomorphic up to
unit rescaling to
\[
 [\Z_\ell\xrightarrow{\ \ell^{c_\ell(f)}\ }\Z_\ell].
\]
Hence the resonance divisor records exactly the primary lengths of the integral Moore defect.
\end{theorem}

\begin{proof}
Multiplication by the nonzero integer $M(f)$ is injective on $\Z$, so the cokernel is
$\Z/M(f)\Z$. The displayed direct sum is the canonical primary decomposition of a finite
cyclic group. Writing $M(f)=\ell^{c_\ell(f)}u_\ell$ with
$u_\ell\in\Z_\ell^\times$ gives the local two-term model after multiplying one basis vector by
$u_\ell^{-1}$.
\end{proof}

\begin{corollary}[Resonance divisor is additive]\label{P1C-cor:resonance-divisor-additive}
For composable confluence morphisms $f,g$,
\[
 \boxed{
 \operatorname{Res}(g\circ f)
 =\operatorname{Res}(f)+\operatorname{Res}(g).}
\]
In particular $f$ is an integral unit confluence if and only if its resonance divisor is zero.
For a fixed total multiplicity $m$, every prime in the support of $\operatorname{Res}(f)$ is at most
$m$.
\end{corollary}

\begin{proof}
By \cref{P1C-thm:carry-functorial-length}, $M(g\circ f)=M(g)M(f)$. Taking every prime
valuation gives the divisor identity. Every confluence scalar is a product of multinomial
coefficients formed from integers at most $m$, so no prime larger than $m$ occurs.
\end{proof}


\begin{corollary}[Divisor-valued complete scalar invariant]
\label{P1C-cor:resonance-divisor-complete}
For every confluence morphism $f$,
\[
 \boxed{
 M(f)=\prod_{\ell\le m}\ell^{c_\ell(f)},
 \qquad
 \operatorname{Res}(f)=\sum_{\ell\le m}c_\ell(f)[\ell].}
\]
Thus the resonance divisor is a complete invariant of the positive integral scalar $M(f)$.
Equivalently, on morphisms of $\mathsf{Conf}(m)$ there is a canonical additive map
\[
 \operatorname{Res}:\operatorname{Mor}\mathsf{Conf}(m)
 \longrightarrow
 \bigoplus_{\ell\le m}\N[\ell]
\]
sending composition to addition; the coordinate at $q$ is precisely the carry length $\ell_q$.
\end{corollary}

\begin{proof}
This is unique prime factorization together with
\cref{P1C-cor:resonance-divisor-additive}.
\end{proof}

\begin{theorem}[Serial composition exact sequence]\label{P1C-thm:serial-defect-extension}
Let
\[
 \lambda\xrightarrow{f}\mu\xrightarrow{g}
u
\]
be composable confluence morphisms and put $A=M(f)$, $B=M(g)$. Then there is a canonical
short exact sequence
\[
 \boxed{
 0\longrightarrow \mathsf D(f)
 \xrightarrow{\ \iota\ }
 \mathsf D(g\circ f)
 \xrightarrow{\ \pi\ }
 \mathsf D(g)
 \longrightarrow0,}
\]
where
\[
 \iota([x]_A)=[Bx]_{AB},
 \qquad
 \pi([y]_{AB})=[y]_B.
\]
At a prime $q$, if
\[
 a=\ell_q(f),\qquad b=\ell_q(g),
\]
the $q$-primary part is
\[
 \boxed{
 0\to\Z/q^a\to\Z/q^{a+b}\to\Z/q^b\to0.}
\]
If $a,b>0$ this primary extension is nonsplit. Thus repeated carry at one prime stacks
\emph{serially} into one longer cyclic Moore defect rather than a direct sum of the two defects.
\end{theorem}

\begin{proof}
The formulas for $\iota$ and $\pi$ are well defined. The first is injective, the second is
surjective, and the kernel of reduction modulo $B$ is the subgroup of multiples of $B$, exactly the
image of $\iota$. Passing to the canonical $q$-primary subgroups gives the displayed sequence.
If $a,b>0$ and it split, the middle group would be isomorphic to
$\Z/q^a\oplus\Z/q^b$, whose exponent is $q^{\max(a,b)}$, whereas
$\Z/q^{a+b}$ has exponent $q^{a+b}$. Hence no splitting is possible.
\end{proof}

\begin{corollary}[Orthogonal-prime versus serial-prime dichotomy]
\label{P1C-cor:orthogonal-serial-dichotomy}
With notation as in \cref{P1C-thm:serial-defect-extension}:
\begin{enumerate}[label=\textup{(\roman*)}]
\item if $\gcd(A,B)=1$, the exact sequence splits canonically by the Chinese remainder theorem,
so the two defect stages are prime-orthogonal;
\item if a prime $q$ divides both $A$ and $B$, the $q$-primary carry lengths add inside the single
cyclic group $\Z/q^{a+b}$ and the corresponding primary extension is nonsplit.
\end{enumerate}
Thus confluence composition is parallel across distinct primes and serial along repeated occurrences
of the same prime.
\end{corollary}

\begin{proof}
When $(A,B)=1$, CRT gives the unique class $e_B\in\Z/(AB)$ satisfying
$e_B\equiv0\pmod A$ and $e_B\equiv1\pmod B$; sending $[y]_B$ to $[ye_B]_{AB}$ is a section
of $\pi$. The second assertion is the primary statement of
\cref{P1C-thm:serial-defect-extension}.
\end{proof}

\begin{theorem}[Complete singleton integral defect]\label{P1C-thm:singleton-integral-defect}
Let $\mathbf 1^m=(1,\ldots,1)$ be the $m$-singleton profile and let
$f_m:\mathbf1^m\to(m)$ be complete confluence. Then
\[
 \boxed{
 \mathsf D(f_m)\simeq\Z/m!\Z,
 \qquad
 \operatorname{Res}(f_m)
 =\sum_{\ell\le m}v_\ell(m!)[\ell].}
\]
For every prime $q\le m$, its $q$-primary summand is
\[
 \boxed{
 \mathsf D(f_m)_{(q)}\simeq\Z/q^{\delta_q(m)}\Z,
 \qquad
 \delta_q(m)=\frac{m-s_q(m)}{q-1}.}
\]
For odd $q$, the canonical $q$-typical carry algorithm of
\cref{P1C-thm:q-typical-normal-form} realizes this local length as exactly $\delta_q(m)$ elementary
valuation-one carry events.
\end{theorem}

\begin{proof}
The complete singleton scalar is $m!$ by \cref{P1C-thm:merge-tree-coherence}. Apply
\cref{P1C-thm:integral-primary-decomposition} and Legendre's digit formula. The final statement is
\cref{P1C-thm:q-typical-normal-form}.
\end{proof}


\begin{example}[Two successive ternary carries at the prime three]
Consider the coarsening chain
\[
 (2,1,6)\longrightarrow(3,6)\longrightarrow(9).
\]
The two integral confluence scalars are
\[
 \binom31=3,
 \qquad
 \binom93=84.
\]
Both have $3$-adic valuation one.  Hence the composite has $3$-primary length two and the
primary part of \cref{P1C-thm:serial-defect-extension} is
\[
 \boxed{
 0\longrightarrow\Z/3
 \longrightarrow\Z/9
 \longrightarrow\Z/3
 \longrightarrow0.}
\]
The middle term is cyclic of order nine, so the two carry events form one serial depth-two defect.
\end{example}

\begin{example}[Prime-orthogonal composition]
For the chain
\[
 (2,1,2)\longrightarrow(3,2)\longrightarrow(5)
\]
the scalars are
\[
 \binom31=3,
 \qquad
 \binom52=10.
\]
They are coprime.  Thus the composite defect is
\[
 \mathsf D\simeq\Z/30
 \simeq\Z/3\oplus\Z/10,
\]
and \cref{P1C-cor:orthogonal-serial-dichotomy} gives the canonical CRT splitting.  This is the
finite parallel behavior complementary to the serial same-prime example above.
\end{example}

\begin{warning}[The prime two in the integral defect]\label{P1C-warn:two-primary-integral-defect}
The integral scalar and its primary decomposition are valid at the prime $2$ as elementary
arithmetic statements. This does not extend the earlier repeated-input Cartan/partition-Lie
realization to characteristic two: the minimal confluence cell used $2$ invertible, and the binary
primary obstruction sector remains a separate problem. Thus the $2$-primary summand of
$\mathsf D(f)$ is presently a combinatorial Moore defect only.
\end{warning}


\subsection{Carry filtrations and the associated-graded Moore category}

The serial exact sequence has a canonical internal filtration.  This upgrades the numerical carry
length to a filtration on both the defect modules and the confluence category itself.


\begin{definition}[Factorial weight and prime height]\label{P1C-def:factorial-height}
For a multiplicity profile $\lambda=(m_1,\ldots,m_s)$ define its factorial weight
\[
 \boxed{w(\lambda):=\prod_{i=1}^s m_i!}
\]
and, for every prime $q$, its factorial $q$-height
\[
 \boxed{h_q(\lambda):=v_q(w(\lambda)).}
\]
The integral divisor height is
\[
 \Phi(\lambda):=\sum_{\ell\le m}h_\ell(\lambda)[\ell].
\]
\end{definition}

\begin{theorem}[Factorial-potential formula]\label{P1C-thm:factorial-potential}
For every confluence morphism $f:\lambda\to\mu$,
\[
 \boxed{
 M(f)=\frac{w(\mu)}{w(\lambda)},
 \qquad
 \ell_q(f)=h_q(\mu)-h_q(\lambda),
 \qquad
 \operatorname{Res}(f)=\Phi(\mu)-\Phi(\lambda).}
\]
In particular every factorial prime height is nondecreasing under confluence, and the carry length is
an exact height difference rather than merely a path-additive morphism invariant.
\end{theorem}

\begin{proof}
Write $\mu=(n_1,\ldots,n_t)$, where each $n_j$ is the sum of the source multiplicities merged into
the $j$-th target block.  The factorization-independent scalar formula gives
\[
 M(f)=\prod_{j=1}^t\frac{n_j!}{\prod_{i\in B_j}m_i!}
 =\frac{\prod_j n_j!}{\prod_i m_i!}
 =\frac{w(\mu)}{w(\lambda)}.
\]
Taking prime valuations gives the remaining formulas.  Nonnegativity follows because $M(f)$ is a
positive integer.
\end{proof}

\begin{corollary}[Exactness and absence of scalar holonomy]\label{P1C-cor:factorial-no-holonomy}
The multiplicative scalar cocycle $M$ and the additive resonance-divisor cocycle are both exact on
$\mathsf{Conf}(m)$: they are respectively the ratio and difference of object potentials $w$ and
$\Phi$.  Consequently every two directed confluence paths with the same endpoints have identical
integral scalar, prime-primary lengths and resonance divisor.
\end{corollary}

\begin{theorem}[Rational trivialization and integral lattice defect]\label{P1C-thm:rational-lattice-trivialization}
Let $V=\Q\xi$ be a one-dimensional rational vector space and assign to a profile $\lambda$ the lattice
\[
 \boxed{
 \Lambda_\lambda:=w(\lambda)^{-1}\Z\,\xi\subset V.}
\]
For every confluence morphism $f:\lambda\to\mu$ one has a canonical inclusion
\[
 \Lambda_\lambda\subseteq\Lambda_\mu
\]
of index $M(f)$, and
\[
 \boxed{
 \Lambda_\mu/\Lambda_\lambda\simeq\Z/M(f)\Z\simeq\mathsf D(f).}
\]
Equivalently, after the objectwise rational basis change
\[
 e_\lambda\longmapsto w(\lambda)^{-1}\xi,
\]
the rank-one confluence transport $e_\lambda\mapsto M(f)e_\mu$ becomes the identity map of the
common rational line $V$.  Thus the complete scalar confluence system is rationally gauge-trivial;
its nontrivial content is the finite index between the integral lattices.
\end{theorem}

\begin{proof}
By \cref{P1C-thm:factorial-potential}, $w(\mu)=M(f)w(\lambda)$.  Hence
\[
 \Lambda_\mu
 =\frac1{M(f)w(\lambda)}\Z\xi
 \supseteq
 \frac1{w(\lambda)}\Z\xi
 =\Lambda_\lambda,
\]
and the quotient has order and cyclic type $M(f)$.  For the transport formula,
\[
 M(f)e_\mu\longmapsto
 \frac{M(f)}{w(\mu)}\xi
 =\frac1{w(\lambda)}\xi,
\]
which is the image of $e_\lambda$.
\end{proof}

\begin{corollary}[Prime carry as local lattice index]\label{P1C-cor:q-carry-lattice-index}
After tensoring the common rational line with $\Q_q$, the localized lattices satisfy
\[
 \Lambda_{\lambda,q}\subseteq\Lambda_{\mu,q},
 \qquad
 \boxed{
 \operatorname{length}_{\Z_q}
 (\Lambda_{\mu,q}/\Lambda_{\lambda,q})
 =\ell_q(f).}
\]
Thus base-$q$ carry is exactly the local lattice-index length of confluence.
\end{corollary}


\subsection{Universal integralization of rationally trivial confluence}

The factorial lattice construction is not merely one convenient integral model.  Once the split
singleton object is primitively normalized, it is forced by naturality of the rationally trivial scalar
transport.

\begin{definition}[Normalized scalar confluence representation]\label{P1C-def:scalar-confluence-representation}
Let $\mathbf1^m=(1,\ldots,1)$ be the discrete singleton profile.  Define a rank-one rational
representation
\[
 \mathscr E_m:\mathsf{Conf}(m)\longrightarrow \operatorname{Vect}^{(1)}_{\Q}
\]
by
\[
 \mathscr E_m(\lambda)=\Q e_\lambda,
 \qquad
 \mathscr E_m(f)(e_\lambda)=M(f)e_\mu
\]
for $f:\lambda\to\mu$.  Multiplicativity of $M$ makes this a functor.
\end{definition}

\begin{theorem}[Universal integralization theorem]\label{P1C-thm:universal-integralization}
Fix a rational line $V=\Q\xi$ and the primitive normalization
\[
 \tau_{\mathbf1^m}(e_{\mathbf1^m})=\xi.
\]
Then the following hold.
\begin{enumerate}[label=\textup{(\roman*)}]
\item There is a unique natural isomorphism
\[
 \tau:\mathscr E_m\xrightarrow{\ \sim\ }\underline V
\]
to the constant rational-line functor with the displayed normalization.  It is
\[
 \boxed{\tau_\lambda(e_\lambda)=w(\lambda)^{-1}\xi.}
\]
\item The images of the canonical integral lines $\Z e_\lambda$ are therefore forced to be the
factorial lattices
\[
 \boxed{\Lambda_\lambda=w(\lambda)^{-1}\Z\xi.}
\]
For every confluence $f:\lambda\to\mu$, the rational identity restricts to the canonical lattice
inclusion $\Lambda_\lambda\hookrightarrow\Lambda_\mu$.
\item The quotient
\[
 \boxed{\mathsf D(f)=\Lambda_\mu/\Lambda_\lambda\simeq\Z/M(f)\Z}
\]
is the universal cokernel of this integralization defect: every homomorphism
$\alpha:\Lambda_\mu\to T$ satisfying $\alpha(\Lambda_\lambda)=0$ factors uniquely through
$\mathsf D(f)$.  If $T$ is finite and $\ker(\alpha)=\Lambda_\lambda$, then $\mathsf D(f)$ embeds
in $T$, so
\[
 |T|\ge M(f).
\]
Thus $\mathsf D(f)$ is the unique smallest finite faithful detector of the failure of the rational
identity to be an integral lattice isomorphism.
\item Its zeroth Fitting ideal and divisor are
\[
 \boxed{
 \operatorname{Fitt}_0\mathsf D(f)=(M(f)),
 \qquad
 \operatorname{div}\operatorname{Fitt}_0\mathsf D(f)
 =\operatorname{Res}(f).}
\]
Hence the resonance divisor is intrinsic to the integral defect module, not to a chosen
factorization of the confluence arrow.
\end{enumerate}
\end{theorem}

\begin{proof}
For any profile $\lambda$ there is a confluence $f_\lambda:\mathbf1^m\to\lambda$, and
\cref{P1C-thm:factorial-potential} gives $M(f_\lambda)=w(\lambda)$.  Naturality of a normalized
trivialization forces
\[
 \xi
 =\tau_{\mathbf1^m}(e_{\mathbf1^m})
 =\tau_\lambda\bigl(w(\lambda)e_\lambda\bigr),
\]
so necessarily $\tau_\lambda(e_\lambda)=w(\lambda)^{-1}\xi$.  This proves uniqueness, while
\cref{P1C-thm:factorial-potential} verifies naturality and hence existence.  The lattice statement is the
image of $\Z e_\lambda$ under this unique trivialization.

The quotient statement is the universal property of the cokernel of
$\Lambda_\lambda\hookrightarrow\Lambda_\mu$.  If $\ker(\alpha)=\Lambda_\lambda$, the induced map
$\mathsf D(f)\to T$ is injective, giving the cardinality bound and the minimality assertion.
Finally $\mathsf D(f)\simeq\Z/M(f)\Z$ has zeroth Fitting ideal $(M(f))$; unique prime
factorization identifies its divisor with $\operatorname{Res}(f)$.
\end{proof}

\begin{remark}[Universal Moore shadow is not a faithful confluence invariant]\label{P1C-rem:universal-shadow-not-faithful}
The universal property above concerns the rank-one integral scalar shadow.  Since
$M(f)=w(\mu)/w(\lambda)$ depends only on the endpoint profiles, two distinct marked
coarsenings with the same endpoints have the same factorial lattice quotient and the same
universal Moore complex.  The full confluence category retains incidence, marking and higher-coherence
data which are deliberately invisible to $\mathsf K_{\Z}(f)$.  Thus ``universal'' here means universal
for integral torsion detected by the normalized rank-one transport, not a complete invariant of the
underlying Cartan/confluence morphism.
\end{remark}

\begin{theorem}[Universal Moore base-change property]\label{P1C-thm:universal-moore-base-change}
Let $\mathcal C$ be a stable category equipped with an exact action of
$\operatorname{Perf}(\Z)$, and let $X\in\mathcal C$.  For every confluence morphism $f$ there is a
canonical equivalence
\[
 \boxed{
 \mathsf K_{\Z}(f)\otimes^{\mathbf L}_{\Z}X
 \simeq
 \operatorname{cofib}\bigl(M(f)\operatorname{id}_X:X\to X\bigr).}
\]
Consequently $\mathsf K_{\Z}(f)$ is the universal rank-one Moore defect attached to $f$: every
$\Z$-linear stable realization obtains its confluence torsion by exact base change from the same
integral two-cell object.
\end{theorem}

\begin{proof}
By definition $\mathsf K_{\Z}(f)$ is the two-term perfect complex
$[\Z\xrightarrow{M(f)}\Z]$, equivalently the cofiber of multiplication by $M(f)$ on the unit
object of $\operatorname{Perf}(\Z)$.  Exactness of the $\operatorname{Perf}(\Z)$-action carries this
cofiber to the displayed cofiber on $X$.
\end{proof}

\begin{corollary}[Localization and support of resonance]\label{P1C-cor:resonance-support}
For a commutative ring $R$,
\[
 \mathsf K_{\Z}(f)\otimes^{\mathbf L}_{\Z}R
 \simeq [R\xrightarrow{M(f)}R].
\]
It is acyclic exactly when $M(f)$ is a unit in $R$.  In particular its support over
$\operatorname{Spec}\Z$ is precisely the finite set of primes in $\operatorname{Res}(f)$, and at a
prime $q$ the local torsion length is
\[
 \boxed{
 \operatorname{length}_{\Z_q}H^0
 \bigl(\mathsf K_{\Z}(f)\otimes^{\mathbf L}\Z_q\bigr)
 =\ell_q(f).}
\]
Thus the resonance divisor is simultaneously the Fitting divisor, the support-with-multiplicity of
the universal Moore defect, and the primewise base-$q$ carry profile.
\end{corollary}



\section{Finite support--coefficient-depth package}

This section concerns the coefficient profile of the finite carrier.  Its depth
coordinate \(\nu\) must not be confused with the valuation-depth coordinate in the
fibre of \(\tauone:\FullMot\to\RecOneMot\).  A comparison between the two requires a
separate arithmetic realization map.

\begin{definition}[Finite coefficient-depth datum]
A finite coefficient-depth datum is a function
\[
 \nu:\mathbb P_{\mathrm{odd}}\longrightarrow\Z_{\ge0}
\]
with finite support.  Its coefficient modulus and coefficient ring are
\[
 N_\nu:=\prod_p p^{\nu_p},
 \qquad
 R_\nu:=\Z/N_\nu.
\]
We write \(|\nu|:=\sum_p\nu_p\).
\end{definition}

The squarefree truncation of \(\nu\) is
\[
 \nu^{\mathrm{sf}}_p:=\min(\nu_p,1).
\]
The present theory retains \(\nu\) itself rather than replacing it by \(\nu^{\mathrm{sf}}\).

\begin{definition}[Support--coefficient-depth object]
A finite support--coefficient-depth index is a pair
\[
 (I,\nu),
\]
where \(I\) is a finite ordered set and \(\nu\) is a finite depth datum.  The integers \(|I|\) and \(|\nu|\) are independent coordinates.
\end{definition}

\begin{principle}[No arity--multiplicity--coefficient-depth identification]
No relation of the form \(|I|=|\nu|\), \(|I|=\nu_p\), ``confluence multiplicity $m$ equals coefficient exponent $\nu_p$'', or ``the $r$-th support stage is the $p^r$-stage'' is built into the theory.  The pointed $q^2/(2,1)$ alignment of the minimal model is a theorem about that seed, not a global indexing convention.
\end{principle}

For fixed finite sets \(S\) of support labels and \(\Pi\) of odd primes, and a depth ceiling \(\bar\nu\in\Z_{\ge0}^{\Pi}\), let
\[
 \mathcal B(S):=\{T:T\subseteq S\},
 \qquad
 \mathcal D(\bar\nu):=\{\mu:0\le\mu\le\bar\nu\}.
\]
The first is ordered by inclusion.  For the second we use reduction arrows from larger depth to smaller depth.

\begin{definition}[Finite support--coefficient-depth category]
The category \(\SD(S,\bar\nu)\) has objects \((T,\mu)\in\mathcal B(S)\times\mathcal D(\bar\nu)\).  There is a unique morphism
\[
 (T,\mu)\longrightarrow(U,\lambda)
\]
precisely when
\[
 T\subseteq U,
 \qquad
 \mu\ge\lambda
\]
coordinatewise.
\end{definition}

Thus every morphism factors uniquely as
\[
 (T,\mu)\longrightarrow(T,\lambda)
 \longrightarrow(U,\lambda),
\]
first reducing depth and then enlarging support.

\begin{proposition}[Mixed Reedy structure]\label{P1C-prop:mixed-reedy}
Give \((T,\mu)\) degree
\[
 d(T,\mu):=|T|+|\mu|.
\]
Let \(\SD^-\) consist of morphisms that keep support fixed and strictly reduce depth, and let \(\SD^+\) consist of morphisms that keep depth fixed and strictly enlarge support, together with identities.  Then \(\SD(S,\bar\nu)\) is a finite Reedy category: every morphism factors uniquely as an inverse morphism followed by a direct morphism, nonidentity inverse morphisms lower degree, and nonidentity direct morphisms raise degree.
\end{proposition}

\begin{proof}
The displayed factorization is forced by the target support and target depth.  A strict depth reduction lowers \(|\mu|\), while a strict support inclusion raises \(|T|\).  Uniqueness follows from the product-poset description.
\end{proof}

This is the categorical reason for treating support and depth asymmetrically.  Reedy language and its latching/matching formalism are standard; see, for example, Berger--Moerdijk \cite{BergerMoerdijkReedy}.

\subsection{Prime-power invariants, first failure, and exact-depth compression}

\subsubsection{The local prime-power permutation system}

Fix throughout this section a finite set \(I\) with
\[
 |I|=m\ge1,
\]
an odd prime \(p\), and an integer \(r\ge1\).  Put
\[
 R_r:=\Z/p^r.
\]

\begin{definition}[Permutation and augmentation modules]
Define
\[
 P_{I,r}:=R_r[I]=\bigoplus_{i\in I}R_r e_i,
\]
with its natural \(S_I\)-action.  The augmentation map is
\[
 \epsilon_r:P_{I,r}\longrightarrow R_r,
 \qquad
 \epsilon_r\!\left(\sum_i a_i e_i\right)=\sum_i a_i.
\]
Its kernel is
\[
 A_{I,r}:=\ker\epsilon_r.
\]
The diagonal vector is
\[
 \mathbf1_I:=\sum_{i\in I}e_i.
\]
\end{definition}

Reduction modulo \(p^r\) gives equivariant surjections
\[
 P_{I,r+1}\twoheadrightarrow P_{I,r},
 \qquad
 A_{I,r+1}\twoheadrightarrow A_{I,r}.
\]
The second map is surjective because any augmentation-zero vector modulo \(p^r\) may be lifted coordinatewise and then corrected in one coordinate by a multiple of \(p^r\).

\begin{lemma}[One-step depth extension]\label{P1C-lem:one-step-extension}
For every \(r\ge1\) there is a natural short exact sequence of \(S_I\)-modules
\[
 \boxed{
 0\longrightarrow A_{I,1}
 \xrightarrow{\ p^r\ }
 A_{I,r+1}
 \longrightarrow A_{I,r}
 \longrightarrow0.}
\]
Here the first map sends a class \(\bar x\in A_{I,1}\) to \(p^r\tilde x\), where \(\tilde x\) is any lift to \(P_{I,r+1}\).
\end{lemma}

\begin{proof}
The map is well defined because changing \(\tilde x\) by \(p y\) changes \(p^r\tilde x\) by \(p^{r+1}y=0\).  Its image is the kernel of reduction on \(A_{I,r+1}\).  Indeed, an element of the kernel has the form \(p^r x\), and the condition that it lie in \(A_{I,r+1}\) is exactly
\[
 p^r\epsilon_1(\bar x)=0\pmod{p^{r+1}},
\]
that is, \(\bar x\in A_{I,1}\).
\end{proof}

\begin{remark}
The sequence in \cref{P1C-lem:one-step-extension} is the canonical source of the vertical connecting maps used later.  No section \(R_r\to R_{r+1}\) is chosen.
\end{remark}

\subsubsection{Invariant depth profile}

Write
\[
 v:=v_p(m).
\]

\begin{theorem}[Invariant depth formula]\label{P1C-thm:invariant-depth}
For every \(r\ge1\),
\[
 \boxed{
 A_{I,r}^{S_I}
 =\{a\mathbf1_I:ma=0\text{ in }R_r\}
 \simeq \Z/p^{\min(r,v)}.}
\]
More explicitly,
\[
 A_{I,r}^{S_I}=
 \begin{cases}
 0,&v=0,\\[1mm]
 R_r\mathbf1_I,&1\le r\le v,\\[1mm]
 p^{r-v}R_r\mathbf1_I,&r>v.
 \end{cases}
\]
\end{theorem}

\begin{proof}
An \(S_I\)-invariant vector in the natural permutation module has all coordinates equal, hence is \(a\mathbf1_I\).  It belongs to the augmentation kernel exactly when \(ma=0\) in \(R_r\).  Write \(m=p^v u\) with \(u\) a \(p\)-adic unit.  If \(r\le v\), this condition is automatic.  If \(r>v\), it is equivalent to \(a\in p^{r-v}R_r\).  The resulting cyclic group has order \(p^{\min(r,v)}\).
\end{proof}

\begin{corollary}[Valuation reconstruction]\label{P1C-cor:valuation-reconstruction}
The finite invariant tower determines \(v_p(m)\).  In fact,
\[
 \boxed{
 v_p(m)=\max\{r\ge0:\mathbf1_I\in A_{I,r}\},}
\]
with the convention that the maximum is zero if \(p\nmid m\).
\end{corollary}

\begin{proof}
The diagonal lies in \(A_{I,r}\) exactly when \(m=0\) in \(\Z/p^r\), equivalently \(p^r\mid m\).
\end{proof}

\begin{definition}[Primitive diagonal survival depth]
The integer
\[
 d_p(I):=v_p(|I|)
\]
is called the primitive diagonal survival depth at \(p\).
\end{definition}

Although the numerical invariant is an ordinary valuation, the transition structure is not merely numerical.

\begin{theorem}[Classification of invariant transition maps]\label{P1C-thm:transition-classification}
Let
\[
 \rho_r:A_{I,r+1}^{S_I}\longrightarrow A_{I,r}^{S_I}
\]
be reduction.  Assume \(v\ge1\).
\begin{enumerate}[label=\textup{(\roman*)}]
\item If \(r<v\), then \(\rho_r\) is the ordinary surjection
\[
 R_{r+1}\mathbf1_I\twoheadrightarrow R_r\mathbf1_I.
\]
\item If \(r\ge v\), define the canonical generator
\[
 g_r:=p^{r-v}\mathbf1_I\in A_{I,r}^{S_I}.
\]
Then \(g_r\) has order \(p^v\), generates \(A_{I,r}^{S_I}\), and
\[
 \boxed{\rho_r(g_{r+1})=p g_r.}
\]
Thus, after identifying both invariant groups with \(\Z/p^v\) by their canonical generators, \(\rho_r\) is multiplication by \(p\).
\end{enumerate}
\end{theorem}

\begin{proof}
The first assertion is immediate from \cref{P1C-thm:invariant-depth}.  For \(r\ge v\), the element \(g_r\) belongs to the augmentation kernel because
\[
 m p^{r-v}=p^r u=0\quad\text{in }R_r,
\]
and it has order \(p^v\).  Reduction sends
\[
 g_{r+1}=p^{r+1-v}\mathbf1_I
\]
to \(p^{r+1-v}\mathbf1_I=p g_r\).
\end{proof}

\begin{remark}[The break is visible in the transition law]
Before depth \(v\), the invariant tower grows by ordinary coefficient extension.  From the transition \(v+1\to v\) onward, the invariant transition is multiplication by \(p\).  Thus the valuation break is a change in the internal dynamics of the invariant tower, not only a change in group order.
\end{remark}

\subsubsection{The first-failure obstruction}

Assume now that \(p\mid m\), so \(v\ge1\).  Applying \(S_I\)-cohomology to \cref{P1C-lem:one-step-extension} gives a connecting homomorphism
\[
 \delta_r:A_{I,r}^{S_I}\longrightarrow H^1(S_I,A_{I,1}).
\]

\begin{theorem}[The local obstruction space]\label{P1C-thm:h1-one-dimensional}
If \(p\) is odd and \(p\mid m\), then
\[
 \boxed{H^1(S_I,A_{I,1})\simeq\F_p.}
\]
\end{theorem}

\begin{proof}
Consider
\[
 0\longrightarrow A_{I,1}
 \longrightarrow \F_p[I]
 \xrightarrow{\epsilon}\F_p
 \longrightarrow0.
\]
The invariant line \(\F_p[I]^{S_I}=\F_p\mathbf1_I\) maps to zero under augmentation because \(m=0\) in \(\F_p\).  Hence the connecting map injects \(\F_p\) into \(H^1(S_I,A_{I,1})\).

On the other hand,
\[
 \F_p[I]\simeq\operatorname{Ind}_{S_{m-1}}^{S_m}\F_p.
\]
By Shapiro's lemma,
\[
 H^1(S_m,\F_p[I])\simeq H^1(S_{m-1},\F_p).
\]
For odd \(p\), the latter group is zero: it is \(\Hom(S_{m-1}^{\mathrm{ab}},\F_p)\), and the symmetric-group abelianization is trivial or of exponent two.  Exactness therefore forces
\[
 H^1(S_I,A_{I,1})\simeq\F_p.
\]
\end{proof}

\begin{definition}[First-failure class]\label{P1C-def:first-failure}
Let \(v=v_p(m)\ge1\).  Since \(\mathbf1_I\in A_{I,v}^{S_I}\), define
\[
 \boxed{
 \omega_p(I):=\delta_v(\mathbf1_I)
 \in H^1(S_I,A_{I,1}).}
\]
\end{definition}

\begin{theorem}[First-failure theorem]\label{P1C-thm:first-failure}
Let \(p\) be odd and \(v=v_p(m)\ge1\).
\begin{enumerate}[label=\textup{(\roman*)}]
\item For every \(r<v\), the primitive diagonal lifts invariantly from depth \(r\) to depth \(r+1\), and
\[
 \delta_r(\mathbf1_I)=0.
\]
\item At \(r=v\), the primitive diagonal has no \(S_I\)-invariant lift to \(A_{I,v+1}\), and
\[
 \boxed{\omega_p(I)\ne0.}
\]
\item Consequently \(\omega_p(I)\) is the generator of the one-dimensional space in \cref{P1C-thm:h1-one-dimensional}.
\end{enumerate}
\end{theorem}

\begin{proof}
If \(r<v\), then \(p^{r+1}\mid m\), so the same diagonal vector \(\mathbf1_I\) lies in \(A_{I,r+1}^{S_I}\) and reduces to \(\mathbf1_I\).

At \(r=v\), suppose \(a\mathbf1_I\in A_{I,v+1}^{S_I}\) reduced to \(\mathbf1_I\).  Then \(a\equiv1\pmod{p^v}\), so \(a\) is a unit modulo \(p^{v+1}\).  The augmentation condition gives
\[
 am=0\pmod{p^{v+1}},
\]
which would imply \(p^{v+1}\mid m\), contradicting the definition of \(v\).  Thus the lift does not exist.  Exactness identifies this failure with a nonzero connecting class, which must generate the one-dimensional target.
\end{proof}

\begin{theorem}[Explicit first-failure cocycle]\label{P1C-thm:explicit-cocycle}
Write
\[
 m=p^v u,
 \qquad p\nmid u,
\]
and choose \(i_0\in I\).  A lift of \(\mathbf1_I\in A_{I,v}\) to \(A_{I,v+1}\) is
\[
 \widetilde{\mathbf1}_I
 :=\mathbf1_I-p^v u e_{i_0}.
\]
Under the kernel identification in \cref{P1C-lem:one-step-extension}, the connecting cocycle is
\[
 \boxed{
 c_{i_0}(\sigma)
 =-\bar u\bigl(e_{\sigma(i_0)}-e_{i_0}\bigr)
 \in A_{I,1}.}
\]
Its cohomology class is \(\omega_p(I)\) and is independent of the choice of \(i_0\).
\end{theorem}

\begin{proof}
The displayed lift has augmentation
\[
 m-p^v u=0.
\]
For \(\sigma\in S_I\),
\[
 \sigma\widetilde{\mathbf1}_I-\widetilde{\mathbf1}_I
 =-p^v u(e_{\sigma(i_0)}-e_{i_0}).
\]
Dividing by the kernel factor \(p^v\) and reducing modulo \(p\) gives the displayed cocycle.  A different lift changes the cocycle by a coboundary; changing \(i_0\) is a special case.
\end{proof}

\begin{proposition}[Post-break recurring defect]\label{P1C-prop:post-break}
For every \(r\ge v\), the canonical generator \(g_r=p^{r-v}\mathbf1_I\) has no invariant lift to \(A_{I,r+1}\).  Therefore
\[
 \delta_r(g_r)\ne0
\]
and hence generates \(H^1(S_I,A_{I,1})\simeq\F_p\).
\end{proposition}

\begin{proof}
By \cref{P1C-thm:transition-classification}, the image of
\[
 A_{I,r+1}^{S_I}\to A_{I,r}^{S_I}
\]
is \(pA_{I,r}^{S_I}\), whereas \(g_r\) is a generator and does not lie in this image.  Exactness gives the claim.
\end{proof}

\begin{definition}[Derived finite depth profile]
Fix a finite depth ceiling \(R\ge1\).  For odd \(p\mid |I|\), the local derived finite depth profile through depth \(R\) is
\[
 \mathfrak D_p^{\le R}(I)
 :=\left(
 \{A_{I,r}^{S_I}\}_{1\le r\le R},
 \{\rho_r\}_{1\le r<R},
 \omega_p(I)\ \text{when }v_p(|I|)\le R
 \right).
\]
Its numerical break coordinate is \(d_p(I)=v_p(|I|)\), whenever the chosen ceiling reaches the break.
\end{definition}

\subsubsection{Equivariant compression at exact prime-power depth}

The next problem is intrinsic: how much of \(P_{I,r}\) can be quotiented out while preserving the full order \(p^r\) of the diagonal?

\begin{definition}[Exact-depth quotient]
An \(S_I\)-equivariant quotient
\[
 q:P_{I,r}\twoheadrightarrow Q
\]
is \emph{exact-depth preserving} if
\[
 \Ord(q(\mathbf1_I))=p^r.
\]
\end{definition}

\begin{lemma}[Resonant kernels meet the diagonal]\label{P1C-lem:kernel-meets-diagonal}
Assume \(p\mid m\).  If \(0\ne K\subseteq P_{I,r}\) is an \(S_I\)-stable \(R_r\)-submodule, then
\[
 \boxed{K\cap R_r\mathbf1_I\ne0.}
\]
More precisely, \(K\) contains a nonzero multiple of \(p^{r-1}\mathbf1_I\).
\end{lemma}

\begin{proof}
Choose \(0\ne x\in K\), and let \(t\) be the largest integer such that \(x\in p^tP_{I,r}\).  Thus \(0\le t<r\), and write \(x=p^t y\) with \(\bar y\in\F_p[I]\) nonzero.

If \(\bar y\) is a nonzero scalar diagonal, then multiplying \(x\) by \(p^{r-1-t}\) gives a nonzero scalar multiple of \(p^{r-1}\mathbf1_I\) in \(K\).

If \(\bar y\) is not diagonal, choose \(i,j\) with \(y_i-y_j\) a unit modulo \(p\).  Then
\[
 x-(ij)x=p^t(y_i-y_j)(e_i-e_j).
\]
After multiplication by a unit, \(p^t(e_i-e_j)\in K\).  Conjugating shows that \(K\) contains \(p^t(e_a-e_b)\) for all \(a,b\).  Multiplying by \(p^{r-1-t}\) shows that \(K\) contains the full difference space \(p^{r-1}A_{I,1}\).  Because \(p\mid m\), the mod-\(p\) diagonal \(\mathbf1_I\) lies in \(A_{I,1}\).  Hence \(p^{r-1}\mathbf1_I\in K\).
\end{proof}

\begin{theorem}[Local exact-depth compression dichotomy]\label{P1C-thm:local-compression}
Let \(p\) be odd and \(r\ge1\).
\begin{enumerate}[label=\textup{(\roman*)}]
\item If \(p\nmid m\), the augmentation quotient
\[
 \epsilon_r:P_{I,r}\twoheadrightarrow R_r
\]
is exact-depth preserving and has minimal possible cardinality among all exact-depth preserving equivariant quotients.
\item If \(p\mid m\), every exact-depth preserving equivariant quotient has zero kernel.  Hence the identity quotient
\[
 P_{I,r}\twoheadrightarrow P_{I,r}
\]
is the unique quotient up to isomorphism with minimal kernel, and no nontrivial equivariant compression preserves full depth.
\end{enumerate}
\end{theorem}

\begin{proof}
If \(p\nmid m\), then \(m\) is a unit in \(R_r\), so \(\epsilon_r(\mathbf1_I)=m\) has order \(p^r\).  Any finite quotient containing an element of order \(p^r\) has at least \(p^r\) elements, and the augmentation quotient has exactly that cardinality.

If \(p\mid m\) and \(K\ne0\) is the kernel, \cref{P1C-lem:kernel-meets-diagonal} gives a nonzero diagonal element in \(K\).  Therefore the image of \(\mathbf1_I\) has order strictly smaller than \(p^r\), contradiction.
\end{proof}

\begin{remark}[Depth does not soften full resonant rigidity]
The module lattice over \(\Z/p^r\) is much richer than over \(\F_p\), but full exact-order preservation remains binary: the nonresonant component compresses to one cyclic line, while the resonant component cannot be compressed at all.
\end{remark}

\subsection{Finite CRT and Moore depth transport}

\subsubsection{The multi-prime CRT hybrid}

Let \(\nu\) be a finite odd-prime depth datum and
\[
 N=N_\nu=\prod_{p\in\Pi}p^{\nu_p},
 \qquad
 R_N=\Z/N,
\]
where \(\Pi=\operatorname{supp}\nu\).  By the finite Chinese remainder theorem,
\[
 \boxed{
 R_N\simeq\prod_{p\in\Pi}R_{p^{\nu_p}},
 \qquad
 R_{p^{\nu_p}}:=\Z/p^{\nu_p}.}
\]
Let \(e_p\in R_N\) be the corresponding idempotents.

Define
\[
 P_{I,N}:=R_N[I].
\]
Then
\[
 P_{I,N}\simeq\prod_{p\in\Pi}P_{I,p^{\nu_p}}.
\]

\begin{definition}[Resonant and nonresonant depth primes]
For \(m=|I|\), put
\[
 \Pi^{\mathrm{res}}(m):=\{p\in\Pi:p\mid m\},
 \qquad
 \Pi^{\mathrm{nr}}(m):=\Pi\setminus\Pi^{\mathrm{res}}(m).
\]
\end{definition}

\begin{definition}[Finite depth CRT hybrid quotient]\label{P1C-def:depth-crt-hybrid}
Set
\[
 \boxed{
 Q_{I,\nu}^{\mathrm{hyb}}
 :=
 \prod_{p\in\Pi^{\mathrm{nr}}(m)}R_{p^{\nu_p}}
 \times
 \prod_{p\in\Pi^{\mathrm{res}}(m)}R_{p^{\nu_p}}[I].}
\]
The quotient map \(q_{I,\nu}^{\mathrm{hyb}}:P_{I,N}\twoheadrightarrow Q_{I,\nu}^{\mathrm{hyb}}\) is augmentation on nonresonant components and the identity on resonant components.
\end{definition}

\begin{theorem}[Prime-power CRT exact-order preservation]\label{P1C-thm:crt-preservation}
The image of \(\mathbf1_I\) under \(q_{I,\nu}^{\mathrm{hyb}}\) has exact additive order
\[
 \boxed{N=\prod_{p\in\Pi}p^{\nu_p}.}
\]
\end{theorem}

\begin{proof}
On a nonresonant \(p\)-component the image is \(m\in R_{p^{\nu_p}}\), a unit multiple of \(1\), and therefore has order \(p^{\nu_p}\).  On a resonant component the full diagonal is retained and also has order \(p^{\nu_p}\).  The order in the finite product is the product of these pairwise coprime prime-power orders.
\end{proof}

\begin{theorem}[Universal minimality of the finite depth CRT hybrid]\label{P1C-thm:crt-minimality}
Among all \(S_I\)-equivariant quotients of \(P_{I,N}\) in which the diagonal has exact order \(N\), the quotient \(Q_{I,\nu}^{\mathrm{hyb}}\) has minimal cardinality.  Primewise, its component is forced to have cardinality at least
\[
 \begin{cases}
 p^{\nu_p},&p\nmid m,\\
 p^{m\nu_p},&p\mid m,
 \end{cases}
\]
and these bounds are attained simultaneously by \cref{P1C-def:depth-crt-hybrid}.
\end{theorem}

\begin{proof}
Every \(R_N\)-module and every equivariant quotient splits canonically by the CRT idempotents.  Exact order \(N\) of the diagonal is equivalent to exact order \(p^{\nu_p}\) in every prime-power component.  Apply \cref{P1C-thm:local-compression} componentwise.  Minimal cardinalities multiply.
\end{proof}

\begin{definition}[Primewise compression-length vector]
The exact-depth compression-length vector is
\[
 \boxed{
 \lambda_{p}(I,\nu):=
 \begin{cases}
 \nu_p,&p\nmid |I|,\\
 |I|\nu_p,&p\mid |I|.
 \end{cases}}
\]
Thus
\[
 |Q_{I,\nu}^{\mathrm{hyb}}|
 =\prod_{p\in\Pi}p^{\lambda_p(I,\nu)}.
\]
\end{definition}

\begin{remark}
The numerical valuation \(v_p(|I|)\) controls the invariant depth profile, while the binary resonance predicate \(p\mid |I|\) controls whether full exact-order equivariant compression is possible.  These are related but distinct outputs of the finite theory.
\end{remark}

\subsubsection{Moore seeds and diagonal depth transport}

The preceding results are module-theoretic.  They become Moore-theoretic whenever an exact-order coefficient class is supplied.

\begin{definition}[Finite Moore seed]
Let \(\mathcal C\) be a dg category or a stable category with an integral scalar action.  An \(N\)-Moore seed is a pair of homogeneous classes represented by
\[
 b,\quad c
\]
with
\[
 db=0,
 \qquad
 dc=Nb,
\]
and such that the cohomology class \([b]\) has exact additive order \(N\).
\end{definition}

For a finite support \(I\), define the diagonal Moore class
\[
 \boldsymbol\beta_{I,N}:=[b]\otimes\mathbf1_I.
\]
It lives in the coefficient-permutation extension of the chosen carrier.

\begin{proposition}[Coefficient reduction of Moore depth]\label{P1C-prop:moore-depth-reduction}
Let \(N=p^r\) and suppose \([b]\) has exact order \(p^r\).  Under the scalar mean detector \(\epsilon_r\), the diagonal Moore class has order
\[
 \boxed{
 p^{\max(r-v_p(|I|),0)}.}
\]
\end{proposition}

\begin{proof}
The mean detector sends \(\boldsymbol\beta_{I,p^r}\) to \(m[b]\), where \(m=|I|\).  Multiplication by an integer of valuation \(v\) lowers the order of an element of exact order \(p^r\) to \(p^{\max(r-v,0)}\).
\end{proof}

\begin{corollary}[Squarefree visibility versus finite depth]\label{P1C-cor:squarefree-vs-depth}
At depth one, the mean detector records only whether \(p\mid |I|\).  Across finite depths \(r\), the sequence of detected orders recovers the entire valuation \(v_p(|I|)\).
\end{corollary}

\begin{remark}
No assertion is made here that a class of coefficient depth \(p^r\) has motivic, cohomological, or geometric ``level \(r\).''  Prime-power depth is an arithmetic filtration coordinate.  Any additional interpretation requires a separate realization theorem.
\end{remark}

\subsection{Support--coefficient-depth Moore--Reedy control geometry}

We now place the finite algebra into the mixed Reedy index of \cref{P1C-prop:mixed-reedy}.  Let \(\mathcal C\) be a category with the finite colimits and limits required below, and let
\[
 X:\SD(S,\bar\nu)\longrightarrow\mathcal C
\]
be a diagram.

\subsubsection{Support latching and depth matching}

Because the direct subcategory changes only support, the latching object at \((T,\mu)\) is
\[
 \boxed{
 L^{\mathrm{supp}}_{T,\mu}X
 :=\colim_{U\subsetneq T}X_{U,\mu}.}
\]
Because the inverse subcategory changes only depth, the matching object is
\[
 \boxed{
 M^{\mathrm{depth}}_{T,\mu}X
 :=\limop_{\lambda<\mu}X_{T,\lambda},}
\]
where \(\lambda<\mu\) means coordinatewise \(\lambda\le\mu\) and \(\lambda\ne\mu\).

\begin{proposition}[Separation of latching and matching directions]\label{P1C-prop:latch-match-separation}
For the mixed Reedy structure on \(\SD(S,\bar\nu)\), the Reedy latching category at \((T,\mu)\) is canonically the proper support-face category of \(T\) at fixed depth \(\mu\), and the Reedy matching category is canonically the proper lower-depth category below \(\mu\) at fixed support \(T\).
\end{proposition}

\begin{proof}
A direct map into \((T,\mu)\) must keep depth fixed and strictly enlarge support, hence starts at \((U,\mu)\) with \(U\subsetneq T\).  An inverse map out of \((T,\mu)\) must keep support fixed and strictly reduce depth, hence ends at \((T,\lambda)\) with \(\lambda<\mu\).
\end{proof}

This separation is the basic structural advantage of the mixed variance.

\subsubsection{Elementary mixed squares}

For \(i\notin T\) and a prime \(p\) with \(\mu_p+1\le\bar\nu_p\), functoriality gives a commutative square
\[
 \begin{CD}
 X_{T,\mu+e_p} @>>> X_{T,\mu}\\
 @VVV @VVV\\
 X_{T\cup\{i\},\mu+e_p} @>>> X_{T\cup\{i\},\mu}.
 \end{CD}
\]
The horizontal arrows are depth reductions and the vertical arrows are support inclusions.

\begin{definition}[Elementary mixed defect]
Assume \(\mathcal C\) is stable.  The elementary support--depth defect is the total cofiber
\[
 \boxed{
 \Theta_{i,p}(T,\mu;X)
 :=\tcofib
 \begin{pmatrix}
 X_{T,\mu+e_p} &\longrightarrow& X_{T,\mu}\\
 \downarrow &&\downarrow\\
 X_{T\cup\{i\},\mu+e_p}&\longrightarrow&X_{T\cup\{i\},\mu}
 \end{pmatrix}.}
\]
Equivalently, up to the standard suspension convention in a stable category, one may use the total fiber.
\end{definition}

\begin{proposition}[Functoriality of the mixed defect]\label{P1C-prop:mixed-defect-functorial}
A natural transformation \(X\to Y\) induces canonical maps
\[
 \Theta_{i,p}(T,\mu;X)
 \longrightarrow
 \Theta_{i,p}(T,\mu;Y).
\]
If the elementary mixed square of \(X\) is bicartesian, then \(\Theta_{i,p}(T,\mu;X)\simeq0\).
\end{proposition}

\begin{proof}
Total cofiber is functorial in commutative squares.  In a stable category a bicartesian square has zero total cofiber.
\end{proof}

\begin{definition}[Support frontier and depth defect]
In a stable category define
\[
 C^{\mathrm{supp}}_{T,\mu}(X)
 :=\cofib\bigl(L^{\mathrm{supp}}_{T,\mu}X\to X_{T,\mu}\bigr)
\]
and
\[
 F^{\mathrm{depth}}_{T,\mu}(X)
 :=\fib\bigl(X_{T,\mu}\to M^{\mathrm{depth}}_{T,\mu}X\bigr).
\]
These are respectively the exact-support frontier at fixed depth and the exact-depth matching defect at fixed support.
\end{definition}

\begin{principle}[Mixed Moore--Reedy reading]
The support frontier, coefficient-depth defect, and elementary mixed defect are
distinct invariants.  A theory that records only support latching cannot reconstruct
coefficient-depth matching data without extra input, while a pure coefficient-depth
tower does not encode support frontiers.  These finite matching data are not the
valuation-depth fibre of the full motive tower without an additional realization.
\end{principle}

\subsubsection{The canonical permutation diagram}

The local modules assemble into an explicit mixed Reedy diagram.  Fix a finite ambient support \(S\), a prime \(p\), and a depth ceiling \(R\).

For \(T\subseteq S\) and \(1\le r\le R\), put
\[
 \mathcal P_{T,r}:=(\Z/p^r)[T],
 \qquad
 \mathcal A_{T,r}:=\ker\bigl(\epsilon:(\Z/p^r)[T]\to\Z/p^r\bigr).
\]
Support inclusions extend a vector by zero, and depth morphisms reduce coefficients.

\begin{proposition}[Canonical support--depth diagrams]\label{P1C-prop:canonical-diagrams}
The assignments \((T,r)\mapsto\mathcal P_{T,r}\) and \((T,r)\mapsto\mathcal A_{T,r}\) define functors on the mixed support--depth Reedy category.  Every elementary support-inclusion/depth-reduction square commutes strictly.
\end{proposition}

\begin{proof}
Extension by zero commutes with coefficient reduction and preserves the augmentation-zero condition.
\end{proof}

\begin{remark}[Why the augmentation-kernel diagram is the nontrivial one]
The free permutation diagram \(\mathcal P\) is coefficientwise generated by independent basis vectors.  The augmentation-kernel diagram \(\mathcal A\) couples those basis directions through one global linear relation.  The invariant-depth break and the class \(\omega_p(I)\) are features of this coupled diagram.
\end{remark}

\begin{proposition}[Depth extension is natural in support]\label{P1C-prop:extension-natural-support}
For every support inclusion \(T\subseteq U\), the one-step exact sequences
\[
 0\to A_{T,1}\to A_{T,r+1}\to A_{T,r}\to0
\]
and
\[
 0\to A_{U,1}\to A_{U,r+1}\to A_{U,r}\to0
\]
form a commutative diagram of short exact sequences under extension by zero.
\end{proposition}

\begin{proof}
All three vertical maps are induced by the same basis inclusion, and multiplication by \(p^r\) as well as coefficient reduction commute with that inclusion.
\end{proof}

This naturality is the finite algebraic seed for any later support-dependent comparison of depth obstruction classes.  No such external realization is needed for the present theory.

\subsubsection{Classification package and finite invariants}

We collect the local and finite multi-prime information into a single finite package.

\begin{definition}[Finite support--depth signature]
For a finite set \(I\) and finite odd-prime depth datum \(\nu\), define
\[
 \mathfrak S(I,\nu)
 :=\left(
 \{\mathfrak D_p^{\le\nu_p}(I)\}_{p\in\operatorname{supp}\nu},
 Q_{I,\nu}^{\mathrm{hyb}},
 \{\lambda_p(I,\nu)\}_p
 \right),
\]
with every primewise profile already truncated at the finite ceiling \(r\le\nu_p\).
\end{definition}

\begin{theorem}[Finite classification of the diagonal depth sector]\label{P1C-thm:finite-classification}
For the canonical symmetric permutation/augmentation system, the diagonal depth sector at \((I,\nu)\) is determined by the following finite data:
\begin{enumerate}[label=\textup{(\roman*)}]
\item for each \(p\mid N_\nu\), the valuation \(v_p(|I|)\);
\item the transition law of \cref{P1C-thm:transition-classification};
\item when \(p\mid |I|\), the unique nonzero first-failure line generated by \(\omega_p(I)\);
\item the primewise exact-depth compression choice: scalar for \(p\nmid |I|\), full permutation module for \(p\mid |I|\).
\end{enumerate}
Conversely, these data reconstruct the invariant tower and the minimal exact-depth CRT hybrid quotient through the prescribed finite depth datum \(\nu\).
\end{theorem}

\begin{proof}
Items (i)--(iii) reconstruct the invariant groups and their transition maps by \cref{P1C-thm:invariant-depth,P1C-thm:transition-classification,P1C-thm:first-failure,P1C-prop:post-break}.  Item (iv) reconstructs the minimal quotient primewise by \cref{P1C-thm:local-compression}; the finite CRT product gives \cref{P1C-def:depth-crt-hybrid}.  The converse follows from the same formulas.
\end{proof}

\begin{remark}[What is classified]
The theorem classifies the canonical diagonal depth sector, not all \(S_I\)-submodules of \((\Z/p^r)[I]\), and not all possible mixed Reedy diagrams.  The distinction is important: the full modular representation theory of symmetric groups over \(\Z/p^r\) is much larger than the sector required for exact diagonal depth.
\end{remark}

\subsubsection{Examples}

\begin{example}[A first break at depth one]
Let \(m=p u\) with \(p\nmid u\).  Then
\[
 A_{I,1}^{S_I}=\F_p\mathbf1_I,
\]
but
\[
 A_{I,2}^{S_I}=pR_2\mathbf1_I.
\]
The reduction map sends the generator \(p\mathbf1_I\) to zero modulo \(p\), so the primitive diagonal at depth one has no invariant lift.  The class \(\omega_p(I)\) is the generator of \(H^1(S_I,A_{I,1})\).
\end{example}

\begin{example}[Break at depth three]
Suppose \(v_p(m)=3\).  Then
\[
 A_{I,1}^{S_I}\simeq\Z/p,
 \quad
 A_{I,2}^{S_I}\simeq\Z/p^2,
 \quad
 A_{I,3}^{S_I}\simeq\Z/p^3,
\]
and the primitive diagonal lifts invariantly through these three levels.  At the transition \(4\to3\),
\[
 A_{I,4}^{S_I}=pR_4\mathbf1_I\simeq\Z/p^3,
\]
and its image is \(pA_{I,3}^{S_I}\).  Thus the first missing invariant lift is one-dimensional modulo \(p\).
\end{example}

\begin{example}[A mixed finite modulus]
Let
\[
 N=p^2q^3s,
\]
where \(p,q,s\) are distinct odd primes, and suppose
\[
 p\mid m,
 \qquad
 q\nmid m,
 \qquad
 s\mid m.
\]
Then
\[
 Q_{I,\nu}^{\mathrm{hyb}}
 \simeq
 (\Z/p^2)[I]
 \times
 \Z/q^3
 \times
 (\Z/s)[I].
\]
Its diagonal has exact order \(p^2q^3s\), and no smaller equivariant quotient has that property.  The compression-length vector is
\[
 (2m,3,m)
\]
in the \((p,q,s)\)-components.
\end{example}

\begin{example}[Mean loss versus retained hidden depth]
Assume \(p\parallel m\) and let \([b]\) have exact order \(p^2\).  The scalar mean sends the diagonal class to
\[
 m[b]=pu[b],
\]
which has order \(p\), so one power of \(p\) is lost.  The full permutation component retains order \(p^2\), while the augmentation-kernel tower records the first-failure class \(\omega_p(I)\).  These are different readouts of the same finite depth phenomenon.
\end{example}

\subsubsection{Scope boundaries}

The source-level theory stops at finite support and finite coefficient depth.  The
realization parts below add the full-motive and stack levels without turning
coefficient depth into valuation depth.  Coefficient depth is not identified
with geometric degree, confluence multiplicity, or valuation depth.  The
binary-primary case requires a separate obstruction calculation.  These exclusions
do not affect the finite diagonal-depth and confluence results proved above.
\section{Branch calibrations and finite source closure}
\label{P1L-sec:hierarchy}

\begin{definition}[Standing branch inputs]
\label{P1C-def:imported-hypotheses}
The numbered branch inputs used below are:
\begin{enumerate}[label=\textup{(I\arabic*)}]
\item an odd repeated prime \(q\) together with the marked old and repeated
      direction slots of the selected Qi branch;
\item normalization of the repeated cubic coefficient to one in the chosen
      primitive basis;
\item\label{P1C-input:I3} a free rank-one line
      \(\widetilde W_E=\Z_q\widetilde\kappa\), a nonzero residual point
      \(\bar\kappa\), a distinguished nonzero carry vector \(\gamma E\), and
      its valuation \(s_{\rm car}\);
\item for \(K=\Q(\sqrt{-519})\) and \(p=5\), the exact-\(\lambda=3\)
      calibration, internal up to the finite scalar \(\tau_{R,5}\).
\end{enumerate}
Thus input \textup{(I3)} is exactly the primitive filtered repeated package;
it is not an existence assertion for any higher arithmetic jet.
\end{definition}

\begin{theorem}[Carry-normalized repeated line]
\label{P1C-thm:qi-carry-comparison-package}
Write
\[
 \gamma E=q^{s_{\rm car}}a_0\widetilde\kappa,
 \qquad a_0\in\Z_q^\times .
\]
There is a unique primitive lift of the residual point for which the unit
part of the carry is Teichm\"uller:
\[
 \boxed{
 \kappa_{\rm Qi}^{\rm can}
 =\frac{a_0}{[\bar a_0]}\widetilde\kappa,
 \qquad
 \gamma E=q^{s_{\rm car}}[\bar a_0]\kappa_{\rm Qi}^{\rm can}.}
\]
For every finite depth \(q^r\), sending this generator to the normalized
repeated Moore generator defines the unique pointed cyclic comparison.  Thus
the distinguished carry removes the principal-unit ambiguity.
\end{theorem}

\begin{theorem}[Canonical finite Qi third-jet relation]
\label{P1C-thm:finite-qi-bifiltered-relation}
After the odd-unit PD normalization and reduction modulo \(q^r\), the marked
root third jet satisfies
\[
 \boxed{
 d\eta_{c,r}
 =\bigl(\alpha_{q^r}e_pe_q+e_p\gamma_2(e_q)\bigr)
   \kappa_{\rm Qi}^{\rm can}.}
\]
The first summand records coefficient filtration and the second records
parameter degree three.  Modulo \(q\), the first vanishes and one recovers the
normalized repeated secondary cell
\[
 d\bar\eta_c=e_p\gamma_2(e_q)\bar\kappa .
\]
\end{theorem}

\begin{theorem}[Exact-\(\lambda=3\) finite Artin gate]
\label{P1C-thm:lambda3-heisenberg-central-gate}
\label{P1C-cor:lambda3-quartic-handoff}
For the explicit \(K=\Q(\sqrt{-519})\), \(p=5\) branch, let
\(L_{\rm ab}/K\) be the degree-\(25\) elementary abelian extension cut out by
\((x,y)\), put
\[
 A=\Gal(L_{\rm ab}/K),\qquad V=\Cl(L_{\rm ab})/5,
\]
and choose \(w_0\mid\mathfrak P_0\).  Then
\[
 \boxed{
 \dim_{\F_5}V_A=1,\qquad
 \tau_{R,5}=0\Longleftrightarrow[w_0]=0\text{ in }V_A.}
\]
These conditions are equivalent to the nonzero quartic root projection and
to the existence of the marked residual coefficient
\(\mathfrak o^{\rm rec}_{3,1}\).  When they hold, the residual class has a
primitive lift to the free rank-one \(5\)-adic root lattice; the
carry-compatible attachment supplies the corresponding relative integral
quartic cell.
\end{theorem}

\begin{definition}[Finite PD Moore carrier]
\label{P1C-def:abstract-finite-pd-carrier}
For \(N>1\), a finite marked direction set \(S\), and \(m\ge0\), put
\[
 \boxed{
 \mathbb W^{\mathrm{PD}}_{N;S,m}
 :=
 [\Z C\xrightarrow{\,N\,}\Z B]\otimes_{\Z}\Gamma^m(D_S).}
\]
For primitive \(v\in\Gamma^m(D_S)\), the class of \(B\otimes v\) has exact
order \(N\).  Direction-linear maps act functorially through \(\Gamma^m\).
\end{definition}

\begin{theorem}[Finite Confluent source package]
\label{P1C-thm:finite-confluent-interface}


For finite support \(S\), fixed order \(m\), and finite odd-prime depth datum
\(\nu\), the constructions above form the finite package
\[
 \boxed{
 \mathcal I^{\rm Conf}_{S,m,\nu}
 =
 \bigl(
 D_S,\Gamma^m(D_S),R_\nu,
 \mathsf{Conf}(m),\mathbb D_m,\operatorname{Res},
 \mathbb W^{\mathrm{PD}}_{N_\nu;S,m}
 \bigr).}
\]
Support deletion, relabelling, and direction-linear maps act through
\(\Gamma^m\); confluence coarsenings act through the scalar \(M(f)\), the
torsion-defect functor \(\mathbb D_m\), and the resonance divisor
\(\operatorname{Res}(f)\).  The higher-support coherence ledger extends this
package without changing its finite carrier type.  Support, multiplicity,
coefficient depth, obstruction height, and realization remain distinct
coordinates.
\end{theorem}

\medskip
The support-level source package is now complete.  We next construct its
motivic realization and full-motive thickening.  No special arithmetic
calibration is built into the realization functors.

\part{From the support tower to low motivic realization}

\begin{remark}[Status ledger for Part~III]
\label{P2R-rem:part-three-status-ledger}
The theorem-like statements in this compressed bridge have the following
status.
\begin{enumerate}[label=\textup{(\alph*)}]
\item \emph{Imported constructional inputs:}
      \cref{P2L-thm:v48r19-rooted-zero-motive,P2R-thm:v92-relative-cylinder-presents-lift}.
\item \emph{Recalled from the source parts:}
      \cref{P2M-thm:imported-fsccomp-interface,P2M-thm:v33-imported-confluent-jet-interface}.
\item \emph{Internally proved root and formal consequences:}
      \cref{P2L-thm:v48r4-motivic-seed,P2M-prop:appendix-root-moore-complex,P2R-thm:v42-no-segal,P2R-thm:v42-moore-matching-completion,P2R-prop:v42-same-rank-commute,P2R-thm:v42-single-dg-globalization,P2R-thm:v42-stack-globalization,P2M-prop:cobar-rees-normal-form}.
\end{enumerate}
The stack handle itself is a fixed-presentation construction.  Its comparison
under a pointed equivalence preserving the complete marking is proved at the
end of Part~VI.
\end{remark}

\section{Low seed and exact-order Moore realization}

\begin{theorem}[Rooted \(p\to pq\) zero-motive]
\label{P2L-thm:v48r19-rooted-zero-motive}
Let \(\mathcal F_p\) be a pointed one-direction root object and
\(\mathfrak C^G_{pq}\) a flat pointed \(pq\) Cartan cell.  A chosen pointed
equivalence
\[
 \iota_p:\mathcal F_p\xrightarrow{\sim}
 \operatorname{Face}_p(\mathfrak C^G_{pq})
\]
defines the rooted enlargement
\[
 \mathcal M^{(0)}_{p,q}:=[\mathfrak C^G_{pq},\iota_p].
\]
It is invariant under compatible finite Fox--Spencer gauge replacement and
recovers \(\mathcal F_p\) on its \(p\)-face.
\end{theorem}

\begin{theorem}[Primitive motivic Moore seed]
\label{P2L-thm:v48r4-motivic-seed}
Let \(N\) be odd and invertible on the regular conductor base \(U_d\), and let
\(\mathcal L_{w,d}\) be the fixed invertible primitive integral coefficient
line.  The pure-weight cyclotomic root localization supplies classes
\[
 db_{\rm mot}=0,\qquad
 dc_{\rm mot}=Nb_{\rm mot},\qquad
 \operatorname{ord}[b_{\rm mot}]=N.
\]
The antecedent \(c_{\rm mot}\) is represented by the invariant root coordinate
\(t=u^N\), hence is an actual root-localization correspondence rather than a
freely adjoined Moore generator.  The two-sided bar over the constant
controller compares this motivic pair with the residual arithmetic Moore line.
\end{theorem}

\begin{proof}
The quotient-stack localization calculation, including exact order and the
geometric antecedent, is proved in
\cref{P2M-prop:appendix-root-moore-complex}.  Applying that proposition to the
fixed coefficient line gives the displayed pair.  The two-sided-bar comparison
is then the ordinary functorial bar construction over the constant controller.
\end{proof}

\begin{remark}[Canonical source line]
The realization uses only the canonical theta--Moore line and the marked
reversal splitting supplied by the source construction above.
\end{remark}

\section{Moore-corrected support globalization}

\begin{proposition}[Proper-face no-Segal obstruction]
\label{P2R-thm:v42-no-segal}
For an exact support \(I\), proper-face restriction kills the top cyclic
complex
\[
 \mathbf Z\Gamma_I^{\rm top}\xrightarrow{\,N\,}
 \mathbf Z\Omega_I^{\rm top},
 \qquad
 H^\bullet_{\rm top}\cong\mathbf Z/N.
\]
Hence ordinary proper-face Segal reconstruction cannot recover the fresh
Moore line.
\end{proposition}

\begin{proof}
Every generator in the displayed complex has exact support \(I\), so its
restriction to \(J\subsetneq I\) is zero.  The proper-face limit therefore
contains no copy of its nonzero \(\mathbf Z/N\)-cohomology line.
\end{proof}

\begin{theorem}[Moore-corrected matching completion]
\label{P2R-thm:v42-moore-matching-completion}
Adjoining the pointed exact-support root/Moore pair and its two-sided-bar
comparison to the ordinary proper-face matching datum gives the completed
support-\(I\) stage.
\end{theorem}

\begin{proof}
The new pair supplies precisely the missing cyclic cofiber detected by
\cref{P2R-thm:v42-no-segal}; the semi-free attachment is initial among
extensions carrying the prescribed antecedent and pointed comparison.
\end{proof}

\begin{proposition}[Same-rank attachment commutation]
\label{P2R-prop:v42-same-rank-commute}
Attachments of distinct exact supports having the same cardinality commute.
\end{proposition}

\begin{proof}
Their differentials factor through the common lower-cardinality stage and
their fresh exact-support summands are disjoint.  The two iterated semi-free
pushouts are therefore canonically isomorphic.
\end{proof}

\begin{theorem}[Single-dg support globalization]
\label{P2R-thm:v42-single-dg-globalization}
For a fixed compatible semi-free presentation, cardinality-wise completion
produces one support-indexed dg category
\(\mathfrak G^{\rm rec}_{S,N}\).  Its restriction to every exact support is
the Moore-corrected matching completion of
\cref{P2R-thm:v42-moore-matching-completion}.
\end{theorem}

\begin{proof}
Attach exact supports in increasing cardinality.  The preceding proposition
makes the result independent of the order chosen within a fixed cardinality,
and strict support deletion follows by removing precisely the non-surviving
exact-support summands.
\end{proof}

\begin{proposition}[Derived stack of a fixed dg presentation]
\label{P2R-thm:v42-stack-globalization}
For the dg category \(\mathfrak G^{\rm rec}_{S,N}\) obtained from the fixed
presentation, the moduli of pseudo-perfect modules
\[
 \mathfrak M^{\rm rec}_{S,N}
 :=\mathcal M_{\mathfrak G^{\rm rec}_{S,N}}
\]
is the associated derived stack.  Independence under replacement by another
dg presentation is not asserted here.
\end{proposition}

\section{Line-valued Moore completion}

\begin{definition}[Relative Moore class stack]
\label{P2R-def:v91-relative-moore-line}
For a derived stack \(X\), an invertible line \(\mathscr L\), and an integer
\(d\), put
\[
 \mathfrak C^{\rm all}_{N,X}(\mathscr L)[d]
 :=\mathbf V_X\!\left(
 \operatorname{cofib}(N:\mathscr L[-d]\to\mathscr L[-d])
 \right),
\]
and let \(\mathfrak C_{N,X}(\mathscr L)[d]\) be its generator open.
\end{definition}

\begin{definition}[Root object twisted by the detector line]
\label{P2R-def:v92-relative-root-object}
Let \(\mathfrak R_N[d]\) be the universal shifted cyclotomic root object with
\(d\Gamma_N[d]=N\Omega_N[d]\), equipped with its weight-one
\(\mathbf G_m\)-action.  Define
\[
 \mathfrak R_{N,X}(\mathscr L)[d]
 :=\operatorname{Fr}_X(\mathscr L)
   \times^{\mathbf G_m}\mathfrak R_N[d],
\]
with the descended class map
\[
 \rho_{N,X,\mathscr L}[d]:
 \mathfrak R_{N,X}(\mathscr L)[d]
 \longrightarrow\mathfrak C_{N,X}(\mathscr L)[d].
\]
\end{definition}

\begin{definition}[Line-valued Moore lift]
\label{P2R-def:v92-line-valued-lift}
For a generating section
\(\alpha:X\to\mathfrak C_{N,X}(\mathscr L)[d]\), put
\[
 \operatorname{Lift}_{N,X}(\mathscr L,\alpha)[d]
 :=X\times^h_{\mathfrak C_{N,X}(\mathscr L)[d]}
      \mathfrak R_{N,X}(\mathscr L)[d].
\]
\end{definition}

\begin{theorem}[Cylinder presentation]
\label{P2R-thm:v92-relative-cylinder-presents-lift}
The detector-line twist of the finite Moore cylinder is a semi-free
path-object presentation of
\(\operatorname{Lift}_{N,X}(\mathscr L,\alpha)[d]\).  Hence the arithmetic
detector line and the root top line are compared without choosing a global
frame.
\end{theorem}

\section{Source controller and realization normalizations}

\begin{proposition}[Source controller at the realization level]
\label{P2M-thm:imported-fsccomp-interface}
\label{P2M-thm:v33-imported-confluent-jet-interface}
The finite package of \cref{P1C-thm:finite-confluent-interface} supplies closed
singletons \(D_i\), binary cells \(H_{ij}\) with
\(dH_{ij}=[D_i,D_j]\), and fresh Moore cells \(K_I\) with
\(dK_I=N\Omega_I^{\rm rec}\) for \(|I|\ge3\).  Its finite confluence ledger is
strict under ordered deletion and natural under relabelling and marked
confluence.  Support, confluence multiplicity, coefficient depth, and jet
height remain distinct indices.
\end{proposition}

\begin{proposition}[Cobar--Rees normalization]
\label{P2M-prop:cobar-rees-normal-form}
\label{P2M-prop:eta-unique-primitive}
Put \(\eta(A)=\max\{|A|-2,0\}\).  In the associative envelope,
\[
 dX_A=
 \sum_{B\sqcup C=A}^{\rm ord}
 \sigma(B,C)\,
 N^{\eta(A)-\eta(B)-\eta(C)}X_BX_C,
\]
where \(B,C\ne\varnothing\).  After inverting \(N\) and setting
\(Y_A=N^{-\eta(A)}X_A\), this is the coefficient-free cobar differential.
The singleton/nonsingleton split forces \(\eta(m)=m-2\) for \(m\ge2\).
\end{proposition}

\begin{proof}
Substitution of \(X_A=N^{\eta(A)}Y_A\) cancels the displayed exponent in
every ordered split term.  The singleton/nonsingleton initial values give
\(\eta(1)=\eta(2)=0\), and the split recurrence then gives
\(\eta(m)=m-2\) for every \(m\ge2\).
\end{proof}

\begin{proposition}[Rank-one root localization]
\label{P2M-prop:appendix-root-moore-complex}
Let \(U_d\) be regular, let \(N\in\mathcal O_{U_d}^{\times}\), and work in an
integral symmetric monoidal stable dg enhancement of motives on tame quotient
stacks, with localization, purity, and finite-etale transfers.  Let
\(\mathcal L_{w,d}\) be an invertible primitive integral coefficient line.
For
\[
 \mathfrak X_N=[\mathbf A^1_{U_d,u}/\mu_N],\qquad
 \mathfrak D_N=B\mu_N,\qquad
 \mathfrak U_N=[\mathbf G_{m,U_d,u}/\mu_N]\simeq\mathbf G_{m,U_d,t},
 \qquad t=u^N,
\]
the Gysin boundary on the invariant coordinate line is multiplication by
\(N\).  After tensoring with \(\mathcal L_{w,d}\) and Gysin totalization, its
cyclic summand is
\[
 \Sigma^{-2,\tau_{\rm root}}\mathcal L_{w,d}
 \xrightarrow{\,N\,}
 \Sigma^{-1,\tau_{\rm root}}\mathcal L_{w,d},
\]
Writing its cocycle and geometric antecedent as \(b^{\Mot}\) and \(c^{\Mot}\)
gives
\[
 db^{\Mot}=0,\qquad dc^{\Mot}=Nb^{\Mot},\qquad
 \operatorname{ord}[b^{\Mot}]=N.
\]
\end{proposition}

\begin{proof}
Because \(N\) is invertible, \(\mu_N\) is finite etale and tame; the displayed
quotient stacks therefore lie in the equivariant six-operation formalism
\cite{HoyoisEquivariantSix,KhanRaviStacks}.  The closed zero section
\(\mathfrak D_N\hookrightarrow\mathfrak X_N\) is a regular immersion of
codimension one with complement \(\mathfrak U_N\), so purity gives its Gysin
localization triangle.

The invariant function \(t=u^N\) has principal divisor
\[
 \operatorname{div}_{\mathfrak X_N}(t)=N\mathfrak D_N.
\]
The motivic residue--divisor identity consequently gives
\[
 \partial[t]=N[\mathfrak D_N].
\]
Equivalently, the ranked boundary on the invariant coordinate line is
\(N\operatorname{id}\).  Tensoring the resulting cofiber with the invertible
line \(\mathcal L_{w,d}\) preserves this boundary and yields the displayed
two-term complex in the common Gysin-normalized Tate sector.

It remains to check that the order is exactly \(N\), rather than merely a
divisor of \(N\).  Restrict to a geometric point of \(U_d\).  Integral motivic
cohomology of the quotient stack agrees with its equivariant higher Chow group
\cite{ChoudhuryDeshmukhHogadi}; the standard equivariant Chow computation
\cite{TotaroChowBG} gives, for the tautological character line
\(\lambda\) on \(B\mu_N\),
\[
 \CH^*(B\mu_N)=\Z[\xi]/(N\xi),\qquad \xi=c_1(\lambda),
\]
on the rank-one coordinate under consideration.  Thus \(\xi\), and hence the
restricted top class, has exact order \(N\).  No proper divisor of \(N\) can
annihilate the original class.  Invertibility and primitivity of
\(\mathcal L_{w,d}\) preserve that exact order.

Finally, in the fixed dg enhancement the localization triangle is represented
by its two-term Gysin complex.  Its closed generator is \(b^{\Mot}\), while the
nullhomotopy supplied by the invariant coordinate \(t=u^N\) is
\(c^{\Mot}\); hence \(dc^{\Mot}=Nb^{\Mot}\).  Both are represented by the
zero-section/Gysin and principal-divisor correspondences, so the antecedent is
geometric.
\end{proof}
\part{Motivic Moore--Reedy structural core}
\section{The pure-weight motivic root target}
\label{P2M-sec:v32-motivic-root-target}

Fix an odd squarefree coefficient order $N$, a pure weight $w\ge2$, and a
conductor parameter $d>1$, and put
\begin{equation}
 \boxed{U_d:=\Spec(\mathcal O_d[1/N]).}
\end{equation}
All root-stack motives in this section are taken in
\begin{equation}
 \boxed{\DM_d:=\DM(U_d;\Z)}
\end{equation}
with integral coefficients.  Write $\DM_d^{(w)}\subseteq\DM_d$ for the
chosen pure-weight-\(w\) source sector.  For
bigraded realization we use its Tate-saturated stable orbit
\begin{equation}\label{P2M-eq:v42-tate-saturated-weight-orbit}
 \boxed{
 \DM_d^{\langle w\rangle}
 :=\left\langle M(b)[a]:M\in\DM_d^{(w)},\ a,b\in\Z\right\rangle_{\rm thick}
 \subseteq\DM_d.}
\end{equation}
The coefficient provenance remains in $\DM_d^{(w)}$; allowing the orbit only
makes the unavoidable Gysin/Tate shifts honest target objects.

\begin{warning}[Tate saturation is not off-weight provenance]
\label{P2M-warn:tate-saturation-not-new-provenance}
The passage $\DM_d^{(w)}\subset\DM_d^{\langle w\rangle}$ closes the chosen
pure-weight coefficient object under the shifts already required by motivic
localization and by the bigrading.  It does not adjoin arbitrary weight-mixed
Quillen--$K$ correction classes, and no shifted copy is treated as an
independent arithmetic or motivic provenance source.
\end{warning}

We use the
explicit motivic bidegree convention
\begin{equation}\label{P2M-eq:v32-bisuspension-convention}
 \boxed{
 \Sigma^{a,b}M:=M(b)[a],
 \qquad
 \REnd^{a,b}(M):=
 \operatorname{Hom}_{\DM_d^{\langle w\rangle}}(M,\Sigma^{a,b}M).}
\end{equation}
Thus the first coordinate is the dg/cohomological degree and the second is the
Tate twist.  The recursive direction labels below are independent of both
$w$ and $d$: $w$ fixes the pure-weight provenance orbit, whereas $d$ fixes the
conductor base $U_d$.

Consider, over $U_d$, the cyclotomic root localization diagram
\[
 \mathfrak X^{\Mot}_{N,d}
 =[\mathbf A^1_{U_d,u}/\mu_N]
   \times_{U_d}\mathbf G_{m,U_d,x},
 \qquad
 \mathfrak U^{\Mot}_{N,d}
 =[\mathbf G_{m,U_d,u}/\mu_N]
   \times_{U_d}\mathbf G_{m,U_d,x},
\]
with closed stratum
\[
 \mathfrak D^{\Mot}_{N,d}
 =B\mu_N\times_{U_d}\mathbf G_{m,U_d,x}.
\]
The codimension-one Gysin term is shifted by
\begin{equation}\label{P2M-eq:v32-gysin-bidegree}
 \boxed{[2](1)=\Sigma^{2,1}.}
\end{equation}
Let $\delta_{N;w,d}^{\Mot}$ denote the resulting ranked equivariant
localization boundary after restriction to $\DM_d^{(w)}$, and put
\begin{equation}
 \boxed{
 \mathbb C^{\Mot}_{w,d}:=\cofib(\delta_{N;w,d}^{\Mot}).}
\end{equation}
For a fixed three-element coefficient-calibration set $P$, define the
pure-weight Boolean root object by
\begin{equation}
 \boxed{
 \mathbb W^{\Mot}_{P,w,d}
 :=\bigoplus_{\varnothing\ne T\subseteq P}
   \Sigma^{|T|-1,0}\mathbb C^{\Mot}_{w,d}.}
\end{equation}
The second coordinate is deliberately zero here: Boolean support bookkeeping
adds only the bar/dg suspension and does not alter the fixed Tate weight.
When the supplied calibration is the original Three-$pqr$ realization, one
has $P=\{p,q,r\}$ and $N=pqr$.

\begin{theorem}[Gysin-normalized common Tate sector]
\label{P2M-thm:root-common-tate-sector}
Let the pure-weight root construction be written in the preceding convention
with an integral conductor class
\[
 \overline c^{\Mot}_{r,d}\in H^1_{\Mot}(U_d,\Z(r)).
\]
For the unit classes $[x],[t]\in H^1_{\Mot}(\mathbf G_m,\Z(1))$, put
\[
 \Omega^{\Mot}_{r,d}=\overline c^{\Mot}_{r,d}\smile[x],
 \qquad
 \Gamma^{\Mot}_{r,d}=\overline c^{\Mot}_{r,d}\smile[x]\smile[t].
\]
Then
\[
 \Omega^{\Mot}_{r,d}\in H^2_{\Mot}(-,\Z(r+1)),
 \qquad
 \Gamma^{\Mot}_{r,d}\in H^3_{\Mot}(-,\Z(r+2)).
\]
The codimension-one localization residue has bidegree $(-1,-1)$ on motivic
cohomology, equivalently the closed term of the Gysin triangle carries the
shift $[2](1)$.  Hence the identity
\[
 \delta_N^{\Mot}(\Gamma^{\Mot}_{r,d})=N\Omega^{\Mot}_{r,d}
\]
is homogeneous after Gysin totalization: the antecedent and obstruction lie in
one common Tate sector.  Denote that sector by
\[
 \boxed{\tau_{\rm root}.}
\]
In the displayed preceding convention its absolute twist is $r+1$.  No
identification between the preceding weight symbol $r$ and the bookkeeping label
$w$ used for the ambient pure-weight sector is required here: $\tau_{\rm root}$
is the Tate sector actually selected by the cyclic root object.  Its
value is fixed by the root coefficient object and is independent
of Boolean support, exact-support channel, insertion history and Confluent jet
level.
\end{theorem}

\begin{proof}
The explicit classes and localization identity are the cup-product suspension
of the internal rank-one root calculation in
\cref{P2M-prop:appendix-root-moore-complex}: the invariant coordinate has
residue \(N\) on the regular root divisor.
The unit classes $[x]$ and $[t]$ each add one Tate twist.  For a codimension-one
closed immersion the motivic Gysin triangle shifts the closed term by
$(1)[2]$ \cite{DegliseGysin2011}, so the associated residue lowers
cohomological degree and Tate twist by one.  Therefore $\Gamma^{\Mot}_{r,d}$, of absolute bidegree $(3,r+2)$,
has Gysin-normalized image in bidegree $(2,r+1)$, exactly the bidegree of
$\Omega^{\Mot}_{r,d}$.  All support, history and jet operations used below are
Tate-neutral on the coefficient object, so this common sector is unchanged by
those operations.
\end{proof}

\begin{theorem}[Pure-weight motivic Moore interface]
\label{P2M-thm:v19-pure-weight-root-seed}
For the Three-$pqr$ calibration there is a fixed root carrier
\[
 W_{\rootc}^{\Mot}=\mathbb W^{\Mot}_{P,w,d}\in\DM_d^{\langle w\rangle}
\]
and a distinguished pair of bihomogeneous endomorphisms in the common root
sector of \cref{P2M-thm:root-common-tate-sector},
\begin{equation}
 \boxed{
 b^{\Mot}\in Z^{-1,\tau_{\rm root}}\REnd(W_{\rootc}^{\Mot}),
 \qquad
 c^{\Mot}\in \REnd^{-2,\tau_{\rm root}}(W_{\rootc}^{\Mot}).}
\end{equation}
They satisfy
\begin{equation}
 \boxed{
 db^{\Mot}=0,\qquad
 dc^{\Mot}=Nb^{\Mot},\qquad
 \ord\!\left([b^{\Mot}]\in
 H^{-1,\tau_{\rm root}}\REnd(W_{\rootc}^{\Mot})\right)=N.}
\end{equation}
The class $c^{\Mot}$ is the actual equivariant root-stack localization
antecedent determined by the invariant coordinate $t=u^N$, not a freely
adjoined Moore generator.  The finite-etale cyclotomic correspondence
\begin{equation}
 \boxed{
 \mathsf T_N:=\pi_{N,*}\pi_N^*=N\,\mathrm{id},
 \qquad \bideg(\mathsf T_N)=(0,0),}
\end{equation}
is part of the same geometric interface.

Thus the codimension-one Gysin shift $(2,1)$ is not erased.  It is absorbed
only in the precise sense of \cref{P2M-thm:root-common-tate-sector}: the two Moore
terms acquire the same Gysin-normalized Tate sector $\tau_{\rm root}$.  No
claim that the absolute Tate degree is zero is needed anywhere below.
\end{theorem}

\begin{proof}
The one-graded relation, exact order, and actual localization antecedent are
\cref{P2L-thm:v48r4-motivic-seed,P2M-prop:appendix-root-moore-complex}.  The
common Tate-sector refinement is
\cref{P2M-thm:root-common-tate-sector}.  Together these identify the cyclic
localization summand with the corresponding two-term Moore complex and yield
the displayed bihomogeneous endomorphisms and exact-order class.
\end{proof}

\paragraph{Root-relative versus absolute Tate grading.}
If one writes $q_{\rm rel}:=q-\tau_{\rm root}$ on the cyclic root Moore line,
then the pair has relative bidegrees $(-1,0)$ and $(-2,0)$.  This is only a
local bookkeeping normalization of the fixed root sector; the absolute
motivic Tate coordinate remains $\tau_{\rm root}$ and is not asserted to
vanish.

\section{Boolean support carrier and strict deletion}


The fixed root carrier is now separated from the growing recursive direction
set.  For a finite ordered support $S$, define
\begin{equation}
 \boxed{
 \mathbb B_S^{\Mot}
 :=\bigoplus_{T\subseteq S}(W_{\rootc}^{\Mot})_T,}
\end{equation}
where the subscript is a support tag and does not alter the internal root
geometry.  For $i\in S$ define the bidegree-$(0,0)$ support shift
\begin{equation}
 \mathsf U_i:(W_{\rootc}^{\Mot})_T
 \longrightarrow(W_{\rootc}^{\Mot})_{T\cup\{i\}}
\end{equation}
to be the identity of $W_{\rootc}^{\Mot}$ when $i\notin T$ and zero when
$i\in T$.  Then
\[
 \mathsf U_i^2=0,
 \qquad
 \mathsf U_i\mathsf U_j=\mathsf U_j\mathsf U_i.
\]
For $R\subseteq S$, let
\[
 p_R^S:\mathbb B_S^{\Mot}\to\mathbb B_R^{\Mot}
\]
kill all support tags not contained in $R$, and let $i_R^S$ be the evident
inclusion.

\begin{proposition}[Strict Boolean motivic deletion]
\label{P2M-prop:v19-boolean-motivic-deletion}
For $Q\subseteq R\subseteq S$,
\[
 p_Q^Rp_R^S=p_Q^S,
 \qquad
 i_R^Si_Q^R=i_Q^S,
 \qquad
 p_R^Si_R^S=\mathrm{id}.
\]
Moreover every surviving support shift is intertwined by deletion, while a
shift in a deleted direction is killed.  Thus
$S\mapsto\mathbb B_S^{\Mot}$ is a strict Boolean support object built from
one genuine motivic root carrier.
\end{proposition}

\begin{proof}
All maps are defined on the support-tagged direct-sum basis.  A summand
survives a composite projection exactly when its tag lies in the smallest
support, and the support shifts add one label whenever that label is present
in the target support.
\end{proof}

\begin{proposition}[Boolean direction-skeleton realization]
\label{P2M-prop:v19-boolean-skeleton-realization}
There is a strict deletion-compatible dg map
\begin{equation}
 \boxed{
 \rho_S^{\Mot,\mathrm{Bool}}:
 \mathscr U^{\rm rec}_{S,N}
 \longrightarrow
 \REnd(\mathbb B_S^{\Mot})}
\end{equation}
defined by
\[
 D_i\longmapsto\mathsf U_i,
 \qquad
 H_{ij}\longmapsto0,
 \qquad
 K_I\longmapsto0\quad(|I|\ge3).
\]
It remembers the full recursive support skeleton but intentionally carries no
higher Moore detector by itself.
\end{proposition}

\begin{proof}
The pair relation maps to zero because the $\mathsf U_i$ commute.  For
$|I|\ge3$, every ordered decomposition $I=A\sqcup B$ contains at least one
block of cardinality at least two; the corresponding source cell maps to
zero.  Hence $\rho_S^{\Mot,\mathrm{Bool}}(\Omega_I^{\rm rec})=0$, so
$dK_I=N\Omega_I^{\rm rec}$ is respected.  Strict deletion follows from
\cref{P2M-prop:v19-boolean-motivic-deletion}.
\end{proof}

\paragraph{Target split.}
The Boolean carrier records support and strict restriction; exact-support
Moore curvature is carried by the channel summands.


\section{Cyclotomic-transfer realization of the Rees support law}

The marked FSCComp support tower already supplies the noncommutative support cells and their Moore--Reedy provenance.  The question here is narrower: can its Rees coefficient profile be realized by genuine geometric correspondences?  The answer is yes, directly from the cyclotomic cover carried by the chosen motivic Moore realization.

Let
\[
 \pi_N:\mathbf G_{m,u}\longrightarrow\mathbf G_{m,t},
 \qquad t=u^N,
\]
be the degree-\(N\) finite etale cover on the locus where \(N\) is invertible.
In any motivic category with finite-etale transfers, pullback and transfer
satisfy
\[
 \boxed{\pi_{N,*}\pi_N^*=N\,\mathrm{id}.}
\]
This is the geometric source of the same integer \(N\) that appears in the
root-stack localization relation \(d\Gamma^{\rootc}=N\Omega^{\rootc}\).

Retain the function \(\eta\) from
\cref{P2M-thm:imported-fsccomp-interface} and put
\[
 \delta(a,b):=\eta(a+b)-\eta(a)-\eta(b).
\]
For disjoint nonempty supports, \(\delta(a,b)\) equals \(0,1,2\) in the
singleton--singleton, singleton--nonsingleton, and
nonsingleton--nonsingleton cases, respectively.

\begin{definition}[Cyclotomic-transfer squarefree support object]
\label{P2M-def:transfer-rees-object}
For a finite ordered support \(S\), let
\[
 \mathbb E^{\Mot,\mathrm{tr}}_{S,N}
 :=\bigoplus_{J\subseteq S}
    \Sigma^{|J|}\mathbf 1\otimes\operatorname{Or}(J),
\]
where \(\mathbf1\) is the motivic unit over the pure-weight root-stack base
and \(\operatorname{Or}(J)\) is the rank-one orientation line generated by
the increasing word on \(J\).  For disjoint \(A,B\), define multiplication
from the \((A,B)\)-summand to the \(A\sqcup B\)-summand by
\[
 \boxed{
 \mu_{A,B}
 :=\epsilon(A,B)
   (\pi_{N,*}\pi_N^*)^{\delta(|A|,|B|)},}
\]
using the canonical suspension identification.  If \(A\cap B\ne\varnothing\),
the squarefree product is zero.
\end{definition}

\begin{theorem}[Cyclotomic-transfer Rees profile]
\label{P2M-thm:transfer-rees-profile}
The object \(\mathbb E^{\Mot,\mathrm{tr}}_{S,N}\) is an associative
support-graded algebra object in motivic correspondences and is strictly
functorial under ordered support deletion.  In the oriented support basis,
for disjoint \(A,B\),
\[
 \boxed{
 \tau_A\tau_B
 =\epsilon(A,B)
   N^{\eta(|A|+|B|)-\eta(|A|)-\eta(|B|)}
   \tau_{A\sqcup B}.}
\]
Hence its nonzero multiplication profile is exactly
\[
 \boxed{1,\qquad N,\qquad N^2,}
\]
the fixed-\(N\) Rees profile of Recursive Secondary.
\end{theorem}

\paragraph{Raw versus full versus normalized coefficients.}
The three coefficient conventions used in the paper are related by one
outer Moore factor.  For a split \(A=B\sqcup C\), write
\(\delta=\eta(A)-\eta(B)-\eta(C)\).  Then
\[
\begin{array}{c|ccc}
\text{split type}
 &1+1&1+\text{nonsingleton}
 &\text{nonsingleton}+\text{nonsingleton}\\ \hline
\delta&0&1&2\\
\text{curvature }\Omega_A
 &\text{--}&1&N\\
\text{full differential }dX_A
 &1&N&N^2\\
\text{normalized cobar }dY_A
 &1&1&1\\
\text{motivic transfer product}
 &1&N&N^2
\end{array}
\]
Thus \cref{P2M-thm:imported-fsccomp-interface} records the raw curvature
coefficient \(N^{\delta-1}\) for \(|A|\ge3\), while the present theorem
realizes the coefficient \(N^\delta\) of the \emph{full} differential after
the outer equation \(dK_A=N\Omega_A\) is restored.


\begin{proof}
The transfer identity identifies each occurrence of
\(\pi_{N,*}\pi_N^*\) with multiplication by \(N\).  The exponent
\(\delta(a,b)\) is the coboundary of the single function \(\eta\), so for
three pairwise disjoint supports both parenthesizations have total exponent
\[
 \eta(a+b+c)-\eta(a)-\eta(b)-\eta(c).
\]
The orientation signs are associative exterior signs.  This proves
associativity and the displayed profile.  For deletion \(T\subseteq S\),
kill the summands indexed by \(J\not\subseteq T\) and keep all others.  The
transfer correspondence is independent of the recursive direction label, so
this projection commutes strictly with multiplication and with iterated
deletion.
\end{proof}

\begin{corollary}[Geometric realization of the intrinsic Rees weights]
\label{P2M-cor:rees-weights-motivic}
The coefficients \(1,N,N^2\) in the recursive support curvature need not be
introduced as formal scalar weights: they are the correspondence degrees of
the genuine cyclotomic cover underlying the pure-weight root stack.  Thus the support-weight part of the marked FSCComp/Secondary Cartan Rees law admits a genuine cyclotomic-transfer realization.
\end{corollary}


\section{Channel-complete transfer--Rees refinement}
The intrinsic source already fixes the quadratic exact-support latching
polynomial.  We realize each ordered bipartition in its own orthogonal
motivic channel block.

For a nonempty support $A$ write
\[
 \widehat K_A:=
 \begin{cases}
  D_i,&A=\{i\},\\
  H_{ij},&A=\{i,j\},\\
  K_A,&|A|\ge3.
 \end{cases}
\]
For $I$ with $|I|=m\ge3$, let
\[
 \mathfrak P^{\to}(I)
 :=\{(A,B):A\sqcup B=I,\ A,B\ne\varnothing\}
\]
be the ordered bipartition set.  Put
\[
 \delta(A,B)
 :=\eta(|I|)-\eta(|A|)-\eta(|B|).
\]
Because $m\ge3$, one has
\[
 \delta(A,B)=
 \begin{cases}
  1,&\min(|A|,|B|)=1,\\
  2,&|A|,|B|\ge2.
 \end{cases}
\]
Let
\[
 \mathsf T_N:=\pi_{N,*}\pi_N^*=N\,\mathrm{id}
\]
be the genuine cyclotomic transfer correspondence of
\cref{P2M-thm:transfer-rees-profile}.  Finally let
$\chi_I(A,B)\in\{\pm1\}$ denote the orientation/Koszul sign with which the
ordered monomial $\widehat K_A\widehat K_B$ occurs in the normalized support
curvature $\Omega_I^{\mathrm{rec}}$.  Thus its scalar magnitude is
$N^{\delta(A,B)-1}$.

\begin{definition}[Ordered transfer channel block]
\label{P2M-def:ordered-transfer-channel}
Fix $(A,B)\in\mathfrak P^{\to}(I)$.  Let
$\mathbb V_{I;A,B}^{\Mot,\mathrm{ch}}$ be the three-object graded matrix
correspondence block on
\[
 V_1=W_{\rootc}^{\Mot},\qquad
 V_2=\Sigma^{1-|A|,0}W_{\rootc}^{\Mot},\qquad
 V_3=\Sigma^{3-|I|,-\tau_{\rm root}}W_{\rootc}^{\Mot}.
\]
With the standard convention $E_{ij}\in\operatorname{Hom}(V_j,V_i)$, the
bihomogeneous matrix units satisfy
\begin{equation}
 \boxed{
 \begin{aligned}
 \bideg(E_{12}^{A,B})&=(1-|A|,0),\\
 \bideg(E_{23}^{A,B})&=(2-|B|,-\tau_{\rm root}),\\
 \bideg(E_{13}^{A,B})&=(3-|I|,-\tau_{\rm root}).
 \end{aligned}}
\end{equation}
Their product is bihomogeneous because
\begin{equation}
 \boxed{
 (1-|A|,0)+(2-|B|,-\tau_{\rm root})
 =(3-|I|,-\tau_{\rm root}),}
\end{equation}
and the matrix relations are
\[
 E_{12}^{A,B}E_{23}^{A,B}=E_{13}^{A,B},
 \qquad
 E_{23}^{A,B}E_{12}^{A,B}=0.
\]
Put
\begin{equation}\label{P2M-eq:channel-shift-sign}
 \boxed{\sigma_I^{\rm sh}:=(-1)^{|I|-3}.}
\end{equation}
The canonical shift matrix units are closed.  With the standard Hom-complex
differential and $|E_{13}^{A,B}|=3-|I|$, one therefore has
\begin{equation}\label{P2M-eq:channel-hom-leibniz}
 \boxed{
 d\bigl(E_{13}^{A,B}x\bigr)
 =\sigma_I^{\rm sh}E_{13}^{A,B}dx}
\end{equation}
for every homogeneous coefficient $x$ when $dE_{13}^{A,B}=0$.  This shift
sign is part of the dg bookkeeping and is separate from the orientation sign
$\chi_I(A,B)$.
Since \(\bideg(b^{\Mot})=(-1,\tau_{\rm root})\),
\(\bideg(c^{\Mot})=(-2,\tau_{\rm root})\), and
\(\bideg(\mathsf T_N)=(0,0)\), the compensating matrix twists cancel the
common root sector and the three images below have respectively the source
bidegrees
\[
 (1-|A|,0),\qquad(1-|B|,0),\qquad(1-|I|,0).
\]
Thus both the dg degree and the Tate twist are checked before imposing the
matrix composition law.  On this exact-support-$I$ block give the source
cells the future components
\begin{align}
 \widehat K_A
 &\longmapsto E_{12}^{A,B},
 \label{P2M-eq:channel-KA}\\
 \widehat K_B
 &\longmapsto E_{23}^{A,B}b^{\Mot},
 \label{P2M-eq:channel-KB}\\
 K_I
 &\longmapsto
   \sigma_I^{\rm sh}\chi_I(A,B)\,
   E_{13}^{A,B}\mathsf T_N^{\delta(A,B)-1}c^{\Mot}.
 \label{P2M-eq:channel-KI}
\end{align}
Every other proper-support source cell has zero component in this block.  These are exact-support future components appended to the already constructed proper-support realization data; the entire block is deleted on every proper restriction of $I$.

The matrix block is the coefficient factor of the Rees channel.  When the raw
flatness equation is read, the first two entries are tensored with the support
basis vectors $\tau_A,\tau_B$ of
$\mathbb E^{\Mot,\mathrm{tr}}_{I,N}$ and the carry entry with $\tau_I$.
Thus the transfer factor belongs to support multiplication; it is not applied
a second time to the normalized controller polynomial.
\end{definition}

\begin{lemma}[Proper-support equations remain strict]
\label{P2M-lem:channel-proper-equations}
The future components \eqref{P2M-eq:channel-KA}--\eqref{P2M-eq:channel-KB} preserve
every already imposed proper-support differential equation.
\end{lemma}

\begin{proof}
Both displayed future components are closed: the first uses the motivic
identity and the second uses $db^{\Mot}=0$.  The prescribed differential of
$\widehat K_A$ involves only source cells on proper subsets of $A$; all of
those have zero component in the $(I;A,B)$ block.  For $|A|=2$ this says that
the two singleton directions have zero component there, so the commutator
image is zero; for $|A|\ge3$ the lower latching polynomial has zero image for
the same support reason.  The argument for $B$ is identical.  Hence adding
the future components changes neither lower differential equation.
\end{proof}

\begin{theorem}[Exact ordered-channel realization]
\label{P2M-thm:ordered-channel-realization}
In the block $\mathbb V_{I;A,B}^{\Mot,\mathrm{ch}}$ one has
\begin{equation}
 \boxed{
 \rho_{I;A,B}(\Omega_I^{\mathrm{rec}})
 =\chi_I(A,B)\,
  E_{13}^{A,B}\mathsf T_N^{\delta(A,B)-1}b^{\Mot}.}
 \label{P2M-eq:channel-latching-image}
\end{equation}
Consequently
\begin{equation}
 \boxed{
 d\rho_{I;A,B}(K_I)
 =N\rho_{I;A,B}(\Omega_I^{\mathrm{rec}}).}
 \label{P2M-eq:channel-carry}
\end{equation}
Moreover the same identity is the exact-support-$I$ coefficient of the raw
transfer--Rees flatness equation: the product channel contributes
$\mathsf T_N^{\delta(A,B)}b^{\Mot}$, while the differential of the carry
contributes the same correspondence with the opposite orientation sign.
\end{theorem}

\begin{proof}
Only $\widehat K_A$ and $\widehat K_B$ are nonzero among the proper-support
cells in this block.  The ordered monomial
$\widehat K_A\widehat K_B$ therefore maps to
$E_{13}^{A,B}b^{\Mot}$, whereas the reverse ordered monomial maps to zero
because $E_{23}^{A,B}E_{12}^{A,B}=0$.  By definition of
$\chi_I(A,B)$ and by the normalized recursive curvature formula, the
coefficient of this ordered monomial is the signed operator
$\chi_I(A,B)N^{\delta(A,B)-1}$.  The equality
$\mathsf T_N=N\,\mathrm{id}$ turns that coefficient into the geometric
transfer power in \eqref{P2M-eq:channel-latching-image}.

Now use \eqref{P2M-eq:channel-hom-leibniz} and
$dc^{\Mot}=Nb^{\Mot}$:
\[
 d\bigl(
   \sigma_I^{\rm sh}E_{13}^{A,B}
   \mathsf T_N^{\delta-1}c^{\Mot}
  \bigr)
 =N E_{13}^{A,B}\mathsf T_N^{\delta-1}b^{\Mot}.
\]
This is \eqref{P2M-eq:channel-carry}.  To read the same calculation as raw
Rees flatness, tensor the two proper-support entries with
$\tau_A,\tau_B$ and the carry with $\tau_I$.  Then
\cref{P2M-thm:transfer-rees-profile} gives the product coefficient
$\mathsf T_N^{\delta}$, while
\[
 \mathsf T_N^{\delta-1}dc^{\Mot}
 =\mathsf T_N^{\delta}b^{\Mot}.
\]
There is therefore exactly one transfer contribution at the raw Rees level:
after dividing the support-$I$ equation by its common factor $N$, it becomes
the normalized coefficient $\mathsf T_N^{\delta-1}$ in
\eqref{P2M-eq:channel-latching-image}.  No transfer factor is double counted.
\end{proof}

\begin{definition}[Support-functorial motivic Moore--Reedy target]
\label{P2M-def:genuine-motivic-rees-cartan}
For a channel block $\mathbb V_{I;A,B}^{\Mot,\mathrm{ch}}$, write
$|\mathbb V_{I;A,B}^{\Mot,\mathrm{ch}}|$ for the direct sum of its three
shifted root-motive factors, forgetting only the displayed matrix notation.
For a finite ordered support $S$, define the total motivic target object
\begin{equation}
 \boxed{
 \mathbb M_{S,N}^{\Mot,\MR}
 :=\mathbb B_S^{\Mot}
 \oplus
 \bigoplus_{\substack{I\subseteq S\\|I|\ge3}}
 \ \bigoplus_{(A,B)\in\mathfrak P^{\to}(I)}
 |\mathbb V_{I;A,B}^{\Mot,\mathrm{ch}}|.}
\end{equation}
Let
\begin{equation}
 \boxed{
 \mathscr M_{S,N}^{\Mot,\MR}
 \subseteq
 \REnd(\mathbb M_{S,N}^{\Mot,\MR})}
\end{equation}
be the dg correspondence subalgebra generated by the Boolean support
projections and shifts $\mathsf U_i$, the channel matrix units, the root Moore
pair $(b^{\Mot},c^{\Mot})$, and the cyclotomic transfer correspondence
$\mathsf T_N$ on the relevant root factors.

For later comparison write
\[
 \mathscr E_S^{\Mot,\mathrm{ch}}
 :=\prod_{\substack{I\subseteq S\\|I|\ge3}}
   \ \prod_{(A,B)\in\mathfrak P^{\to}(I)}
   \REnd\bigl(\mathbb V_{I;A,B}^{\Mot,\mathrm{ch}}\bigr)
\]
for the exact-support channel factor of this target.
A support deletion $T\subseteq S$ defines
\begin{equation}\label{P2M-eq:v19-total-target-deletion-map}
 P_T^S:\mathbb M_{S,N}^{\Mot,\MR}\longrightarrow
 \mathbb M_{T,N}^{\Mot,\MR}
\end{equation}
by $p_T^S$ on the Boolean carrier, by zero on every channel block whose exact
support is not contained in $T$, and by the identity on all remaining channel
blocks.
\end{definition}

\begin{proposition}[Strict deletion of the total motivic target]
\label{P2M-prop:v19-total-target-deletion}
For $R\subseteq T\subseteq S$ the preceding deletion maps satisfy
\begin{equation}
 \boxed{P_R^T P_T^S=P_R^S,}
\end{equation}
and they preserve the generating dg subalgebras
$\mathscr M_{-,N}^{\Mot,\MR}$.  Hence
$S\mapsto\mathscr M_{S,N}^{\Mot,\MR}$ is strict under ordered support
deletion.
\end{proposition}

\begin{proof}
The Boolean summand is strict by
\cref{P2M-prop:v19-boolean-motivic-deletion}.  Channel blocks are indexed by exact
supports and are either retained identically or removed.  Every algebra
generator is therefore retained or killed by the same support rule, so the
restricted generator multiplication and differential agree with the smaller
target.
\end{proof}

\begin{theorem}[Geometric support-functorial Moore--Reedy correspondence realization]
\label{P2M-thm:channel-complete-realization}
For every finite ordered $S$ and odd squarefree $N$ for which the genuine
pure-weight root pair $(b^{\Mot},c^{\Mot})$ is fixed, the Boolean skeleton and
all exact-support channels combine to a dg realization
\begin{equation}\label{P2M-eq:v19-total-motivic-realization}
 \boxed{
 \rho_S^{\Mot,\MR}:
 \mathscr U^{\mathrm{rec}}_{S,N}
 \longrightarrow
 \mathscr M_{S,N}^{\Mot,\MR}.}
\end{equation}
Every generating target morphism is represented by a composite of identities on shifted motivic summands, zero-section/Gysin correspondences, and graph/transpose-graph correspondences of the cyclotomic cover.  In this structural sense the realization is geometric.  Its Boolean projection is $\rho_S^{\Mot,\mathrm{Bool}}$ and its exact-support
channel projection is the previously defined map
\[
 \rho_S^{\Mot,\mathrm{ch}}:
 \mathscr U^{\mathrm{rec}}_{S,N}
 \longrightarrow
 \mathscr E_S^{\Mot,\mathrm{ch}}.
\]
It has the following properties.
\begin{enumerate}[label=\textup{(\roman*)}]
\item The total target and the realization are strict under ordered support deletion.
\item Every ordered bipartition monomial in every higher latching polynomial is
realized in a distinct geometric channel; no singleton or genuine partition
channel is discarded.
\item The coefficients of the corresponding raw flatness law are supplied by
the cyclotomic transfer powers $\mathsf T_N$ and hence reproduce the complete
$1,N,N^2$ Rees multiplication profile.
\item For every $I$, $|I|\ge3$, the class
$[\rho_S^{\Mot,\mathrm{ch}}(\Omega_I^{\mathrm{rec}})]$ has exact additive
order $N$.
\item The image of $K_I$ is a genuine motivic antecedent assembled from matrix
placements of the original root-stack localization class $c^{\Mot}$; no new
abstract Moore complex is introduced at higher support.
\item Attachments of the same support cardinality are strictly independent and
therefore satisfy the same-rank Fubini property.
\item Relabelling a finite support by a permutation carries the block
$(I;A,B)$ to $(\sigma I;\sigma A,\sigma B)$ and transports the orientation
line by the Koszul sign.  These maps give a coherent symmetric-group action on
the tower.  Thus the construction has no distinguished support selector.
\end{enumerate}
Thus the intrinsic FSC Moore--Reedy controller has a single support-functorial
motivic target.  The Boolean summand retains recursive support geometry, while
the channel summands realize every nontrivial exact-support latching class and
its genuine motivic Moore antecedent before any arithmetic base change.
\end{theorem}

\begin{proof}
The Boolean component is a dg map by
\cref{P2M-prop:v19-boolean-skeleton-realization}.  The nontrivial dg identities in
the exact-support channel summands hold blockwise by
\cref{P2M-lem:channel-proper-equations,P2M-thm:ordered-channel-realization}; orthogonality
of the direct summands therefore gives the total dg map
\eqref{P2M-eq:v19-total-motivic-realization}.  Strict deletion is
\cref{P2M-prop:v19-total-target-deletion}.  Every ordered
bipartition occurs exactly once as a channel index, proving (ii), and (iii) is
the raw-flatness statement of \cref{P2M-thm:ordered-channel-realization}.

For (iv), every channel coordinate is a transfer multiple of
$[b^{\Mot}]$, so the full vector is annihilated by $N$.  Fix a highest-face
channel $(A,B)=(\{i\},I\setminus\{i\})$.  Here $\delta(A,B)=1$, so
projection to this block sends the normalized latching class to
$\pm E_{13}^{A,B}b^{\Mot}$.  The matrix coordinate is split and
$[b^{\Mot}]$ has exact order $N$; hence the full channel vector has exact
order $N$.  Statement (v) follows from \eqref{P2M-eq:channel-KI}: when both
blocks of a partition are nonsingleton the additional factor is the genuine
transfer correspondence $\mathsf T_N$, and otherwise it is the identity.
In either case the only Moore antecedent used is $c^{\Mot}$ itself.

Different exact supports label orthogonal products, so same-rank additions
commute strictly, proving (vi).  Finally a permutation relabels the Boolean support tags and carries $\mathsf U_i$ to $\mathsf U_{\sigma i}$; on the exact-support summands it bijects ordered
bipartitions and transports the orientation generator.  The identity and
composition laws are those of these support relabellings together with the
multiplicativity of the Koszul orientation sign.  This proves (vii).

For geometricity, the root pair is represented by the regular zero section and
the invariant principal-divisor correspondence by
\cref{P2M-prop:appendix-root-moore-complex}.  The transfer
\(\mathsf T_N=\pi_{N,*}\pi_N^*\) is the composite of the transpose and ordinary
graphs of the finite-etale cyclotomic cover.  Boolean shifts and channel matrix
units are identity correspondences placed between finite shifted direct
summands.  Finite sums, compositions, shifts, and cones remain inside the
stable additive category of motivic correspondences.  Hence every generator
and every differential used by \(\rho_S^{\Mot,\MR}\) has the asserted
geometric correspondence representative.
\end{proof}

\begin{remark}[Scope of the geometric realization]
\label{P2M-rem:geometric-mr-scope}
The theorem realizes the structural root, transfer, latching, support-deletion,
and Reedy data by motivic correspondences.  Transporting a Confluent jet or
higher arithmetic operation as a marking on these carriers does not by itself
construct an independent algebraic correspondence whose operation law is that
higher Massey product.  Such operation-level provenance remains a separate
theorem.
\end{remark}

\subsection{Role-separated objectification of the support action}

The Boolean shifts are morphisms on one direct-sum carrier.  They must not be
renamed as independent controller objects.  The following finite additive
enlargement performs the required objectification while retaining the original
correspondences as the evaluation maps.

\begin{definition}[Role-separated Boolean objectification]
\label{P2M-def:role-separated-objectification}
Write the finite ordered support as
\(S=\{q_1<\cdots<q_r\}\), put
\(S_{<j}:=\{q_1,\ldots,q_{j-1}\}\), and denote the Boolean summand
\((W_{\rootc}^{\Mot})_T\) by \(B_T\).  Superscripts on a summand below denote
labelled copies, not new motives.  In the finite additive Karoubi envelope of
the summands of \(\mathbb M_{S,N}^{\Mot,\MR}\), define
\begin{align}
 d_{q_i}^{\rm sel}(S)
 &:={B_{\{q_i\}}}^{\rm core}
   \oplus\bigoplus_{j>i}
      B_{S_{<j}\setminus\{q_i\}}^{(i,j)},
 \label{P2M-eq:objectified-controller}\\
 e_{q_j}^{\rm sel}(S)
 &:=B_{S_{<j}}^{\rm birth}
 \oplus
 \bigoplus_{\substack{I\subseteq S_{<j}\\|I|\ge3}}
 \ \bigoplus_{(A,B)\in\mathfrak P^{\to}(I)}
 |\mathbb V_{I;A,B}^{\Mot,\mathrm{ch}}|^{(j)}.
 \label{P2M-eq:objectified-birth}
\end{align}
For \(i<j\), let
\begin{equation}
 \boxed{
 \lambda_{ij}^{\rm sel}:
 d_{q_i}^{\rm sel}(S)\longrightarrow e_{q_j}^{\rm sel}(S)}
 \label{P2M-eq:objectified-control-edge}
\end{equation}
be zero on every summand except the \((i,j)\)-summand and there the restriction
\[
 \mathsf U_{q_i}:
 B_{S_{<j}\setminus\{q_i\}}^{(i,j)}
 \longrightarrow B_{S_{<j}}^{\rm birth}
 \lhookrightarrow e_{q_j}^{\rm sel}(S).
\]
Let \(\mathcal O_{S,N}^{\rm sel}\) be the full pretriangulated idempotent-complete
dg subcategory generated by the individual Boolean and exact-support channel
summands of \(\mathbb M_{S,N}^{\Mot,\MR}\), the displayed role objects, the
maps \(\lambda_{ij}^{\rm sel}\), and the channel markings supplied by
\(\rho_S^{\Mot,\MR}\).
\end{definition}

The core summand keeps the newest controller nonzero even before it has an
outgoing edge.  The labelled edge summands keep different future evaluations
independent.  Thus neither a controller nor a birth is being identified with
the endomorphism \(\mathsf U_{q_i}\).

\begin{theorem}[Object-level birth--control realization]
\label{P2M-thm:role-separated-objectification}
The preceding construction has the following properties.
\begin{enumerate}[label=\textup{(\roman*)}]
\item The objects \(d_{q_i}^{\rm sel}(S)\) and
      \(e_{q_j}^{\rm sel}(S)\) are role-separated objects, and
      a distinguished nonzero closed degree-zero control morphism
      \(\lambda_{ij}^{\rm sel}\) is marked exactly for \(i<j\).
\item Each control edge is the indicated objectwise evaluation of the genuine
      Boolean correspondence \(\mathsf U_{q_i}\); it is not obtained by treating
      that endomorphism as an object.
\item For every chosen order ideal of exact supports on \(S_{<j}\), the
      ordered-channel cells of
      \cref{P2M-thm:channel-complete-realization} restrict to the corresponding
      marked higher comparison cells landing at the \(j\)-th birth profile.
      All their faces and Moore antecedents are retained.
\item Under \(S_r\subset S_{r+1}\), every old birth object and old edge is
      retained, each old controller acquires precisely its new
      \((i,r+1)\)-summand, and the new controller has only its core summand.
      Projection back to \(S_r\) deletes exactly these new data.  These maps
      are strict and compose under further cutoff deletion.
\item Adjoining the role objects does not change the Morita type.  More
      precisely, the inclusion of the dg category generated by the individual
      Boolean and channel summands into \(\mathcal O_{S,N}^{\rm sel}\) is a
      Morita equivalence.
\end{enumerate}
Consequently the finite motivic Moore--Reedy tower contains independent
controller objects, independent birth objects, and object-level maps
\(d_{q_i}^{\rm sel}\to e_{q_j}^{\rm sel}\), while its original Boolean
endomorphisms remain correctly typed as correspondences on the carrier.
\end{theorem}

\begin{proof}
The Boolean idempotents split in the additive Karoubi envelope, so every
summand and every finite labelled copy in
\eqref{P2M-eq:objectified-controller}--\eqref{P2M-eq:objectified-birth} is an
object.  Projection to the \((i,j)\)-copy followed by the restriction of
\(\mathsf U_{q_i}\) is therefore a closed morphism.  It exists in the prescribed
triangular profile because \(q_i\in S_{<j}\) exactly when \(i<j\).  This proves
\textup{(i)} and \textup{(ii)}.

Channel completeness supplies one orthogonal block for every ordered
bipartition monomial and its Moore antecedent.  Restricting to an order ideal
keeps exactly the selected blocks, and strict support deletion preserves or
kills each block on the nose.  The same direct-sum rule proves the cutoff
description of the role objects, hence \textup{(iii)} and \textup{(iv)}.

Finally, every adjoined role object is a finite direct sum of labelled copies
of Boolean and channel summands already in the perfect envelope.  Conversely the original
Boolean and channel summands remain among the generators of
\(\mathcal O_{S,N}^{\rm sel}\).  The two dg subcategories therefore have the
same idempotent-complete pretriangulated hull, which is precisely the Morita
criterion.  This proves \textup{(v)}.
\end{proof}

\begin{remark}[External selected carriers]
\label{P2M-rem:objectification-external-carrier-scope}
The theorem constructs the canonical role-separated presentation attached to
the Moore--Reedy carrier.  To identify these objects with independently chosen
objects in another arithmetic or motivic carrier, one must supply a pointed
marked Morita equivalence carrying the displayed role, edge, and higher-cell
markings.  The presentation-independence theorem below then performs the
transport; bare equality of support labels is insufficient.
\end{remark}


\begin{theorem}[Finite Confluent jet decoration reuses the existing motivic Moore--Reedy recursion]
\label{P2M-thm:v33-confluent-jet-decoration}
Fix an odd squarefree coefficient order $N$, a finite ordered support $S$ and a
finite jet ceiling $M$.  Equip the intrinsic source with the finite ledger
$\mathfrak R^{\rm Conf}_{S,\le M}$ of
\cref{P2M-thm:v33-imported-confluent-jet-interface}.  Define the jet-decorated
motivic stage by
\begin{equation}
 \boxed{
 \mathbb M_{S,N,\le M}^{\Mot,\MR;\rm jet}
 :=
 \bigl(\mathbb M_{S,N}^{\Mot,\MR},\,
       \mathfrak R^{\rm Conf}_{S,\le M}\bigr),}
 \label{P2M-eq:v33-jet-decorated-motivic-stage}
\end{equation}
where the second entry is a marking ledger and is not an additional direct
summand.  Then the realization
$\rho_S^{\Mot,\MR}$ of \cref{P2M-thm:channel-complete-realization} extends
canonically as a marked realization with the following properties.
\begin{enumerate}[label=\textup{(\roman*)}]
\item The underlying motivic object $\mathbb M_{S,N}^{\Mot,\MR}$, dg
correspondence algebra $\mathscr M_{S,N}^{\Mot,\MR}$, Boolean carrier,
proper-face latching object, exact-support channel blocks and structural Moore
equations $dK_I=N\Omega_I^{\rm rec}$ are unchanged.
\item Ordered support deletion and support permutation act diagonally: by the
existing strict maps on $\mathbb M_{S,N}^{\Mot,\MR}$ and by the internal natural
maps on $\mathfrak R^{\rm Conf}_{S,\le M}$.  Hence the jet-decorated stages are
strict under deletion and permutation-coherent.
\item If the Confluent ledger selects a primitive repeated/root line on a marked
exact-support channel, its motivic realization uses the already existing
primitive $b^{\Mot}$ line in that channel and the same root-localization
antecedent $c^{\Mot}$.  No new Moore pair and no duplicate structural
exact-support cell is introduced.  Distinct jet levels remain distinct as
markings even when they reuse the same underlying root carrier.
\item The vanishing/nonvanishing, effective-carry and closed-lift status
recorded by $\mathfrak o_n^{\rm cc}$ and $\beta_n^{\rm eff}$ is transported as
marked source data.  The present theorem does not reprove those arithmetic
status statements and does not identify jet order with support size,
confluence multiplicity or coefficient depth.
\item The reduced insertion/frontier-history pseudofunctor of
\cref{P2M-thm:v21-bicategorical-motivic-realization} is unchanged.  Jet decorations
are carried through its relabelling and survival rules without adding a second
history-completion mechanism.
\end{enumerate}
Thus finite motivic jet recursion is an enrichment of the existing motivic
Moore--Reedy support recursion, not a second tower.
\end{theorem}

\begin{proof}
By \cref{P2M-thm:channel-complete-realization}, every marked exact-support latching
channel already has a primitive motivic top line and a genuine antecedent built
from the fixed root pair $(b^{\Mot},c^{\Mot})$, and the whole target is strict
under ordered deletion and coherent under permutations.  By
\cref{P2M-thm:v33-imported-confluent-jet-interface}, the finite jet ledger has the
same deletion/permutation/confluence naturality and changes only the marking
status attached to the finite support data.  Taking the ordered pair in
\eqref{P2M-eq:v33-jet-decorated-motivic-stage} therefore requires no change to the
dg object or its latching maps.  A selected primitive repeated/root line is
placed on the marked exact-support channel already specified by the Confluent
interface; exact order and the motivic antecedent then follow from parts
\textup{(iv)--(v)} of \cref{P2M-thm:channel-complete-realization}.  Finally the
history realization acts by relabelling and survival on marked generators, so
adding the jet labels preserves the same bicategorical coherence.
\end{proof}


\begin{definition}[Finite support--jet index category]
\label{P2M-def:v34-support-jet-index}
Fix a finite ordered support $S$ and a finite jet ceiling $M\ge3$.  Let
$\mathsf J_{S,M}$ be the finite category whose objects are pairs $(I,n)$ with
$\varnothing\ne I\subseteq S$ and $3\le n\le M$.  It is generated by two
kinds of arrows:
\begin{enumerate}[label=\textup{(\alph*)}]
\item a support deletion
\[
 (I,n)\longrightarrow(J,n),\qquad \varnothing\ne J\subseteq I,
\]
\item a jet successor
\[
 (I,n)\longrightarrow(I,n+1).
\]
\end{enumerate}
The only mixed relations are the evident interchange squares obtained by
performing one support deletion and one jet successor in either order.
\end{definition}

\begin{definition}[Marked motivic stage category]
\label{P2M-def:v34-marked-motivic-stage-category}
Let $\mathbf{MotMR}_N$ denote the category generated by the finite motivic
Moore--Reedy stages $\mathbb M_{I,N}^{\Mot,\MR}$ and the strict support maps
already constructed in this paper.  Write $\mathbf{MotMRJet}^{\rm mk}_N$ for
the category whose objects are pairs $(\mathbb M,\mathfrak R)$ consisting of
an object of $\mathbf{MotMR}_N$ and a finite Confluent ledger marking.  A
morphism is an underlying motivic dg map together with a compatible map of
marking ledgers.  Forgetting the second coordinate defines
\[
 \operatorname{Und}:\mathbf{MotMRJet}^{\rm mk}_N
 \longrightarrow\mathbf{MotMR}_N.
\]
No new motivic direct summand is implicit in this notation.
\end{definition}

\begin{theorem}[Finite support--jet motivic Moore--Reedy tower]
\label{P2M-thm:v34-support-jet-tower}
Fix odd squarefree $N$, finite ordered $S$ and $M\ge3$.  For
$(I,n)\in\mathsf J_{S,M}$ put
\begin{equation}
 \boxed{
 \mathbb T^{\Mot,\MR}_{N}(I,n)
 :=
 \bigl(\mathbb M_{I,N}^{\Mot,\MR},
       \mathfrak R^{\rm Conf}_{I,\le n}\bigr).}
 \label{P2M-eq:v34-two-axis-stage}
\end{equation}
Then the following data define a strict functor
\begin{equation}
 \boxed{
 \mathbb T^{\Mot,\MR}_{N;S,M}:
 \mathsf J_{S,M}\longrightarrow
 \mathbf{MotMRJet}^{\rm mk}_N.}
 \label{P2M-eq:v34-two-axis-functor}
\end{equation}
\begin{enumerate}[label=\textup{(\roman*)}]
\item For $J\subseteq I$, the support arrow is the existing strict deletion
$P_J^I$ on $\mathbb M_{I,N}^{\Mot,\MR}$ together with the internal support
restriction on $\mathfrak R^{\rm Conf}_{I,\le n}$.
\item The jet successor
\[
 \jmath_{I,n}:\mathbb T_N^{\Mot,\MR}(I,n)
 \longrightarrow\mathbb T_N^{\Mot,\MR}(I,n+1)
\]
has underlying motivic map $\operatorname{id}_{\mathbb M_{I,N}^{\Mot,\MR}}$
and enlarges only the ledger from level $n$ to level $n+1$.
\item Support deletion and jet successor commute strictly:
\begin{equation}
 \boxed{
 (P_J^I,\mathrm{res})\,\jmath_{I,n}
 =
 \jmath_{J,n}\,(P_J^I,\mathrm{res}).}
 \label{P2M-eq:v34-support-jet-interchange}
\end{equation}
Hence every mixed path in the finite grid with the same endpoints induces the
same marked motivic map.
\item The underlying motivic tower is constant in the jet direction:
\begin{equation}
 \boxed{
 \operatorname{Und}\bigl(\mathbb T_N^{\Mot,\MR}(I,n)\bigr)
 =\mathbb M_{I,N}^{\Mot,\MR},
 \qquad
 \operatorname{Und}(\jmath_{I,n})=\operatorname{id}.}
 \label{P2M-eq:v34-constant-underlying-jet-axis}
\end{equation}
Thus increasing the jet ceiling introduces no duplicate Boolean block,
exact-support channel, Moore pair or frontier-history block.
\item Confluence acts only on the marking ledger unless an independent
motivic correspondence realizing that confluence map has already been fixed.
Thus confluence naturality is retained without promoting it to an unproved
third structural tower axis.
\end{enumerate}
\end{theorem}

\begin{proof}
The support maps and their strict composition law are
\cref{P2M-prop:v19-total-target-deletion}.  The ledger restrictions and their
compatibility with support deletion are imported in
\cref{P2M-thm:v33-imported-confluent-jet-interface}.  Jet truncation gives the
canonical inclusion of finite ledgers
$\mathfrak R^{\rm Conf}_{I,\le n}\hookrightarrow
\mathfrak R^{\rm Conf}_{I,\le n+1}$ while leaving the motivic object fixed.
Naturality of every level-$n$ ledger component under support deletion gives
\eqref{P2M-eq:v34-support-jet-interchange}.  The remaining assertions are then
formal from the definitions and from
\cref{P2M-thm:v33-confluent-jet-decoration}.
\end{proof}

\begin{theorem}[Jetwise Reedy reconstruction and latching invariance]
\label{P2M-thm:v34-jetwise-reedy-reconstruction}
Let $|I|\ge3$ and $3\le n\le M$.  At jet level $n$, the proper-face motivic
latching object is still exactly
\[
 \mathbb L_I^{\Mot}
 =\operatorname*{colim}_{J\in\partial I}
   \mathbb M_{J,N}^{\Mot,\MR},
\]
and the Boolean-saturated splitting remains
\begin{equation}
 \boxed{
 \mathbb M_{I,N}^{\Mot,\MR}
 \cong
 \widehat{\mathbb L}_I^{\Mot}
 \oplus\mathbb Q_I^{\Mot,\MR}.}
 \label{P2M-eq:v34-jetwise-reedy-splitting}
\end{equation}
The level-$n$ decoration of the proper-face datum is simply the compatible
family
\[
 \bigl\{\mathfrak R^{\rm Conf}_{J,\le n}\bigr\}_{J\in\partial I};
\]
it does not alter the object colimit.  Consequently:
\begin{enumerate}[label=\textup{(\roman*)}]
\item the underlying latching map and its cofiber are independent of $n$;
\item the fresh exact-support motivic correction is still precisely
$(W_{\rootc}^{\Mot})_I\oplus\mathbb Q_I^{\Mot,\MR}$ before Boolean saturation,
and $\mathbb Q_I^{\Mot,\MR}$ after saturation;
\item the structural antecedent remains $\gamma_I^{\Mot}$ with
$d\gamma_I^{\Mot}=N\beta_I^{\Mot}$;
\item the jet successor from $n$ to $n+1$ commutes with reconstruction: it
changes only the proper-face and exact-support marking ledgers and acts as the
identity on every object in the split formula
\eqref{P2M-eq:v34-jetwise-reedy-splitting}.
\end{enumerate}
Thus Reedy support reconstruction can be performed before or after any finite
number of jet-successor steps with the same underlying motivic result.
\end{theorem}

\begin{proof}
The object latching theorem
\cref{P2M-thm:v20-explicit-motivic-latching} and reconstruction theorem
\cref{P2M-thm:v20-motivic-moore-reedy-reconstruction} depend only on the support
face diagram and the fixed exact-support channel package.  By
\cref{P2M-thm:v34-support-jet-tower}, all jet successors are identities on these
underlying motivic objects and all support restrictions of the ledger commute
with jet truncation.  Hence the face colimit, Boolean saturation, exact-support
cofiber and antecedent are unchanged at every finite jet level.
\end{proof}

\begin{corollary}[No structural cell growth along the jet axis]
\label{P2M-cor:v34-no-jet-structural-growth}
For fixed support $I$ and squarefree $N$, the number of structural motivic
Moore packages in the stages
\[
 \mathbb T_N^{\Mot,\MR}(I,3),\ldots,
 \mathbb T_N^{\Mot,\MR}(I,M)
\]
is independent of $M$.  A higher Confluent obstruction may change from
undetected to detected, from nonprimitive to primitive, or from obstructed to
filled in the marking ledger, but it does not force another copy of
$(b^{\Mot},c^{\Mot})$ or another exact-support channel package.  Any stronger
claim that a higher arithmetic operation has an independent geometric
correspondence realization remains a separate provenance question.
\end{corollary}

\begin{warning}[Layer-I baseline and primitive prime-power input]
\label{P2M-warn:v33-squarefree-jet-scope}
The unconditional all-support baseline of Layer I is proved for odd squarefree
$N_{\rm red}$.  Layer II may carry an arbitrary finite coefficient-depth datum
$\nu$, but the label $\nu$ alone neither creates a primitive exact-depth source
line nor changes support or obstruction height.  When the imported Confluent
package supplies compatible pointed primitive lines of exact orders
$p^{\nu_p}$, the coefficient-faithful comparison theorem
\cref{P2M-thm:prime-power-all-support-comparison} promotes them to a genuine
all-support order-$N_\nu$ motivic Moore--Reedy realization.  Thus the remaining
gate is the arithmetic primitive-line and pointing input, not additional
support geometry.  No coefficient change is identified with $n\mapsto n+1$.
\end{warning}

\begin{corollary}[The tetrahedral realization gate is closed]
\label{P2M-cor:remaining-geometric-gate}
At support four the channel-complete construction simultaneously realizes all
eight singleton--triangle ordered channels and all six ordered $2+2$
channels.  The singleton channels carry the primitive root class, while the
$2+2$ channels carry its transfer multiple.  Hence the tetrahedral Pfaffian
correction and all four highest-face detectors occur in one
permutation-coherent motivic stage.  No separate four-direction assignment to
the three coefficient roles is required.
\end{corollary}


\section{Main Theorem A: motivic Moore--Reedy reconstruction}

Proper faces supply the old motivic data, the Boolean skeleton adds the
full-support tag, and one exact-support Moore package carries the new
obstruction and antecedent.

\subsection{Split face inclusions and the ordinary motivic latching object}

For $J\subseteq I$, let
\begin{equation}\label{P2M-eq:v20-total-target-face-inclusion}
 \iota_J^I:
 \mathbb M_{J,N}^{\Mot,\MR}\longrightarrow
 \mathbb M_{I,N}^{\Mot,\MR}
\end{equation}
be the inclusion which is $i_J^I$ on the Boolean carrier and the identity on
every channel block whose exact support is contained in $J$.  It is split by
$P_J^I$ of \eqref{P2M-eq:v19-total-target-deletion-map}.

\begin{proposition}[Split face system of the motivic target]
\label{P2M-prop:v20-split-face-system}
For $K\subseteq J\subseteq I$,
\[
 \iota_J^I\iota_K^J=\iota_K^I,
 \qquad
 P_K^JP_J^I=P_K^I,
 \qquad
 P_J^I\iota_J^I=\operatorname{id}.
\]
Thus the proper-face diagram of $\mathbb M_{I,N}^{\Mot,\MR}$ is a diagram of
literal split motivic summands.
\end{proposition}

\begin{proof}
The assertion is coordinatewise on the Boolean support tags and on the
exact-support channel blocks.  A channel block occurs in the $J$-stage
exactly when its exact support is contained in $J$.
\end{proof}

For $|I|\ge3$ let $\partial I$ denote the poset of nonempty proper subsets of
$I$, ordered by inclusion.  In the idempotent-complete additive/stable
category generated by the chosen root motive and its finite shifted sums,
define the ordinary motivic latching object
\begin{equation}
 \boxed{
 \mathbb L_I^{\Mot}
 :=
 \operatorname*{colim}_{J\in\partial I}
 \mathbb M_{J,N}^{\Mot,\MR}}
\end{equation}
along the maps \eqref{P2M-eq:v20-total-target-face-inclusion}.  Put
\[
 \mathbb B_I^{<I}
 :=
 \bigoplus_{T\subsetneq I}(W_{\rootc}^{\Mot})_T.
\]

\begin{theorem}[Explicit motivic latching object]
\label{P2M-thm:v20-explicit-motivic-latching}
The split face diagram identifies
\begin{equation}
 \boxed{
 \mathbb L_I^{\Mot}
 \cong
 \mathbb B_I^{<I}
 \oplus
 \bigoplus_{\substack{K\subsetneq I\\|K|\ge3}}
 \ \bigoplus_{(A,B)\in\mathfrak P^\to(K)}
 |\mathbb V_{K;A,B}^{\Mot,\mathrm{ch}}|.}
\end{equation}
The canonical latching map
$\lambda_I^{\rm obj}:\mathbb L_I^{\Mot}\to
\mathbb M_{I,N}^{\Mot,\MR}$ is the inclusion of these summands.  Its cofiber
is
\begin{equation}\label{P2M-eq:v20-latching-cofiber}
 \boxed{
 \cofib(\lambda_I^{\rm obj})
 \cong
 (W_{\rootc}^{\Mot})_I
 \oplus
 \mathbb Q_I^{\Mot,\MR},}
\end{equation}
where
\begin{equation}
 \boxed{
 \mathbb Q_I^{\Mot,\MR}
 :=
 \bigoplus_{(A,B)\in\mathfrak P^\to(I)}
 |\mathbb V_{I;A,B}^{\Mot,\mathrm{ch}}|.}
\end{equation}
\end{theorem}

\begin{proof}
Because every face map is a split coordinate inclusion, the colimit is the
union of the summands appearing on at least one proper face.  A Boolean tag
$T$ appears precisely when $T\subsetneq I$, and a channel block appears
precisely when its exact support $K$ is proper in $I$.  The only remaining
summands of the $I$-stage are the full Boolean tag and the exact-support-$I$
channel package, giving \eqref{P2M-eq:v20-latching-cofiber}.
\end{proof}

\begin{definition}[Boolean-saturated motivic latching object]
\label{P2M-def:v20-boolean-saturated-latching}
Define
\begin{equation}
 \boxed{
 \widehat{\mathbb L}_I^{\Mot}
 :=
 \mathbb L_I^{\Mot}\oplus(W_{\rootc}^{\Mot})_I.}
\end{equation}
This is the ordinary proper-face latching object after the purely skeletal
operation of adding the new full-support Boolean root tag.  It carries no new
Moore detector class.
\end{definition}

Hence objectwise
\begin{equation}\label{P2M-eq:v20-object-reedy-splitting}
 \boxed{
 \mathbb M_{I,N}^{\Mot,\MR}
 \cong
 \widehat{\mathbb L}_I^{\Mot}
 \oplus
 \mathbb Q_I^{\Mot,\MR}.}
\end{equation}
The second summand is where the genuinely new exact-support differential
lives.

\subsection{The latching action on the fresh motivic package}

Let $\mathscr U_{\partial I,N}^{\rm rec}\subset
\mathscr U_{I,N}^{\rm rec}$ be the dg subalgebra generated by all cells
$\widehat K_A$ with $\varnothing\ne A\subsetneq I$.  The support-$I$ source
stage is the semi-free extension
\begin{equation}\label{P2M-eq:v20-source-moore-attachment}
 \mathscr U_{I,N}^{\rm rec}
 =
 \mathscr U_{\partial I,N}^{\rm rec}
 \langle K_I;\ dK_I=N\Omega_I^{\rm rec}\rangle.
\end{equation}

\begin{definition}[Exact-support motivic latching action]
\label{P2M-def:v20-exact-support-latching-action}
On the block indexed by $(A,B)\in\mathfrak P^\to(I)$ define
\[
 \widehat K_A\longmapsto E_{12}^{A,B},
 \qquad
 \widehat K_B\longmapsto E_{23}^{A,B}b^{\Mot},
\]
and send every other proper-support generator to zero.  Taking the direct sum
over all ordered bipartitions defines
\begin{equation}\label{P2M-eq:v20-exact-support-latching-action}
 \boxed{
 \Lambda_I^{\Mot}:
 \mathscr U_{\partial I,N}^{\rm rec}
 \longrightarrow
 \REnd(\mathbb Q_I^{\Mot,\MR}).}
\end{equation}
\end{definition}

\begin{proposition}[The future-channel latching action is dg]
\label{P2M-prop:v20-latching-action-dg}
The map $\Lambda_I^{\Mot}$ is a dg map.  Its value on the support-$I$
curvature polynomial is
\begin{equation}
 \boxed{
 \beta_I^{\Mot}
 :=
 \Lambda_I^{\Mot}(\Omega_I^{\rm rec})
 =
 \bigoplus_{(A,B)\in\mathfrak P^\to(I)}
 \chi_I(A,B)
 E_{13}^{A,B}\mathsf T_N^{\delta(A,B)-1}b^{\Mot}.}
\end{equation}
\end{proposition}

\begin{proof}
Blockwise this is exactly the proper-support calculation of
\cref{P2M-lem:channel-proper-equations,P2M-thm:ordered-channel-realization}.  The
proper-support differential of a nonzero channel entry uses only smaller
cells, all of which have zero component in that exact-support block.  Taking
the orthogonal direct sum therefore preserves every proper-support equation.
The curvature formula is \eqref{P2M-eq:channel-latching-image} in every block.
\end{proof}

Define the exact-support antecedent
\begin{equation}
 \boxed{
 \gamma_I^{\Mot}
 :=
 \bigoplus_{(A,B)\in\mathfrak P^\to(I)}
 \sigma_I^{\rm sh}\chi_I(A,B)
 E_{13}^{A,B}\mathsf T_N^{\delta(A,B)-1}c^{\Mot}.}
\end{equation}

\begin{theorem}[Exact-support motivic Moore package]
\label{P2M-thm:v20-exact-support-motivic-moore-package}
For every finite $I$ with $|I|\ge3$,
\begin{equation}
 \boxed{
 d\gamma_I^{\Mot}=N\beta_I^{\Mot},
 \qquad
 \ord[\beta_I^{\Mot}]=N.}
\end{equation}
Moreover every proper support deletion kills
$\mathbb Q_I^{\Mot,\MR}$ and hence kills both
$\beta_I^{\Mot}$ and $\gamma_I^{\Mot}$.
Thus $\mathbb Q_I^{\Mot,\MR}$ is an exact-support motivic Moore package for
the latching action \eqref{P2M-eq:v20-exact-support-latching-action}.
\end{theorem}

\begin{proof}
The differential identity follows blockwise from
$dc^{\Mot}=Nb^{\Mot}$.  Every component is killed by $N$.  Projection to any
highest-face channel
$(\{i\},I\setminus\{i\})$ gives
$\pm E_{13}b^{\Mot}$ because $\delta=1$ there.  Since
$[b^{\Mot}]$ has exact order $N$, the full vector has exact order $N$.
Exact-support deletion is part of the definition of
$P_J^I$ for $J\subsetneq I$.
\end{proof}

\subsection{The target-side Moore--Reedy theorem}

\begin{theorem}[Motivic Moore--Reedy latching reconstruction]
\label{P2M-thm:v20-motivic-moore-reedy-reconstruction}
Let $I$ be a finite ordered support with $|I|\ge3$.  The support-$I$ motivic
stage and realization are reconstructed from the following data:
\begin{enumerate}[label=\textup{(\roman*)}]
\item the ordinary proper-face motivic latching object
      $\mathbb L_I^{\Mot}$;
\item the full-support Boolean root tag
      $(W_{\rootc}^{\Mot})_I$, giving
      $\widehat{\mathbb L}_I^{\Mot}$;
\item the single exact-support channel package
      $\mathbb Q_I^{\Mot,\MR}$ with its latching action
      $\Lambda_I^{\Mot}$;
\item the antecedent $\gamma_I^{\Mot}$ satisfying
      $d\gamma_I^{\Mot}=N\Lambda_I^{\Mot}(\Omega_I^{\rm rec})$.
\end{enumerate}
More precisely, the object reconstruction is the canonical split formula
\eqref{P2M-eq:v20-object-reedy-splitting}, and the dg realization on
$\mathscr U_{\partial I,N}^{\rm rec}$ is the direct sum of the previously
constructed proper-face realization and the future-channel action
$\Lambda_I^{\Mot}$.  Extending this map across
\eqref{P2M-eq:v20-source-moore-attachment} by
\[
 K_I\longmapsto\gamma_I^{\Mot}
\]
recovers exactly the support-$I$ component of
$\rho_I^{\Mot,\MR}$.

Conversely, restriction of the completed support-$I$ stage to every proper
face removes $\mathbb Q_I^{\Mot,\MR}$ and the full Boolean tag, recovering
the original proper-face diagram.  Hence
\[
 \boxed{
 \begin{gathered}
 \text{proper faces}
 +\text{Boolean support saturation}\\
 {}+\text{one exact-support motivic Moore package}\\
 =\text{the support-}I\text{ motivic stage}
 \end{gathered}}
\]
\end{theorem}

\begin{proof}
The object statement is
\cref{P2M-thm:v20-explicit-motivic-latching}.  On old summands the realization is
the face-colimit of the already constructed stages.  On the fresh
exact-support package the proper-support part is
$\Lambda_I^{\Mot}$ by definition.  Its obstruction is
$\beta_I^{\Mot}$, and
\cref{P2M-thm:v20-exact-support-motivic-moore-package} supplies the genuine
antecedent $\gamma_I^{\Mot}$ required by the source differential
$dK_I=N\Omega_I^{\rm rec}$.  Thus the semi-free extension is a dg map and is
the blockwise map already used in
\cref{P2M-thm:channel-complete-realization}.  Exact-support deletion removes the
new package and the full support tag, proving the converse.
\end{proof}

\begin{corollary}[Motivic no-Segal and exact Moore correction]
\label{P2M-cor:v20-motivic-no-segal}
The ordinary proper-face latching object alone does not reconstruct the
support-$I$ motivic target: its cofiber contains the nonzero exact-support
package \eqref{P2M-eq:v20-latching-cofiber}.  After the harmless Boolean
saturation, the only remaining correction is
$\mathbb Q_I^{\Mot,\MR}$, whose latching class has exact order $N$.
Thus the failure of ordinary face reconstruction in the motivic target is
precisely repaired by the exact-support Moore package constructed above.
\end{corollary}


\section{Main Theorem B: motivic insertion and frontier history}

Strict deletion does not remember an insertion path.  The reduced source
does: each genuine intermediate frontier is a free closed degree-\(-1\)
history line, while structural Moore completion occurs only once.  We add
one motivic history block for each such line.

\begin{proposition}[The unmarked terminal target forgets insertion history]
\label{P2M-prop:v21-unmarked-history-collapse}
Let two structurally reduced insertion histories have the same terminal ordered
support $I$.  If a downstream correspondence model is determined only by the
terminal object $\mathbb M_{I,N}^{\Mot,\MR}$, its split face inclusions and its
exact-support channel packages, then the two histories have the same underlying
terminal motivic correspondence.  In particular such a model cannot be faithful
on the free closed frontier-history lattice of $\mathcal E_{\rm Moore}$.
\end{proposition}

\begin{proof}
All summands of $\mathbb M_{I,N}^{\Mot,\MR}$ are indexed by the final Boolean
supports and final exact-support channel labels.  None is indexed by an
intermediate compressed support occurring only in a chosen insertion history.
The preceding reduced normal form, on the other hand, retains one free closed
generator for every such intermediate frontier.  Forgetting those labels before
passing to the terminal object therefore identifies histories with the same
endpoint.  The statement is only a history-separation no-go; it does not affect
the deletion/Reedy realization already proved.
\end{proof}

\begin{definition}[Elementary motivic history block]
\label{P2M-def:v21-motivic-history-block}
For an oriented frontier-history label $h$, let
$\mathbb H_h^{\Mot}$ be a two-object graded matrix correspondence block on two
shifts of the genuine root carrier $W_{\rootc}^{\Mot}$, with the shifts chosen
so that its distinguished matrix unit
\begin{equation}
 \boxed{\varepsilon_h\in
 \REnd^{-1}(\mathbb H_h^{\Mot})}
\end{equation}
has degree $-1$.  Give the block zero internal differential on this matrix unit:
\begin{equation}\label{P2M-eq:v21-history-closed}
 \boxed{d\varepsilon_h=0.}
\end{equation}
Distinct history blocks are orthogonal.  Reversing the orientation of the
frontier sends $\varepsilon_h$ to $-\varepsilon_h$.
\end{definition}

The matrix unit is free, matching the source history line; using the
\(N\)-torsion class \(b^{\Mot}\) would have the wrong integral type.

\begin{definition}[Reduced motivic history apex]
\label{P2M-def:v21-reduced-motivic-history-apex}
Let $\gamma$ be a marked structurally reduced insertion history with terminal
support $I_m$ and preceding closed-history normal form
\[
 \mathscr C_{I_m}^{\rm all}
 \left\langle
   \delta_{I_1},\ldots,\delta_{I_{m-1}};
   \ d\delta_{I_t}=0
 \right\rangle .
\]
Define its motivic history apex by
\begin{equation}
 \boxed{
 \mathbb A_\gamma^{\Mot,\MR}
 :=
 \mathbb M_{I_m,N}^{\Mot,\MR}
 \oplus
 \bigoplus_{t=1}^{m-1}\mathbb H_{I_t}^{\Mot}.}
\end{equation}
Let $\mathscr A_\gamma^{\Mot,\MR}$ be the dg subalgebra of its derived
endomorphism algebra generated by the terminal structural algebra
$\mathscr M_{I_m,N}^{\Mot,\MR}$ and the history matrix units
$\varepsilon_{I_t}$.
\end{definition}

\begin{proposition}[Reduced-apex dg realization]
\label{P2M-prop:v21-reduced-apex-dg-realization}
There is a canonical marked dg map on the reduced normal form
\begin{equation}
 \boxed{
 \varrho_\gamma^{\Mot,\mathrm{hist}}:
 \mathscr C_{I_m}^{\rm all}
 \left\langle\delta_{I_t};\ d\delta_{I_t}=0\right\rangle
 \longrightarrow
 \mathscr A_\gamma^{\Mot,\MR}}
\end{equation}
characterized by
\[
 K_J\longmapsto \rho_{I_m}^{\Mot,\MR}(K_J),
 \qquad
 \delta_{I_t}\longmapsto\varepsilon_{I_t}.
\]
It is injective on the free abelian history lattice
$\bigoplus_t\Z\delta_{I_t}$.
\end{proposition}

\begin{proof}
The structural generators satisfy their prescribed differentials because
$\rho_{I_m}^{\Mot,\MR}$ is already a dg realization.  Every history generator
is closed preceding and maps to the closed matrix unit
\eqref{P2M-eq:v21-history-closed}.  The reduced preceding theorem says that these
history generators are inert for the structural Moore recursion, so no extra
compatibility equation with $dK_J=N\Omega_J$ has to be imposed.  Distinct
$\varepsilon_{I_t}$ occupy distinct orthogonal matrix blocks and are therefore
linearly independent, proving injectivity on the marked history lattice.
\end{proof}


\begin{definition}[History-marked motivic cospan bicategory]
\label{P2M-def:v21-motivic-history-bicategory}
Let $\mathbf{MotHistCosp}_{N}^{\MR}$ be the smallest marked reduced cospan bicategory generated as
follows.
\begin{enumerate}[label=\textup{(\roman*)}]
\item Its support-$I$ object is $\mathbb M_{I,N}^{\Mot,\MR}$.
\item A strict deletion $I\to J$ is represented by the graph cospan
      \[
       \mathbb M_{I,N}^{\Mot,\MR}
       \xrightarrow{\ P_J^I\ }
       \mathbb M_{J,N}^{\Mot,\MR}
       \xleftarrow{\mathrm{id}}
       \mathbb M_{J,N}^{\Mot,\MR}.
      \]
\item A reduced insertion history $\gamma:J\rightsquigarrow I$ has apex
      $\mathbb A_\gamma^{\Mot,\MR}$ together with the decoration
      $\varrho_\gamma^{\Mot,\mathrm{hist}}$; its source leg is the canonical
      split face inclusion into the terminal structural summand and its target
      leg is the inclusion of $\mathbb M_{I,N}^{\Mot,\MR}$ into the apex.
\item A reordering is represented by the common terminal support object with
      one leg the corresponding permutation relabelling and the other the
      identity; oriented history blocks are relabelled with the same sign rule.
\item Horizontal composition is \emph{reduced}: for a composite history one
      uses the terminal structural object exactly once and adds precisely the
      history blocks indexed by the closed generators in the preceding reduced
      normal form.  One does not push out already completed structural motivic
      apices.
\item Marked $2$-morphisms are blockwise motivic isomorphisms induced by the
      preceding marked structural $2$-cells and oriented transport of the
      history labels.
\end{enumerate}
Associators and unitors are the canonical reindexing isomorphisms of the
terminal structural summand and the finite direct sum of history blocks.
\end{definition}

\begin{theorem}[Bicategorical history-marked motivic Moore--Reedy realization]
\label{P2M-thm:v21-bicategorical-motivic-realization}
For odd squarefree $N$, the source-side structurally reduced pseudofunctor of
the preceding FSCComp theorem admits a marked motivic realization
\begin{equation}
 \boxed{
 \mathcal R_{N}^{\Mot,\MR,\mathrm{hist}}:
 \FSCComp^{\rm mk}
 \longrightarrow
 \mathbf{MotHistCosp}_{N}^{\MR}.}
\end{equation}
It has the following properties.
\begin{enumerate}[label=\textup{(\roman*)}]
\item On objects and strict deletions, forgetting history blocks recovers the
      support-functorial target $\mathbb M_{I,N}^{\Mot,\MR}$ and the strict
      deletion maps of the Moore--Reedy reconstruction.
\item On a reduced insertion chain
      $I_0\subset I_1\subset\cdots\subset I_m$, its apex is exactly
      \[
       \mathbb M_{I_m,N}^{\Mot,\MR}
       \oplus\bigoplus_{t=1}^{m-1}\mathbb H_{I_t}^{\Mot};
      \]
      in particular every genuine intermediate frontier remains visible.
\item Structural higher-support cells are realized only in the terminal
      $\mathbb M_{I_m,N}^{\Mot,\MR}$ factor.  Hence no exact-support motivic
      Moore package is duplicated under composition.
\item Reordering $2$-cells transport the history blocks with their orientation
      signs; mixed deletion--insertion $2$-cells retain a block exactly when
      both labels of its frontier survive and otherwise send it to zero.
\item The induced map on the free marked history lattice is injective.
\item Forgetting the history summands gives the deletion/Reedy motivic
      realization $\rho^{\Mot,\MR}$ and its latching reconstruction theorem.
\end{enumerate}
Thus the constructed motivic target admits a bicategorical insertion lift after
adding only the minimal free history blocks required by the preceding reduced
normal form.
\end{theorem}

\begin{proof}
Apply the preceding reduced pseudofunctor $\mathcal E_{\rm Moore}$ and use its
canonical closed-history normal form.  Send its terminal structural controller
to $\mathbb M_{I_m,N}^{\Mot,\MR}$ by the realization of
\cref{P2M-thm:v20-motivic-moore-reedy-reconstruction}, and send each closed history
generator to the matrix unit of
\cref{P2M-def:v21-motivic-history-block}.  Proposition
\ref{P2M-prop:v21-reduced-apex-dg-realization} gives a dg decoration on every apex.

Upstream reduced composition first composes the pair/frontier history and then
performs structural Moore completion once.  Definition
\ref{P2M-def:v21-motivic-history-bicategory} uses the identical order of operations:
the terminal motivic structural target occurs once and the extra summands are
only history blocks.  Hence composition does not duplicate any
$\mathbb Q_J^{\Mot,\MR}$.  The preceding associators, unitors, reordering maps
and mixed deletion--insertion maps act only by the corresponding relabelling,
orientation and survival rules on the closed history generators.  The same
rules define the blockwise motivic isomorphisms, so the pentagon and triangle
identities reduce to those already satisfied preceding together with canonical
associativity of finite direct sums.  This proves pseudofunctoriality and
\textup{(i)--(vi)}.
\end{proof}

\begin{corollary}[Two-step insertion versus direct insertion remains visible]
\label{P2M-cor:v21-two-step-history-block}
For $J\subset I\subset K$, the direct reduced insertion has motivic apex
$\mathbb M_{K,N}^{\Mot,\MR}$, whereas the two-step insertion has apex
\begin{equation}
 \boxed{
 \mathbb M_{K,N}^{\Mot,\MR}\oplus\mathbb H_I^{\Mot}.}
\end{equation}
Thus the two histories are not identified by the marked motivic realization;
forgetting $\mathbb H_I^{\Mot}$ is precisely the downstream history
contraction which would identify them.
\end{corollary}

\begin{corollary}[Insertion compatibility of motivic latching packages]
\label{P2M-cor:v21-insertion-latching-compatibility}
Under $\mathcal R_N^{\Mot,\MR,\mathrm{hist}}$, the support-$I$ exact-support
package $\mathbb Q_I^{\Mot,\MR}$ is attached exactly once, at the terminal
structural Moore completion where support $I$ first becomes structural.  On any
later insertion it is transported as a proper-support summand of the terminal
motivic target; the only new data contributed by an intermediate compressed
stage are the corresponding free history blocks.

Consequently the Moore relation
\[
 d\gamma_I^{\Mot}=N\beta_I^{\Mot}
\]
and the exact-order-$N$ latching detector are invariant under reduced insertion
composition and reordering.
\end{corollary}



\subsection{Jet-order lift of the frontier-history pseudofunctor}
\label{P2M-subsec:v35-history-jet-interchange}

The two-axis tower of \cref{P2M-thm:v34-support-jet-tower} is objectwise in the
support--jet grid.  We now check that the same jet successor is compatible
with the genuinely bicategorical insertion history of
\cref{P2M-thm:v21-bicategorical-motivic-realization}.  No new history block or
structural Moore summand is introduced.

\begin{definition}[Jet-decorated reduced history apex]
\label{P2M-def:v35-jet-decorated-history-apex}
Let
\[
 \gamma:I_0\rightsquigarrow I_m
\]
be a structurally reduced marked insertion history and let
$\operatorname{Hist}(\gamma)$ be its finite set of genuine intermediate
frontier labels.  For $3\le n\le M$, define the finite history-typed ledger
\begin{equation}
 \boxed{
 \mathfrak R^{\rm Conf,hist}_{\gamma,\le n}
 :=
 \left(
   \mathfrak R^{\rm Conf}_{I_m,\le n};
   \bigl\{
    \operatorname{res}_{e(h)}
    \mathfrak R^{\rm Conf}_{I_m,\le n}
   \bigr\}_{h\in\operatorname{Hist}(\gamma)}
 \right).}
 \label{P2M-eq:v35-history-typed-ledger}
\end{equation}
Here $e(h)$ is the marked frontier support typing the closed history generator
$\delta_h$, and $\operatorname{res}_{e(h)}$ denotes the support/frontier
restriction supplied by
\cref{P2M-thm:v33-imported-confluent-jet-interface}.  The corresponding
jet-decorated motivic apex is
\begin{equation}
 \boxed{
 \mathbb A^{\Mot,\MR;\mathrm{jet}}_{\gamma,n}
 :=
 \bigl(
   \mathbb A^{\Mot,\MR}_{\gamma},
   \mathfrak R^{\rm Conf,hist}_{\gamma,\le n}
 \bigr).}
 \label{P2M-eq:v35-jet-decorated-history-apex}
\end{equation}
The second component is marking data only: the underlying motivic apex and all
history summands $\mathbb H_h^{\Mot}$ are exactly those of
\cref{P2M-def:v21-motivic-history-bicategory}.
\end{definition}

\begin{definition}[Jet-decorated history cospan bicategory]
\label{P2M-def:v35-history-jet-bicategory}
Fix a finite jet ceiling $M$.  Let
$\mathbf{MotHistJetCosp}^{\MR}_{N,\le M}$ be the smallest marked bicategory
containing, for every $3\le n\le M$, the objects
$(\mathbb M_{I,N}^{\Mot,\MR},\mathfrak R^{\rm Conf}_{I,\le n})$ and the
reduced insertion apices of
\eqref{P2M-eq:v35-jet-decorated-history-apex}.  At a fixed level $n$ the
underlying cospans, reduced horizontal composition, history blocks,
associators, unitors, reorderings and mixed deletion--insertion $2$-morphisms
are exactly those of $\mathbf{MotHistCosp}^{\MR}_{N}$.  Between consecutive
levels it also contains the jet arrows whose underlying motivic maps are
identities and whose marking maps are the canonical ledger inclusions.
Forgetting the jet level and every ledger defines a strict homomorphism
\[
 \operatorname{Und}^{\rm hist}:
 \mathbf{MotHistJetCosp}^{\MR}_{N,\le M}
 \longrightarrow
 \mathbf{MotHistCosp}^{\MR}_{N}.
\]
No new motivic direct summand is part of this bicategory.
\end{definition}

\begin{theorem}[Bicategorical history--jet interchange tower]
\label{P2M-thm:v35-history-jet-pseudonatural-tower}
Fix odd squarefree $N$ and a finite jet ceiling $M\ge3$.  For every
$3\le n\le M$, the history-marked realization of
\cref{P2M-thm:v21-bicategorical-motivic-realization} has a jet-decorated lift
\begin{equation}
 \boxed{
 \mathcal R^{\Mot,\MR,\mathrm{hist};\mathrm{jet}}_{N,n}:
 \FSCComp^{\rm mk}
 \longrightarrow
 \mathbf{MotHistJetCosp}^{\MR}_{N,\le M}.}
 \label{P2M-eq:v35-history-jet-pseudofunctor}
\end{equation}
These lifts satisfy the following properties.
\begin{enumerate}[label=\textup{(\roman*)}]
\item Forgetting jet markings recovers the same bicategorical motivic
      pseudofunctor at every level:
      \[
       \operatorname{Und}^{\rm hist}
       \mathcal R^{\Mot,\MR,\mathrm{hist};\mathrm{jet}}_{N,n}
       =
       \mathcal R_N^{\Mot,\MR,\mathrm{hist}}.
      \]
\item For $n<M$ there is a pseudonatural jet-successor transformation
      \begin{equation}
       \boxed{
       \mathbf j_n^{\rm hist}:
       \mathcal R^{\Mot,\MR,\mathrm{hist};\mathrm{jet}}_{N,n}
       \Longrightarrow
       \mathcal R^{\Mot,\MR,\mathrm{hist};\mathrm{jet}}_{N,n+1}.}
       \label{P2M-eq:v35-history-jet-successor}
      \end{equation}
      On every support object, reduced insertion apex and history block its
      underlying motivic map is the identity; it only applies the ledger
      inclusions
      $\mathfrak R^{\rm Conf}_{-,\le n}\hookrightarrow
       \mathfrak R^{\rm Conf}_{-,\le n+1}$.
\item If $\gamma_1$ and $\gamma_2$ are composable reduced insertion histories,
      then reduced horizontal composition commutes with jet successor.  More
      explicitly, whether one first forms the reduced composite and then
      applies $\mathbf j_n^{\rm hist}$, or first applies the two jet successors
      and then composes at level $n+1$, the resulting decorated apex has the
      same terminal structural summand and the same finite set of history
      blocks, and the two marking ledgers are equal.
\item The associators, unitors, reordering $2$-cells and mixed
      deletion--insertion $2$-cells of
      \cref{P2M-def:v21-motivic-history-bicategory} commute with
      $\mathbf j_n^{\rm hist}$.  Hence the pentagon and triangle identities
      hold simultaneously at every finite jet level.
\item Along the entire jet axis, the underlying history apex is constant:
      \begin{equation}
       \boxed{
       \operatorname{Und}^{\rm hist}
       \bigl(\mathbb A^{\Mot,\MR;\mathrm{jet}}_{\gamma,n}\bigr)
       =
       \mathbb A^{\Mot,\MR}_{\gamma}
       =
       \mathbb M^{\Mot,\MR}_{I_m,N}
        \oplus
        \bigoplus_{h\in\operatorname{Hist}(\gamma)}
        \mathbb H_h^{\Mot}.}
       \label{P2M-eq:v35-history-apex-constant-in-jet}
      \end{equation}
      In particular no new frontier-history block and no duplicate structural
      Moore package is created by raising jet order.
\end{enumerate}
\end{theorem}

\begin{proof}
At level $n$, decorate the existing object, deletion and reduced insertion
cospans by the terminal and frontier restrictions of the Confluent ledger.
The support/frontier restrictions commute with jet truncation by
\cref{P2M-thm:v33-imported-confluent-jet-interface,P2M-thm:v34-support-jet-tower}.
The underlying history pseudofunctor uses only the terminal structural object,
closed orthogonal history blocks, relabelling, orientation and
survival/deletion rules.  None of these depends on jet level.

For a reduced composite, the preceding normal form takes the terminal structural
support exactly once and takes the union of the surviving genuine frontier
labels.  Restricting a ledger to those labels and then enlarging the jet
ceiling gives the same tuple as enlarging first and restricting afterwards.
This proves the equality in \textup{(iii)}.  The same observation applies to
reorderings and mixed deletion--insertion maps, while associators and unitors
are only canonical reindexings of the terminal summand and finite direct sum of
history blocks.  Therefore the jet inclusions define the pseudonatural
transformations in \eqref{P2M-eq:v35-history-jet-successor} and commute with all
structural $2$-cells.  The pentagon and triangle identities are inherited from
\cref{P2M-thm:v21-bicategorical-motivic-realization}.  Finally every underlying
map of $\mathbf j_n^{\rm hist}$ is the identity, giving
\eqref{P2M-eq:v35-history-apex-constant-in-jet}.
\end{proof}

\begin{corollary}[Finite support--history--jet interchange]
\label{P2M-cor:v35-support-history-jet-interchange}
Within the odd-squarefree all-support scope, any finite diagram obtained from
ordered support deletions, reduced insertion-history composition or reordering,
and jet-successor maps has the following invariance: two paths with the same
source and target and related by the support--jet interchange squares or the
history associator/unitor/reordering relations induce the same underlying
motivic map and the same final Confluent marking ledger.

Thus the three structures are compatible but play different categorical
roles:
\[
 \boxed{
 \begin{array}{ccl}
 \text{support} &:& \text{Reedy object axis},\\
 \text{history} &:& \text{bicategorical morphism/2-cell layer},\\
 \text{jet order} &:& \text{finite marking axis}.
 \end{array}}
\]
No third structural motivic object tower is introduced by the history or jet
coordinates.
\end{corollary}

\begin{corollary}[No history growth along the jet axis]
\label{P2M-cor:v35-no-history-growth-along-jet-axis}
For a fixed reduced insertion history $\gamma$, the number and motivic types of
frontier-history blocks in
$\mathbb A^{\Mot,\MR;\mathrm{jet}}_{\gamma,n}$ are independent of $n$.
Changing the jet level can change only the obstruction, carry, detector,
primitive-line, filler and pointing status recorded on those marked supports.
Any claim that a higher jet class requires a new geometric history
correspondence is therefore an additional provenance theorem, not part of the
finite Moore--Reedy tower itself.
\end{corollary}

\section{Main Theorem C: secondary bar geometry and all-support transgression}

\subsection{Jetwise constancy of the diagonal bar and secondary pages}

The history--jet theorem above says that increasing jet order changes no
terminal motivic apex and no frontier-history block.  The bar construction is
functorial in precisely those dg data.  We therefore record explicitly that
the secondary-transgression machinery below is not a fourth recursive axis.

\begin{definition}[Jet-decorated diagonal-bar stage]
\label{P2M-def:v36-jet-decorated-diagonal-bar}
Let $\gamma:J\rightsquigarrow I$ be a reduced insertion history and let
$3\le n\le M$.  Write
\[
 \mathfrak R^{\rm Conf}_{\gamma,\le n}
 :=\left(
   \mathfrak R^{\rm Conf}_{I,\le n};
   \{\mathfrak R^{\rm Conf}_{e(h),\le n}\}_{h\in\operatorname{Hist}(\gamma)}
 \right)
\]
for the terminal and surviving frontier restrictions of the finite Confluent
ledger.  Define the marked bar stage
\begin{equation}
 \boxed{
 \mathbb{DBar}_{\gamma,n}^{\Mot,\mathrm{jet}}
 :=\bigl(
       \mathbb{DBar}_{\gamma}^{\Mot},
       \mathfrak R^{\rm Conf}_{\gamma,\le n}
     \bigr).}
 \label{P2M-eq:v36-jet-decorated-diagonal-bar}
\end{equation}
The underlying normalized Hochschild/bar complex, coefficient bimodule and
bar-length filtration are exactly those of
\cref{P2M-def:v23-diagonal-bar-complex,P2M-eq:v24-bar-length-filtration}; the second
coordinate is a marking only.
\end{definition}

\begin{theorem}[Bar--jet interchange and spectral-sequence constancy]
\label{P2M-thm:v36-bar-jet-interchange}
Fix odd squarefree $N$, a finite reduced history $\gamma$ and a finite jet
ceiling $M$.  For every $3\le n<M$, jet successor induces a marked filtered
chain map
\begin{equation}
 \boxed{
 \mathbf j^{\rm bar}_{\gamma,n}:
 \mathbb{DBar}_{\gamma,n}^{\Mot,\mathrm{jet}}
 \longrightarrow
 \mathbb{DBar}_{\gamma,n+1}^{\Mot,\mathrm{jet}}}
 \label{P2M-eq:v36-bar-jet-successor}
\end{equation}
whose underlying filtered chain map is the identity on
$\mathbb{DBar}_{\gamma}^{\Mot}$.  Consequently:
\begin{enumerate}[label=\textup{(\roman*)}]
\item $\mathbf j^{\rm bar}_{\gamma,n}$ commutes strictly with the internal
      differential, the length-lowering Hochschild composition differential
      and the bar-length filtration;
\item on every page on which the bar-length spectral sequence is defined, the
      induced underlying map is the identity and
      \begin{equation}
       \boxed{
       d_r^{\Mot}\,\mathbf j^{\rm bar}_{\gamma,n}
       =
       \mathbf j^{\rm bar}_{\gamma,n}\,d_r^{\Mot}.}
       \label{P2M-eq:v36-jet-bar-page-interchange}
      \end{equation}
\item the native squarefree symbol carrier
      $\mathbb S_\gamma^{\Mot,\rm sq}$ and the symbol isomorphism
      $\sigma_\gamma^{\Mot,\rm sq}$ are independent of $n$ on the underlying
      motivic side; jet successor only extends the ledger marking on their
      support labels;
\item the maps in \eqref{P2M-eq:v36-bar-jet-successor} are natural for reduced
      history composition, associators, unitors, oriented reorderings and
      mixed deletion--insertion maps.  Hence they form the bar-level image of
      the pseudonatural history--jet transformations of
      \cref{P2M-thm:v35-history-jet-pseudonatural-tower};
\item for exact support $|I|=m\ge3$, on the history-zero exact-multilinear
      structural sector used in
      \cref{P2M-thm:v28-full-squarefree-secondary-transgression,P2M-cor:v28-full-motivic-secondary-transgression},
      the same Moore-corrected class exists at every finite jet level and
      satisfies
      \begin{equation}
       \boxed{
       d_{m-1,n}^{\Mot}
       \bigl[N^{m-3}\mathfrak J_I^{\rm bar}\bigr]
       =
       \varepsilon_I
       [1_{\mathscr A}\otimes s\Omega_I^{\rm rec}]
       \ne0,
       \qquad
       \ord[1_{\mathscr A}\otimes s\Omega_I^{\rm rec}]=N.}
       \label{P2M-eq:v36-jetwise-secondary-transgression}
      \end{equation}
      Here the subscript $n$ records only the jet decoration; the underlying
      spectral differential is the original $d_{m-1}^{\Mot}$.
\end{enumerate}
Thus raising jet order does not create a new bar page, alter the Rees
coefficient, or change the motivic Moore target of the secondary
transgression.
\end{theorem}

\begin{proof}
By \cref{P2M-thm:v35-history-jet-pseudonatural-tower}, the jet successor is the
identity on the terminal structural summand, every history block and every dg
map used to make the history apex into an
$\mathscr O_\gamma^{\rm full}$-bimodule.  The normalized Hochschild/bar
construction is functorial in a dg algebra and a compatible bimodule map.
Hence the induced map on
$\mathscr A_\gamma^{\Mot,\MR}\otimes
T^c(s\overline{\mathscr O}_\gamma^{\rm full})$ is literally the identity.
It therefore preserves each bar-length summand and commutes separately with
both parts of the bar differential, proving \textup{(i)}.  The spectral
sequence functoriality of a filtered identity gives \textup{(ii)}.

The squarefree representative lattice and coefficient-unit length-one symbols
depend only on the same fixed source operation algebra and coefficient apex,
so \textup{(iii)} follows from
\cref{P2M-thm:v23-native-bar-symbol-fidelity}.  Naturality in \textup{(iv)} is the
bar functor applied to the pseudonatural squares of
\cref{P2M-thm:v35-history-jet-pseudonatural-tower}; all underlying squares already
commute, while the Confluent ledger restrictions commute with jet truncation.
Finally the staircase chains, no-short-cross-talk theorem and filtered
comparison proving the full motivic transgression are chains and maps in this
unchanged filtered complex.  Applying the identity at level $n$ gives exactly
\eqref{P2M-eq:v36-jetwise-secondary-transgression}, including exact order $N$.
\end{proof}

\begin{corollary}[Jet-invariance of the Moore page, lift length and transfer exponent]
\label{P2M-cor:v36-three-index-jet-invariance}
For $|I|=m\ge3$ and every finite jet level $n$ in the squarefree all-support
tower,
\[
 \boxed{
 \begin{array}{c}
 \text{first full diagonal-bar Moore page}=m-1,\\[1mm]
 \text{minimal total Moore-lift bar length}=m,\\[1mm]
 \text{Rees/transfer exponent}=m-2
 \end{array}}
\]
is independent of $n$.  The support cardinality controls the cobar/bar
indices, the Rees profile controls the transfer exponent, and jet order only
records the independent Confluent obstruction/filler status.  No equality
between these indices is introduced by the finite motivic tower.
\end{corollary}

\begin{warning}[Bar constancy is not independent higher-operation provenance]
\label{P2M-warn:v36-bar-jet-provenance-firewall}
The theorem says that an arithmetic higher-jet class selected by the Confluent
ledger can be tracked through the already constructed motivic bar and Moore
channels without changing their structural differential.  It does not assert
that every new higher arithmetic operation has an independently constructed
geometric correspondence realizing its full operation law.  Such a
correspondence remains additional provenance data, exactly as in
\cref{P2M-cor:v34-no-jet-structural-growth}.
\end{warning}

The structural algebra entering this section is not ad hoc.  By
\cref{P2M-prop:cobar-rees-normal-form}, after forgetting history marks it is the
\(N\)-Rees lattice inside the classical cobar algebra
\[
 \boxed{
 \mathscr U^{\rm rec}_{I,N}
 \subset
 \Cobar(\overline\Lambda_I)\otimes\Z[1/N].}
\]
Thus the bar calculations below are classical bar--cobar geometry with one
new ingredient: tracking the integral Rees lattice.  We use the standard
bar--cobar adjunction and twisting-cochain conventions of
\cite{Adams1956,HMS1974,LodayVallette2012}; the genuinely new conclusions are
the exact Rees coefficient, its motivic realization, and its compatibility
with FSC frontier history.

History blocks preserve frontier labels but are orthogonal to the
structural target, so they kill mixed products that remain nonzero preceding.
We add operation-separated fidelity without replacing the geometric
Moore--Reedy factor.

\subsection{Primary no-go and the finite operation envelope}

Let $\gamma$ be a structurally reduced marked insertion history with terminal
support $I$ and closed-history normal form
\begin{equation}
 \mathscr C_{\gamma}^{\rm red}
 :=\mathscr C_I^{\rm all}
   \left\langle \delta_h\ (h\in\operatorname{Hist}(\gamma));\ d\delta_h=0\right\rangle .
\end{equation}
Here $\operatorname{Hist}(\gamma)$ is the finite set of genuine intermediate frontier labels.
Every structural generator $K_A$ carries support $A$ and every history generator
$\delta_h$ carries the two-label frontier support $e(h)$.

\begin{proposition}[Orthogonal history blocks kill mixed words]
\label{P2M-prop:v22-orthogonal-mixed-word-nogo}
Let $x$ be a structural generator whose motivic image lies in
$\mathscr M_{I,N}^{\Mot,\MR}$ and let $\delta_h$ be a history generator.  In the
history-marked apex
\[
 \mathbb A_{\gamma}^{\Mot,\MR}
 =\mathbb M_{I,N}^{\Mot,\MR}
  \oplus\bigoplus_{h\in\operatorname{Hist}(\gamma)}\mathbb H_h^{\Mot}
\]
one has
\begin{equation}
 \boxed{
 \varrho_{\gamma}^{\Mot,\mathrm{hist}}(x)
 \varrho_{\gamma}^{\Mot,\mathrm{hist}}(\delta_h)=0,
 \qquad
 \varrho_{\gamma}^{\Mot,\mathrm{hist}}(\delta_h)
 \varrho_{\gamma}^{\Mot,\mathrm{hist}}(x)=0.}
\end{equation}
Consequently the history-marked target is faithful on the marked history lattice but not
on any source sector containing a nonzero mixed word $x\delta_h$ or
$\delta_hx$.
\end{proposition}

\begin{proof}
The structural algebra acts block-diagonally on the terminal structural summand,
whereas $\varepsilon_h$ has source and target inside the orthogonal block
$\mathbb H_h^{\Mot}$.  The two off-block composites therefore vanish.  This is
a target relation, not a source relation; the reduced source declares
$\delta_h$ free and closed and imposes no annihilation relation with structural
cells.
\end{proof}

\paragraph{Finite squarefree operation envelope.}

The compact repair uses only the squarefree support sector.  This is the sector
relevant to the support-separated FSC products used throughout the Moore--Reedy
and Hall constructions, and it is finite at each fixed support.

\begin{definition}[Marked squarefree operation envelope]
\label{P2M-def:v22-squarefree-operation-envelope}
Let $\mathcal G_{\gamma}$ be the finite marked generating set consisting of all
structural cells $K_A$ in $\mathscr C_I^{\rm all}$ together with all history
cells $\delta_h$.  Give them supports
\[
 \operatorname{supp}(K_A)=A,
 \qquad
 \operatorname{supp}(\delta_h)=e(h).
\]
Let
\begin{equation}
 \boxed{
 \mathscr O_{\gamma}^{\rm full}:=T_{\Z}(\mathcal G_{\gamma})}
\end{equation}
be the marked associative tensor envelope, with the differential extended as a
graded derivation from the structural Moore equations and $d\delta_h=0$.  The
marked generator realization of the history-marked construction extends multiplicatively to
$\mathscr O_{\gamma}^{\rm full}$.

Let $\mathfrak J_{\gamma}^{\rm rep}$ be the graded ideal generated by every
marked word $g_1\cdots g_r$ for which two factor supports meet, and define
\begin{equation}
 \boxed{
 \mathscr O_{\gamma}^{\rm sq}
 :=\mathscr O_{\gamma}^{\rm full}/\mathfrak J_{\gamma}^{\rm rep}.}
\end{equation}
Write
\[
 q_{\gamma}^{\rm sq}:\mathscr O_{\gamma}^{\rm full}
 \twoheadrightarrow\mathscr O_{\gamma}^{\rm sq}
\]
for the quotient.  Thus the nonzero quotient words have pairwise disjoint
nonempty supports.  Distinct words, orderings and marked history labels remain
distinct basis coordinates.
\end{definition}

\begin{lemma}[Squarefree operation envelope is a finite dg algebra]
\label{P2M-lem:v22-squarefree-operation-dg}
The repeated-support ideal $\mathfrak J_{\gamma}^{\rm rep}$ is stable under the
differential.  Hence $\mathscr O_{\gamma}^{\rm sq}$ is a dg algebra.  Its
underlying graded abelian group is finite free.
\end{lemma}

\begin{proof}
The differential of a structural generator of support $A$ is a signed integral
sum of products whose factor supports partition $A$; the history generators are
closed.  If a word already contains a second factor whose support intersects
$A$, then after replacing $A$ by any such partition that second support still
intersects at least one new factor.  The repeated-support ideal is therefore a
dg ideal.

Every surviving factor has nonempty support and surviving factor supports are
pairwise disjoint subsets of the finite set $I$.  Thus a nonzero word has length
at most $|I|$.  There are only finitely many marked generators on each support
and only finitely many ordered disjoint-support words, proving finite freeness.
\end{proof}

\paragraph{Squarefree scope.}
The squarefree quotient is a finite fidelity sector, not a source relation;
repeated-support words remain in the Ind-completion.


\paragraph{Native mixed-operation symbols.}


The left-regular construction proves multiplicative faithfulness, but its
underlying dg lattice should not be regarded as an unrelated extra motivic
object.  There is a canonical bar-geometric model for exactly that lattice.

Let
\[
 \epsilon_\gamma:\mathscr O_\gamma^{\rm full}\longrightarrow\Z
\]
be the augmentation sending every positive marked generator to zero, and write
\(\overline{\mathscr O}_\gamma^{\rm full}:=\ker\epsilon_\gamma\).
Let
\[
 \widetilde\varrho_{\gamma}^{\Mot,\mathrm{hist}}:
 \mathscr O_\gamma^{\rm full}
 \longrightarrow
 \mathscr A_\gamma^{\Mot,\MR}
\]
be the multiplicative extension of the geometric/history realization from
the history-marked construction.  We use the same map on the left and on the right to make
\(\mathscr A_\gamma^{\Mot,\MR}\) an
\(\mathscr O_\gamma^{\rm full}\)-bimodule.

\begin{definition}[Diagonal motivic bar/Hochschild complex]
\label{P2M-def:v23-diagonal-bar-complex}
The \emph{diagonal motivic bar complex} of the reduced apex is the normalized
Hochschild/bar complex
\begin{equation}
 \boxed{
 \mathbb{DBar}_{\gamma}^{\Mot}
 :=
 C_\bullet^{\rm Hoch}\!\left(
   \mathscr O_\gamma^{\rm full};
   \mathscr A_\gamma^{\Mot,\MR}
 \right)
 =
 \mathscr A_\gamma^{\Mot,\MR}
 \otimes
 T^c\!\left(s\overline{\mathscr O}_\gamma^{\rm full}\right),}
\end{equation}
where the coefficient bimodule structure is induced on both sides by
\(\widetilde\varrho_{\gamma}^{\Mot,\mathrm{hist}}\).
We use the standard normalized suspension convention.  Thus for a homogeneous
\(a\in\overline{\mathscr O}_\gamma^{\rm full}\), the coefficient-unit
length-one symbol satisfies
\begin{equation}\label{P2M-eq:v23-length-one-bar-differential}
 d\!\left(1_{\mathscr A}\otimes sa\right)
 =
 -\,1_{\mathscr A}\otimes s(da)
 +\widetilde\varrho_\gamma(a)
 -\widetilde\varrho_\gamma(a)
 =
 -\,1_{\mathscr A}\otimes s(da).
\end{equation}
\end{definition}

\paragraph{Three-direction specialization.}
With distinct endpoint actions, the same length-one bar formula is precisely
the same arithmetic--motivic bar comparison established in the preceding structural construction; the
current v0.48r13 snapshot is informational only.  The diagonal bar is its
equal-endpoint specialization.  In particular the preceding cells satisfy
\[
 d\beta_B^{\Mot}=b_{\rm Gal}e_{\Mot}-e_{\Mot}b^{\Mot},
 \qquad
 d\beta_C^{\Mot}=-N\beta_B^{\Mot}-e_{\Mot}c^{\Mot}.
\]
The endpoint bridge is Tate-neutral.  The Moore pair lies in the common sector
$\tau_{\rm root}$ of \cref{P2M-thm:root-common-tate-sector}; the comparison cells
are typed in that same sector.  Equivalently, after the root-relative grading
$q_{\rm rel}=q-\tau_{\rm root}$, the displayed comparison equations have the
old relative bidegrees $(-1,0)$ and $(-2,0)$.  Thus the equations are
bihomogeneous without assuming absolute Tate degree zero.

The cancellation in \eqref{P2M-eq:v23-length-one-bar-differential} is the diagonal
analogue of that geometric comparison-bar formula.  There the two endpoint
actions are different and the length-one differential records their
difference; here they are deliberately the same, so the endpoint discrepancy
vanishes and the source operation itself survives as a closedly transported
bar symbol.

Let \(\mathscr W_\gamma^{\rm sq}\subset
\mathscr O_\gamma^{\rm full}\) be the free abelian subcomplex spanned by the
unit and the marked words whose factor supports are pairwise disjoint.  By
\cref{P2M-lem:v22-squarefree-operation-dg}, the quotient map
\(q_\gamma^{\rm sq}\) restricts to an isomorphism of dg lattices
\begin{equation}\label{P2M-eq:v23-squarefree-representative-lattice}
 \mathscr W_\gamma^{\rm sq}
 \xrightarrow{\ \sim\ }
 \mathscr O_\gamma^{\rm sq}.
\end{equation}

\begin{definition}[Finite diagonal-bar symbol carrier]
\label{P2M-def:v23-bar-symbol-carrier}
Let \(\mathbb S_\gamma^{\Mot,\rm sq}\) be the finite motivic twisted complex
whose unit line is \(\mathbf1^{\Mot}\) and whose nonunit basis lines are the
desuspended coefficient-unit length-one symbols
\[
 \left(1_{\mathscr A}\otimes sw\right)[-1],
 \qquad
 0\ne w\in\mathscr W_\gamma^{\rm sq}.
\]
Its differential is induced from
\eqref{P2M-eq:v23-length-one-bar-differential}.  Equivalently it is the finite
subcomplex of the shifted diagonal bar complex generated by the coefficient
unit and the squarefree length-one symbols.
\end{definition}

\begin{theorem}[Native geometric bar-symbol fidelity]
\label{P2M-thm:v23-native-bar-symbol-fidelity}
There is a canonical isomorphism of dg motivic lattices
\begin{equation}
 \boxed{
 \sigma_\gamma^{\Mot,\rm sq}:
 \mathbf1^{\Mot}\otimes_\Z\mathscr O_\gamma^{\rm sq}
 \xrightarrow{\ \sim\ }
 \mathbb S_\gamma^{\Mot,\rm sq}.}
\end{equation}
It sends the unit word to the coefficient unit and every nonunit squarefree
word \(w\) to
\[
 \sigma_\gamma^{\Mot,\rm sq}(w)
 =
 \left(1_{\mathscr A}\otimes s\widetilde w\right)[-1],
\]
where \(\widetilde w\in\mathscr W_\gamma^{\rm sq}\) is its unique squarefree
representative.  Consequently distinct squarefree mixed
history/structural words remain distinct already in a canonical
bar-correspondence symbol complex.
\end{theorem}

\begin{proof}
Equation \eqref{P2M-eq:v23-squarefree-representative-lattice} gives a unique
squarefree representative for every basis word in
\(\mathscr O_\gamma^{\rm sq}\).  The coefficient unit of
\(\mathscr A_\gamma^{\Mot,\MR}\) is a primitive integral line, so distinct
length-one tensors \(1_{\mathscr A}\otimes s\widetilde w\) are linearly
independent.  After shifting the length-one sector by \([-1]\),
\eqref{P2M-eq:v23-length-one-bar-differential} becomes
\[
 d\,\sigma_\gamma^{\Mot,\rm sq}(w)
 =
 \sigma_\gamma^{\Mot,\rm sq}(dw).
\]
Thus the displayed map is a chain isomorphism on the indicated finite basis.
No injectivity property of the endpoint action
\(\widetilde\varrho_\gamma\) is used: even when a mixed word acts trivially on
the orthogonal geometric apex, its marked source word remains a distinct bar
tensor.
\end{proof}

\begin{corollary}[Bar-native mixed-word separation]
\label{P2M-cor:v23-bar-native-mixed-word-separation}
For a structural marked word \(x\) and history generator \(\delta_h\) with
pairwise disjoint factor supports, any two distinct nonzero squarefree words
among
\[
 x\delta_h,\qquad \delta_hx
\]
have distinct nonzero images in
\(\mathbb S_\gamma^{\Mot,\rm sq}\), even if both products vanish in the
orthogonal geometric endomorphism algebra of the history-marked construction.
\end{corollary}

\begin{proposition}[Naturality of diagonal-bar symbols]
\label{P2M-prop:v23-bar-symbol-naturality}
The symbol complexes and maps
\(\sigma_\gamma^{\Mot,\rm sq}\) are natural for the reduced insertion
associators and unitors, oriented reorderings, and mixed
deletion--insertion maps of the preceding marked Moore-history pseudofunctor.
A surviving marked word is transported to the correspondingly relabelled bar
symbol; a word whose support loses a required label is sent to zero.
\end{proposition}

\begin{proof}
All preceding structural \(2\)-cells act by marked dg maps on the reduced
generating set and hence on
\(\mathscr O_\gamma^{\rm full}\).  They also act on the motivic coefficient
apex by the history-marked pseudofunctor.  Normalized Hochschild/bar chains are functorial
for a dg algebra map together with a compatible bimodule map.  The coefficient
unit is preserved, so the length-one symbol sector is carried to itself with
the same orientation sign and survival rule.  The preceding pentagon and
triangle identities therefore induce the corresponding identities on symbols.
\end{proof}

\paragraph{Bar scope.}
Bar symbols separate words at chain level; total bar homology also imposes
composition differentials, beginning with \(d_1\).


\subsection{The bar-length spectral sequence and first composition transgression}


Filter the diagonal complex of
\cref{P2M-def:v23-diagonal-bar-complex} by bar length:
\begin{equation}\label{P2M-eq:v24-bar-length-filtration}
 F_p\mathbb{DBar}_{\gamma}^{\Mot}
 :=
 \bigoplus_{0\le r\le p}
 \mathscr A_\gamma^{\Mot,\MR}
 \otimes
 \left(s\overline{\mathscr O}_\gamma^{\rm full}\right)^{\otimes r}.
\end{equation}
The internal differential preserves length, while every Hochschild
multiplication or endpoint-action term lowers length by one.  Thus
\eqref{P2M-eq:v24-bar-length-filtration} is an increasing dg filtration and gives
the bar-length spectral sequence
\[
 E_r^{p,*}(\gamma)\Longrightarrow
 H^*\!\left(\mathbb{DBar}_{\gamma}^{\Mot}\right)
\]
whenever the usual convergence hypotheses are imposed.  Only the first two
pages will be used below, so no global convergence assertion is needed.

\begin{proposition}[First two bar pages]
\label{P2M-prop:v24-first-two-bar-pages}
The \(E_0\)-differential is induced only by the internal differentials of
\(\mathscr A_\gamma^{\Mot,\MR}\) and
\(\mathscr O_\gamma^{\rm full}\).  The first spectral differential
\[
 d_1:E_1^{p,*}\longrightarrow E_1^{p-1,*}
\]
is induced by the length-lowering Hochschild composition differential:
left endpoint action, adjacent multiplication of bar letters, and right
endpoint action.
For a homogeneous internally closed element
\(a\in\overline{\mathscr O}_\gamma^{\rm full}\), the coefficient-unit
length-one class
\[
 [\,1_{\mathscr A}\otimes sa\,]\in E_1^{1,*}
\]
is a \(d_1\)-cycle.
\end{proposition}

\begin{proof}
The decomposition of the normalized Hochschild differential by bar-length
change is standard.  The length-one endpoint formula is
\eqref{P2M-eq:v23-length-one-bar-differential}; after passing to internal
homology the remaining left and right endpoint terms are equal with opposite
Hochschild signs, so the induced \(d_1\) vanishes on the coefficient-unit
length-one symbol.
\end{proof}

For closed homogeneous \(a,b\), let
\[
 \operatorname{Alt}_{\rm Hoch}(a,b)
 =
 1_{\mathscr A}\otimes(sa|sb)
 -
 \varepsilon_{\rm Hoch}(a,b)\,
 1_{\mathscr A}\otimes(sb|sa),
\]
where \(\varepsilon_{\rm Hoch}(a,b)\in\{\pm1\}\) is the standard suspended
Koszul sign which makes the two pairs of endpoint-action terms cancel.  The
remaining sign below depends only on the fixed suspension convention.

\begin{theorem}[First composition differential is graded cyclicization]
\label{P2M-thm:v24-first-composition-cyclicization}
Let \(a,b\in\overline{\mathscr O}_\gamma^{\rm full}\) be homogeneous internal
cycles.  On the \(E_1\)-page of the bar-length filtration,
\begin{equation}\label{P2M-eq:v24-commutator-transgression}
 \boxed{
 d_1\!\left[\operatorname{Alt}_{\rm Hoch}(a,b)\right]
 =
 \pm\,
 \left[
   1_{\mathscr A}\otimes
   s\!\left(
     ab-(-1)^{|a||b|}ba
   \right)
 \right].}
\end{equation}
Consequently every graded commutator class represented by internally closed
positive operations maps to zero on the coefficient-unit length-one
\(E_2\)-edge.
\end{theorem}

\begin{proof}
Apply the length-lowering part of the normalized graded Hochschild
differential to the two length-two tensors.  By the defining choice of
\(\varepsilon_{\rm Hoch}(a,b)\), the left endpoint terms cancel pairwise and
the right endpoint terms cancel pairwise.  The only surviving contributions
are the two adjacent-multiplication terms.  Their suspended Koszul signs are
exactly the graded-commutator sign, up to one common overall orientation unit.
This gives \eqref{P2M-eq:v24-commutator-transgression}.
\end{proof}

\begin{corollary}[The \(E_2\) symbol edge factors through graded abelianization]
\label{P2M-cor:v24-E2-graded-abelianization}
Let \(H_+(\mathscr O_\gamma)\) denote the positive part of the internal
homology algebra.  The coefficient-unit length-one symbol map to the
\(E_2\)-page factors through the graded abelianization
\begin{equation}
 \boxed{
 H_+(\mathscr O_\gamma)
 \longrightarrow
 H_+(\mathscr O_\gamma)_{\rm ab}
 :=
 \frac{H_+(\mathscr O_\gamma)}
      {[H_+(\mathscr O_\gamma),H_+(\mathscr O_\gamma)]_{\rm gr}}
 \longrightarrow
 E_2^{1,*}(\gamma).}
\end{equation}
No equality with the full graded abelianization is asserted: additional
first-page boundaries may impose further relations.
\end{corollary}

There is nevertheless a canonical primitive sector which the first
composition differential cannot remove.  Put
\[
 \mathfrak I_\gamma:=\ker\epsilon_\gamma,\qquad
 Q_\gamma^{\rm ind}:=
 \mathfrak I_\gamma/\mathfrak I_\gamma^2.
\]
Every structural or history generator has decomposable differential, so the
induced differential on \(Q_\gamma^{\rm ind}\) is zero.  Let
\[
 p_{\mathbf1}:
 \mathscr A_\gamma^{\Mot,\MR}\longrightarrow
 \Z\,1_{\mathscr A}
\]
be the marked coefficient-unit projection, killing every positive structural
or history operator.

\begin{theorem}[Primitive marked lines survive the first composition page]
\label{P2M-thm:v24-primitive-survival}
Let \(g\) be an internally closed marked generator whose image in
\(Q_\gamma^{\rm ind}\) is nonzero.  Then
\begin{equation}\label{P2M-eq:v24-primitive-E2-survival}
 \boxed{
 [\,1_{\mathscr A}\otimes sg\,]
 \ne0
 \quad\text{in }E_2^{1,*}(\gamma).}
\end{equation}
In particular every singleton direction generator \(D_i\) and every free
closed reduced-history generator \(\delta_h\) survives the first composition
differential.
\end{theorem}

\begin{proof}
The class is nonzero on \(E_1\): if \(g\) were an internal boundary, its image
under the chain map to the zero-differential indecomposable quotient
\(Q_\gamma^{\rm ind}\) would vanish.  Suppose its coefficient-unit bar symbol
were a \(d_1\)-boundary.  Apply \(p_{\mathbf1}\) to the coefficient and then
project the remaining length-one operation to \(Q_\gamma^{\rm ind}\).
Endpoint-action terms have positive coefficient and are killed by
\(p_{\mathbf1}\); adjacent-multiplication terms land in
\(\mathfrak I_\gamma^2\) and are killed by the indecomposable projection.
Hence every \(d_1\)-boundary maps to zero, whereas the symbol of \(g\) maps to
its assumed nonzero indecomposable class.  This contradiction proves
\eqref{P2M-eq:v24-primitive-E2-survival}.
\end{proof}


\paragraph{Primitive retract and permanent frontier classes.}


The \(E_2\)-argument above used the coefficient-unit projection only against
the first composition differential.  In fact the same idea gives a chain map
on the entire diagonal bar complex and therefore rules out every possible
higher zigzag on the primitive marked sector.

\begin{lemma}[Unit-coefficient firewall]
\label{P2M-lem:v25-unit-coefficient-firewall}
The marked coefficient-unit projection
\[
 p_{\mathbf1}:
 \mathscr A_\gamma^{\Mot,\MR}\longrightarrow \Z\,1_{\mathscr A}
\]
is a chain projection.  If
\(a\in\mathfrak I_\gamma=\ker\epsilon_\gamma\) is a positive source operation
and \(m\in\mathscr A_\gamma^{\Mot,\MR}\), then
\begin{equation}\label{P2M-eq:v25-unit-coefficient-firewall}
 \boxed{
 p_{\mathbf1}\!\left(
   m\,\widetilde\varrho_\gamma(a)
 \right)
 =
 p_{\mathbf1}\!\left(
   \widetilde\varrho_\gamma(a)\,m
 \right)
 =0.}
\end{equation}
\end{lemma}

\begin{proof}
Every image of a positive source operation has positive marked support/path
degree: singleton directions are support shifts, higher structural cells are
placed in positive channel entries, and frontier-history generators lie in
their marked history factors.  The differentials preserve the direct sum of
positive marked operators because every source attaching differential is
decomposable and has no scalar term.  The explicitly generated target contains
no inverse positive path which could multiply one of these images back to the
global identity correspondence.  Thus the identity line is a dg direct
summand and \eqref{P2M-eq:v25-unit-coefficient-firewall} follows.
\end{proof}

Recall that
\[
 Q_\gamma^{\rm ind}
 =
 \mathfrak I_\gamma/\mathfrak I_\gamma^2
\]
has zero induced differential: every marked generator has either zero
differential or a decomposable attaching differential.  Write
\(sQ_\gamma^{\rm ind}\) for the same suspension used in the normalized bar.

\begin{definition}[Total-bar indecomposable detector]
\label{P2M-def:v25-total-bar-ind-detector}
Define
\begin{equation}\label{P2M-eq:v25-total-bar-ind-detector}
 \boxed{
 \Pi_\gamma^{\rm ind}:
 \mathbb{DBar}_\gamma^{\Mot}
 \longrightarrow
 sQ_\gamma^{\rm ind}}
\end{equation}
on a homogeneous normalized bar tensor by
\[
 \Pi_\gamma^{\rm ind}
 \bigl(
 m\otimes(sa_1|\cdots|sa_r)
 \bigr)
 :=
 \begin{cases}
 p_{\mathbf1}(m)\,s\overline{a_1},&r=1,\\[1mm]
 0,&r\ne1.
 \end{cases}
\]
Give \(sQ_\gamma^{\rm ind}\) the zero differential and the filtration
concentrated in bar length one.
\end{definition}

\begin{theorem}[The indecomposable detector is a filtered chain map]
\label{P2M-thm:v25-ind-detector-chain-map}
The map \(\Pi_\gamma^{\rm ind}\) of
\eqref{P2M-eq:v25-total-bar-ind-detector} is a filtered chain map.
\end{theorem}

\begin{proof}
For the internal differential on a length-one tensor, the coefficient
contribution has zero unit projection because \(p_{\mathbf1}\) is a chain
projection, while the operation contribution maps to zero in
\(Q_\gamma^{\rm ind}\) because every generator differential is decomposable.
For the length-lowering Hochschild differential, only a length-two tensor can
contribute to the detector.  Its two endpoint-action terms have zero unit
coefficient by \cref{P2M-lem:v25-unit-coefficient-firewall}, and its adjacent
multiplication term lies in \(\mathfrak I_\gamma^2\).  Tensors of larger
length still have length at least two after one Hochschild multiplication and
are sent to zero.  Finally a length-one tensor is sent by the endpoint
differential to length zero, where the detector is zero.  Hence every component
of the total differential commutes with \(\Pi_\gamma^{\rm ind}\).
Filtration compatibility is immediate.
\end{proof}

Let \(\mathscr P_\gamma\subset Q_\gamma^{\rm ind}\) denote the marked primitive
sublattice
\begin{equation}
 \boxed{
 \mathscr P_\gamma
 :=
 \bigoplus_{i\in I_m}\Z\,\overline{D_i}
 \ \oplus\
 \bigoplus_{h\in\operatorname{Hist}(\gamma)}
 \Z\,\overline{\delta_h}.}
\end{equation}
Because the operation envelope is free on the marked generators before
decomposable attaching relations are imposed, the displayed lines are
independent summands of \(Q_\gamma^{\rm ind}\).  Let
\(\operatorname{pr}_{\mathscr P}:Q_\gamma^{\rm ind}\to\mathscr P_\gamma\)
be the corresponding marked projection and put
\[
 \Pi_\gamma^{\mathscr P}
 :=
 s\operatorname{pr}_{\mathscr P}\circ\Pi_\gamma^{\rm ind}.
\]

\begin{definition}[Primitive bar section]
\label{P2M-def:v25-primitive-bar-section}
Define
\begin{equation}
 \jmath_\gamma^{\mathscr P}:
 s\mathscr P_\gamma
 \longrightarrow
 \mathbb{DBar}_\gamma^{\Mot}
\end{equation}
by
\[
 s\overline{D_i}\longmapsto1_{\mathscr A}\otimes sD_i,
 \qquad
 s\overline{\delta_h}\longmapsto1_{\mathscr A}\otimes s\delta_h.
\]
\end{definition}

\begin{theorem}[Primitive diagonal-bar retract]
\label{P2M-thm:v25-primitive-bar-retract}
The maps
\(\jmath_\gamma^{\mathscr P}\) and
\(\Pi_\gamma^{\mathscr P}\) are filtered chain maps and satisfy
\begin{equation}
 \boxed{
 \Pi_\gamma^{\mathscr P}\,
 \jmath_\gamma^{\mathscr P}
 =
 \operatorname{id}_{s\mathscr P_\gamma}.}
\end{equation}
Consequently there is a filtered chain decomposition
\begin{equation}\label{P2M-eq:v25-total-bar-splitting}
 \boxed{
 \mathbb{DBar}_\gamma^{\Mot}
 \cong
 \jmath_\gamma^{\mathscr P}(s\mathscr P_\gamma)
 \oplus
 \ker\Pi_\gamma^{\mathscr P}.}
\end{equation}
\end{theorem}

\begin{proof}
Both \(D_i\) and \(\delta_h\) are internally closed.  Their diagonal
coefficient-unit length-one bar tensors have vanishing endpoint differential,
so \(\jmath_\gamma^{\mathscr P}\) is a chain map.  It is filtered in length
one.  The detector is a filtered chain map by
\cref{P2M-thm:v25-ind-detector-chain-map}, and the marked projection to
\(\mathscr P_\gamma\) is a chain map because the indecomposable differential is
zero.  The composite is visibly the identity on the displayed marked basis,
which gives the direct-sum decomposition.
\end{proof}

\begin{corollary}[Native motivic bar-homology history classes]
\label{P2M-cor:v25-native-bar-history-classes}
Every singleton direction and every reduced frontier-history line defines a
nonzero independent class in the homology of the total native diagonal bar:
\begin{equation}
 \boxed{
 \bigoplus_{i\in I_m}\Z\,
 [1_{\mathscr A}\otimes sD_i]
 \ \oplus\
 \bigoplus_{h\in\operatorname{Hist}(\gamma)}\Z\,
 [1_{\mathscr A}\otimes s\delta_h]
 \hookrightarrow
 H^*\!\left(\mathbb{DBar}_\gamma^{\Mot}\right).}
\end{equation}
The image is a free direct summand.
\end{corollary}

\begin{proof}
Take homology in the filtered chain splitting
\eqref{P2M-eq:v25-total-bar-splitting}.  The first summand has zero differential,
so its homology is itself.
\end{proof}

\begin{corollary}[No higher bar zigzag can erase primitive history]
\label{P2M-cor:v25-no-higher-zigzag}
On every page on which the bar-length spectral sequence is defined, the
classes represented by
\[
 1_{\mathscr A}\otimes sD_i,
 \qquad
 1_{\mathscr A}\otimes s\delta_h
\]
form a direct summand concentrated in filtration one.  They support no
nonzero \(d_r\), and no \(d_r\) can hit them, for any \(r\ge1\).
In particular the \(E_2\to E_3\) differential vanishes identically on the
primitive direction/history sector.
\end{corollary}

\begin{proof}
The splitting \eqref{P2M-eq:v25-total-bar-splitting} is filtered, so the associated
spectral sequence splits pagewise into the spectral sequence of the zero
differential complex \(s\mathscr P_\gamma\), concentrated in filtration one,
and that of \(\ker\Pi_\gamma^{\mathscr P}\).  The first summand has no
differentials on any page and cannot receive a differential from the second.
\end{proof}

\begin{corollary}[All higher bar transgressions are decomposable-side phenomena]
\label{P2M-cor:v25-higher-transgression-kernel}
Every nonzero higher bar-length differential \(d_r\), \(r\ge2\), is supported
inside the complementary filtered complex
\[
 \boxed{\ker\Pi_\gamma^{\mathscr P}.}
\]
Thus the next unresolved bar-homology problem concerns only decomposable,
cyclicized, and higher-relation sectors; primitive support and insertion
frontier history is already permanently realized.
\end{corollary}

\begin{proposition}[Naturality of the primitive bar retract]
\label{P2M-prop:v25-primitive-retract-naturality}
The primitive retract is natural for the reduced insertion associators and
unitors, oriented reorderings, and mixed deletion--insertion maps.  A direction
class is relabelled or deleted with its support label; a history class is
transported with its orientation sign exactly when its frontier survives and
is sent to zero when a frontier label is deleted.
\end{proposition}

\begin{proof}
The preceding reduced pseudofunctor transports \(D_i\) and the closed
\(\delta_h\) by precisely these marked rules.  The history-marked motivic insertion
realization uses the same rules on their target factors.  Both the
coefficient-unit projection and the indecomposable marked projection commute
with relabelling and survival/deletion.  Hence the section and retraction
commute with all indicated structural \(2\)-cells.
\end{proof}


\subsection{The first genuine secondary bar differential}


The primitive retract rules out every higher differential on
\(D_i\) and \(\delta_h\), but it does not make the complementary
decomposable sector spectrally trivial.  The first nonzero secondary
transgression already occurs on three singleton directions.

Write
\[
 q_\gamma^{\rm sq}:
 \mathscr O_\gamma^{\rm full}
 \longrightarrow
 \mathscr O_\gamma^{\rm sq}
\]
for the squarefree dg-algebra quotient and retain the coefficient-unit
projection \(p_{\mathbf1}\) of
\cref{P2M-lem:v25-unit-coefficient-firewall}.

\begin{definition}[Squarefree augmentation bar]
\label{P2M-def:v26-squarefree-augmentation-bar}
The \emph{squarefree augmentation bar} is
\begin{equation}
 \boxed{
 \mathbb B_\gamma^{\rm sq}
 :=
 C_\bullet^{\rm Hoch}\!\left(
   \mathscr O_\gamma^{\rm sq};\Z_\epsilon
 \right)
 =
 \Z\otimes
 T^c\!\left(s\overline{\mathscr O}_\gamma^{\rm sq}\right),}
\end{equation}
where both endpoint actions on \(\Z_\epsilon\) are given by the augmentation.
It is filtered by bar length.
\end{definition}

\begin{proposition}[Motivic-to-squarefree bar comparison]
\label{P2M-prop:v26-motivic-squarefree-bar-comparison}
The formula
\begin{equation}
 \boxed{
 \mathfrak q_\gamma^{\rm bar}
 \bigl(
   m\otimes(sa_1|\cdots|sa_r)
 \bigr)
 :=
 p_{\mathbf1}(m)\,
 \bigl(
   s q_\gamma^{\rm sq}(a_1)|\cdots|
   s q_\gamma^{\rm sq}(a_r)
 \bigr)}
\end{equation}
defines a filtered chain map
\[
 \mathfrak q_\gamma^{\rm bar}:
 \mathbb{DBar}_\gamma^{\Mot}
 \longrightarrow
 \mathbb B_\gamma^{\rm sq}.
\]
Hence it induces a morphism of bar-length spectral sequences.
\end{proposition}

\begin{proof}
The internal differentials commute with both
\(p_{\mathbf1}\) and \(q_\gamma^{\rm sq}\).
For an endpoint-action term with a positive bar letter,
\cref{P2M-lem:v25-unit-coefficient-firewall} gives zero after
\(p_{\mathbf1}\), exactly matching the trivial endpoint action in
\(\Z_\epsilon\).  Adjacent multiplication commutes with the dg-algebra
quotient \(q_\gamma^{\rm sq}\).  Thus every normalized Hochschild
differential component is preserved.  Bar length is unchanged by the map.
\end{proof}

Now fix an oriented three-element support
\[
 I=\{i,j,k\},\qquad i<j<k.
\]
Let \(\mathscr O_{I}^{\rm sq,ex}\) be the exact-support-\(I\) part of the
squarefree operation algebra.

\begin{lemma}[The exact-support-three squarefree hexagon]
\label{P2M-lem:v26-three-support-hexagon}
In cohomological degrees \(0,-1,-2\), the exact-support-three squarefree
internal complex is
\begin{equation}
 \boxed{
 \Z K_{ijk}
 \xrightarrow{\ d\ }
 \Z^{6}\!
 \left\langle
   H_{ab}D_c,\,
   D_cH_{ab}
 \right\rangle_{\{a,b,c\}=I}
 \xrightarrow{\ d\ }
 \Z^{6}\!
 \left\langle
   D_{\sigma(i)}D_{\sigma(j)}D_{\sigma(k)}
 \right\rangle_{\sigma\in S_3}.}
\end{equation}
The second differential is the oriented incidence map of a hexagon.  Its
kernel is the primitive rank-one line
\begin{equation}
 \boxed{
 \ker d|_{-1}
 =
 \Z\,\Omega_{ijk},}
\end{equation}
where
\[
 \Omega_{ijk}
 =
 [D_i,H_{jk}]
 +[D_j,H_{ki}]
 +[D_k,H_{ij}]
\]
with the fixed support orientation.  Since
\[
 dK_{ijk}=N\Omega_{ijk},
\]
one has
\begin{equation}
 \boxed{
 H^{-1}\!\left(\mathscr O_I^{\rm sq,ex}\right)
 \cong
 \Z/N,
 \qquad
 [\Omega_{ijk}]\text{ a generator}.}
\end{equation}
\end{lemma}

\begin{proof}
Squarefreeness leaves exactly the six ordered singleton words in degree zero.
A degree-\(-1\) exact-support word must contain one pair filler and the
remaining singleton, in either order, giving the six displayed generators.
Their differentials are
\[
 d(H_{ab}D_c)
 =
 (D_aD_b-D_bD_a)D_c,
 \qquad
 d(D_cH_{ab})
 =
 D_c(D_aD_b-D_bD_a),
\]
so every degree-\(-1\) generator joins two adjacent permutations of
\((i,j,k)\).  The resulting graph is the Cayley hexagon of \(S_3\), whose
integral incidence kernel is one primitive cycle.  Expanding that oriented
cycle gives precisely the displayed \(\Omega_{ijk}\).
There is no other squarefree degree-\(-2\) exact-support generator besides
\(K_{ijk}\), and its differential is \(N\Omega_{ijk}\).  Quotienting the
primitive cycle line by this image gives \(\Z/N\).
\end{proof}

For the spectral-sequence calculation, write
\(\partial_0\) for the internal differential and
\(\partial_1\) for the length-lowering Hochschild composition differential.
The standard zigzag formula says that if
\[
 \partial_0x=0,
 \qquad
 \partial_1x=-\partial_0y,
\]
then
\begin{equation}\label{P2M-eq:v26-d2-zigzag-formula}
 d_2[x]=[\partial_1y].
\end{equation}

Let
\[
 \mathfrak J_{ijk}^{\rm bar}
 :=
 \operatorname{Alt}^{(3)}_{\rm Hoch}(D_i,D_j,D_k)
\]
be the normalized fully antisymmetrized length-three singleton bar chain,
using the standard suspended Koszul signs.  Define
\(Y_{ijk}^{\rm bar}\) to be the canonical oriented length-two correction
obtained from
\(\partial_1\mathfrak J_{ijk}^{\rm bar}\) by replacing each occurrence of
\([D_a,D_b]\) with its chosen pair filler \(H_{ab}\).

\begin{lemma}[Three-singleton bar zigzag]
\label{P2M-lem:v26-three-singleton-zigzag}
With the preceding orientation conventions,
\begin{equation}\label{P2M-eq:v26-three-singleton-first-zigzag}
 \boxed{
 \partial_1\mathfrak J_{ijk}^{\rm bar}
 =
 -\partial_0Y_{ijk}^{\rm bar},}
\end{equation}
and
\begin{equation}\label{P2M-eq:v26-three-singleton-second-zigzag}
 \boxed{
 \partial_1Y_{ijk}^{\rm bar}
 =
 \varepsilon_{ijk}\,
 s\Omega_{ijk},
 \qquad
 \varepsilon_{ijk}\in\{\pm1\}.}
\end{equation}
\end{lemma}

\begin{proof}
Expand the six permutations in the Hochschild antisymmetrizer.
At the first composition step, adjacent singleton multiplications group
pairwise into the three graded commutators
\([D_i,D_j]\), \([D_j,D_k]\), and \([D_k,D_i]\), each tensored on the
appropriate side with the remaining singleton.  Replacing each commutator by
\(dH_{ab}\) gives
\eqref{P2M-eq:v26-three-singleton-first-zigzag}.
Apply \(\partial_1\) once more.  The endpoint terms vanish in the augmentation
bar, and the remaining adjacent products are the oriented commutators
\([D_i,H_{jk}]\), \([D_j,H_{ki}]\), and \([D_k,H_{ij}]\).
Their common suspension sign is \(\varepsilon_{ijk}\), giving
\eqref{P2M-eq:v26-three-singleton-second-zigzag}.
\end{proof}

\begin{theorem}[First genuine secondary bar transgression]
\label{P2M-thm:v26-first-secondary-bar-transgression}
In the bar-length spectral sequence of
\(\mathbb B_\gamma^{\rm sq}\), the three-singleton class satisfies
\begin{equation}
 \boxed{
 d_2
 \bigl[
   \mathfrak J_{ijk}^{\rm bar}
 \bigr]
 =
 \varepsilon_{ijk}\,
 [s\Omega_{ijk}]
 \ne0.}
\end{equation}
The image has exact additive order \(N\).
Consequently the corresponding \(d_2\) in the full motivic diagonal-bar
spectral sequence is nonzero whenever the three-support block survives under
\(\mathfrak q_\gamma^{\rm bar}\).
\end{theorem}

\begin{proof}
The first two identities of
\cref{P2M-lem:v26-three-singleton-zigzag} give the spectral-sequence zigzag
\eqref{P2M-eq:v26-d2-zigzag-formula}.  By
\cref{P2M-lem:v26-three-support-hexagon},
\([s\Omega_{ijk}]\) is the suspended generator of the exact-order-\(N\)
internal homology line.  It cannot be a first-page bar boundary in the
augmentation bar: after projection to internal homology indecomposables in
exact support three and degree \(-1\), every product boundary vanishes,
whereas \([\Omega_{ijk}]\) generates the unique \(\Z/N\) line.
Thus the \(d_2\)-image is nonzero of exact order \(N\).
Naturality of the comparison map
\(\mathfrak q_\gamma^{\rm bar}\) implies that a zero \(d_2\) upstairs would
map to zero downstairs, so the motivic differential is nonzero whenever this
exact-support quotient is present.
\end{proof}

\begin{corollary}[Secondary Moore class is a bar-Massey transgression]
\label{P2M-cor:v26-moore-as-bar-massey}
At support three, the intrinsic Moore relation
\[
 dK_{ijk}=N\Omega_{ijk}
\]
has the following downstream motivic bar interpretation:
the pair fillers \(H_{ij},H_{jk},H_{ki}\) null-homotope the first composition
boundary of the three-singleton antisymmetrizer, and the resulting secondary
bar differential is the exact-order-\(N\) Moore obstruction
\([\Omega_{ijk}]\).
Thus the first genuinely secondary Moore line is recovered by native bar
geometry without using the left-regular fidelity factor.
\end{corollary}

\paragraph{Secondary versus primary commutators.}
Primary cyclicization applies to commutators of internal cycles; the pair
fillers \(H_{ab}\) are not internal cycles, so it does not kill
\([s\Omega_{ijk}]\).


\subsection{Integral Moore staircases and the canonical total cycle}


The support-three computation has a dimension-independent origin, but one must
state it with the correct integral normalization.  Starting at support four,
the raw singleton antisymmetrizer already has lower-support secondary
transgressions.  The class which reaches the top Moore obstruction is the
\emph{Moore-corrected} singleton class obtained by pulling back the normalized
Stasheff partition staircase.

For a nonempty block \(A\), put
\[
 \eta(A):=\eta(|A|)=\max\{|A|-2,0\}.
\]
This is the preceding Moore support weight rewritten in the present namespace;
it is unrelated to the Layer-II coefficient-depth vector $\nu$.
Let \(\kappa_A\) denote the coefficient-normalized Stasheff generator of the
latest intrinsic source and let
\[
 X_A:=
 \begin{cases}
 D_i,&A=\{i\},\\
 H_{ij},&A=\{i,j\},\\
 K_A,&|A|\ge3
 \end{cases}
\]
denote its integral FSC counterpart.  Upstream v0.39 supplies the injective
normalization map
\begin{equation}
 \boxed{
 \Phi_N(X_A)=N^{\eta(A)}\kappa_A.}
\end{equation}
Its differential is
\begin{equation}
 dX_J
 =
 \sum_{\{A,B\}\in\operatorname{Bip}(J)}
 \epsilon(A,B)
 N^{\eta(J)-\eta(A)-\eta(B)}
 [X_A,X_B].
\end{equation}

Fix an oriented support \(I\), \(|I|=m\ge3\).  For \(1\le b\le m\), let
\(\operatorname{Part}_b(I)\) be the ordered set partitions
\[
 \pi=(A_1|\cdots|A_b)
\]
of \(I\) into nonempty blocks.  Write \(\epsilon(\pi)\) for the Hall/Stasheff
orientation transported from the fixed order on \(I\), and put
\[
 \eta(\pi):=\sum_{r=1}^b\eta(A_r).
\]

\begin{definition}[Normalized partition-bar staircase]
\label{P2M-def:v27-normalized-partition-staircase}
In the augmentation bar of the normalized Stasheff controller define
\begin{equation}
 \boxed{
 \mathfrak Z^{\rm nor}_{I,b}
 :=
 \sum_{\pi=(A_1|\cdots|A_b)\in\operatorname{Part}_b(I)}
 \epsilon(\pi)\,
 (s\kappa_{A_1}|\cdots|s\kappa_{A_b}).}
\end{equation}
Thus
\[
 \mathfrak Z^{\rm nor}_{I,m}
 =
 \mathfrak J_I^{\rm bar}
\]
is the fully antisymmetrized singleton bar chain and
\[
 \mathfrak Z^{\rm nor}_{I,1}=s\kappa_I.
\]
Let \(\partial_0\) denote the internal Stasheff differential and
\(\partial_1\) the length-lowering augmentation-bar composition differential.
\end{definition}

\begin{lemma}[Partition staircase identities]
\label{P2M-lem:v27-normalized-staircase-identities}
For \(2\le b\le m\),
\begin{equation}
 \boxed{
 \partial_1\mathfrak Z^{\rm nor}_{I,b}
 =
 -\,\partial_0\mathfrak Z^{\rm nor}_{I,b-1}.}
\end{equation}
\end{lemma}

\begin{proof}
A term of \(\partial_1\mathfrak Z^{\rm nor}_{I,b}\) merges two adjacent
blocks of an ordered \(b\)-partition.  The same oriented incidence appears in
\(\partial_0\mathfrak Z^{\rm nor}_{I,b-1}\) by applying the normalized
Stasheff differential to the merged block and splitting it back into those
two blocks.  The bar suspension and Hall orientation give opposite incidence
signs.  Every adjacent refinement/coarsening pair occurs exactly once on each
side.  This is precisely the bar form of the normalized twisting-cochain
Maurer--Cartan equation.
\end{proof}

The point is now integral.  For \(b\ge2\), every ordered \(b\)-partition
satisfies
\begin{equation}\label{P2M-eq:v27-partition-weight-bound}
 \boxed{
 \eta(\pi)\le m-3.}
\end{equation}
Indeed, if \(q\) blocks are nonsingletons then
\[
 \eta(\pi)=m-b-q,
\]
unless all blocks are singletons, in which case it is zero.

\begin{definition}[Integral Moore-corrected partition chains]
\label{P2M-def:v27-integral-partition-chains}
For \(2\le b\le m\), define
\begin{equation}\label{P2M-eq:v27-integral-partition-chain}
 \boxed{
 \mathfrak Z^{\rm int}_{I,b}
 :=
 \sum_{\pi=(A_1|\cdots|A_b)}
 \epsilon(\pi)\,
 N^{m-3-\eta(\pi)}
 (sX_{A_1}|\cdots|sX_{A_b}).}
\end{equation}
All exponents are nonnegative by
\eqref{P2M-eq:v27-partition-weight-bound}.  The induced bar normalization satisfies
\begin{equation}
 \boxed{
 \Phi_N^{\rm bar}
 \bigl(\mathfrak Z^{\rm int}_{I,b}\bigr)
 =
 N^{m-3}\mathfrak Z^{\rm nor}_{I,b}.}
\end{equation}
\end{definition}

\begin{proposition}[Integral staircase and terminal saturation defect]
\label{P2M-prop:v27-integral-staircase}
For \(3\le b\le m\),
\begin{equation}
 \boxed{
 \partial_1\mathfrak Z^{\rm int}_{I,b}
 =
 -\,\partial_0\mathfrak Z^{\rm int}_{I,b-1}.}
\end{equation}
At the final step,
\begin{equation}\label{P2M-eq:v27-integral-terminal-step}
 \boxed{
 \partial_1\mathfrak Z^{\rm int}_{I,2}
 =
 \varepsilon_I\,s\Omega_I^{\rm rec},
 \qquad
 \varepsilon_I\in\{\pm1\}.}
\end{equation}
Moreover
\begin{equation}\label{P2M-eq:v27-top-singleton-prefactor}
 \boxed{
 \mathfrak Z^{\rm int}_{I,m}
 =
 N^{m-3}\mathfrak J_I^{\rm bar}.}
\end{equation}
\end{proposition}

\begin{proof}
For \(b\ge3\), apply the injective bar extension of \(\Phi_N\) and use
\cref{P2M-lem:v27-normalized-staircase-identities}; both sides map to
\(N^{m-3}\) times the same normalized incidence chain.

For \(b=2\), the normalized identity gives
\[
 \Phi_N^{\rm bar}
 \bigl(\partial_1\mathfrak Z^{\rm int}_{I,2}\bigr)
 =
 \pm N^{m-3}s\,d\kappa_I.
\]
On the other hand,
\[
 \Phi_N(K_I)=N^{m-2}\kappa_I,
 \qquad
 dK_I=N\Omega_I^{\rm rec},
\]
so
\[
 \Phi_N(\Omega_I^{\rm rec})
 =
 N^{m-3}d\kappa_I.
\]
Injectivity gives \eqref{P2M-eq:v27-integral-terminal-step}.  Finally every block
of the unique singleton partition has Moore weight zero, so
\eqref{P2M-eq:v27-integral-partition-chain} gives
\eqref{P2M-eq:v27-top-singleton-prefactor}.
\end{proof}

The two-block correction makes the transfer profile visible without any
further calculation.  If \(I=A\sqcup B\), then
\[
 m-3-\eta(A)-\eta(B)
 =
 \begin{cases}
 0,&\min(|A|,|B|)=1,\\
 1,&|A|,|B|\ge2.
 \end{cases}
\]
Thus \(\mathfrak Z^{\rm int}_{I,2}\) contains singleton--highest-face terms
with coefficient one and genuine nonsingleton--nonsingleton terms with
coefficient \(N\), exactly the normalized latching pattern of the motivic
channel realization.

\begin{definition}[Canonical integral Moore staircase]
\label{P2M-def:v27-moore-staircase-subcomplex}
For \(2\le b\le m\) put
\[
 E_b:=\mathfrak Z^{\rm int}_{I,b},
 \qquad
 F_b:=\partial_0E_b,
\]
set
\[
 F_1:=s\Omega_I^{\rm rec},
 \qquad
 G_1:=sK_I,
\]
and let
\(\mathbb S_I^{\rm Moore,bar}\) be the filtered subcomplex of the squarefree
augmentation bar generated by
\[
 G_1,\ F_1,\ E_2,F_2,\ldots,E_m,F_m.
\]
Here \(F_m=0\) because the singleton generators are internally closed.
\end{definition}

\begin{lemma}[Staircase subcomplex is closed]
\label{P2M-lem:v27-moore-staircase-closed}
The displayed span is stable under the total bar differential.  More
precisely,
\[
 dG_1=NF_1,\qquad dF_b=0,
\]
and, up to the fixed orientation signs,
\[
 dE_b=F_b-F_{b-1}
 \qquad(2\le b\le m).
\]
\end{lemma}

\begin{proof}
Use \(dK_I=N\Omega_I^{\rm rec}\),
\(\partial_0^2=0\), the anticommutation
\(\partial_0\partial_1+\partial_1\partial_0=0\), and
\cref{P2M-prop:v27-integral-staircase}.  The formula for \(dE_b\) is just the sum
of its internal and composition components.
\end{proof}

\begin{theorem}[All-support Moore-staircase differential]
\label{P2M-thm:v27-all-support-moore-staircase}
The bar-length spectral sequence of
\(\mathbb S_I^{\rm Moore,bar}\) has
\[
 d_r\bigl[N^{m-3}\mathfrak J_I^{\rm bar}\bigr]=0
 \qquad(1\le r<m-1),
\]
and its first nonzero outgoing differential is
\begin{equation}
 \boxed{
 d_{m-1}
 \bigl[
   N^{m-3}\mathfrak J_I^{\rm bar}
 \bigr]
 =
 \varepsilon_I
 [s\Omega_I^{\rm rec}].}
\end{equation}
The target is a cyclic line of exact order \(N\).
\end{theorem}

\begin{proof}
The successive corrections
\[
 E_{m-1},E_{m-2},\ldots,E_2
\]
give the standard length-filtration zigzag using
\cref{P2M-prop:v27-integral-staircase}.  Hence all earlier outgoing
differentials vanish and the \(d_{m-1}\)-representative is
\(\partial_1E_2=\varepsilon_Is\Omega_I^{\rm rec}\).

Inside the staircase the only length-one antecedent for \(F_1\) is \(G_1\),
with \(dG_1=NF_1\); therefore the target line is \(\Z/N\).  Equivalently,
the exact-order statement is certified by the motivic Moore--Reedy
realization of \cref{P2M-thm:channel-complete-realization}, which sends
\(\Omega_I^{\rm rec}\) to an exact-order-\(N\) highest-face detector class.
\end{proof}

\begin{corollary}[Rees budget of the integral staircase]
\label{P2M-cor:v27-bar-transfer-exponent}
At support \(m\ge3\), the corrected partition-bar staircase uses
\(N^{m-3}\) before the terminal Moore step, and
\(dK_I=N\Omega_I\) supplies one final factor:
\begin{equation}
 \boxed{N^{m-3}\cdot N=N^{m-2}.}
\end{equation}
This is the same exponent
\(\eta(m)=m-2\) carried by every pointed weight-zero Hall operation tree.
The equality is a Rees-weight identity; the separate bar-length value \(m\)
comes from the cobar conilpotence of \(e_I\).
\end{corollary}

\begin{proof}
The first factor is
\eqref{P2M-eq:v27-top-singleton-prefactor}, the terminal factor is
\(dK_I=N\Omega_I\), and
\(\eta(m)=m-2\) is \cref{P2M-prop:eta-unique-primitive}.  The Hall realization
uses the same Rees exponent by
\cref{P2M-cor:v14-weight-zero-common-transfer}.
\end{proof}

\begin{corollary}[Explicit tetrahedral corrected staircase]
\label{P2M-cor:v27-tetrahedral-staircase}
For \(I=\{1,2,3,4\}\), the top class is
\[
 E_4=N\mathfrak J_{1234}^{\rm bar}.
\]
The length-three correction is \(N\) times the oriented sum of the
\(H_{ab}|D_c|D_d\) partition tensors.  The length-two correction has the form
\begin{equation}
 \boxed{
 E_2
 =
 \sum_{i=1}^4
  \varepsilon_i
  \bigl(
    sD_i|sK_{I\setminus\{i\}}
    \ \pm\
    sK_{I\setminus\{i\}}|sD_i
  \bigr)
 +
 N\!\!\sum_{\substack{I=A\sqcup B\\|A|=|B|=2}}
  \varepsilon_{A,B}
  \bigl(
    sH_A|sH_B
    \ \pm\
    sH_B|sH_A
  \bigr),}
\end{equation}
with the Hall/Koszul orientation signs.  Consequently the canonical
tetrahedral staircase has
\[
 \boxed{
 d_3[N\mathfrak J_{1234}^{\rm bar}]
 =
 \pm[s\Omega_{1234}^{\rm rec}]}
\]
and the target line has order \(N\).
\end{corollary}

\paragraph{Corrected singleton class.}
For \(m\ge4\) the raw singleton antisymmetrizer has proper-support
transgressions; the all-support class is the Moore-corrected
\(N^{m-3}\mathfrak J_I^{\rm bar}\).


\paragraph{Exact-multilinear bar degree and Moore lifts.}


Give every structural generator \(X_A\) its support multidegree
\[
 \deg_{\rm supp}(X_A)=\mathbf1_A\in\N^I
\]
and every history generator an additional history-count degree one.  Both
degrees are additive under products and bar tensors.

\begin{lemma}[Exact-multilinear history-zero bar summand]
\label{P2M-lem:v28-exact-multilinear-summand}
The total squarefree augmentation bar decomposes by total support multidegree
and history count.  In particular the component
\begin{equation}
 \boxed{
 \mathbb B_I^{\rm ml}
 :=
 \left(
 \mathbb B_\gamma^{\rm sq}
 \right)_{\deg_{\rm supp}=\mathbf1_I,\ \deg_{\rm hist}=0}}
\end{equation}
is a filtered dg direct summand.  The entire Moore staircase of the preceding staircase construction lies in
\(\mathbb B_I^{\rm ml}\).
\end{lemma}

\begin{proof}
Every structural differential replaces one support \(A\) by products whose
factor supports partition \(A\), so total support multidegree is preserved.
History generators are closed and no structural differential creates a
history generator.  Hochschild composition only concatenates adjacent words
and therefore preserves both gradings.  The direct-sum decomposition follows.
\end{proof}

\begin{definition}[Integral Moore lift]
\label{P2M-def:v28-integral-moore-lift}
Here the bar complex is the normalized bar of the associative envelope
$\mathscr U^{\rm rec}_{I,N}$; the internal multilinear Lie controller enters
through its canonical commutator inclusion.  An \emph{integral Moore lift of
\(K_I\)} is a total cycle
\[
 \Theta\in Z^*\!\left(\mathbb B_I^{\rm ml}\right)
\]
whose length-one indecomposable component contains \(sK_I\) with coefficient
\(1\) and contains no other exact-support-\(I\) indecomposable generator.  Write
\[
 \ell_{\rm bar}(\Theta)
\]
for its maximal nonzero bar length.
\end{definition}

For \(I=\{i,j\}\), the familiar pair lift is
\[
 sH_{ij}
 +\text{the oriented two-singleton bar chain},
\]
so its maximal length is two and its top singleton coefficient is one.


\begin{proposition}[Bar--cobar weight and the unique top singleton line]
\label{P2M-prop:v31-barcobar-weight}
Let \(I\) have cardinality \(m\), and work first in the untwisted associative
envelope \(\Cobar(\overline\Lambda_I)\).  The iterated reduced coproduct of the
exact-support generator satisfies
\begin{equation}\label{P2M-eq:v31-iterated-reduced-coproduct}
 \overline\Delta^{(m-1)}(e_I)
 =
 \sum_{\sigma\in S_I}
 \varepsilon(\sigma)\,
 e_{\sigma(i_1)}\otimes\cdots\otimes e_{\sigma(i_m)}
 \ne0,
 \qquad
 \overline\Delta^{(m)}(e_I)=0.
\end{equation}
Hence the bar--cobar unit
\[
 \eta_{\rm bc}:
 \overline\Lambda_I
 \longrightarrow
 B\Cobar(\overline\Lambda_I)
\]
has, on \(e_I\), maximal bar length exactly \(m\), and its top component is
the primitive antisymmetric singleton line
\[
 \pm\mathfrak J_I^{\rm bar}.
\]
Moreover any exact-multilinear bar cycle whose bar--cobar counit is
\(e_I\) has maximal bar length at least \(m\); if the maximal length is
exactly \(m\), its top singleton coefficient is uniquely \(\pm1\), with the
sign fixed by the chosen orientation of \(I\).
\end{proposition}

\begin{proof}
Formula \eqref{P2M-eq:v31-iterated-reduced-coproduct} is the ordered-partition
expansion of the reduced exterior coproduct.  It is nonzero at depth
\(m-1\) because the singleton tensors are distinct basis vectors, and one
more reduced coproduct is impossible.  The classical bar--cobar unit is the
sum of these iterated reduced coproducts, with the standard suspensions; see
\cite{Adams1956,HMS1974,LodayVallette2012}.  Thus its top component is the
oriented all-singleton antisymmetrizer.

For the minimality statement, filter \(B\Cobar(\overline\Lambda_I)\) by bar
length.  In exact support \(I\), bar length \(m\) has only singleton
letters, so its internal differential is zero, and there is no
exact-multilinear bar length \(m+1\) from which a composition boundary could
enter.  Hence the top singleton orientation coefficient is invariant under
changing the representative of the bar--cobar class.  A representative of
length \(<m\) would have zero image under the \((m-1)\)-fold reduced
coproduct detector, contradicting
\eqref{P2M-eq:v31-iterated-reduced-coproduct}.  The exact-multilinear top singleton orientation line in this associative
envelope is rank one, so its coefficient is forced to be \(\pm1\).  This is
the rank-one statement used by \cref{P2M-def:v28-integral-moore-lift}; no separate
rank-one claim for the preceding free Lie complex is needed.
\end{proof}

\begin{lemma}[Highest-face branch extraction]
\label{P2M-lem:v28-highest-face-branch}
Let \(I\) have cardinality \(m\ge3\), fix \(i\in I\), and put
\(J=I\setminus\{i\}\).  If \(\Theta\) is an integral Moore lift of \(K_I\),
then its length-two component contains the oriented
singleton/complement branch
\begin{equation}
 \boxed{
 \pm N\,(sD_i|sK_J)
 \ \pm\
 N\,(sK_J|sD_i)}
\end{equation}
with the Hall/Koszul signs forced by the \(I\)-orientation.  After fixing the
left or right \(D_i\)-slot and factoring it off, the higher-length terms
needed to cancel the internal differential of this branch form \(N\) times an
integral Moore lift of \(K_J\).
\end{lemma}

\begin{proof}
The bar-length-one part of the cycle equation is
\[
 \partial_0(sK_I)+\partial_1(\Theta_2)=0,
\]
hence
\[
 \partial_1(\Theta_2)=-N\,s\Omega_I.
\]
In the full-support curvature, the proper cell \(K_J\) occurs only in the
singleton/complement term associated with \(\{i\}\sqcup J\).  Upstream v0.39
records this as the universal highest-face formula
\[
 \pi_{J,K}\partial(\Omega_I)
 =
 \pm D_i\otimes sK_J,
\]
and proves that no genuine nonsingleton--nonsingleton partition contributes
to the same \(K_J\)-coordinate.  In the free associative word basis the only
bar products which yield \(D_iK_J\) or \(K_JD_i\) are therefore the two
displayed singleton/complement tensors, forcing coefficient \(N\).

Now apply the internal differential.  Since
\[
 dK_J=N\Omega_J,
\]
the forced branch produces \(N^2\) times the corresponding
\(D_i|\Omega_J\) or \(\Omega_J|D_i\) term.  Exact multidegree prevents any
support outside \(J\) from entering that branch.  Factoring off the fixed
\(D_i\)-slot, the remaining cycle equations are exactly the defining
equations for a Moore lift of \(K_J\), scaled by the already forced outer
factor \(N\).  Finally, the \(S_I\)-equivariance of the channel realization
transports this forced branch through all choices of \(i\) with the Hall
orientation signs; consequently the extremal all-singleton contributions
assemble on the single antisymmetric line
\(\mathfrak J_I^{\rm bar}\), rather than on unrelated branch coordinates.
\end{proof}

\begin{theorem}[Minimal bar--cobar/Rees lift theorem]
\label{P2M-thm:v28-minimal-moore-lift}
Let \(I\) be an ordered support of cardinality \(m\ge2\).  Every integral
Moore lift \(\Theta\) of \(K_I\) satisfies
\begin{equation}
 \boxed{\ell_{\rm bar}(\Theta)\ge m.}
\end{equation}
This lower bound is independent of \(N\): it remains true after every base
change in which \(N\) is invertible.

If equality holds, then the top antisymmetric singleton coefficient is
uniquely determined:
\begin{equation}
 \boxed{
 \operatorname{coeff}_{\mathfrak J_I^{\rm bar}}(\Theta)
 =
 \pm N^{m-2},}
\end{equation}
where the sign is the chosen Hall/bar orientation.  Both statements are
sharp.
\end{theorem}

\begin{proof}
Localize at \(N\) and use
\cref{P2M-prop:cobar-rees-normal-form}.  A lift whose length-one component is
\(X_I\) with coefficient \(1\) becomes a bar cycle for the untwisted cobar
algebra whose counit is
\[
 N^{\eta(I)}e_I=N^{m-2}e_I.
\]
By \cref{P2M-prop:v31-barcobar-weight}, no representative of this nonzero
bar--cobar class can have maximal length \(<m\).  Since localization does not
increase bar length, the same lower bound holds integrally.

If the maximal length is exactly \(m\), the top exact-multilinear length-\(m\)
line is rank one and \cref{P2M-prop:v31-barcobar-weight} forces coefficient
\(\pm1\) for a lift of \(e_I\).  Multiplying the counit class by
\(N^{\eta(I)}\) therefore forces the unique coefficient
\[
 \pm N^{\eta(I)}
 =
 \pm N^{m-2}
\]
for a lift of \(X_I\).  No additional integral solution can change this
coefficient: after localization it would change the unique rank-one top
associated-graded component.  The canonical cycle below attains the bound.
\end{proof}

\paragraph{The canonical total Moore bar cycle.}


\begin{definition}[Canonical total Moore bar cycle]
\label{P2M-def:v28-total-moore-bar-cycle}
For \(|I|=m\ge3\), define
\begin{equation}
 \boxed{
 \Theta_I^{\rm bar}
 :=
 sK_I
 +
 N\sum_{b=2}^{m}
 \mathfrak Z_{I,b}^{\rm int}.}
\end{equation}
For \(m=2\), use the pair lift.
\end{definition}

\begin{theorem}[Total-cycle identity and sharp transfer coefficient]
\label{P2M-thm:v28-total-moore-cycle}
The element \(\Theta_I^{\rm bar}\) is a total cycle in
\(\mathbb B_I^{\rm ml}\):
\begin{equation}
 \boxed{d\Theta_I^{\rm bar}=0.}
\end{equation}
Its maximal bar length is \(m\), and its top singleton component is
\begin{equation}
 \boxed{
 \left(\Theta_I^{\rm bar}\right)_{\rm top}
 =
 N^{m-2}\mathfrak J_I^{\rm bar}.}
\end{equation}
Thus it attains both sharp bounds of
\cref{P2M-thm:v28-minimal-moore-lift}.
\end{theorem}

\begin{proof}
Use the notation of
\cref{P2M-def:v27-moore-staircase-subcomplex}.  There
\[
 d(sK_I)=NF_1,
 \qquad
 dE_b=F_b-F_{b-1},
 \qquad
 F_m=0.
\]
Therefore
\[
 d\Theta_I^{\rm bar}
 =
 NF_1+
 N\sum_{b=2}^{m}(F_b-F_{b-1})
 =
 NF_m=0.
\]
Finally
\[
 N E_m
 =
 N\cdot N^{m-3}\mathfrak J_I^{\rm bar}
 =
 N^{m-2}\mathfrak J_I^{\rm bar}.
\]
\end{proof}

\begin{corollary}[Transfer exponent is the Rees top coefficient]
\label{P2M-cor:v28-transfer-minimal-saturation}
For a minimal integral Moore lift on support \(|I|=m\), the top singleton
coefficient is not merely bounded below: it is identically
\[
 \boxed{\pm N^{m-2}.}
\]
Equivalently,
\[
 \boxed{
 \text{cyclotomic transfer exponent}
 =
 \eta(m)
 =
 m-2.}
\]
The exponent is the Rees weight of the one-block generator; the separate
length statement \(\ell_{\rm bar}=m\) is the coefficient-free cobar weight.
\end{corollary}

\subsection{No-short-cross-talk and full-bar promotion}


\begin{theorem}[No short incoming bar differential to the Moore line]
\label{P2M-thm:v28-no-short-cross-talk}
Let \(|I|=m\ge3\).  In the full exact-multilinear history-zero squarefree
augmentation-bar spectral sequence, the class
\[
 [s\Omega_I^{\rm rec}]
\]
survives to the \(E_{m-1}\)-page.  Equivalently, it is not hit by any
\(d_r\) with \(r<m-1\).
\end{theorem}

\begin{proof}
Suppose \(s\Omega_I^{\rm rec}\) were hit on some page \(r<m-1\).
A defining system for an incoming differential to filtration one may be
chosen with source components in filtrations \(2,\ldots,r+1\); in particular
it has no length-one indecomposable component.  Thus there is a bar chain
\(Z\) of maximal length at most \(r+1\le m-1\), with zero length-one
indecomposable projection, such that
\[
 dZ=s\Omega_I^{\rm rec}.
\]
This choice is also forced by exact multidegree: in history degree zero the
only length-one indecomposable of exact support \(I\) is \(K_I\), whereas
\(dZ=s\Omega_I^{\rm rec}\) has only decomposable proper-support products in
length one.

Consequently
\[
 sK_I-NZ
\]
is a total cycle whose length-one indecomposable component is exactly
\(sK_I\) with coefficient one and whose maximal bar length is at most \(m-1\).  It is therefore an integral
Moore lift of \(K_I\), contradicting
\cref{P2M-thm:v28-minimal-moore-lift}.
\end{proof}

\begin{theorem}[Full squarefree all-support secondary transgression]
\label{P2M-thm:v28-full-squarefree-secondary-transgression}
Let \(|I|=m\ge3\).  The Moore-corrected class with leading term
\[
 N^{m-3}\mathfrak J_I^{\rm bar}
\]
is a well-defined nonzero class on the \(E_{m-1}\)-page of the
\emph{full} exact-multilinear squarefree augmentation-bar spectral sequence,
and
\begin{equation}
 \boxed{
 d_{m-1}
 \bigl[
   N^{m-3}\mathfrak J_I^{\rm bar}
 \bigr]
 =
 \varepsilon_I
 [s\Omega_I^{\rm rec}]
 \ne0.}
\end{equation}
The image has exact additive order \(N\).
\end{theorem}

\begin{proof}
The defining chains
\[
 \mathfrak Z_{I,m-1}^{\rm int},
 \ldots,
 \mathfrak Z_{I,2}^{\rm int}
\]
of the preceding staircase construction are actual chains in the full bar, so the staircase identities prove
that the displayed leading term survives every earlier outgoing
\(d_r\), \(r<m-1\), and that its \(d_{m-1}\)-representative is
\(\varepsilon_Is\Omega_I^{\rm rec}\).

By \cref{P2M-thm:v28-no-short-cross-talk}, that target has not been killed by any
earlier incoming differential, so the displayed \(d_{m-1}\) is genuinely
nonzero in the full spectral sequence.  Its order divides \(N\) because
\(dK_I=N\Omega_I\).  It is exactly \(N\): the channel-complete motivic
realization sends \([\Omega_I]\) to the highest-face exact-order-\(N\)
Moore detector, so no smaller positive multiple can vanish already in the
source internal homology or on the surviving target line.
\end{proof}

\begin{corollary}[Full motivic diagonal-bar all-support transgression]
\label{P2M-cor:v28-full-motivic-secondary-transgression}
The coefficient-unit copy of the same Moore-corrected defining system in the
full motivic diagonal bar satisfies
\begin{equation}
 \boxed{
 d_{m-1}^{\Mot}
 \bigl[
   N^{m-3}\mathfrak J_I^{\rm bar}
 \bigr]
 =
 \varepsilon_I
 [1_{\mathscr A}\otimes s\Omega_I^{\rm rec}]
 \ne0,}
\end{equation}
and its image has exact order \(N\).
\end{corollary}

\begin{proof}
Apply the filtered comparison
\(\mathfrak q_\gamma^{\rm bar}\) of
\cref{P2M-prop:v26-motivic-squarefree-bar-comparison}.  It sends the motivic
coefficient-unit defining system to the full squarefree one and the target to
\([s\Omega_I]\).  If the motivic differential vanished, its squarefree image
would vanish, contradicting
\cref{P2M-thm:v28-full-squarefree-secondary-transgression}.  The Moore relation
gives \(N\)-torsion upstairs, while the squarefree image has exact order \(N\),
so the motivic target has exact order \(N\) as well.
\end{proof}

\begin{corollary}[Cobar/Rees separation of the three indices]
\label{P2M-cor:v28-closed-bar-transfer-triangle}
For every support cardinality \(m\ge3\),
\[
 \boxed{
 \begin{array}{c}
 \text{first full diagonal-bar Moore page}=m-1,\\[1mm]
 \text{minimal total Moore-lift bar length}=m,\\[1mm]
 \text{Rees/transfer exponent}=\eta(m)=m-2.
 \end{array}}
\]
The first two numbers come from the bar filtration and the conilpotence
length of the exterior-coalgebra generator \(e_I\).  The third is logically
different: it is the unique coboundary primitive of the imported
\(1,N,N^2\) coefficient profile.  Their numerical relation is a compatibility
between cobar weight and Rees normalization, not a derivation of all three
from one recurrence.
\end{corollary}

\begin{corollary}[Explicit mixed history transgression]
\label{P2M-cor:v24-mixed-history-transgression}
Let \(h\) be a reduced frontier-history label with pair support
\(\supp(h)\), and choose \(i\notin\supp(h)\).  Then
\(D_i\) and \(\delta_h\) are closed and have disjoint supports.  Their
squarefree graded commutator
\[
 [D_i,\delta_h]_{\rm gr}
 =
 D_i\delta_h-
 (-1)^{|D_i||\delta_h|}\delta_hD_i
\]
is a nonzero operation symbol on the symbol page, but
\begin{equation}
 \boxed{
 [\,1_{\mathscr A}\otimes
 s[D_i,\delta_h]_{\rm gr}\,]
 =0
 \quad\text{in }E_2^{1,*}(\gamma).}
\end{equation}
Thus the first composition differential already detects a genuine distinction
between chain-level word separation and homological operation-history
fidelity.
\end{corollary}

\begin{proof}
Squarefree nonvanishing follows from the free marked operation envelope:
neither ordering repeats a support label, and the two ordered words are
distinct basis words.  Since \(D_i\) and \(\delta_h\) are internal cycles,
\cref{P2M-thm:v24-first-composition-cyclicization} makes their graded commutator
a \(d_1\)-boundary.
\end{proof}

\paragraph{Filtration firewall.}
The preceding Moore-weight spectral sequence has \(E_2=0\); the downstream
bar-length filtration has a permanent primitive summand.


\begin{corollary}[Why multiplicative strictification is genuinely stronger]
\label{P2M-cor:v24-regular-strictification-necessary}
The native diagonal bar supplies canonical operation symbols, but its first
composition differential necessarily forgets graded-commutator information.
Therefore any downstream target required to remain faithful on the full
noncommutative marked operation algebra needs additional multiplicative
structure beyond total diagonal-bar homology.  The left regular action of
The left-regular construction supplies exactly such a strictification on the squarefree sector.
\end{corollary}


\subsection{Strictification and operation fidelity}

The diagonal bar gives native geometric \emph{separation} of the squarefree
operation ledger, but \cref{P2M-thm:v24-first-composition-cyclicization} shows that
its first composition differential already kills graded commutators.  Thus if
one wants a faithful \emph{multiplicative} action of the full noncommutative
squarefree operation algebra, the following strictification is not merely
cosmetic: left regular multiplication restores information which diagonal-bar
homology is designed to cyclicize.  Because the target category is $\Z$-linear and stable, the
finite free dg lattice may be tensored with the integral motivic unit
$\mathbf1^{\Mot}$, normalized by
$\operatorname{End}(\mathbf1^{\Mot})\cong\Z$.  The object below is therefore not a new
geometric carrier: by \cref{P2M-thm:v23-native-bar-symbol-fidelity} its underlying
dg motivic lattice is canonically the diagonal-bar symbol carrier.  The only
additional datum is its left regular endomorphism action.

\begin{definition}[Squarefree motivic operation carrier]
\label{P2M-def:v22-squarefree-motivic-operation-carrier}
Define
\begin{equation}
 \boxed{
 \mathbb T_{\gamma}^{\Mot,\rm sq}
 :=\mathbf 1^{\Mot}\otimes_{\Z}\mathscr O_{\gamma}^{\rm sq}.}
\end{equation}
The differential is the tensor differential.  Left multiplication on the
second factor defines
\begin{equation}
 \boxed{
 \lambda_{\gamma}^{\Mot,\rm sq}:
 \mathscr O_{\gamma}^{\rm sq}
 \longrightarrow
 \REnd\!\left(\mathbb T_{\gamma}^{\Mot,\rm sq}\right).}
\end{equation}
\end{definition}

\begin{corollary}[Regular carrier is a multiplicative bar-symbol strictification]
\label{P2M-cor:v23-regular-is-bar-strictification}
There is a canonical isomorphism of dg motivic objects
\[
 \boxed{
 \Phi_\gamma:
 \mathbb T_\gamma^{\Mot,\rm sq}
 \xrightarrow{\ \sim\ }
 \mathbb S_\gamma^{\Mot,\rm sq}}
\]
given by \(\Phi_\gamma=\sigma_\gamma^{\Mot,\rm sq}\) on the operation-lattice
factor.  Under this identification the left regular action
\(\lambda_\gamma^{\Mot,\rm sq}\) is additional multiplicative structure on the
bar-symbol object; it is not needed merely to separate marked words.
\end{corollary}

\begin{theorem}[Finite squarefree multiplicative local-faithfulness theorem]
\label{P2M-thm:v22-squarefree-local-faithfulness}
For every finite reduced history $\gamma$, the map
$\lambda_{\gamma}^{\Mot,\rm sq}$ is an injective dg-algebra map.  Hence every
nonzero squarefree marked word, including every squarefree mixed
history/structural word, survives in the motivic operation carrier.
\end{theorem}

\begin{proof}
The left regular action is a dg action because the differential on
$\mathscr O_{\gamma}^{\rm sq}$ is a derivation.  If
$\lambda_{\gamma}^{\Mot,\rm sq}(a)=0$, evaluate the regular operator on the unit
word $1\in\mathscr O_{\gamma}^{\rm sq}$.  The result is $a$.  Hence $a=0$.
The integral motivic unit has faithful scalar endomorphism line $\Z$, so the
copy indexed by the unit word maps to the copy indexed by $a$ with exactly the
integral coefficient of $a$.  No torsion property of the root Moore carrier is
used in this argument.
\end{proof}

\begin{corollary}[Mixed history/structural words survive]
\label{P2M-cor:v22-mixed-words-survive}
Let $x$ be a structural marked word and $\delta_h$ a history generator such
that the supports of all factors are pairwise disjoint.  If
$x\delta_h\ne0$ and $\delta_hx\ne0$ in the marked source envelope, then
\[
 \lambda_{\gamma}^{\Mot,\rm sq}(x\delta_h)\ne0,
 \qquad
 \lambda_{\gamma}^{\Mot,\rm sq}(\delta_hx)\ne0.
\]
Moreover the two operators remain distinct whenever the two marked source words
are distinct.
\end{corollary}

\paragraph{Geometric and operation-fidelity factors.}

The multiplicatively strictified bar-symbol carrier is not used in place of the
geometric target.  It is added as a second block so that one projection retains
the genuine root-localization and Moore geometry while the other detects
operation words by the left regular action on a dg object already identified
with the native diagonal-bar symbol sector.

\begin{definition}[Operation-faithful motivic history apex]
\label{P2M-def:v22-operation-faithful-apex}
For a reduced history $\gamma$ define
\begin{equation}
 \boxed{
 \mathbb A_{\gamma}^{\Mot,\rm op}
 :=\mathbb A_{\gamma}^{\Mot,\MR}
   \oplus\mathbb T_{\gamma}^{\Mot,\rm sq}.}
\end{equation}
Let
$\widetilde\varrho_{\gamma}^{\Mot,\rm hist}:
 \mathscr O_{\gamma}^{\rm full}\to\mathscr A_{\gamma}^{\Mot,\MR}$
be the multiplicative extension of the history-marked generator realization.
Define
\begin{equation}
 \boxed{
 \varrho_{\gamma}^{\Mot,\rm op}(a)
 :=\widetilde\varrho_{\gamma}^{\Mot,\rm hist}(a)
   \oplus\lambda_{\gamma}^{\Mot,\rm sq}
          \bigl(q_{\gamma}^{\rm sq}(a)\bigr),
 \qquad a\in\mathscr O_{\gamma}^{\rm full}.}
\end{equation}
The first component is the geometric Moore--Reedy/history realization on the
full marked associative envelope; the second is the squarefree
operation-separated regular detector.
\end{definition}

\begin{theorem}[Geometric and operation fidelity coexist]
\label{P2M-thm:v22-geometric-operation-coexistence}
The map $\varrho_{\gamma}^{\Mot,\rm op}$ is injective on the free abelian
submodule of $\mathscr O_{\gamma}^{\rm full}$ spanned by squarefree marked
words.  Projection to the first factor recovers all
history-marked geometric conclusions: root localization, exact-support motivic Moore
packages, latching reconstruction, transfer--Rees weights and free frontier
history blocks.  Projection to the second factor is faithful on every
squarefree marked associative operation word.  In particular adding operation
fidelity introduces no duplicate structural Moore package.
\end{theorem}

\begin{proof}
On the squarefree-word submodule the quotient map
$q_{\gamma}^{\rm sq}$ is injective, so injectivity follows from the injective
regular second component of \cref{P2M-thm:v22-squarefree-local-faithfulness}.  The first projection is exactly
the history-marked apex and therefore preserves the already established geometric
realization.  The second summand is a detector block built from shifts of the integral
motivic unit; it carries no new structural generator $K_I$ and does not modify
the differential in the first summand.
\end{proof}

\paragraph{Reduced bicategorical compatibility.}

The preceding reduced theorem transports the closed history generators by oriented relabelling,
preserves them under mixed deletion--insertion exactly when their frontier
labels survive, and recomputes structural Moore completion once at the terminal
support.  These operations act functorially on marked words and hence on the
regular carrier.

\begin{definition}[Operation-faithful history cospans]
\label{P2M-def:v22-operation-history-bicategory}
Let $\mathbf{MotOpHistCosp}_{N}^{\MR,\rm sq}$ be the reduced marked motivic
cospan bicategory obtained from
$\mathbf{MotHistCosp}_{N}^{\MR}$ by adjoining to every reduced insertion apex
its carrier $\mathbb T_{\gamma}^{\Mot,\rm sq}$.  Horizontal composition first
uses the preceding reduced history normal form and then rebuilds the squarefree
operation carrier of that composite normal form.  It never pushes out
already-completed operation carriers.  Marked $2$-cells act by the induced
transport of generators and words.
\end{definition}

\begin{theorem}[Squarefree operation-faithful bicategorical enhancement]
\label{P2M-thm:v22-operation-faithful-bicategorical-enhancement}
For odd squarefree $N$, the history-marked motivic pseudofunctor admits a canonical
refinement
\begin{equation}
 \boxed{
 \mathcal R_N^{\Mot,\MR,\rm op}:
 \FSCComp^{\rm mk}
 \longrightarrow
 \mathbf{MotOpHistCosp}_{N}^{\MR,\rm sq}.}
\end{equation}
For every reduced insertion apex its dg decoration is faithful on the marked
squarefree operation envelope.  Forgetting the regular operation carrier
recovers $\mathcal R_N^{\Mot,\MR,\mathrm{hist}}$ exactly.
\end{theorem}

\begin{proof}
Apply the preceding structurally reduced normal form first.  Reordering and mixed
base-change maps act on the finite marked generating set and preserve support
intersection, so they induce dg maps of the squarefree operation envelopes and
therefore maps of the regular twisted complexes.  The preceding associators and
unitors identify the same composite reduced normal form; the induced word
reindexings satisfy the same pentagon and triangle relations.  Reduced
composition rebuilds the operation carrier only after the structural/history
normal form is known, so no structural Moore or operation detector factor is
duplicated.  Apexwise faithfulness is
\cref{P2M-thm:v22-squarefree-local-faithfulness}.
\end{proof}

\paragraph{Local-faithfulness scope.}
Faithfulness is on the marked squarefree operation algebra at each reduced
apex; local fullness of the whole motivic cospan bicategory is not claimed.


\paragraph{Regular and Hall fidelity.}
The regular action preserves associative words.  The Hall/root--Bass factor
retains the pointed Hall basis and adjacent-refinement differential.


\paragraph{Regular action after Ind-completion.}

Repeated-support words require the formal Ind-completion; the same
bar-symbol and regular construction then applies without the squarefree
quotient.

\begin{definition}[Full Ind-motivic bar-symbol and regular carrier]
\label{P2M-def:v22-ind-regular-carrier}
Let $\mathscr O_{\gamma}^{\infty}:=\mathscr O_{\gamma}^{\rm full}$ with no
repeated-support quotient.  In
$\operatorname{Ind}\langle \mathbf 1^{\Mot}\rangle$ put
\begin{equation}
 \boxed{
 \mathbb T_{\gamma}^{\Mot,\infty}
 :=\mathbf 1^{\Mot}\otimes_{\Z}\mathscr O_{\gamma}^{\infty}.}
\end{equation}
\end{definition}

\begin{corollary}[Full marked bar-symbol and multiplicative fidelity in the Ind-motivic envelope]
\label{P2M-cor:v22-ind-full-word-faithfulness}
The coefficient-unit length-one sector of the unrestricted diagonal bar gives
an injective chain-level symbol map on
\(\mathscr O_\gamma^\infty\), and left multiplication gives an injective
dg-algebra map
\[
 \mathscr O_{\gamma}^{\infty}
 \hookrightarrow
 \REnd\!\left(\mathbb T_{\gamma}^{\Mot,\infty}\right).
\]
Thus every finite marked associative word in structural and history generators,
including repeated-support words, has both a native diagonal-bar symbol and a
faithful multiplicative regular action after formal Ind-completion.
\end{corollary}

\begin{proof}
The unrestricted tensor envelope is a countably free dg lattice, so its tensor
with $\mathbf1^{\Mot}$ is an object of the Ind-completion.  The diagonal-bar
length-one calculation is identical to
\eqref{P2M-eq:v23-length-one-bar-differential}, and unit-word evaluation in the
regular model proves multiplicative injectivity exactly as in
\cref{P2M-thm:v22-squarefree-local-faithfulness}.
\end{proof}


\section{Sharpness and operation-history fidelity}

\subsection{Tetrahedral coefficient-framing obstruction}

The Three-\(pqr\) motivic target contains more than a Moore coefficient line:
its marked direction cells are tied to three actual coefficient roles.  That
framing is harmless at arity three but cannot itself serve as the direction labelling of a faithful all-support realization.

Let \(P=\{p,q,r\}\).  If \(J\) is a three-element recursive direction set, a
\emph{literal Three-\(pqr\) face framing} is a bijection
\[
 \lambda_J:J\xrightarrow{\sim}P
\]
used to identify the three recursive directions with the three coefficient
roles of the old pure-weight target.

\begin{definition}[Overlap-compatible tetrahedral framing]
\label{P2M-def:compatible-triad-framing}
Let \(I\) be a four-element direction set.  A family of literal face framings
\(\{\lambda_J\}_{|J|=3}\) is \emph{overlap compatible} if for every pair of
three-faces \(J,J'\subset I\),
\[
 \lambda_J|_{J\cap J'}=\lambda_{J'}|_{J\cap J'}.
\]
This is exactly the condition required if the framed face targets are to
restrict to the same coefficient-labelled object on every common edge while
the recursive direction labels themselves are held fixed.
\end{definition}

\begin{remark}[Tetrahedral coefficient-framing obstruction]
\label{P2M-thm:triad-framing-nogo}
No overlap-compatible tetrahedral framing exists.
\end{remark}

\begin{proof}
Suppose \(\{\lambda_J\}\) were compatible.  For a vertex \(i\in I\), choose
any three-face \(J\ni i\) and define \(\lambda(i):=\lambda_J(i)\).  Pairwise
overlap compatibility makes this independent of \(J\), so it defines a global
map
\[
 \lambda:I\longrightarrow P.
\]
Every three-element subset \(J\subset I\) must restrict bijectively to \(P\),
because \(\lambda|_J=\lambda_J\).  But a map from four elements to three
colours has two vertices of the same colour.  Any three-face containing that
pair fails to be injective, a contradiction.
\end{proof}

\begin{corollary}[Literal facewise reuse of the framed seed is not Reedy]
\label{P2M-cor:literal-face-reuse-nogo}
There is no support-labelled proper-face diagram on a four-element direction
set obtained by taking the old numerically framed Three-\(pqr\) target on every
three-face and identifying common lower faces by label-preserving restriction.
In particular a faithful realization of the higher FSCComp support tower cannot be produced merely by reusing the three coefficient roles locally.
\end{corollary}

\begin{proof}
Such a diagram would induce the compatible framings forbidden by
\cref{P2M-thm:triad-framing-nogo}.
\end{proof}

\paragraph{Permutation quotient.}
Relabelling also transports the numerical power weights, so quotienting
direction labels would erase support data rather than repair the framing
obstruction.


The conclusion is a realization constraint rather than a provenance obstruction: the intrinsic direction geometry is supplied by FSCComp, while the coefficient realization must remain separate before arity four.

\subsection{Finite-scalar realization no-go at the 3-primary place}

The coefficient-framing obstruction rules out literal reuse of three fixed
roles.  One might instead try to assign each recursive direction an arbitrary
power-correspondence scalar and each pair an arbitrary scalar multiple of the
edge loop.  For unbounded support this also fails once the detector has been
scalarized at the prime three.

\begin{definition}[Scalar triangular detector atlas]
\label{P2M-def:scalar-detector-atlas}
Let \(S\) be a finite direction set.  A scalar triangular detector atlas over
\(\F_3\) consists of vertex weights \(\lambda_i\in\F_3\) and symmetric edge
weights \(h_{ij}=h_{ji}\in\F_3\) such that the projected detector on a
three-support \(\{i,j,k\}\) is
\[
 \boxed{
 b_{ijk}
 :=\lambda_i h_{jk}+\lambda_j h_{ik}+\lambda_k h_{ij}.}
\]
This is the form obtained after a one-dimensional root line has forgotten all
non-scalar information in the direction and pair correspondences.
\end{definition}

\begin{theorem}[Sharp characteristic-three scalar-atlas bound]
\label{P2M-thm:finite-scalar-atlas-nogo}
Let \(S\) carry a scalar triangular detector atlas over \(\F_3\) with
\[
 b_{ijk}
e0
 \qquad\text{for every }\{i,j,k\}\subseteq S.
\]
Then
\begin{equation}
 \boxed{|S|\le22.}
\end{equation}
More sharply, on any subset \(T\subseteq S\) on which
\(\lambda_i=\lambda\in\F_3^\times\) is constant,
\begin{equation}
 \boxed{|T|\le10,}
\end{equation}
and the bound \(10\) is attained.
\end{theorem}

\begin{proof}
Assume first that \(\lambda_i=\lambda\ne0\) on \(T\).  Then
\[
 b_{ijk}
e0
 \quad\Longleftrightarrow\quad
 h_{ij}+h_{ik}+h_{jk}
e0
 \quad\text{in }\F_3.
\]
For three elements of \(\F_3\), the sum is zero exactly when they are either
all equal or all distinct.  Thus the edge-coloring
\[
 h:\binom{T}{2}\longrightarrow\F_3
\]
has neither a monochromatic triangle nor a rainbow triangle.  Equivalently it
is a Gallai three-coloring with no monochromatic \(K_3\).  The classical
Gallai--Ramsey value
\[
 \operatorname{gr}_3(K_3:K_3)=11
\]
therefore gives \(|T|\le10\); see
\cite{ChungGraham1983,GyarfasSimonyi2004}.

The bound is sharp.  Partition ten vertices into five pairs
\(V_0,\ldots,V_4\).  Color the edge inside every \(V_a\) by \(0\), color
edges between \(V_a,V_b\) by \(1\) when \(ab\) is an edge of a five-cycle,
and by \(2\) otherwise.  The two inter-pair color graphs are complementary
five-cycles, hence triangle-free; every triangle using two vertices from one
pair has colors \(0,c,c\).  Thus no triangle is monochromatic or rainbow.

For arbitrary vertex weights, the zero-weight class has size at most two,
since any three vertices with \(\lambda=0\) give \(b_{ijk}=0\).  Each of the
two nonzero weight classes has size at most ten by the preceding paragraph.
Hence \(|S|\le2+10+10=22\).
\end{proof}

\begin{corollary}[Non-scalar geometry is mandatory]
\label{P2M-cor:nonscalar-mandatory}
On a coefficient package containing the prime three, an all-support detector
cannot factor through one-dimensional scalar direction/pair data: every such
atlas has \(|S|\le22\).  Non-scalar correspondence data must therefore be
retained before the triangular Moore detector is taken.
\end{corollary}

\begin{warning}[Characteristic three is special here]
\label{P2M-warn:scalar-prime-three}
The sharp Gallai--Ramsey reduction uses
\(3\lambda h=0\) and is specific to the triangular detector in
characteristic three.  For \(\ell\ne3\), a monochromatic triangle does not
automatically annihilate
\(\lambda_i h_{jk}+\lambda_jh_{ik}+\lambda_kh_{ij}\); an
\(\ell\)-primary sharpening would require an appropriate zero-sum Ramsey
invariant rather than the ordinary Gallai--Ramsey number.
\end{warning}


\subsection{Operation-history fidelity and Hall completion}

Channels retain the quadratic latching polynomial but not arbitrary
rooted grafting history.  We compare them with the pointed Hall/root--Bass model constructed below.

Throughout this section fix an exact ordered support $I$, a distinguished label
$i_0\in I$, and the Hall orientation convention of the intrinsic source.  Write
$W_0K_I^{\mathrm{St}}$ for the normalized weight-zero pointed Hall complex and
$\mathscr H_{I,i_0}^{\mathrm{RB}}$ for its operation-separated root/Bass model.
The preceding theorem gives an integral chain isomorphism
\begin{equation}\label{P2M-eq:v14-preceding-hall-iso}
 \operatorname{ev}_{\mathrm{RB}}:
 W_0K_I^{\mathrm{St}}
 \xrightarrow{\ \sim\ }
 \mathscr H_{I,i_0}^{\mathrm{RB}}.
\end{equation}
A Hall basis vector is indexed by a pointed ordered partition
\[
 \pi=(A_1|\cdots|A_b),\qquad i_0\in A_1,
\]
and has the right-normed form
\begin{equation}\label{P2M-eq:v14-hall-word}
 h_\pi=[K_{A_b},[K_{A_{b-1}},\ldots,[K_{A_2},K_{A_1}]\ldots]].
\end{equation}
In the weight-zero sector every block $A_j$ is a singleton or a pair.  Distinct
$\pi$ have distinct leading operation-tree lines in
$\mathscr H_{I,i_0}^{\mathrm{RB}}$.

\subsubsection{Why a three-step channel cannot retain full Hall history}

For an ordered transfer channel let $\mathfrak n_{I;A,B}$ be the ideal in the
three-object matrix block generated by the positive-path entries
$E_{12}^{A,B}$, $E_{23}^{A,B}$ and $E_{13}^{A,B}$.  The defining matrix
relations give
\begin{equation}\label{P2M-eq:v14-channel-radical-cube-zero}
 \boxed{\mathfrak n_{I;A,B}^{\,3}=0.}
\end{equation}

\begin{theorem}[Three-step Hall-fidelity no-go]
\label{P2M-thm:v14-three-step-hall-nogo}
Let $h_\pi$ be a pointed Hall word with $b\ge3$ blocks.  Under the ordered
channel-complete transfer--Rees realization,
\begin{equation}
 \boxed{\rho_I^{\Mot,\mathrm{ch}}(h_\pi)=0.}
\end{equation}
Consequently the channel-complete realization is not faithful on the full
pointed weight-zero Hall operation-history complex.  This failure already occurs in
support three, where the all-singleton Hall word has three blocks.
\end{theorem}

\begin{proof}
Expanding the iterated graded commutator \eqref{P2M-eq:v14-hall-word} gives a signed
sum of associative words, every one of length $b$ in the positive-support
source cells.  In a fixed ordered channel block only two designated
proper-support cells have nonzero positive-path components.  Every associative
word of length at least three therefore lies in
$\mathfrak n_{I;A,B}^{3}$ and vanishes by
\eqref{P2M-eq:v14-channel-radical-cube-zero}.  The argument is blockwise, hence the
word vanishes in the product of all ordered channels.

The vanishing is a genuine loss of operation history rather than a source
relation.  By the pointed root/Bass Hall theorem of
, the Hall words have distinct leading tagged
operation trees and are integrally linearly independent.  In particular the
three-singleton word at support three is nonzero preceding.
\end{proof}

\begin{corollary}[Fixed-depth channel architectures cannot be all-support Hall faithful]
\label{P2M-cor:v14-fixed-depth-hall-nogo}
Suppose a realization architecture has a positive-operation ideal $\mathfrak n$
with $\mathfrak n^L=0$ uniformly in the support cardinality.  Then every Hall
word with at least $L$ blocks is killed.  Since the all-singleton pointed Hall
word on an $m$-element support has $m$ blocks, a realization faithful to Hall
history at all finite supports must have operation depth growing with $m$.
In particular no fixed-size upper-triangular matrix channel can be uniformly
Hall faithful.
\end{corollary}

\paragraph{Binary scope.}
Channels are complete for quadratic latching curvature; full Hall history is
strictly stronger.


\subsubsection{The transfer cocycle telescopes on operation trees}

The geometric transfer profile has precisely the property needed to restore
operation history without introducing parenthesization-dependent coefficients.
Let $T$ be a rooted binary operation tree whose leaves are pairwise disjoint
nonempty supports $A_1,\ldots,A_b$ with union $I$.  At an internal vertex $v$
write $L_v,R_v$ for the supports of its two incoming subtrees and put
\[
 \delta_v:=\eta(|L_v\sqcup R_v|)-\eta(|L_v|)-\eta(|R_v|).
\]

\begin{lemma}[Treewise telescoping of cyclotomic transfer weight]
\label{P2M-lem:v14-tree-transfer-telescoping}
For every rooted binary operation tree $T$,
\begin{equation}
 \boxed{
 \sum_{v\in\operatorname{Int}(T)}\delta_v
 =\eta(|I|)-\sum_{j=1}^{b}\eta(|A_j|).}
\end{equation}
Hence the total transfer exponent depends only on the leaf supports and not on
parenthesization, grafting order, or the chosen binary tree.
\end{lemma}

\begin{proof}
Every internal subtree support appears once with positive sign as
$\eta(|L_v\sqcup R_v|)$ at its own vertex and once with negative sign as one
of the two incoming terms at its parent.  These contributions cancel along
every internal edge.  Only the root contribution $\eta(|I|)$ and the negative
leaf contributions remain.
\end{proof}

\begin{corollary}[Common transfer weight on the weight-zero Hall sector]
\label{P2M-cor:v14-weight-zero-common-transfer}
Let $|I|=m\ge3$.  For every pointed weight-zero Hall partition $\pi$ of $I$,
all blocks are singletons or pairs and therefore have $\eta(|A_j|)=0$.  Thus
any binary expansion of the corresponding operation tree carries the common
geometric transfer factor
\begin{equation}
 \boxed{\mathsf T_N^{\,m-2}.}
\end{equation}
In particular the transfer weight cannot identify two different pointed Hall
trees.
\end{corollary}

\subsubsection{Hall-completed transfer realization}

The previous corollary lets us decorate the intrinsic operation-separated
ledger rather than replace it by a bounded-depth matrix shadow.  Let
$\tau_I$ denote the exact-support basis vector in the cyclotomic transfer--Rees
support algebra.  Define the pointed transfer-decorated Hall target
\begin{equation}
 \boxed{
 \mathscr H_{I,i_0;N}^{\mathrm{RB,tr}}
 :=\mathscr H_{I,i_0}^{\mathrm{RB}}\otimes_{\Z}\Z\tau_I.}
\end{equation}
The exact-support basis line $\Z\tau_I$ is the corresponding line of the transfer--Rees support algebra.  We use the equality
\[
 \mathsf T_N^{m-2}=N^{m-2}\operatorname{id}
\]
to regard the scalar coefficient below as the geometric iterate of the pull--push correspondence.

\begin{theorem}[Pointed Hall-faithful transfer completion]
\label{P2M-thm:v14-hall-faithful-transfer}
For every pointed exact support $(I,i_0)$ with $|I|=m\ge3$, the formula
\begin{equation}
 \boxed{
 \Real_{I,i_0}^{\mathrm{Hall,tr}}(h_\pi)
 :=N^{m-2}\,
   \operatorname{ev}_{\mathrm{RB}}(h_\pi)\otimes\tau_I}
\end{equation}
defines an integral injective chain map
\[
 \Real_{I,i_0}^{\mathrm{Hall,tr}}:
 W_0K_I^{\mathrm{St}}
 \hookrightarrow
 \mathscr H_{I,i_0;N}^{\mathrm{RB,tr}}.
\]
It has the following properties.
\begin{enumerate}[label=\textup{(\roman*)}]
\item Distinct pointed Hall words remain distinct; in particular every tagged
      leading operation tree survives.
\item The adjacent-refinement differential is preserved exactly.
\item The map is natural under order-preserving relabellings of the exact support which preserve the distinguished label.  In a larger ambient support, strict deletion keeps this exact-support factor unchanged when $I$ survives and drops the whole factor otherwise.
\item The only new coefficient is the genuine cyclotomic transfer
      correspondence; no new Hall differential or operation-tree relation is
      imposed in the target.
\end{enumerate}
\end{theorem}

\begin{proof}
By \cref{P2M-cor:v14-weight-zero-common-transfer}, every basis vector in the
exact-support weight-zero Hall complex acquires the same transfer factor
$N^{m-2}$, geometrically represented by $\mathsf T_N^{m-2}$.  The preceding evaluation
\eqref{P2M-eq:v14-preceding-hall-iso} is a chain isomorphism, so multiplying every
basis vector by one common closed central correspondence preserves the
adjacent-refinement differential.

On the integral operation-tree lattice the correspondence acts as
$N^{m-2}$ times the identity.  Multiplication by the nonzero integer
$N^{m-2}$ on a free abelian group is injective.  Since the preceding Hall
transition matrix to tagged operation-tree coordinates is unitriangular,
distinct Hall words remain linearly independent after the common transfer
factor.  Pointed relabelling naturality is inherited from the preceding root/Bass
realization.  Exact-support deletion is the same block rule used throughout this paper: a surviving support-$I$ factor is unchanged and a non-surviving factor is removed.
\end{proof}

\paragraph{Integral scope.}
The Hall transfer map is injective over \(\Z\), not a saturated splitting for
\(m>2\).


\paragraph{Free/torsion separation.}
The Hall ledger is free, whereas the Moore line is \(N\)-torsion.


To remove the auxiliary choice of root from the realization data, keep all
pointings simultaneously:
\begin{equation}
 \boxed{
 \mathscr H_{I,N}^{\mathrm{RB,tr},\bullet}
 :=\prod_{i_0\in I}\mathscr H_{I,i_0;N}^{\mathrm{RB,tr}}.}
\end{equation}
Permutations relabel the pointing factors.  For an ambient support deletion $T\subseteq S$, an exact-support factor indexed by $I$ is kept, with all of its pointings, exactly when $I\subseteq T$; otherwise the whole factor is dropped.

\begin{corollary}[Hall-completed transfer--Rees realization package]
\label{P2M-cor:v14-hall-completed-package}
The product realization data
\begin{equation}\label{P2M-eq:v14-hall-completed-package}
 \boxed{
 \mathscr E_S^{\mathrm{fid}}
 :=\mathscr E_S^{\Mot,\mathrm{ch}}
 \times
 \prod_{\substack{I\subseteq S\\|I|\ge3}}
 \mathscr H_{I,N}^{\mathrm{RB,tr},\bullet}}
\end{equation}
retain simultaneously:
\begin{enumerate}[label=\textup{(\roman*)}]
\item every ordered quadratic FSCComp latching monomial and the full
      $1,N,N^2$ transfer--Rees profile;
\item the genuine exact-order-$N$ Moore class and its motivic antecedent;
\item every pointed weight-zero Hall operation tree and its
      adjacent-refinement differential.
\end{enumerate}
The first factor is the dg motivic channel realization of
\cref{P2M-thm:channel-complete-realization}; the second is chain-level
operation-history data.  No claim is made that their product is already a
single irreducible motivic dg algebra or a full realization of the
bicategorical insertion history.
\end{corollary}

\subsubsection{A fidelity ledger for the successive compressions}

\begin{summary}
For the realization constructions of this paper, the retained information is
strictly nested as follows.
\begin{enumerate}[label=\textup{(\alph*)}]
\item The ordered channel-complete model is faithful to every ordered quadratic
      bipartition monomial, the raw transfer--Rees coefficients, and the Moore
      carry, but not to Hall words with three or more blocks.
\item Paired diamonds forget ordered reversal but retain every unordered
      commutator/bipartition channel and strict Rees flatness.
\item The face-minimal model forgets genuine partition homotopies and retains
      the full non-null highest-face Moore detector module.
\item The CRT-hybrid model retains exactly the primewise-minimal global
      exact-order detector; at nonresonant primes it need not recover every
      individual face coordinate.
\item The Hall-completed package
      \eqref{P2M-eq:v14-hall-completed-package} restores the full pointed
      weight-zero operation-tree history without changing the Moore or Rees
      conclusions of the ordered channel factor.
\end{enumerate}
Thus detector minimality and Hall fidelity are different optimization
problems: compressing to the minimal Moore detector necessarily discards data
which are essential for operation-history fidelity.
\end{summary}


\subsection{Source explanation of the transfer exponent}

\subsubsection{Bicategorical frontier shadow and one transfer factor}

The transfer coefficient is already visible in structural span geometry:
successive insertions retain intermediate \(A_N\)-frontiers absent from direct
insertion.  A finite-span shadow turns each such frontier into one factor
\(N\) after linearization.

\paragraph{A pullback-preserving finite-span shadow.}

Let $U_{\rm grp}$ and $U_{\rm mod}$ denote the underlying-set functors on
finite groups and finite $A_N$-modules.  Put
\[
 \mathbf{FinPair}:=\mathbf{FinSet}\times\mathbf{FinSet}.
\]
For a typed FSC envelope state
\[
 \mathbf X=(G_X,S_X)\in\mathbf{FSCEnv}_{A_N}
\]
define
\begin{equation}
 \boxed{
 \mathcal U_{\rm str}(\mathbf X)
 :=\bigl(U_{\rm grp}(G_X),U_{\rm mod}(S_X)\bigr).}
\end{equation}
On structural morphisms use the underlying pair of set maps.

\begin{proposition}[Faithful pullback shadow of the typed envelope]
\label{P2M-prop:v15-underlying-pullback-shadow}
The functor
\begin{equation}\label{P2M-eq:v15-underlying-pullback-functor}
 \boxed{
 \mathcal U_{\rm str}:
 \mathbf{FSCEnv}_{A_N}\longrightarrow\mathbf{FinPair}}
\end{equation}
is faithful and preserves every pullback used to compose structural FSC
spans.  Consequently it induces a canonical pseudofunctor of span
bicategories
\begin{equation}\label{P2M-eq:v15-span-shadow-pseudofunctor}
 \boxed{
 \mathfrak S_N^{\rm str}:
 \FSCComp\longrightarrow\operatorname{Span}(\mathbf{FinPair}).}
\end{equation}
The induced map on each groupoid of structural 2-morphisms is faithful.
\end{proposition}

\begin{proof}
The typed envelope category used preceding is a full finite-pullback-closed
subcategory of
$\mathbf{Grp}\times A_N\text{-}\mathbf{Mod}$, and its structural pullbacks
are computed componentwise.  The underlying-set functor on groups and the
underlying-set functor on modules both create limits, hence send each such
pullback to the corresponding pullback of finite sets.  They are faithful,
so their product is faithful.

A pullback-preserving functor induces the standard span pseudofunctor:
objects and both legs are sent componentwise, while the compositor is the
canonical comparison between the image of a source pullback and the target
pullback.  Here that comparison is an isomorphism because
\eqref{P2M-eq:v15-underlying-pullback-functor} preserves the pullback.  The source
associators and unitors therefore map to the ordinary span associators and
unitors in $\operatorname{Span}(\mathbf{FinPair})$.  Finally a source
2-morphism is an isomorphism of typed middle objects commuting with both legs;
its underlying pair of bijections determines the original group/module map,
so two distinct source 2-morphisms cannot acquire the same image.
\end{proof}

\paragraph{Shadow scope.}
The finite-set span functor is faithful on actual structural \(2\)-cells but
does not reflect arbitrary span isomorphisms.


\paragraph{The intrinsic frontier remains visible.}

Let
\[
 \mathfrak i_{r,r+2}^{(2)}
\]
denote the composite of the two terminal insertion spans
$[r]\rightsquigarrow[r+1]\rightsquigarrow[r+2]$, and let
$\mathfrak i_{r,r+2}^{\rm dir}$ denote the direct insertion span with the same
endpoints.  Upstream, the controller projection of the composite middle
object is canonically
\begin{equation}\label{P2M-eq:v15-source-frontier-factor}
 \mathcal H_{r+2}^{(2)}\times A_N
 \longrightarrow
 \mathcal H_{r+2}^{(2)},
\end{equation}
where the second factor is invisible to both endpoint legs.

\begin{theorem}[Frontier history survives the finite-span realization]
\label{P2M-thm:v15-frontier-survival}
Assume $N>1$.  The two 1-morphisms
\[
 \mathfrak S_N^{\rm str}(\mathfrak i_{r,r+2}^{(2)}),
 \qquad
 \mathfrak S_N^{\rm str}(\mathfrak i_{r,r+2}^{\rm dir})
\]
are not 2-isomorphic in $\operatorname{Span}(\mathbf{FinPair})$.
Already on the controller component, the middle set of the two-step history
has cardinality exactly $N$ times the middle set of direct insertion.
Consequently the elementary FSC frontier history is not lost by
\eqref{P2M-eq:v15-span-shadow-pseudofunctor}.
\end{theorem}

\begin{proof}
By the preceding frontier-excess theorem, the controller middle object of the
two-step insertion is
$\mathcal H_{r+2}^{(2)}\times A_N$ over the direct middle object
$\mathcal H_{r+2}^{(2)}$.  Taking underlying sets preserves this product, and
$|A_N|=N$.  A 2-isomorphism of finite spans would in particular give a
bijection between their controller middle sets commuting with both legs.
Their cardinalities differ by the factor $N>1$, so no such bijection exists.
\end{proof}


\paragraph{Finite-span linearization and cyclotomic transfer.}

For a finite span
\[
 X\xleftarrow{p}M\xrightarrow{q}Y
\]
define its integral incidence linearization
\begin{equation}\label{P2M-eq:v15-span-linearization}
 \Lin_{\Z}(M):\Z[X]\longrightarrow\Z[Y],
 \qquad
 e_x\longmapsto
 \sum_{m\in p^{-1}(x)}e_{q(m)}.
\end{equation}
Equivalently the $(y,x)$ matrix entry is
$|\{m\in M:p(m)=x,\ q(m)=y\}|$.  Pullback composition of finite spans is sent
exactly to matrix composition.  We apply this construction only to the
controller component of \eqref{P2M-eq:v15-span-shadow-pseudofunctor}.

Let $W_{\rootc}^{\Mot}$ be the fixed root coefficient object used throughout
this paper.  Tensoring \eqref{P2M-eq:v15-span-linearization} with
$\operatorname{id}_{W_{\rootc}^{\Mot}}$ gives the additive root linearization
$\Lin_{\rootc}^{\Mot}$.  Scalar multiplication by $N$ on the root factor is
geometrically represented by the degree-$N$ cyclotomic pull--push operator
\[
 \mathsf T_N=\pi_{N,*}\pi_N^*=N\,\operatorname{id}.
\]

\begin{theorem}[One intrinsic frontier decategorifies to one cyclotomic transfer]
\label{P2M-thm:v15-frontier-transfer-identity}
For the two-step and direct insertion histories above, controller-span
linearization satisfies
\begin{equation}\label{P2M-eq:v15-frontier-linearization-N}
 \boxed{
 \Lin_{\Z}\bigl(
   \mathfrak S_{N,\rm ctl}^{\rm str}(\mathfrak i_{r,r+2}^{(2)})
 \bigr)
 =N\,
 \Lin_{\Z}\bigl(
   \mathfrak S_{N,\rm ctl}^{\rm str}(\mathfrak i_{r,r+2}^{\rm dir})
 \bigr).}
\end{equation}
After root-motive decoration this becomes
\begin{equation}\label{P2M-eq:v15-frontier-transfer-motivic}
 \boxed{
 \Lin_{\rootc}^{\Mot}\bigl(
   \mathfrak i_{r,r+2}^{(2)}
 \bigr)
 =\mathsf T_N\circ
 \Lin_{\rootc}^{\Mot}\bigl(
   \mathfrak i_{r,r+2}^{\rm dir}
 \bigr).}
\end{equation}
Thus the basic transfer factor used in the transfer--Rees realization is the
additive/cyclotomic decategorification of one intrinsic FSC insertion
frontier.
\end{theorem}

\begin{proof}
In \eqref{P2M-eq:v15-source-frontier-factor} the $A_N$ coordinate is invisible to
both endpoint legs.  Hence each point of the direct middle span is replaced by
exactly $N$ points with identical incidence data in the two-step span.  Every
entry of the incidence matrix is therefore multiplied by $N$, proving
\eqref{P2M-eq:v15-frontier-linearization-N}.  Tensor with the root coefficient
object and use $\pi_{N,*}\pi_N^*=N\operatorname{id}$ to obtain
\eqref{P2M-eq:v15-frontier-transfer-motivic}.
\end{proof}

\begin{corollary}[Transfer has a source-side history interpretation]
\label{P2M-cor:v15-transfer-source-history}
The appearance of a single $\mathsf T_N$ in the realization is not merely a
coefficient normalization.  At the elementary two-step insertion level it is
forced by the cardinality of the intrinsic frontier retained by the source
span bicategory.  The span remembers the frontier before linearization; the
transfer operator is its additive geometric readout afterwards.
\end{corollary}


\subsubsection{Binary block frontiers and the full transfer exponent}


The elementary insertion calculation above counts intermediate compressed
\emph{stages}.  That is not the same combinatorial statistic as the transfer
exponent in a binary latching monomial.  Indeed a long sequential insertion
history may carry arbitrarily many intermediate frontiers, whereas
$\delta(A,B)$ is always either one or two.  The correct controller-envelope shadow is
instead blockwise: before the two blocks are grafted, restore the hidden outer
edge of each block which is itself nonsingleton.

For a nonempty ordered block $C$, define the auxiliary block compression
\begin{equation}
 \mathcal G_C^{\rm blk}:=
 \begin{cases}
  \mathcal H_C^{(2)},&|C|=1,\\
  \mathcal H_C^{(2)}/\langle z_{\min C,\max C}\rangle,&|C|\ge2,
 \end{cases}
 \qquad
 q_C^{\rm blk}:\mathcal H_C^{(2)}\twoheadrightarrow
 \mathcal G_C^{\rm blk}.
\end{equation}
Put
\begin{equation}
 F_C:=\ker q_C^{\rm blk}.
\end{equation}
Thus
\begin{equation}\label{P2M-eq:v16-block-frontier-rank}
 \boxed{
 F_C\cong
 \begin{cases}
  0,&|C|=1,\\
  A_N,&|C|\ge2.
 \end{cases}}
\end{equation}
For $|C|\ge3$ this is the controller component of the ordinary FSC static
compression; for $|C|=2$ it is the same outer-edge quotient used only as an
auxiliary binary-block compression.  No new FSC provenance object is being
postulated.

\begin{lemma}[Transfer exponent counts nonsingleton blocks]
\label{P2M-lem:v16-delta-block-count}
For every ordered bipartition $I=A\sqcup B$ with $|I|\ge3$,
\begin{equation}
 \boxed{
 \delta(A,B)=
 \mathbf 1_{\{|A|\ge2\}}
 +\mathbf 1_{\{|B|\ge2\}}.}
\end{equation}
\end{lemma}

\begin{proof}
If one block is a singleton, the other has cardinality at least two and the
right-hand side is one, equal to the defining value of $\delta$.  If both
blocks have cardinality at least two, the right-hand side is two, again the
defining value of $\delta$.
\end{proof}

Let
\begin{equation}
 \alpha_{A,B}:\mathcal H_I^{(2)}\longrightarrow
 \mathcal G_A^{\rm blk}\times\mathcal G_B^{\rm blk},
 \qquad
 h\longmapsto
 \bigl(q_A^{\rm blk}d_{I,A}h,
       q_B^{\rm blk}d_{I,B}h\bigr).
\end{equation}
The direct controller binary-grafting shadow is the span
\begin{equation}
 \mathfrak g^{\rm dir}_{I;A,B}:
 \mathcal G_A^{\rm blk}\times\mathcal G_B^{\rm blk}
 \xleftarrow{\ \alpha_{A,B}\ }
 \mathcal H_I^{(2)}
 \xrightarrow{\ q_I\ }
 \mathcal G_I^{\triangle}.
\end{equation}
Now restore the local complete-envelope lift on each block and define
\begin{equation}\label{P2M-eq:v16-resolved-middle}
 \boxed{
 \mathcal B_{I;A,B}:=
 \mathcal H_I^{(2)}
 \times_{\mathcal G_A^{\rm blk}\times\mathcal G_B^{\rm blk}}
 \bigl(\mathcal H_A^{(2)}\times\mathcal H_B^{(2)}\bigr),}
\end{equation}
where the second map to the base is
$q_A^{\rm blk}\times q_B^{\rm blk}$.  Its resolved span is
\begin{equation}\label{P2M-eq:v16-resolved-binary-span}
 \mathfrak g^{\rm res}_{I;A,B}:
 \mathcal G_A^{\rm blk}\times\mathcal G_B^{\rm blk}
 \xleftarrow{\ q_A^{\rm blk}\times q_B^{\rm blk}\ }
 \mathcal B_{I;A,B}
 \xrightarrow{\ q_I\operatorname{pr}_I\ }
 \mathcal G_I^{\triangle}.
\end{equation}

\begin{theorem}[Binary block-frontier factorization]
\label{P2M-thm:v16-binary-frontier-factorization}
For every ordered bipartition $I=A\sqcup B$, the projection
$\operatorname{pr}_I:\mathcal B_{I;A,B}\to\mathcal H_I^{(2)}$ has a canonical
section
\begin{equation}
 s_{A,B}(h)=\bigl(h,d_{I,A}h,d_{I,B}h\bigr).
\end{equation}
Moreover there is a canonical isomorphism of controller groups over
$\mathcal H_I^{(2)}$
\begin{equation}
 \boxed{
 \mathcal B_{I;A,B}
 \cong
 \mathcal H_I^{(2)}\times F_A\times F_B
 \cong
 \mathcal H_I^{(2)}\times A_N^{\delta(A,B)}.}
\end{equation}
Under this isomorphism the two endpoint maps in
\eqref{P2M-eq:v16-resolved-binary-span} are independent of the frontier
coordinates.
\end{theorem}

\begin{proof}
For fixed $h\in\mathcal H_I^{(2)}$, a point of the pullback
\eqref{P2M-eq:v16-resolved-middle} is a pair of lifts $h_A,h_B$ of
$q_A^{\rm blk}d_{I,A}h$ and $q_B^{\rm blk}d_{I,B}h$.  The restrictions
$d_{I,A}h,d_{I,B}h$ give the canonical base lifts.  Hence every other pair is
uniquely
\[
 h_A=d_{I,A}h\,f_A,
 \qquad
 h_B=d_{I,B}h\,f_B,
 \qquad
 (f_A,f_B)\in F_A\times F_B.
\]
The frontier groups are central outer-edge lines in the corresponding
class-two envelopes, so
\[
 (h,f_A,f_B)\longmapsto
 \bigl(h,d_{I,A}h\,f_A,d_{I,B}h\,f_B\bigr)
\]
is a group isomorphism, not merely a bijection of fibres.  The final
identification follows from \cref{P2M-eq:v16-block-frontier-rank,P2M-lem:v16-delta-block-count}.
Both endpoint maps factor through the relevant block quotients or through
$q_I\operatorname{pr}_I$, so they forget $f_A,f_B$.
\end{proof}

\begin{corollary}[The tetrahedral genuine partition has two frontiers]
\label{P2M-cor:v16-two-two-frontier}
For $I=\{1,2,3,4\}$ and the $2+2$ partition
$A=\{1,2\}$, $B=\{3,4\}$,
\begin{equation}
 \boxed{
 \mathcal B_{1234;12,34}
 \cong
 \mathcal H_{1234}^{(2)}\times A_N^2.}
\end{equation}
Thus the first genuine partition channel carries two independent block
compression frontiers at the controller-resolution level.
\end{corollary}

\begin{theorem}[Full raw transfer exponent from block frontiers]
\label{P2M-thm:v16-full-transfer-frontier-readout}
Apply finite-span incidence linearization to the direct and resolved controller
spans.  Then
\begin{equation}\label{P2M-eq:v16-full-transfer-linearization}
 \boxed{
 \Lin_{\Z}(\mathfrak g^{\rm res}_{I;A,B})
 =N^{\delta(A,B)}
  \Lin_{\Z}(\mathfrak g^{\rm dir}_{I;A,B}).}
\end{equation}
After tensoring with the fixed root coefficient object and using
$\mathsf T_N=\pi_{N,*}\pi_N^*=N\operatorname{id}$, this becomes
\begin{equation}\label{P2M-eq:v16-full-transfer-cyclotomic}
 \boxed{
 \Lin_{\rootc}^{\Mot}(\mathfrak g^{\rm res}_{I;A,B})
 =\mathsf T_N^{\delta(A,B)}\circ
  \Lin_{\rootc}^{\Mot}(\mathfrak g^{\rm dir}_{I;A,B}).}
\end{equation}
Hence the complete raw binary transfer exponent in the channel law is the
cyclotomic decategorification of the two blockwise compression frontiers.
The normalized latching coefficient is one transfer power lower because the
Moore equation $dK_I=N\Omega_I^{\rm rec}$ removes the common factor $N$.
\end{theorem}

\begin{proof}
By \cref{P2M-thm:v16-binary-frontier-factorization}, every incidence point of the
direct middle span is replaced by exactly
$|A_N|^{\delta(A,B)}=N^{\delta(A,B)}$ points with identical endpoint data in
the resolved span.  Thus every entry of the incidence matrix is multiplied by
that number, proving \eqref{P2M-eq:v16-full-transfer-linearization}.  The
cyclotomic identity replaces each scalar factor $N$ by the genuine transfer
correspondence $\mathsf T_N$, proving
\eqref{P2M-eq:v16-full-transfer-cyclotomic}.  The final normalization statement is
exactly the passage from the raw transfer coefficient
$\mathsf T_N^{\delta}$ to the normalized channel coefficient
$\mathsf T_N^{\delta-1}$ in \cref{P2M-thm:ordered-channel-realization}.
\end{proof}

\begin{proposition}[Sequential frontier length is not the transfer exponent]
\label{P2M-prop:v16-sequential-not-delta}
Iterating the terminal insertion span through $k\ge2$ successive insertions
produces $k-1$ intermediate compression frontiers relative to direct
insertion.  This count depends on the length of the chosen insertion history
and therefore cannot equal the binary invariant $\delta(A,B)$ in general.
For example a four-step terminal enlargement carries three intermediate
$A_N$-frontiers, whereas every binary partition with both blocks
nonsingleton has $\delta=2$.
\end{proposition}

\begin{proof}
The case $k=2$ is the preceding frontier-excess theorem.  Composing one more
insertion span introduces one additional pullback over the new intermediate
compressed controller, whose fibre is the one-dimensional kernel of its
static quotient.  Induction gives one factor for each intermediate stage,
hence $k-1$.  Since $\delta$ depends only on whether the two final blocks are
singletons or nonsingletons, the two counts have different dependence on the
data.
\end{proof}

\subsubsection{Pair-normal binary history and the transfer exponent}


The controller resolution of the preceding section counts one hidden outer
edge for each nonsingleton child block.  The preceding pair-normal cospan
theorem provides an integral marked dg version of the same count.  Upstream
the internal source additionally supplies a structurally reduced all-support higher-Moore
history lift, but the transfer exponent already occurs at pair-normal level
and does not depend on that stronger formal coherence theorem.

For a nonempty ordered block $C$, write
\[
 e(C):=\{\min C,\max C\}
\]
when $|C|\ge2$, and let
\[
 \alpha_C:=[D_{\min C},D_{\max C}].
\]
Inside the parent complete pair seed $\mathfrak P_I^{(2),\Z}$ denote the
canonical filler of this pair by $H_{e(C)}^I$.

\begin{definition}[Binary pair-history presentation]
\label{P2M-def:v18-binary-pair-history}
For an ordered bipartition $I=A\sqcup B$, form the semi-free multi-pushout
\begin{equation}
 \mathfrak P_{I;A|B}^{(2),\rm hist}
 :=
 \mathfrak P_I^{(2),\Z}
 \underset{\mathfrak P_{A,<}^{(2),\rm vis}}{\amalg}
 \mathfrak P_A^{(2),\Z}
 \underset{\mathfrak P_{B,<}^{(2),\rm vis}}{\amalg}
 \mathfrak P_B^{(2),\Z},
\end{equation}
where a singleton block contributes the identity pushout because its visible
and complete pair presentations coincide.  By the same semi-free pushout
calculation as the preceding two-step pair-history theorem, this has the
explicit presentation
\begin{equation}
 \boxed{
 \mathfrak P_{I;A|B}^{(2),\rm hist}
 \cong\mathfrak P_I^{(2),\Z}
 \left\langle
   \eta_C\ (C\in\{A,B\},\ |C|\ge2);\
   d\eta_C=\alpha_C
 \right\rangle.}
\end{equation}
Define the marked contraction
\begin{equation}
 q_{A|B}^{\rm pair}:
 \mathfrak P_{I;A|B}^{(2),\rm hist}
 \longrightarrow
 \mathfrak P_I^{(2),\Z}
\end{equation}
by fixing the parent complete seed and sending
$\eta_C\mapsto H_{e(C)}^I$.  Put
\begin{equation}
 \theta_C:=\eta_C-H_{e(C)}^I.
\end{equation}
The \emph{binary frontier-line ledger} is the free marked submodule
\begin{equation}
 \mathsf{Fr}_{A|B}:=
 \bigoplus_{\substack{C\in\{A,B\}\\ |C|\ge2}}
 \Z\theta_C.
\end{equation}
\end{definition}

\paragraph{Type firewall.}
Binary latching readouts, structural insertion histories and the reduced
Moore-history pseudofunctor are different objects.  The transfer character
records only the binary frontier ledger.


\begin{theorem}[Pair-normal source formula for the full transfer exponent]
\label{P2M-thm:v18-pair-history-delta}
For every ordered bipartition $I=A\sqcup B$:
\begin{enumerate}[label=\textup{(\roman*)}]
\item each $\theta_C$ is closed;
\item the marked lines $\theta_C$ are independent and
\begin{equation}\label{P2M-eq:v18-frontier-ledger-rank}
 \boxed{
 \mathsf{Fr}_{A|B}\cong\Z^{\delta(A,B)},
 \qquad
 \delta(A,B)=
 \mathbf1_{\{|A|\ge2\}}+
 \mathbf1_{\{|B|\ge2\}};}
\end{equation}
\item after reduction modulo $N$ the frontier-line shadow is
\begin{equation}\label{P2M-eq:v18-frontier-ledger-modN}
 \boxed{
 \mathsf{Fr}_{A|B}\otimes A_N
 \cong A_N^{\delta(A,B)};}
\end{equation}
\item finite incidence/cardinality readout contributes the factor
$N^{\delta(A,B)}$, and cyclotomic decoration contributes
$\mathsf T_N^{\delta(A,B)}$.
\end{enumerate}
Hence the entire raw $1,N,N^2$ transfer profile is already the
cyclotomic readout of pair-normal marked frontier history.  No
higher-Moore insertion pseudofunctor is used in this derivation.
\end{theorem}

\begin{proof}
For every nonsingleton child $C$, the preceding pair-normal insertion cospan uses the complete child seed $\mathfrak P_C^{(2),\Z}$ and the labelled inclusion of its visible chart into the parent complete seed.  Its pushout therefore retains a second copy $\eta_C$ of the child hidden filler with the same boundary $\alpha_C$ as the canonical parent copy $H_{e(C)}^I$.  Hence $d\theta_C=0$.  Because $A$ and $B$ are disjoint, their hidden pair labels are distinct, and semi-free pushout keeps the corresponding extra fillers as independent marked generators.  Their differences therefore give independent marked history lines, one for each nonsingleton child.  The number
of such children is exactly \cref{P2M-lem:v16-delta-block-count}, proving
\eqref{P2M-eq:v18-frontier-ledger-rank}.  Reduction modulo $N$ gives
\eqref{P2M-eq:v18-frontier-ledger-modN}.  Each $A_N$ line has $N$ elements, so
forgetting the dg operations and retaining only the marked frontier
coordinates multiplies incidence multiplicity by $N^{\delta(A,B)}$.
Finally use $\mathsf T_N=\pi_{N,*}\pi_N^*=N\operatorname{id}$.
\end{proof}

\begin{corollary}[The $2+2$ Pfaffian has two pair-history lines]
\label{P2M-cor:v18-two-two-pair-history}
For $I=\{1,2,3,4\}$ and $A=\{1,2\}$, $B=\{3,4\}$,
\[
 \boxed{
 \mathsf{Fr}_{12|34}
 =\Z\theta_{12}\oplus\Z\theta_{34},
 \qquad
 \mathsf{Fr}_{12|34}\otimes A_N\cong A_N^2.}
\]
Its finite/cyclotomic readouts are therefore $N^2$ and $\mathsf T_N^2$.
\end{corollary}

\begin{proposition}[Compatibility with the controller block-frontier shadow]
\label{P2M-prop:v18-pair-controller-agreement}
The pair-history ledger of \cref{P2M-thm:v18-pair-history-delta} and the controller
resolution of \cref{P2M-thm:v16-binary-frontier-factorization} have the same
frontier index set: one marked coordinate for each nonsingleton child block.
After reduction modulo $N$ and forgetting brackets/products, both yield
$A_N^{\delta(A,B)}$ and hence the same transfer readout.

This is an agreement of marked frontier shadows, not an identification of the
full semi-free dg cospan apex with the finite controller pullback group.
\end{proposition}


\subsubsection{Transfer as a binary support-operation contraction character}


The preceding pair cospan enhancement retains sequential pair-normal insertion history as closed frontier-difference lines in semi-free pushout apices, and the reduced all-support theorem extends this formally to all supports by a structurally reduced higher-Moore pseudofunctor.  The transfer--Rees realization nevertheless uses a different combinatorial object: a binary support-operation tree in the Cartan--Moore latching/Hall language.  The pair-history theorem of the preceding section assigns each binary node a rigorously typed frontier multiplicity, and the resulting transfer character is a decategorified readout on that binary shadow, not a replacement for the formal history pseudofunctor.


\begin{definition}[Transfer frontier count and character]
\label{P2M-def:v17-transfer-history-character}
Let $\tau$ be a rooted binary support-operation tree in the latching/Hall language.  For an internal vertex
$v$, let $A_v,B_v$ be the disjoint child supports grafted at $v$ and put
\[
 \kappa_v:=\delta(|A_v|,|B_v|)
 =\mathbf 1_{\{|A_v|\ge2\}}
  +\mathbf 1_{\{|B_v|\ge2\}}.
\]
Define
\begin{equation}\label{P2M-eq:v17-history-kappa}
 \boxed{
 \kappa_{\rm tr}(\tau):=
 \sum_{v\in\operatorname{Int}(\tau)}\kappa_v,}
\end{equation}
and its scalar and cyclotomic transfer characters by
\begin{equation}\label{P2M-eq:v17-history-transfer-character}
 \boxed{
 \Theta_N^{\Z}(\tau):=N^{\kappa_{\rm tr}(\tau)},
 \qquad
 \Theta_N^{\Mot}(\tau):=
 \mathsf T_N^{\kappa_{\rm tr}(\tau)}.}
\end{equation}
The leaves may themselves be nontrivial support blocks; no singleton-leaf
assumption is made.
\end{definition}

\begin{theorem}[Grafting multiplicativity of the transfer character]
\label{P2M-thm:v17-history-character-multiplicative}
Suppose $\tau$ is obtained by grafting rooted histories $\tau_A,\tau_B$ below
a new binary root with child supports $A,B$.  Then
\begin{equation}\label{P2M-eq:v17-kappa-grafting}
 \boxed{
 \kappa_{\rm tr}(\tau)
 =\kappa_{\rm tr}(\tau_A)
  +\kappa_{\rm tr}(\tau_B)
  +\delta(|A|,|B|).}
\end{equation}
Consequently
\begin{equation}\label{P2M-eq:v17-theta-grafting}
 \boxed{
 \Theta_N^{\Mot}(\tau)
 =\Theta_N^{\Mot}(\tau_A)
  \Theta_N^{\Mot}(\tau_B)
  \mathsf T_N^{\delta(|A|,|B|)},}
\end{equation}
with the analogous scalar identity for $\Theta_N^{\Z}$.
\end{theorem}

\begin{proof}
The internal vertices of $\tau$ are the disjoint union of the internal
vertices of $\tau_A$, those of $\tau_B$, and the new root vertex.  Summing
\eqref{P2M-eq:v17-history-kappa} gives \eqref{P2M-eq:v17-kappa-grafting}; exponentiating
and using composition of powers of $\mathsf T_N$ gives
\eqref{P2M-eq:v17-theta-grafting}.
\end{proof}

\begin{theorem}[Binary frontier resolution realizes the local contraction factor]
\label{P2M-thm:v17-local-history-contraction-factor}
For a one-vertex binary graft $A|B$, the pair-normal frontier ledger of
\cref{P2M-thm:v18-pair-history-delta} and the auxiliary controller resolution of
\cref{P2M-thm:v16-binary-frontier-factorization} have the same mod-$N$ invisible frontier shadow
\[
 F_A\times F_B\cong A_N^{\delta(A,B)}.
\]
Hence its finite-span linearization is multiplication by
$\Theta_N^{\Z}(A|B)=N^{\delta(|A|,|B|)}$, and its root-motive decoration is
multiplication by
$\Theta_N^{\Mot}(A|B)=\mathsf T_N^{\delta(|A|,|B|)}$.

Recursively applying the local resolution at every internal vertex of a
binary history $\tau$ produces an auxiliary iterated frontier resolution whose
invisible finite fibre has cardinality
\begin{equation}\label{P2M-eq:v17-tree-frontier-cardinality}
 \boxed{N^{\kappa_{\rm tr}(\tau)}.}
\end{equation}
Thus the transfer character is the finite/cyclotomic decategorification of the
block-frontier part of the uncontracted history tree.
\end{theorem}

\begin{proof}
The one-vertex statement is
\cref{P2M-thm:v18-pair-history-delta,P2M-prop:v18-pair-controller-agreement,P2M-thm:v16-full-transfer-frontier-readout}.
For a tree, form the auxiliary resolution recursively from the leaves upward.
At each new vertex the local frontier coordinates are independent of the
endpoint incidence data and contribute the product factor
$A_N^{\delta(A_v,B_v)}$.  Their cardinalities therefore multiply.  Taking the
product over all internal vertices gives \eqref{P2M-eq:v17-tree-frontier-cardinality};
finite-span linearization and cyclotomic decoration give
\eqref{P2M-eq:v17-history-transfer-character}.
\end{proof}

\begin{corollary}[The contraction character telescopes]
\label{P2M-cor:v17-history-character-telescopes}
Let the leaf supports of $\tau$ be $L_1,\ldots,L_s$ and let its root support
be $I$.  Then
\begin{equation}
 \boxed{
 \kappa_{\rm tr}(\tau)
 =\eta(|I|)-\sum_{a=1}^s\eta(|L_a|).}
\end{equation}
In particular $\Theta_N^{\Mot}(\tau)$ depends only on the root and leaf
support cardinalities, not on the parenthesization or the internal grafting
history.  On the pointed weight-zero Hall sector, whose leaves are singletons
or pairs, this becomes
\[
 \kappa_{\rm tr}(\tau)=|I|-2,
 \qquad
 \Theta_N^{\Mot}(\tau)=\mathsf T_N^{|I|-2}.
\]
\end{corollary}

\begin{proof}
The first formula is exactly the treewise coboundary telescoping identity of
\cref{P2M-lem:v14-tree-transfer-telescoping}.  Since $\eta(1)=\eta(2)=0$, the
weight-zero specialization follows.
\end{proof}

\begin{theorem}[Transfer contraction is necessarily history-forgetting]
\label{P2M-thm:v17-transfer-not-history-faithful}
The character $\Theta_N^{\Mot}$ is not faithful on binary operation-history
shadows.  Any two distinct rooted binary support-operation trees with the same
root support and the same multiset of leaf support cardinalities have the
same transfer character.  In particular different parenthesizations of a
fixed singleton/pair Hall partition are indistinguishable after transfer
contraction.  A fortiori this character cannot recover the richer structural
history retained preceding.
\end{theorem}

\begin{proof}
By \cref{P2M-cor:v17-history-character-telescopes}, the exponent is determined by $|I|$ and the leaf cardinalities alone.  Already at pair-normal level, sequential insertion cospans retain distinct closed frontier-history lines for different intermediate stages, while the binary transfer character depends only on the final support-operation tree through its telescoping exponent.  Therefore the transfer readout is necessarily many-to-one relative to the structural history retained by the source cospan calculus.
\end{proof}

\begin{proposition}[Relation with the preceding reduced higher-Moore history lift]
\label{P2M-prop:v20-relation-to-preceding-moore-lift}
On the rigorously proved pair-normal sector, the transfer readout factors as
\[
 \begin{gathered}
 \FSCComp_{\le2}^{\rm mk}
 \xrightarrow{\ \mathcal E_{\rm pair}\ }
 \mathbf{PairCosp}_{\Z}^{\rm mk}
 \longrightarrow
 \{\text{marked binary frontier ledgers}\}\\
 {}\xrightarrow{\ \operatorname{rk}\ }
 \mathbf N
 \xrightarrow{\ k\mapsto\mathsf T_N^k\ }
 \langle\mathsf T_N\rangle
 \end{gathered}
\]
The second arrow keeps only the closed frontier-difference lines associated to
the two child blocks of the selected binary latching node.  Its rank is
$\delta(A,B)$ by \cref{P2M-thm:v18-pair-history-delta}.

For higher support, the internal source supplies the structurally reduced formal
history pseudofunctor $\mathcal E_{\rm Moore}$.  The displayed transfer readout
is compatible with its pair-normal forgetful shadow, but it factors through a
much coarser binary frontier ledger and therefore forgets most formal history.
No claim is made that this scalar/cyclotomic character itself is the
pseudofunctor $\mathcal E_{\rm Moore}$ or a downstream motivic insertion
realization.
\end{proposition}


\paragraph{The combined fidelity package.}

\begin{summary}
For every finite ordered support system, the constructions separate six
levels of information:
\begin{enumerate}[label=\textup{(\roman*)}]
\item the structural correspondence bicategory $\FSCComp$;
\item its proved pair-normal marked cospan enhancement
      $\mathcal E_{\rm pair}$, retaining hidden pair frontiers and sequential
      pair-history lines;
\item the formal all-support Cartan--Moore controller together with the
      structurally reduced history pseudofunctor $\mathcal E_{\rm Moore}$;
\item the binary-operation transfer character $\Theta_N^{\Mot}$, whose local
      exponent is the rank of the pair-normal binary frontier ledger;
\item the motivic channel realization $\rho_S^{\Mot,\mathrm{ch}}$, retaining
      the complete quadratic latching polynomial, the $1,N,N^2$ transfer--Rees
      law, and the genuine exact-order Moore pair;
\item the Hall-completed transfer realization, retaining every pointed
      weight-zero operation tree and its adjacent-refinement differential.
\end{enumerate}
Levels \textup{(ii)} and \textup{(iv)} meet at each binary node by
\cref{P2M-thm:v18-pair-history-delta}; levels \textup{(iv)} and \textup{(v)} agree
on the corresponding transfer exponent.  The identifications at levels \textup{(ii)}--\textup{(iv)} are coarse
readouts of level \textup{(iii)}; they do not replace its formal history
data and do not by themselves construct a downstream motivic insertion
pseudofunctor.
\end{summary}


\section{Canonical detector compression}

The complete model stores both orders of every bipartition.  Folding this
involution preserves transfer--Rees flatness; the derived Moore sector then
admits a sharper detector compression.

For an unordered bipartition
\[
 P=\{A,B\},\qquad A\sqcup B=I,
\]
write $\delta(P):=\delta(A,B)$; this is independent of the ordering.  Let
\[
 \mathfrak P(I)
 :=\mathfrak P^{\to}(I)/(A,B)\sim(B,A),
 \qquad
 \#\mathfrak P(I)=2^{|I|-1}-1.
\]

\begin{definition}[Paired diamond channel]
\label{P2M-def:paired-diamond-channel}
Fix $P=\{A,B\}\in\mathfrak P(I)$.  The paired diamond
$\mathbb V_{I;P}^{\Mot,\diamond}$ has four graded objects
$s,v_A,v_B,t$ built from shifted copies of $W_{\rootc}^{\Mot}$.  Its only
matrix paths used below are
\[
 s\longrightarrow v_A\longrightarrow t,
 \qquad
 s\longrightarrow v_B\longrightarrow t.
\]
Choose the shifts so that the first edge on the $A$-path has degree
$1-|A|$, the first edge on the $B$-path has degree $1-|B|$, and the two
second edges, after multiplication by $b^{\Mot}$, have the complementary
source-cell degrees.  Denote the corresponding matrix units by
$E_{sA},E_{At},E_{sB},E_{Bt}$ and the common top unit by $E_{st}$; in
particular $|E_{st}|=3-|I|$ and $dE_{st}=0$.

Give the two proper-support cells the future components
\begin{align}
 \widehat K_A
 &\longmapsto
 E_{sA}
 +\chi_I(B,A)E_{Bt}b^{\Mot},
 \\
 \widehat K_B
 &\longmapsto
 E_{sB}
 +\chi_I(A,B)E_{At}b^{\Mot},
\end{align}
and put
\begin{equation}
 \boxed{
 K_I\longmapsto
 2\sigma_I^{\rm sh}E_{st}
 \mathsf T_N^{\delta(P)-1}c^{\Mot}.}
 \label{P2M-eq:paired-diamond-carry}
\end{equation}
All other proper-support cells have zero component in this block.  The block
is exact-support $I$ and disappears under every proper support deletion.
The absence of a transfer factor on the two future edges is essential: the
ordered monomial in $\Omega_I^{\mathrm{rec}}$ already has coefficient
$N^{\delta(P)-1}$, represented motivically by
$\mathsf T_N^{\delta(P)-1}$.  Adding that factor to an edge would count it
twice; for $\delta(P)=2$ it would incorrectly produce $2N^2b^{\Mot}$ in
place of $2Nb^{\Mot}$.
\end{definition}

The two paths are deliberately weighted by the signs of the \emph{opposite}
ordered monomials.  This makes the two curvature contributions agree rather
than cancel.

\begin{theorem}[Ordered-pair folding preserves strict Rees flatness]
\label{P2M-thm:ordered-pair-folding}
For every $P=\{A,B\}\in\mathfrak P(I)$, the paired diamond satisfies
\begin{equation}
 \boxed{
 \rho_{I;P}^{\diamond}(\Omega_I^{\mathrm{rec}})
 =2E_{st}\mathsf T_N^{\delta(P)-1}b^{\Mot},}
 \label{P2M-eq:paired-diamond-latching}
\end{equation}
and
\begin{equation}
 \boxed{
 d\rho_{I;P}^{\diamond}(K_I)
 =N\rho_{I;P}^{\diamond}(\Omega_I^{\mathrm{rec}}).}
 \label{P2M-eq:paired-diamond-dg}
\end{equation}
At the raw transfer--Rees level the two ordered paths contribute together
$2E_{st}\mathsf T_N^{\delta(P)}b^{\Mot}$, exactly the differential of
\eqref{P2M-eq:paired-diamond-carry}.  Hence replacing the two ordered blocks
$(A,B)$ and $(B,A)$ by one paired diamond preserves all chain-level Rees data.
\end{theorem}

\begin{proof}
The product corresponding to $\widehat K_A\widehat K_B$ uses the path
$s\to v_A\to t$ and equals
\[
 \chi_I(A,B)E_{st}b^{\Mot}.
\]
Its curvature coefficient is
$\chi_I(A,B)\mathsf T_N^{\delta-1}$, so its contribution is
$+E_{st}\mathsf T_N^{\delta-1}b^{\Mot}$.  Similarly the reverse ordered
monomial uses $s\to v_B\to t$ and contributes the same element after its own
curvature coefficient is applied.  No cross path composes.  The two
contributions therefore
sum to \eqref{P2M-eq:paired-diamond-latching}.

Since $dc^{\Mot}=Nb^{\Mot}$,
\[
 d\bigl(2\sigma_I^{\rm sh}E_{st}
 \mathsf T_N^{\delta-1}c^{\Mot}\bigr)
 =2N E_{st}\mathsf T_N^{\delta-1}b^{\Mot},
\]
which proves \eqref{P2M-eq:paired-diamond-dg}.  Tensoring the two proper-support
edges with $\tau_A,\tau_B$ restores the one outer Moore factor; hence the raw product is
$2\mathsf T_N^\delta b^{\Mot}$ and agrees with the same differential.
\end{proof}

\begin{corollary}[Strict channel count drops by a factor of two]
\label{P2M-cor:strict-unordered-channel-count}
There is a strict deletion-functorial, permutation-coherent transfer--Rees
realization with one paired diamond for each unordered bipartition.  Thus the
number of strict exact-support channel blocks at cardinality $m$ drops from
\[
 2^m-2
 \qquad\text{to}\qquad
 \boxed{2^{m-1}-1}
\]
without discarding any commutator or genuine partition term of the recursive
curvature.
\end{corollary}

\begin{proof}
Take the finite product of the paired diamonds.  Support deletion and
permutation relabelling act on unordered bipartitions, and
\cref{P2M-thm:ordered-pair-folding} proves the blockwise dg and raw Rees identities.
\end{proof}

\subsection{The derived Moore class sees only highest-face channels}

The previous compression is still a strict chain-level Rees statement.  In
the derived Moore sector one can compress much further because the genuine
partition channels carry an extra transfer factor.

For $|I|=m$ let
\begin{equation}
 \mathcal M_N[I]
 :=\bigl[\Z K_I\xrightarrow{\ N\ }\Z\Omega_I\bigr]
 \label{P2M-eq:support-moore-two-cell}
\end{equation}
with $K_I$ in degree $1-m$ and $\Omega_I$ in degree $2-m$.  Up to shift this
is a projective resolution of $\Z/N$, so
\begin{equation}
 \boxed{
 \operatorname{End}_{D(\Z)}(\mathcal M_N[I])\simeq\Z/N.}
 \label{P2M-eq:moore-derived-end-ring}
\end{equation}

\begin{proposition}[Derived class of one paired channel]
\label{P2M-prop:paired-channel-derived-class}
Restrict the paired-diamond realization to the exact-support Moore two-cell
\eqref{P2M-eq:support-moore-two-cell}.  Under
\eqref{P2M-eq:moore-derived-end-ring}, its derived endomorphism class is
\begin{equation}
 \boxed{
 [f_P]=2N^{\delta(P)-1}\pmod N.}
 \label{P2M-eq:paired-derived-scalar}
\end{equation}
Consequently:
\begin{enumerate}[label=\textup{(\roman*)}]
\item if $P=\{\{i\},I\setminus\{i\}\}$ is a singleton--highest-face
      partition, then $[f_P]=2\in(\Z/N)^\times$;
\item if both blocks of $P$ have cardinality at least two, then $[f_P]=0$.
\end{enumerate}
In the second case the map is explicitly null-homotopic.
\end{proposition}

\begin{proof}
The paired top class and carry are respectively
$2\mathsf T_N^{\delta-1}b^{\Mot}$ and
$2\sigma_I^{\rm sh}\mathsf T_N^{\delta-1}c^{\Mot}$.  Identify the shifted
root Moore pair with \eqref{P2M-eq:support-moore-two-cell} by taking
$\sigma_I^{\rm sh}E_{st}c^{\Mot}$ as its carry basis and
$E_{st}b^{\Mot}$ as its closed basis.  Then the induced chain map
is multiplication by $2N^{\delta-1}$.  Since
$\mathcal M_N[I]\simeq(\Z/N)[m-2]$, its derived endomorphism ring is
$\Z/N$, proving \eqref{P2M-eq:paired-derived-scalar}.  The intrinsic shift
$[m-2]$ here places the cohomology of $\Z/N$ in degree $2-m$; it is the
inverse of the external regrading $[2-m]$ used later to move that cohomology
to degree zero in $\mathfrak D_I^{\rm face}$.

For $\delta=2$ the map is multiplication by $2N$.  An explicit chain homotopy
$H$ is
\[
 H(\Omega_I)=2K_I,
 \qquad
 H(K_I)=0.
\]
Indeed $dH+Hd=2N\,\mathrm{id}$.  For $\delta=1$ the class is multiplication
by $2$, which is a unit because $N$ is odd.
\end{proof}

\begin{corollary}[Genuine partition channels are Moore-null]
\label{P2M-cor:genuine-partition-moore-null}
Every paired channel with $|A|,|B|\ge2$ is essential for strict raw
transfer--Rees flatness but contributes zero to the derived exact-support
Moore class.  Its role is a chain-level partition homotopy, not a new torsion
detector line.
\end{corollary}

This distinction is exactly the one hidden by a raw channel count.  The
exponential family of genuine partition blocks records a strict flatness
presentation; the nontrivial torsion detector is supported only on the
codimension-one faces.

\subsection{The permutation module of highest-face detectors}

Let
\[
 \mathfrak F(I)
 :=\{P_i=\{\{i\},I\setminus\{i\}\}:i\in I\}
 \subset\mathfrak P(I).
\]
It has cardinality $m$ and carries the evident permutation action of $S_I$.

\begin{definition}[Face-minimal detector stage]
\label{P2M-def:face-minimal-detector-stage}
The face-minimal exact-support target is the product of the paired diamonds
indexed only by $\mathfrak F(I)$:
\[
 \boxed{
 \mathscr E_I^{\Mot,\mathrm{face}}
 :=\prod_{i\in I}
   \REnd\bigl(\mathbb V_{I;P_i}^{\Mot,\diamond}\bigr).}
\]
All genuine partition channels are omitted.  The source map is the projection
of the paired strict-Rees realization to these $m$ blocks.  Equivalently, it
keeps every singleton--highest-face commutator and sends every genuine
partition monomial to zero.
\end{definition}

\begin{theorem}[Detector-minimal permutation-coherent tower]
\label{P2M-thm:detector-minimal-tower}
The face-minimal stages assemble over all finite ordered supports into a strict
support-deletion functorial and permutation-coherent dg realization
\[
 \rho_S^{\Mot,\mathrm{face}}:
 \mathscr U^{\mathrm{rec}}_{S,N}
 \longrightarrow
 \mathscr E_S^{\Mot,\mathrm{face}}.
\]
For every exact support $I$:
\begin{enumerate}[label=\textup{(\roman*)}]
\item all $m=|I|$ highest-face detectors survive and each is a unit multiple
      of the original pure-weight class $[b^{\Mot}]$;
\item the support-$I$ latching class has exact order $N$;
\item the induced map of the exact-support Moore two-cell is, in the derived
      category, the projection of the paired strict-Rees map onto its entire
      non-null summand;
\item the omitted genuine partition components are precisely the
      null-homotopic components of
      \cref{P2M-prop:paired-channel-derived-class}.
\end{enumerate}
Thus the strict Rees model and the face-minimal model define the same
nontrivial exact-support Moore detector class after passage to the derived
mapping category, although only the former remembers every raw partition
homotopy on the nose.
\end{theorem}

\begin{proof}
The paired-diamond proof remains valid after projection to singleton--face
partitions, because all omitted proper-support components simply become zero.
The source differential is still respected: the target image of the full
curvature is its singleton--face part, and the image of $K_I$ is the sum of
the corresponding paired carries.  Exact-support deletion and permutation of
faces are strict on the indexing set $\mathfrak F(I)$.

By \cref{P2M-prop:paired-channel-derived-class}, every retained Moore component is
multiplication by the unit $2$ in $\Z/N$, while every omitted genuine
partition component is null-homotopic.  This proves all four statements.
\end{proof}

To state the minimality precisely we retain the full highest-face detector
ledger, not merely nonvanishing of one projection.

\begin{definition}[Split detector-faithful compression]
\label{P2M-def:split-detector-faithful}
A derived exact-support realization is \emph{split detector-faithful} if its
Moore cohomology carries $m$ specified cyclic order-$N$ quotients, indexed by
$i\in I$, such that the $i$th highest-face projection lands by a unit in the
$i$th quotient and the other highest-face coordinate projections vanish
there.  Equivalently the detector ledger contains a split quotient
$(\Z/N)^I$ with its natural permutation action.
\end{definition}

\begin{theorem}[Moore-detector minimality]
\label{P2M-thm:moore-detector-minimality}
Every split detector-faithful compression at support $I$ requires at least
$m=|I|$ cyclic order-$N$ detector summands.  The face-minimal construction of
\cref{P2M-thm:detector-minimal-tower} has exactly $m$ and is therefore minimal in
this sense.
\end{theorem}

\begin{proof}
A split quotient $(\Z/N)^I$ has minimal number of generators $m$ over
$\Z/N$.  Any Moore cohomology group admitting such a quotient therefore
requires at least $m$ cyclic generators.  The face-minimal target has one
paired highest-face block for each $i\in I$, and
\cref{P2M-prop:paired-channel-derived-class} identifies its detector quotient
with $(\Z/N)^I$ by multiplication by the unit $2$ on every coordinate.
\end{proof}

\begin{corollary}[Tetrahedral compression ledger]
\label{P2M-cor:tetrahedral-compression-ledger}
At support four the exact-support channel counts reduce as
\[
 \boxed{
 \begin{gathered}
 14\ \text{ordered channels}
 \longrightarrow
 7\ \text{paired strict-Rees diamonds}\\
 {}\longrightarrow
 4\ \text{derived detector-minimal face blocks}
 \end{gathered}}
\]
The three paired $2+2$ Pfaffian channels are strict Rees homotopies and are
null in the derived Moore detector sector; the four singleton--triangle
channels are precisely the nontrivial order-$N$ detector coordinates.
\end{corollary}

\paragraph{Minimality scope.}
Detector minimality concerns the exact-support Moore detector, not the full
ambient motivic endomorphism category.


\subsection{Representation-valued compression and cardinality resonance}

The highest-face quotients assemble into the natural permutation module.
This exposes the resonance distinction without changing rank.

Put
\[
 R_N:=\Z/N,
 \qquad
 P_I:=R_N[I]=\bigoplus_{i\in I}R_N e_i,
 \qquad m:=|I|.
\]
The symmetric group $S_I$ acts by permuting the basis.  Write
\begin{equation}
 \epsilon_I:P_I\longrightarrow R_N,
 \qquad
 \epsilon_I\!\left(\sum_i x_ie_i\right)=\sum_i x_i,
\end{equation}
and
\begin{equation}
 \Delta_I:R_N\longrightarrow P_I,
 \qquad
 \Delta_I(1)=\mathbf 1_I:=\sum_{i\in I}e_i.
\end{equation}
Then $\epsilon_I\Delta_I=m$.

\begin{definition}[Permutation-valued Moore detector]
\label{P2M-def:permutation-valued-moore-detector}
The representation-valued exact-support Moore detector is
\begin{equation}
 \boxed{
 \mathcal D_I^{\Mot}
 :=\mathcal M_N[I]\otimes_{\Z}\Z[I],}
\end{equation}
with $S_I$ acting on the second tensor factor.  Its top Moore cohomology is
\[
 H^{2-m}(\mathcal D_I^{\Mot})\simeq P_I.
\]
Under the face-minimal realization of
\cref{P2M-thm:detector-minimal-tower}, the normalized highest-face detector vector
is
\begin{equation}\label{P2M-eq:diagonal-detector-vector}
 \boxed{
 \boldsymbol\beta_I
 =2[b^{\Mot}]\otimes\mathbf 1_I.}
\end{equation}
\end{definition}

\begin{theorem}[Equivariant packaging of the face-minimal detector]
\label{P2M-thm:equivariant-face-packaging}
In the derived exact-support Moore detector category, the product of the
$m$ singleton--highest-face blocks of
\cref{P2M-def:face-minimal-detector-stage} is canonically $S_I$-equivalent to the
single vector-valued complex $\mathcal D_I^{\Mot}$.  Under this equivalence:
\begin{enumerate}[label=\textup{(\roman*)}]
\item the $i$th face projection is coefficient extraction along $e_i$;
\item support relabelling acts by the permutation representation on $P_I$;
\item the exact-support latching class is the diagonal vector
      \eqref{P2M-eq:diagonal-detector-vector};
\item the cyclic-generator rank remains $m$.
\end{enumerate}
Thus one may compress $m$ separate detector \emph{blocks} to one equivariant
object, but not the underlying detector rank when all face projections are
retained.
\end{theorem}

\begin{proof}
Every retained diamond induces multiplication by the same unit $2$ on one
copy of the Moore two-cell, by
\cref{P2M-prop:paired-channel-derived-class}.  Taking their finite direct sum gives
$\mathcal M_N[I]\otimes\Z[I]$.  The permutation coherence of
\cref{P2M-thm:detector-minimal-tower} is exactly the permutation action on the
basis $e_i$.  The latching image has the same class $2[b^{\Mot}]$ in every
coordinate, giving \eqref{P2M-eq:diagonal-detector-vector}.  The last statement is
\cref{P2M-thm:moore-detector-minimality} rewritten in representation-valued form.
\end{proof}

Let
\begin{equation}
 A_I:=\ker(\epsilon_I)\subset P_I
\end{equation}
be the augmentation-zero, or standard, permutation submodule.  We use the
word ``standard'' only for this canonical augmentation module; over modular
coefficient fields it need not be irreducible.

\begin{proposition}[Invariant maps and the mean/standard splitting]
\label{P2M-prop:permutation-module-splitting}
Assume $m\ge2$.
\begin{enumerate}[label=\textup{(\roman*)}]
\item The invariant vectors are
\[
 P_I^{S_I}=R_N\mathbf 1_I.
\]
\item Every $S_I$-equivariant map $P_I\to R_N$ to the trivial module is a
      unique scalar multiple of $\epsilon_I$.
\item The exact sequence
\[
 0\longrightarrow A_I\longrightarrow P_I
 \xrightarrow{\ \epsilon_I\ }R_N\longrightarrow0
\]
      splits $S_I$-equivariantly if and only if $m$ is a unit in $R_N$, or
      equivalently $\gcd(m,N)=1$.
\item In the nonresonant case $\gcd(m,N)=1$ there is a canonical decomposition
\begin{equation}\label{P2M-eq:mean-standard-splitting}
 \boxed{
 P_I=R_N\mathbf 1_I\oplus A_I,}
\end{equation}
      with projectors
\[
 p_{\rm mean}=m^{-1}\Delta_I\epsilon_I,
 \qquad
 p_{\rm rel}=1-p_{\rm mean}.
\]
\end{enumerate}
\end{proposition}

\begin{proof}
Transitivity of $S_I$ on $I$ forces every invariant vector to have equal
coordinates, proving (i).  If $f:P_I\to R_N$ is equivariant, then
$f(e_i)=f(e_j)$ for every $i,j$, so $f=a\epsilon_I$ for a unique $a\in R_N$;
this is (ii).

An equivariant section $s:R_N\to P_I$ must send $1$ to an invariant vector
$a\mathbf 1_I$.  The condition $\epsilon_Is=\mathrm{id}$ is therefore
$am=1$ in $R_N$.  Such $a$ exists exactly when $m$ is a unit modulo $N$,
proving (iii).  When it exists, $a=m^{-1}$ gives the displayed projectors and
(iv).
\end{proof}

\begin{theorem}[Support-cardinality resonance for scalar symmetric detectors]
\label{P2M-thm:cardinality-resonance}
Let $q:P_I\to R_N$ be any $S_I$-equivariant scalar compression.  Applied to
the diagonal exact-support class \eqref{P2M-eq:diagonal-detector-vector}, its image
is
\[
 q(\boldsymbol\beta_I)
 =2a m[b^{\Mot}]
\]
for some $a\in R_N$.  Consequently the largest order obtainable by any
$S_I$-equivariant scalar compression is
\begin{equation}\label{P2M-eq:scalar-compression-max-order}
 \boxed{
 \frac{N}{\gcd(N,m)}.}
\end{equation}
In particular a permutation-invariant rank-one detector preserves exact order
$N$ if and only if $\gcd(N,m)=1$.

Equivalently, if $\ell\mid N$ and $\ell\mid m$, then every equivariant scalar
compression kills the $\ell$-primary component of the diagonal latching
class.
\end{theorem}

\begin{proof}
By \cref{P2M-prop:permutation-module-splitting}, $q=a\epsilon_I$.  Hence
\[
 q(\boldsymbol\beta_I)
 =2a\epsilon_I(\mathbf1_I)[b^{\Mot}]
 =2am[b^{\Mot}].
\]
The class $[b^{\Mot}]$ generates a cyclic group of order $N$ and $2$ is a
unit modulo the odd integer $N$.  Choose an integer lift $\widetilde a$ of
$a$.  Multiplication by $am$ therefore gives a class of order
$N/\gcd(N,\widetilde a m)$.  This is maximized by choosing $a$ a unit,
and the maximum is \eqref{P2M-eq:scalar-compression-max-order}.  The final
statement is the same computation in the $\ell$-primary CRT factor.
\end{proof}

\begin{definition}[Resonant support cardinality]
\label{P2M-def:resonant-support-cardinality}
A support cardinality $m$ is \emph{$N$-resonant} if
\[
 \gcd(m,N)>1.
\]
A prime $\ell\mid N$ is resonant at $m$ when $\ell\mid m$.
\end{definition}

\begin{corollary}[Modular resonance filtration]
\label{P2M-cor:modular-resonance-filtration}
Fix a resonant prime $\ell\mid N$ with $\ell\mid m$ and reduce the detector
module modulo $\ell$.  Put
\[
 P_{I,\ell}:=\F_\ell[I],
 \qquad
 A_{I,\ell}:=\ker\!\left(\epsilon_I:P_{I,\ell}\to\F_\ell\right),
 \qquad
 L_{I,\ell}:=\F_\ell\mathbf1_I.
\]
Then
\begin{equation}
 \boxed{
 0\subset L_{I,\ell}\subset A_{I,\ell}\subset P_{I,\ell}.}
\end{equation}
Thus the diagonal latching mode itself lies in the augmentation-zero module
at a resonant prime.  In particular there is no $S_I$-equivariant separation
of ``global mean'' from ``relative face'' modes in that CRT component.
\end{corollary}

\begin{proof}
Since $\ell\mid m$,
\[
 \epsilon_I(\mathbf1_I)=m=0\quad\text{in }\F_\ell,
\]
so $L_{I,\ell}\subset A_{I,\ell}$.  The nonexistence of an equivariant
mean/relative splitting is also the mod-$\ell$ specialization of
\cref{P2M-prop:permutation-module-splitting}.
\end{proof}

\begin{corollary}[What the standard module does and does not retain]
\label{P2M-cor:standard-module-scope}
If $\gcd(m,N)=1$, the decomposition
\eqref{P2M-eq:mean-standard-splitting} separates the detector data into:
\begin{enumerate}[label=\textup{(\roman*)}]
\item one trivial line carrying the entire diagonal exact-support latching
      class, with normalized mean $m^{-1}\epsilon_I(\boldsymbol\beta_I)
      =2[b^{\Mot}]$ of exact order $N$;
\item the augmentation-zero module $A_I$, which records relative differences
      between highest faces and has zero component on the diagonal latching
      vector.
\end{enumerate}
Therefore the standard module alone cannot replace the Moore latching line,
and the trivial line alone cannot recover the individual face projections.
At resonant cardinalities even this canonical splitting ceases to exist.
\end{corollary}

\begin{theorem}[No representation-rank compression with full face recovery]
\label{P2M-thm:no-rank-compression-face-recovery}
Let $Q$ be an $R_N[S_I]$-module equipped with elements $q_i\in Q$ and
$R_N$-linear functionals $\lambda_i:Q\to R_N$, both permuted naturally by
$S_I$, such that for one unit $u\in R_N^\times$,
\begin{equation}\label{P2M-eq:face-recovery-matrix}
 \lambda_i(q_j)=u\,\delta_{ij}.
\end{equation}
Then the map
\[
 P_I\longrightarrow Q,
 \qquad e_i\longmapsto q_i,
\]
is a split injection of $R_N$-modules.  Hence $Q$ requires at least $m$
cyclic generators.  In particular neither a trivial line nor an
$(m-1)$-generator augmentation-standard module can retain all highest-face
projections.
\end{theorem}

\begin{proof}
Define $\Lambda:Q\to P_I$ by
\[
 \Lambda(x):=u^{-1}\sum_i\lambda_i(x)e_i.
\]
Then \eqref{P2M-eq:face-recovery-matrix} gives
$\Lambda(q_j)=e_j$.  Thus $\Lambda$ is a left inverse to the displayed map
$P_I\to Q$.  The generator bound follows because $P_I\simeq R_N^m$.
\end{proof}

\begin{corollary}[Representation-theoretic endpoint of detector compression]
\label{P2M-cor:representation-compression-endpoint}
There are three distinct endpoints for the exact-support detector:
\begin{enumerate}[label=\textup{(\roman*)}]
\item one may package all $m$ face blocks as the single equivariant object
      $\mathcal D_I^{\Mot}$ without changing rank;
\item away from $N$-resonant cardinalities, one may further split this object
      into one exact-order global latching line plus an augmentation-zero
      relative-face module;
\item if all individual highest-face projections are to remain recoverable,
      the detector rank cannot drop below $m$; if only an invariant scalar
      detector is retained, exact order drops to $N/\gcd(N,m)$ at resonant
      cardinalities.
\end{enumerate}
Thus the $m$-face model is not merely coordinate-minimal: it is the uniform
all-cardinality representation-theoretic endpoint compatible with both full
face recovery and exact order.
\end{corollary}


\subsection{CRT-primewise stratification and the canonical hybrid quotient}

Squarefree \(N\) resolves resonance primewise.  Nonresonant components
retain only the mean line, while resonant components force the full face
module; the CRT hybrid is the minimal equivariant diagonal-class quotient.

Let
\[
 \Pi_N:=\{\ell:\ell\mid N\},\qquad m:=|I|,
\]
and split
\begin{align}
 \Pi_N^{\rm nr}(m)&:=\{\ell\mid N:\ell\nmid m\},\\
 \Pi_N^{\rm res}(m)&:=\{\ell\mid N:\ell\mid m\}.
\end{align}
Write \(e_\ell\in R_N\) for the standard CRT idempotent and put
\begin{equation}
 e_{\rm nr}(m):=\sum_{\ell\in\Pi_N^{\rm nr}(m)}e_\ell,
 \qquad
 e_{\rm res}(m):=\sum_{\ell\in\Pi_N^{\rm res}(m)}e_\ell.
\end{equation}
Thus \(e_{\rm nr}+e_{\rm res}=1\).  Set
\[
 P_{I,\ell}:=e_\ell P_I\simeq\mathbf F_\ell[I],
 \qquad
 L_{I,\ell}:=\mathbf F_\ell\mathbf1_I,
 \qquad
 A_{I,\ell}:=\ker(\epsilon_I:P_{I,\ell}\to\mathbf F_\ell).
\]

\begin{lemma}[Submodule lattice of the natural permutation module]
\label{P2M-lem:permutation-submodule-lattice}
Let \(k\) be a field and \(m\ge2\).  For the natural permutation module
\(k[I]\) of \(S_I\), every \(S_I\)-stable submodule is one of
\[
 0,\qquad k\mathbf1_I,\qquad
 \ker\!\left(\sum:k[I]\to k\right),\qquad k[I],
\]
with repetitions allowed when the characteristic forces one of the displayed
inclusions to coincide.  In particular, if \(\operatorname{char}k\mid m\),
then every nonzero \(S_I\)-stable submodule contains the diagonal line
\(k\mathbf1_I\).
\end{lemma}

\begin{proof}
Let \(W\subseteq k[I]\) be stable.  If every vector of \(W\) has all
coordinates equal, then \(W\subseteq k\mathbf1_I\), giving \(0\) or the
diagonal line.  Otherwise choose \(x\in W\) with \(x_i\ne x_j\).  Then
\[
 x-(ij)x=(x_i-x_j)(e_i-e_j)\in W,
\]
so \(e_i-e_j\in W\).  Conjugating by \(S_I\) shows that \(W\) contains all
coordinate differences and hence the whole augmentation kernel.  Therefore
\(W\) is the augmentation kernel or all of \(k[I]\).  If
\(\operatorname{char}k\mid m\), then
\(\sum_i1=m=0\), so \(k\mathbf1_I\) lies inside the augmentation kernel.
\end{proof}

\begin{theorem}[Primewise maximal kernel preserving the diagonal class]
\label{P2M-thm:primewise-maximal-kernel}
Fix \(\ell\mid N\) and let
\(q_\ell:P_{I,\ell}\to Q_\ell\) be an \(S_I\)-equivariant linear map such
that
\begin{equation}\label{P2M-eq:primewise-diagonal-survival}
 q_\ell(\mathbf1_I)\ne0.
\end{equation}
Then:
\begin{enumerate}[label=\textup{(\roman*)}]
\item if \(\ell\nmid m\), the largest possible kernel is
      \(A_{I,\ell}\), and the smallest possible nonzero image has dimension
      one;
\item if \(\ell\mid m\), one necessarily has \(\ker q_\ell=0\).  Hence
      \(q_\ell\) is injective and
      \(\dim_{\mathbf F_\ell}\operatorname{im}q_\ell\ge m\).
\end{enumerate}
Both bounds are attained.
\end{theorem}

\begin{proof}
When \(\ell\nmid m\), the diagonal line is not contained in
\(A_{I,\ell}\), because its augmentation is \(m\ne0\).  Thus the quotient
by \(A_{I,\ell}\) preserves \(\mathbf1_I\), and its image is the one-
dimensional trivial module.  By
\cref{P2M-lem:permutation-submodule-lattice}, no larger proper stable kernel can
avoid the diagonal line.

When \(\ell\mid m\), the same lemma shows that every nonzero stable kernel
contains \(L_{I,\ell}\), hence contains \(\mathbf1_I\), contradicting
\eqref{P2M-eq:primewise-diagonal-survival}.  Therefore the kernel is zero.  The
identity map attains the bound.
\end{proof}

\begin{definition}[Canonical CRT detector kernel and hybrid quotient]
\label{P2M-def:crt-hybrid-quotient}
Define
\begin{equation}\label{P2M-eq:crt-hybrid-kernel}
 \boxed{
 K_I^{\rm CRT}
 :=e_{\rm nr}(m)A_I\subseteq P_I}
\end{equation}
and
\begin{equation}
 \boxed{
 Q_I^{\rm CRT}
 :=P_I/K_I^{\rm CRT}.}
\end{equation}
By CRT there is a canonical \(S_I\)-equivariant identification
\begin{equation}\label{P2M-eq:crt-hybrid-decomposition}
 \boxed{
 Q_I^{\rm CRT}
 \simeq
 e_{\rm nr}(m)R_N
 \oplus
 e_{\rm res}(m)P_I.}
\end{equation}
The first summand is the nonresonant scalar mean sector; the second is the
full resonant face sector.
\end{definition}

On \(e_{\rm nr}(m)R_N\), multiplication by \(m\) is invertible.  Let
\(m_{\rm nr}^{-1}\in e_{\rm nr}(m)R_N\) be its inverse, characterized by
\begin{equation}\label{P2M-eq:partial-mean-inverse}
 m\,m_{\rm nr}^{-1}=e_{\rm nr}(m).
\end{equation}
Under \eqref{P2M-eq:crt-hybrid-decomposition}, the quotient map is
\begin{equation}\label{P2M-eq:crt-hybrid-map}
 \boxed{
 q_I^{\rm CRT}(x)
 =\bigl(m_{\rm nr}^{-1}e_{\rm nr}\epsilon_I(x),\,
        e_{\rm res}x\bigr).}
\end{equation}

\begin{theorem}[Canonical CRT-hybrid detector and exact-order preservation]
\label{P2M-thm:crt-hybrid-exact-order}
The image of the diagonal motivic latching vector
\(\boldsymbol\beta_I=2[b^{\Mot}]\otimes\mathbf1_I\) under
\eqref{P2M-eq:crt-hybrid-map} is
\begin{equation}\label{P2M-eq:crt-hybrid-latching-vector}
 \boxed{
 \boldsymbol\beta_I^{\rm CRT}
 =\bigl(2e_{\rm nr}[b^{\Mot}],\,
        2e_{\rm res}[b^{\Mot}]\otimes\mathbf1_I\bigr).}
\end{equation}
It is nonzero in every CRT component and therefore has exact order \(N\).
Moreover the construction is canonical under relabelling of \(I\).
\end{theorem}

\begin{proof}
Equation \eqref{P2M-eq:partial-mean-inverse} gives
\[
 m_{\rm nr}^{-1}e_{\rm nr}\epsilon_I(\mathbf1_I)
 =m_{\rm nr}^{-1}e_{\rm nr}m=e_{\rm nr},
\]
which yields \eqref{P2M-eq:crt-hybrid-latching-vector}.  For every
\(\ell\in\Pi_N^{\rm nr}(m)\), the \(\ell\)-component is the nonzero scalar
\(2[b^{\Mot}_\ell]\).  For every
\(\ell\in\Pi_N^{\rm res}(m)\), it is the nonzero diagonal vector
\(2[b^{\Mot}_\ell]\mathbf1_I\).  Since \(N\) is squarefree and \(2\) is a
unit at every \(\ell\mid N\), the combined class has exact order \(N\).
All ingredients---augmentation, the diagonal vector and the CRT idempotents---
are invariant under relabelling.
\end{proof}

\begin{theorem}[Universal minimality of the CRT-hybrid quotient]
\label{P2M-thm:crt-hybrid-universal-minimality}
Call an \(S_I\)-equivariant quotient
\(q:P_I\twoheadrightarrow Q\) \emph{CRT-exact} if
\begin{equation}
 e_\ell q(\mathbf1_I)\ne0
 \qquad\text{for every }\ell\mid N.
\end{equation}
Then \(K_I^{\rm CRT}\) is the largest \(S_I\)-stable submodule of \(P_I\)
which can occur as the kernel of a CRT-exact quotient.  Equivalently, every
CRT-exact quotient admits a canonical further quotient
\begin{equation}\label{P2M-eq:any-detector-to-hybrid}
 \operatorname{im}(q)\twoheadrightarrow Q_I^{\rm CRT}.
\end{equation}
Thus \(Q_I^{\rm CRT}\) is the minimal CRT-exact equivariant detector in the
quotient order.
\end{theorem}

\begin{proof}
Reduce prime by prime.  By
\cref{P2M-thm:primewise-maximal-kernel}, at a nonresonant prime the maximal kernel
compatible with diagonal survival is \(A_{I,\ell}\), while at a resonant prime
it is zero.  Taking their CRT product gives exactly
\(e_{\rm nr}A_I=K_I^{\rm CRT}\).  Hence the kernel of any CRT-exact map is
contained in \(K_I^{\rm CRT}\), which is equivalent to the factor quotient
\eqref{P2M-eq:any-detector-to-hybrid}.
\end{proof}

\begin{corollary}[Exact primewise minimal-rank profile]
\label{P2M-cor:crt-primewise-minimal-rank}
For \(\ell\mid N\), put
\begin{equation}\label{P2M-eq:primewise-rank-profile}
 d_\ell(m):=
 \begin{cases}
  1,&\ell\nmid m,\\
  m,&\ell\mid m.
 \end{cases}
\end{equation}
Then
\begin{equation}
 \boxed{
 \dim_{\mathbf F_\ell}
 \bigl(Q_I^{\rm CRT}\otimes_{R_N}\mathbf F_\ell\bigr)
 =d_\ell(m),}
\end{equation}
and this dimension is minimal among all CRT-exact equivariant quotients.
Consequently the minimum number of \(R_N\)-generators is
\begin{equation}\label{P2M-eq:hybrid-global-generator-number}
 \boxed{
 \mu_{R_N}(Q_I^{\rm CRT})
 =\max_{\ell\mid N}d_\ell(m)
 =
 \begin{cases}
  1,&\gcd(m,N)=1,\\
  m,&\gcd(m,N)>1,
 \end{cases}}
\end{equation}
while the total CRT fibre dimension is
\begin{equation}
 \sum_{\ell\mid N}d_\ell(m)
 =\#\Pi_N^{\rm nr}(m)+m\,\#\Pi_N^{\rm res}(m).
\end{equation}
\end{corollary}

\begin{proof}
The field dimensions are immediate from
\eqref{P2M-eq:crt-hybrid-decomposition} and their minimality is
\cref{P2M-thm:primewise-maximal-kernel}.  A module over the finite product of
fields \(R_N\simeq\prod_{\ell\mid N}\mathbf F_\ell\) needs exactly the
maximum of the dimensions of its field components as generators.  This gives
\eqref{P2M-eq:hybrid-global-generator-number}; summing the dimensions gives the
last formula.
\end{proof}

\begin{corollary}[Automatic allocation of face memory]
\label{P2M-cor:automatic-face-memory}
The hybrid quotient retains precisely the face information forced by exact
order and permutation symmetry:
\begin{enumerate}[label=\textup{(\roman*)}]
\item at \(\ell\nmid m\), all augmentation-zero relative-face information is
      discarded and only the normalized global mean survives;
\item at \(\ell\mid m\), no nonzero equivariant face direction can be
      discarded without killing the diagonal \(\ell\)-primary latching mode,
      so the full permutation module survives.
\end{enumerate}
Hence the hybrid model is strictly weaker than full face recovery at
nonresonant primes, but it is already forced to be fully face-faithful at
resonant primes.
\end{corollary}

\begin{definition}[Primewise hybrid motivic detector complex]
\label{P2M-def:primewise-hybrid-motivic-complex}
Let \(\mathcal M_{N,\ell}[I]\) denote the \(e_\ell\)-summand of the
exact-support Moore detector complex in its idempotent-complete derived
category.  Define
\begin{equation}
 \boxed{
 \mathcal D_I^{\Mot,\rm CRT}
 :=
 \bigoplus_{\ell\in\Pi_N^{\rm nr}(m)}
      \mathcal M_{N,\ell}[I]
 \ \oplus\!
 \bigoplus_{\ell\in\Pi_N^{\rm res}(m)}
      \bigl(\mathcal M_{N,\ell}[I]\otimes\Z[I]\bigr).}
\end{equation}
Its top Moore cohomology is canonically \(Q_I^{\rm CRT}\), and its latching
class is \(\boldsymbol\beta_I^{\rm CRT}\).
\end{definition}

\begin{theorem}[Support-functorial CRT-hybrid detector tower]
\label{P2M-thm:crt-hybrid-support-tower}
Replace, for every exact support \(I\) in the face-minimal motivic tower, its
permutation-valued detector object by
\(\mathcal D_I^{\Mot,\rm CRT}\).  The resulting detector package is strictly
functorial under ordered support deletion and permutation coherent.  If
\(T\subseteq S\), an exact-support block indexed by \(I\subseteq S\) is kept
unchanged when \(I\subseteq T\) and is deleted otherwise.  In particular its
cardinality \(|I|\), resonant-prime set and CRT idempotents never change along
a surviving restriction.

The tower retains an exact-order-\(N\) global latching detector on every exact
support and has the primewise minimal-rank profile
\eqref{P2M-eq:primewise-rank-profile} at each layer.
\end{theorem}

\begin{proof}
Exact-support detector blocks in the Moore--Reedy tower are supported on their
own label set \(I\): proper deletion either preserves all of \(I\), in which
case the block and its cardinality are unchanged, or removes the entire block.
Therefore the CRT decomposition and hybrid quotient commute strictly with
support deletion.  Permutation coherence follows from
\cref{P2M-thm:crt-hybrid-exact-order}.  Exact order and primewise minimality are
\cref{P2M-thm:crt-hybrid-exact-order,P2M-cor:crt-primewise-minimal-rank} support by
support.
\end{proof}


\begin{corollary}[Squarefree three-prime calibration]
\label{P2M-cor:three-prime-hybrid-calibration}
If \(N=pqr\) is the fixed squarefree three-prime coefficient order, let
\[
 r_N(m):=\#\{\ell\in\{p,q,r\}:\ell\mid m\}.
\]
Then the hybrid detector has CRT fibre profile
\[
 d_\ell(m)=1+(m-1)\mathbf1_{\ell\mid m},
\]
so its total CRT fibre dimension is
\begin{equation}
 \boxed{3+(m-1)r_N(m).}
\end{equation}
Its \(R_N\)-generator number is one when \((m,N)=1\) and \(m\) otherwise.
The set of cardinalities on which the rank-one symmetric detector is available
has natural density
\begin{equation}
 \boxed{\frac{\varphi(N)}{N}
 =\prod_{\ell\mid N}\left(1-\frac1\ell\right).}
\end{equation}
For the calibration \(N=105=3\cdot5\cdot7\), this density is
\(48/105=16/35\); the complementary density \(19/35\) consists of resonant
layers on which at least one primary component forces the full rank-\(m\)
face module.
\end{corollary}

\begin{proof}
The fibre formula is \eqref{P2M-eq:primewise-rank-profile} summed over the three
CRT factors.  The generator statement is
\eqref{P2M-eq:hybrid-global-generator-number}.  Nonresonant cardinalities are
exactly the integers coprime to \(N\), whose residue classes modulo \(N\)
occupy \(\varphi(N)\) of the \(N\) classes.
\end{proof}

\paragraph{Two minimality notions.}
\(Q_I^{\rm CRT}\) is minimal for the diagonal class primewise; full
highest-face recovery instead forces \(P_I\).


\subsection{Geometric averaging resonance}

The same resonance is geometric.  The highest faces form a finite
constant cover of the root stack; its trace is the all-ones correspondence,
and failure to divide by the covering degree is exactly the CRT resonance.

Keep the cyclotomic root stack \(\mathfrak X^{\Mot}_{N,d}\) and its
pure-weight localization object \(\mathbb C^{\Mot}_{w,d}\) from the
arity-three input.  The latter carries the genuine Moore pair
\((b^{\Mot},c^{\Mot})\) after the marked root correspondence is applied.
For a finite support \(I\), \(|I|=m\), form the finite face cover of the
actual root stack
\begin{equation}
 \boxed{
 \varpi_I:\mathfrak F_I^{\Mot}
 :=\coprod_{i\in I}\mathfrak X^{\Mot}_{N,d}
 \longrightarrow \mathfrak X^{\Mot}_{N,d}.}
\end{equation}
It is the trivial finite \'etale cover of degree \(m\), but its sheets are
canonically indexed by the highest faces \(I\setminus\{i\}\).  Applying
the same pure-weight localization functor on every sheet gives the face object
\[
 \mathcal P_I^{\Mot}
 :=\mathbb C^{\Mot}_{w,d}\otimes_{R_N}P_I
 \simeq \bigoplus_{i\in I}\mathbb C^{\Mot}_{w,d}.
\]
The geometric pullback and transfer along \(\varpi_I\) induce the diagonal
and summation maps on this sheet decomposition and retain the full
\(S_I\)-action.  Write
\begin{equation}
 \boxed{
 \Theta_I:=\varpi_I^*\varpi_{I,*}
 \in\operatorname{End}(\mathcal P_I^{\Mot}).}
\end{equation}
Here \(\varpi_{I,*}\) is summation over the face sheets and
\(\varpi_I^*\) is diagonal pullback.

\begin{theorem}[Finite-face trace relation]
\label{P2M-thm:finite-face-trace-relation}
Under the identification
\(\mathcal P_I^{\Mot}\simeq\mathbb C^{\Mot}_{w,d}\otimes P_I\), one has
\begin{equation}\label{P2M-eq:theta-delta-epsilon}
 \boxed{
 \Theta_I=\operatorname{id}_{\mathbb C^{\Mot}_{w,d}}\otimes
           \Delta_I\epsilon_I,}
\end{equation}
where \(\epsilon_I:P_I\to R_N\) is augmentation and
\(\Delta_I:R_N\to P_I\) is the diagonal map.  Equivalently, the finite
cover satisfies the two pull--push identities
\begin{equation}\label{P2M-eq:face-cover-two-trace-identities}
 \boxed{
 \varpi_{I,*}\varpi_I^*=m\operatorname{id},
 \qquad
 \Theta_I^2=m\Theta_I.}
\end{equation}
The construction is canonical under permutations of \(I\).
\end{theorem}

\begin{proof}
For a tuple \(x=(x_i)_{i\in I}\), transfer gives
\(\varpi_{I,*}(x)=\sum_i x_i\), and pullback returns this sum on every
sheet.  Thus \(\Theta_I(x)=\Delta_I\epsilon_I(x)\), proving
\eqref{P2M-eq:theta-delta-epsilon}.  Since
\(\epsilon_I\Delta_I=m\operatorname{id}\),
\[
 (\Delta_I\epsilon_I)^2
 =\Delta_I(\epsilon_I\Delta_I)\epsilon_I
 =m\Delta_I\epsilon_I.
\]
All maps are defined from the unlabeled finite cover together with its sheet
set, so relabelling only permutes the sheets.
\end{proof}

The second relation in \eqref{P2M-eq:face-cover-two-trace-identities} is the geometric origin of the
mean/relative dichotomy.  It also explains why the resonant case is not simply
a missing choice of normalization: the trace correspondence itself changes
homological type after reduction to a resonant coefficient factor.

\begin{theorem}[Failure-of-averaging correspondence]
\label{P2M-thm:failure-of-averaging-correspondence}
Fix \(\ell\mid N\) and let \(\Theta_{I,\ell}=e_\ell\Theta_I\).  On the
\(\ell\)-primary face detector one has:
\begin{enumerate}[label=\textup{(\roman*)}]
\item if \(\ell\nmid m\), then
\begin{equation}
 \boxed{
 E_{I,\ell}^{\rm mean}:=m^{-1}\Theta_{I,\ell}}
\end{equation}
is an \(S_I\)-equivariant idempotent correspondence with image the diagonal
line \(L_{I,\ell}\) and kernel the augmentation module \(A_{I,\ell}\);
\item if \(\ell\mid m\), then
\begin{equation}
 \boxed{\Theta_{I,\ell}^2=0,}
\end{equation}
while \(\Theta_{I,\ell}
e0\), and
\begin{equation}
 \boxed{
 \operatorname{im}\Theta_{I,\ell}=L_{I,\ell}
 \subset A_{I,\ell}=\ker\Theta_{I,\ell}.}
\end{equation}
Thus the mean correspondence degenerates from a projector to a nonzero
square-zero correspondence precisely at the resonant primes.
\end{enumerate}
\end{theorem}

\begin{proof}
The first statement follows from
\eqref{P2M-eq:face-cover-two-trace-identities} after dividing by the unit \(m\) in
\(\mathbf F_\ell\).  Formula \eqref{P2M-eq:theta-delta-epsilon} identifies its
image and kernel with the diagonal and augmentation modules.

If \(\ell\mid m\), then
\eqref{P2M-eq:face-cover-two-trace-identities} reduces to
\(\Theta_{I,\ell}^2=0\).  The map is nonzero because a basis face vector is
sent to the nonzero diagonal vector.  Its image is therefore the diagonal
line.  Its kernel is the augmentation module by
\eqref{P2M-eq:theta-delta-epsilon}, and the inclusion of the image in the kernel
is exactly the equality \(m=0\) in \(\mathbf F_\ell\).
\end{proof}

\begin{definition}[Geometric resonance coefficient stratum]
\label{P2M-def:geometric-resonance-stratum}
On the coefficient spectrum
\(\Spec R_N=\coprod_{\ell\mid N}\Spec\mathbf F_\ell\), define the open-and-
closed strata
\begin{equation}
 \boxed{
 \mathfrak B_I^{\rm nr}:=\Spec(e_{\rm nr}(m)R_N),
 \qquad
 \mathfrak B_I^{\rm res}:=\Spec(e_{\rm res}(m)R_N).}
\end{equation}
The second is called the \emph{geometric averaging-resonance stratum}.  It is
the maximal coefficient stratum on which the finite-face trace
correspondence \(\Theta_I\) is square-zero rather than normalizable to a mean
projector.
\end{definition}

\begin{corollary}[Resonance is the failure locus of finite-face averaging]
\label{P2M-cor:resonance-is-averaging-failure}
The following are equivalent for a CRT prime \(\ell\mid N\):
\begin{enumerate}[label=\textup{(\roman*)}]
\item \(\ell\mid m\);
\item the \(\ell\)-point lies in \(\mathfrak B_I^{\rm res}\);
\item the degree-\(m\) trace correspondence on the face cover cannot be
      normalized to an idempotent mean projector over \(\mathbf F_\ell\);
\item the trace correspondence is nonzero square-zero on the
      \(\ell\)-primary face motive;
\item the diagonal latching line lies inside the relative-face augmentation
      module.
\end{enumerate}
Hence the support-cardinality resonance of
\cref{P2M-thm:cardinality-resonance} is a geometric failure-of-averaging
phenomenon for the finite highest-face cover.
\end{corollary}

\begin{proof}
The equivalence of \textup{(i)}--\textup{(iv)} is
\cref{P2M-thm:failure-of-averaging-correspondence}; \textup{(i)} is equivalent to
\textup{(v)} because \(\epsilon_I(\mathbf1_I)=m\).
\end{proof}

The preceding theorem can be combined with the identity correspondence on the
resonant coefficient summand.  This produces one idempotent which exactly
realizes the hybrid quotient of \cref{P2M-def:crt-hybrid-quotient}.

\begin{definition}[Geometric CRT-hybrid projector]
\label{P2M-def:geometric-crt-hybrid-projector}
Let \(m_{\rm nr}^{-1}\in e_{\rm nr}(m)R_N\) be as in
\eqref{P2M-eq:partial-mean-inverse}.  Define
\begin{equation}\label{P2M-eq:geometric-crt-hybrid-projector}
 \boxed{
 \Pi_I^{\rm geom}
 :=e_{\rm nr}(m)m_{\rm nr}^{-1}\Theta_I
   +e_{\rm res}(m)\operatorname{id}_{\mathcal P_I^{\Mot}}.}
\end{equation}
\end{definition}

\begin{theorem}[The CRT-hybrid detector is a motivic direct summand]
\label{P2M-thm:crt-hybrid-motivic-projector}
The correspondence \(\Pi_I^{\rm geom}\) is an \(S_I\)-equivariant
idempotent.  In the idempotent-complete category of pure-weight motivic
correspondences with \(R_N\)-coefficients,
\begin{align}
 \ker\Pi_I^{\rm geom}
 &\simeq
 \mathbb C^{\Mot}_{w,d}\otimes e_{\rm nr}(m)A_I,
 \label{P2M-eq:geom-hybrid-kernel}\\
 \operatorname{Im}\Pi_I^{\rm geom}
 &\simeq
 \mathbb C^{\Mot}_{w,d}\otimes
 \bigl(e_{\rm nr}(m)R_N\oplus e_{\rm res}(m)P_I\bigr)
 \simeq
 \mathbb C^{\Mot}_{w,d}\otimes Q_I^{\rm CRT}.
 \label{P2M-eq:geom-hybrid-image}
\end{align}
Under \eqref{P2M-eq:geom-hybrid-image}, the face-coefficient vector of the
projected motivic latching class is exactly
\(\boldsymbol\beta_I^{\rm CRT}\) of
\eqref{P2M-eq:crt-hybrid-latching-vector}.  Restricting the localization object to
the genuine Moore sector generated by \(b^{\Mot},c^{\Mot}\) recovers the
hybrid detector complex \(\mathcal D_I^{\Mot,\rm CRT}\) of
\cref{P2M-def:primewise-hybrid-motivic-complex}.  Thus the canonical
primewise-minimal CRT detector is not merely a module quotient: it is the
face-coefficient part of an actual finite-cover motivic correspondence direct
summand.
\end{theorem}

\begin{proof}
On the nonresonant coefficient summand,
\(m_{\rm nr}^{-1}\Theta_I\) is the mean projector of
\cref{P2M-thm:failure-of-averaging-correspondence}; on the resonant summand the
second term of \eqref{P2M-eq:geometric-crt-hybrid-projector} is the identity.
The two CRT summands are orthogonal, so their sum is idempotent.  Its kernel is
the nonresonant augmentation module and its image is the nonresonant diagonal
line plus the full resonant permutation module.  These are precisely
\eqref{P2M-eq:crt-hybrid-kernel} and \eqref{P2M-eq:crt-hybrid-decomposition}.
Applying the projector to
\(2[b^{\Mot}]\otimes\mathbf1_I\) gives
\eqref{P2M-eq:crt-hybrid-latching-vector}.
\end{proof}

\begin{corollary}[Geometric universal property of the hybrid detector]
\label{P2M-cor:geometric-hybrid-universal-property}
Among \(S_I\)-equivariant direct-summand quotients of the face motive which
preserve the diagonal latching class in every CRT component,
\(\operatorname{Im}\Pi_I^{\rm geom}\) is minimal.  Equivalently,
\(\Pi_I^{\rm geom}\) kills the largest motivically split face summand which
can be discarded without losing a primary component of the exact-order
latching class.
\end{corollary}

\begin{proof}
After factoring out the common pure-weight localization object, the kernel of
\(\Pi_I^{\rm geom}\) is the face coefficient submodule
\(K_I^{\rm CRT}\) by \eqref{P2M-eq:geom-hybrid-kernel}.  Its maximality is
\cref{P2M-thm:crt-hybrid-universal-minimality}.
\end{proof}

\begin{theorem}[Periodic motivic resonance stratification]
\label{P2M-thm:periodic-motivic-resonance-stratification}
The geometric resonance type of a support layer depends only on
\begin{equation}
 \boxed{d_I:=\gcd(|I|,N).}
\end{equation}
In particular:
\begin{enumerate}[label=\textup{(\roman*)}]
\item it is periodic in \(|I|\) with period \(N\);
\item the resonant coefficient stratum is the union of the CRT points
      \(\ell\mid d_I\);
\item the geometric projector \(\Pi_I^{\rm geom}\) has face-coefficient
      rank \(1\) at \(\ell\nmid d_I\) and face-coefficient rank
      \(|I|\) at \(\ell\mid d_I\);
\item for a prescribed subset \(R\subseteq\Pi_N\), the set of support
      cardinalities for which the resonant-prime set is exactly \(R\) has
      natural density
\begin{equation}\label{P2M-eq:exact-resonance-type-density}
 \boxed{
 \prod_{\ell\in R}\frac1\ell
 \prod_{\ell\in\Pi_N\setminus R}
       \left(1-\frac1\ell\right).}
\end{equation}
\end{enumerate}
\end{theorem}

\begin{proof}
For squarefree \(N\), the set of primes dividing both \(|I|\) and \(N\) is
exactly the set of prime divisors of \(d_I\), giving the first three claims.
The resonance condition depends only on \(|I|\bmod N\), so it is periodic.
For the density, the conditions \(\ell\mid |I|\) and
\(\ell\nmid |I|\) are independent under the Chinese remainder theorem over
one period modulo \(N\), giving the product in
\eqref{P2M-eq:exact-resonance-type-density}.
\end{proof}

\paragraph{Geometric scope.}
Resonance is realized by a pull--push correspondence over the root-stack
motive; it is a coefficient-spectrum stratification, not a divisor in the
root stack.


\section{Karoubi decomposition and jetwise constancy}

Singleton--highest-face diamonds carry the Moore detector; genuine
partition diamonds only strictify transfer--Rees flatness.  Their canonical
null-homotopies split these roles already on the designated cellular mapping
complex.

For the Moore two-cell
\[
 \mathcal M_N[I]=[\Z K_I\xrightarrow{N}\Z\Omega_I]
\]
let
\begin{equation}
 \boxed{
 J_I(\Omega_I)=K_I,\qquad J_I(K_I)=0.}
\end{equation}
As a degree-minus-one element of the mapping complex
\(\operatorname{End}^\bullet(\mathcal M_N[I])\), it satisfies
\begin{equation}\label{P2M-eq:universal-moore-homotopy-diff}
 \boxed{d_{\operatorname{End}}J_I=N\operatorname{id}_{\mathcal M_N[I]}.}
\end{equation}
This is the chain-level source of both the surviving Moore torsion and the
contractibility of the genuine partition channels.

Write
\[
 \mathfrak P_{\rm face}(I)
 :=\bigl\{\{\{i\},I\setminus\{i\}\}:i\in I\bigr\},
 \qquad
 \mathfrak P_{\rm gen}(I)
 :=\mathfrak P(I)\setminus\mathfrak P_{\rm face}(I).
\]
Thus
\begin{equation}
 \#\mathfrak P_{\rm gen}(I)=2^{m-1}-1-m,
 \qquad m=|I|.
\end{equation}

\begin{definition}[Cellular exact-support channel mapping complex]
\label{P2M-def:cellular-channel-mapping-complex}
Inside the product of exact-support mapping complexes obtained by restricting
each paired diamond to its designated shifted root Moore pair and identifying
that pair with \(\mathcal M_N[I]\), let
\(\mathfrak C_I^{\rm cell}\) be the subcomplex generated as follows.
For a face partition \(P_i\in\mathfrak P_{\rm face}(I)\), put
\[
 \beta_i:=2\operatorname{id}_{\mathcal M_N[I]},
 \qquad
 \gamma_i:=2J_I,
 \qquad
 d\gamma_i=N\beta_i.
\]
For a genuine partition \(P\in\mathfrak P_{\rm gen}(I)\), put
\[
 \beta_P:=2N\operatorname{id}_{\mathcal M_N[I]},
 \qquad
 h_P:=2J_I,
 \qquad
 dh_P=\beta_P.
\]
Each symbol is placed in its own orthogonal paired channel block.  Define
\begin{align}
 \mathfrak D_I^{\rm face}
 &:={\bigoplus}_{i\in I}
       [\Z\gamma_i\xrightarrow{\ N\ }\Z\beta_i],
 \\
 \mathfrak H_I^{\rm Rees}
 &:={\bigoplus}_{P\in\mathfrak P_{\rm gen}(I)}
       [\Z h_P\xrightarrow{\ 1\ }\Z\beta_P].
\end{align}
The first is the \emph{essential face detector complex}; the second is the
\emph{strict Rees-homotopy complex}.
\end{definition}

\begin{theorem}[Strict cellular detector--homotopy decomposition]
\label{P2M-thm:cellular-detector-homotopy-decomposition}
For every finite exact support \(I\), there is a canonical isomorphism of finite
chain complexes
\begin{equation}\label{P2M-eq:cellular-strict-decomposition}
 \boxed{
 \mathfrak C_I^{\rm cell}
 \cong
 \mathfrak D_I^{\rm face}\oplus\mathfrak H_I^{\rm Rees}.}
\end{equation}
Moreover:
\begin{enumerate}[label=\textup{(\roman*)}]
\item \(\mathfrak H_I^{\rm Rees}\) is contractible, with blockwise contraction
      \(s_P(\beta_P)=h_P\), \(s_P(h_P)=0\);
\item \(H^0(\mathfrak D_I^{\rm face})\simeq(\Z/N)^I\), and all other
      cohomology of \(\mathfrak D_I^{\rm face}\) vanishes;
\item projection onto \(\mathfrak D_I^{\rm face}\) is a strong deformation
      retract of \(\mathfrak C_I^{\rm cell}\);
\item the decomposition is canonical under permutations of \(I\) and is strict
      under support deletion of the exact-support block.
\end{enumerate}
Hence the exponential family of genuine partition channels is a contractible
strictification package at the cellular mapping-complex level; the entire
nontrivial exact-support Moore cohomology is carried by the \(m\) highest-face
cells.
\end{theorem}

\begin{proof}
The paired-channel calculation of
\cref{P2M-prop:paired-channel-derived-class} identifies every face channel with
multiplication by \(2\) and every genuine partition channel with multiplication
by \(2N\) on \(\mathcal M_N[I]\).  Equation
\eqref{P2M-eq:universal-moore-homotopy-diff} therefore gives
\[
 d(2J_I)=N(2\operatorname{id})
 \quad\text{on a face channel},
 \qquad
 d(2J_I)=2N\operatorname{id}
 \quad\text{on a genuine partition channel}.
\]
Because distinct paired diamonds are orthogonal, their cellular mapping
coordinates form the direct sum
\eqref{P2M-eq:cellular-strict-decomposition} on the nose.  Each genuine partition
summand has differential one in the displayed basis, so the stated \(s_P\)
contracts it.  Each face summand is the two-term Moore complex with differential
\(N\), hence has only \(H^0\simeq\Z/N\).  Taking the direct sum over the \(m\)
faces proves \textup{(ii)}.  The direct-sum projection and inclusion, together
with the blockwise contractions, give the strong deformation retract.
Permutation and deletion naturality follow because both
\(\mathfrak P_{\rm face}(I)\) and \(\mathfrak P_{\rm gen}(I)\) are intrinsic
subsets of the unordered bipartition set.
\end{proof}

\begin{corollary}[Cellular Morita reduction of the strict Rees presentation]
\label{P2M-cor:cellular-morita-reduction}
In the derived category of finite exact-support mapping complexes,
\begin{equation}
 \boxed{
 \mathfrak C_I^{\rm cell}\simeq\mathfrak D_I^{\rm face}.}
\end{equation}
The equivalence is represented by an explicit support-functorial strong
deformation retract, not merely by equality of cohomology groups.
\end{corollary}

\begin{remark}[Cellular reduction versus presentation Morita equivalence]
\label{P2M-rem:cellular-versus-presentation-morita}
\Cref{P2M-cor:cellular-morita-reduction} concerns only the displayed finite
exact-support mapping complexes.  It neither objectifies the Boolean shifts
nor compares complete dg presentations.  Those two tasks are performed,
respectively, by \cref{P2M-thm:role-separated-objectification} and
\cref{P2M-thm:marked-morita-presentation-independence}; neither conclusion is
deduced from the cellular deformation retract.
\end{remark}

The geometric averaging projector of
\cref{P2M-thm:crt-hybrid-motivic-projector} now refines the surviving face
complex itself.  Canonically
\[
 \mathfrak D_I^{\rm face}
 \cong
 \bigl(\mathcal M_N[I]\otimes_{\Z}\Z[I]\bigr)[2-m],
 \qquad
 H^0(\mathfrak D_I^{\rm face})\simeq P_I=R_N[I].
\]
Indeed the two terms of $\mathcal M_N[I]$ lie in degrees $1-m$ and
$2-m$, so the common cohomological shift $[2-m]$ places them in degrees
$-1$ and $0$, respectively, matching the degrees of the cellular generators
$\gamma_i$ and $\beta_i$.
The finite-face correspondence \(\Pi_I^{\rm geom}\) geometrically realizes on
this Moore detector the coefficient idempotent
\begin{equation}
 \boxed{
 p_I^{\rm geom}
 :=e_{\rm nr}(m)m_{\rm nr}^{-1}\Delta_I\epsilon_I
   +e_{\rm res}(m)\operatorname{id}_{P_I}}
\end{equation}
in the derived endomorphism ring of the face complex.  Thus the following
splitting is a derived shadow of the actual motivic projector, not an
independently chosen coefficient decomposition.

\begin{definition}[Essential CRT detector and relative-face memory]
\label{P2M-def:essential-crt-relative-face}
In the idempotent-complete derived category define
\begin{align}
 \mathfrak D_I^{\rm CRT}
 &:=\operatorname{Im}(p_I^{\rm geom}),
 \\
 \mathfrak R_I^{\rm rel}
 &:=\operatorname{Im}(1-p_I^{\rm geom}).
\end{align}
The first is the \emph{primewise-essential global detector}; the second records
exactly the nonresonant augmentation-zero face memory which is unnecessary for
global CRT-exact detection but is needed for full face recovery.
\end{definition}

\begin{theorem}[Geometric Karoubi decomposition of the cellular detector]
\label{P2M-thm:geometric-karoubi-cellular-decomposition}
There is a canonical permutation-equivariant decomposition
\begin{equation}\label{P2M-eq:karoubi-face-decomposition}
 \boxed{
 \mathfrak D_I^{\rm face}
 \simeq
 \mathfrak D_I^{\rm CRT}\oplus\mathfrak R_I^{\rm rel}}
\end{equation}
in the idempotent-complete motivic derived category.  Its top Moore cohomology is
\begin{align}
 H^0(\mathfrak D_I^{\rm CRT})&\simeq Q_I^{\rm CRT},
 \\
 H^0(\mathfrak R_I^{\rm rel})&\simeq e_{\rm nr}(m)A_I.
 \label{P2M-eq:relative-face-cohomology}
\end{align}
In particular \(\mathfrak R_I^{\rm rel}=0\) on every resonant CRT summand,
where the primewise-essential detector is already fully face-faithful.
Writing \(\mathfrak C_I^{\rm cell,Kar}\) for the image of the cellular
complex in the idempotent-complete derived category, combining with
\cref{P2M-thm:cellular-detector-homotopy-decomposition} gives the cellular
three-way decomposition
\begin{equation}
 \boxed{
 \mathfrak C_I^{\rm cell,Kar}
 \simeq
 \mathfrak D_I^{\rm CRT}
 \oplus
 \mathfrak R_I^{\rm rel}
 \oplus
 \mathfrak H_I^{\rm Rees}.}
\end{equation}
The last summand is contractible; the middle summand is generally not.
\end{theorem}

\begin{proof}
The two idempotents
\(\Pi_I^{\rm geom}\) and \(1-\Pi_I^{\rm geom}\) are orthogonal and sum to
the identity, so idempotent completion gives
\eqref{P2M-eq:karoubi-face-decomposition}.  By
\cref{P2M-thm:crt-hybrid-motivic-projector}, the image of the first projector
has face-coefficient module \(Q_I^{\rm CRT}\).  The kernel/image computation
of the same theorem identifies the complementary coefficient module with
\(e_{\rm nr}A_I\), proving
\eqref{P2M-eq:relative-face-cohomology}.  On a resonant CRT factor
\(\Pi_I^{\rm geom}=\operatorname{id}\), so the complement vanishes.  Finally
insert \eqref{P2M-eq:karoubi-face-decomposition} into the strict cellular splitting
\eqref{P2M-eq:cellular-strict-decomposition}.
\end{proof}

\begin{corollary}[Two notions of essential detector]
\label{P2M-cor:two-essential-detector-notions}
The decomposition separates two legitimate endpoint requirements:
\begin{enumerate}[label=\textup{(\roman*)}]
\item for full highest-face recovery the essential noncontractible cellular
      object is
      \(\mathfrak D_I^{\rm face}
        =\mathfrak D_I^{\rm CRT}\oplus\mathfrak R_I^{\rm rel}\);
\item for global CRT-exact latching detection the canonical minimal object is
      only \(\mathfrak D_I^{\rm CRT}\).
\end{enumerate}
The Rees-homotopy complex is dispensable in either derived endpoint, but it is
required to retain strict raw transfer--Rees flatness before passage to the
derived mapping category.
\end{corollary}


\subsection{Jetwise constancy of Hall, detector, CRT and Karoubi readouts}
\label{P2M-sec:v37-readout-jet-constancy}

The preceding tower theorems show that jet successor is the identity on the
underlying Moore--Reedy object, on every reduced history apex and on the
normalized diagonal bar complex.  The remaining constructions of the paper
are downstream readouts of those fixed structural objects.  We now record the
corresponding compatibility once, rather than repeat the same observation in
each compression section.

\begin{definition}[Jet-decorated downstream readout package]
\label{P2M-def:v37-jet-readout-package}
Fix odd squarefree $N$, an exact ordered support $I$, a distinguished label
$i_0\in I$, and a finite jet ceiling $M\ge3$.  For $3\le n\le M$ let
\begin{equation}\label{P2M-eq:v37-jet-readout-package}
 \boxed{
 \mathfrak Q_{I,i_0,n}^{\Mot,\mathrm{jet}}
 :=\Bigl(
   \mathscr H_{I,i_0;N}^{\mathrm{RB,tr}},\
   \mathcal D_I^{\Mot},\
   Q_I^{\rm CRT},\
   \mathfrak C_I^{\rm cell,Kar};\
   \mathfrak R^{\rm Conf}_{I,\le n}
 \Bigr).}
\end{equation}
The semicolon is intentional: the first four entries are the already constructed
Hall, detector, CRT, and Karoubi readouts.  The last entry is only the finite
Confluent marking ledger.  No direct sum, product, new motive or new detector
object is asserted by this notation.

The jet successor
\begin{equation}\label{P2M-eq:v37-readout-jet-successor}
 \mathbf j^{\rm rd}_{I,n}:
 \mathfrak Q_{I,i_0,n}^{\Mot,\mathrm{jet}}
 \longrightarrow
 \mathfrak Q_{I,i_0,n+1}^{\Mot,\mathrm{jet}}
\end{equation}
is the identity on the first four entries and the canonical inclusion
$\mathfrak R^{\rm Conf}_{I,\le n}\hookrightarrow
\mathfrak R^{\rm Conf}_{I,\le n+1}$ on the ledger.
\end{definition}

\begin{theorem}[Jetwise constancy of downstream motivic readouts]
\label{P2M-thm:v37-readout-jet-constancy}
With the preceding notation, all downstream readouts commute with jet
successor.  More precisely:
\begin{enumerate}[label=\textup{(\roman*)}]
\item The pointed Hall-faithful map is independent of $n$:
\begin{equation}\label{P2M-eq:v37-hall-jet-constant}
 \Real_{I,i_0,n}^{\mathrm{Hall,tr}}
 =\Real_{I,i_0}^{\mathrm{Hall,tr}}:
 W_0K_I^{\mathrm{St}}
 \hookrightarrow
 \mathscr H_{I,i_0;N}^{\mathrm{RB,tr}}.
\end{equation}
Hence the pointed Hall basis, tagged leading operation trees and
adjacent-refinement differential are unchanged along the jet axis.

\item The face-minimal detector is independent of $n$.  In particular the
permutation-valued Moore detector remains $\mathcal D_I^{\Mot}$ and its
normalized diagonal detector vector remains
\begin{equation}\label{P2M-eq:v37-face-vector-jet-constant}
 \boldsymbol\beta_I=2[b^{\Mot}]\otimes\mathbf1_I.
\end{equation}
No additional highest-face block is introduced by a jet successor.

\item The CRT-hybrid projector and quotient are independent of $n$:
\begin{equation}\label{P2M-eq:v37-crt-projector-jet-constant}
 p_{I,n}^{\rm geom}=p_I^{\rm geom},
 \qquad
 Q_{I,n}^{\rm CRT}=Q_I^{\rm CRT}.
\end{equation}
Thus primewise resonance, the minimal CRT-exact rank profile and exact-order
preservation depend on the coefficient order and support cardinality, not on
jet order.

\item The cellular strong deformation retract and the Karoubi splitting are
independent of $n$.  At every finite jet level one has the same decompositions
\begin{equation}\label{P2M-eq:v37-karoubi-jet-constant}
 \mathfrak C_I^{\rm cell}
 \simeq
 \mathfrak D_I^{\rm face}\oplus\mathfrak H_I^{\rm Rees},
 \qquad
 \mathfrak D_I^{\rm face}
 \simeq
 \mathfrak D_I^{\rm CRT}\oplus\mathfrak R_I^{\rm rel}.
\end{equation}
The Rees-homotopy summand remains contractible and no jet step changes the
primewise-essential or relative-face Karoubi factors.

\item Order-preserving relabellings, support permutations and every strict
support deletion for which the exact-support readout survives commute with
\eqref{P2M-eq:v37-readout-jet-successor}.  In particular all mixed
support--jet squares remain strict after applying any of the readouts above.
\end{enumerate}
Equivalently, if $\operatorname{Read}_I$ denotes any one of the Hall, face,
CRT or Karoubi constructions, then on the finite marked tower
\begin{equation}\label{P2M-eq:v37-readout-interchange}
 \boxed{
 \operatorname{Read}_I\circ\jmath_{I,n}
 =\mathbf j^{\rm rd}_{I,n}\circ\operatorname{Read}_I,}
\end{equation}
where equality means identity on the underlying readout object together with
inclusion of the finite ledger.
\end{theorem}

\begin{proof}
By \cref{P2M-thm:v34-support-jet-tower}, the jet successor is the identity on
$\mathbb M_{I,N}^{\Mot,\MR}$ and changes only the Confluent ledger.  The Hall
completion of \cref{P2M-thm:v14-hall-faithful-transfer} is built from the fixed
preceding Hall complex, the fixed root/Bass evaluation and the transfer scalar
$N^{|I|-2}$, so none of its chain-level data depends on $n$.  This proves
\textup{(i)}.

The face detector of \cref{P2M-thm:equivariant-face-packaging} is obtained from
the same exact-support Moore pair and the permutation module $R_N[I]$;
therefore its diagonal vector is unchanged, proving \textup{(ii)}.  The
CRT-hybrid projector of \cref{P2M-thm:crt-hybrid-motivic-projector} depends only on
the coefficient idempotents and the scalar $m=|I|$ in
$\epsilon_I\Delta_I=m$, so it is independent of jet order.  This proves
\textup{(iii)}.

Finally the cellular complexes and contractions in
\cref{P2M-thm:cellular-detector-homotopy-decomposition} are built blockwise from
the universal Moore homotopy $J_I$, while the Karoubi projector is precisely
the same $p_I^{\rm geom}$.  Hence both splittings are fixed along the jet
axis, proving \textup{(iv)}.  The support naturality in \textup{(v)} follows
from the corresponding deletion/permutation naturality of the underlying
Moore--Reedy stage together with the support naturality of the Confluent
ledger defined in \cref{P2M-thm:v33-imported-confluent-jet-interface}.
\end{proof}

\begin{corollary}[Jet-independent detector resonance and minimal rank]
\label{P2M-cor:v37-jet-independent-resonance}
Let $m=|I|$ and let $\ell\mid N$.  Along every finite jet stage on the fixed
support $I$, the resonance condition remains
\[
 \ell\mid m,
\]
and the primewise minimal CRT-exact detector rank remains the same number
$d_\ell(m)$ of \cref{P2M-cor:crt-primewise-minimal-rank}.  Consequently
\[
 \mu_{R_N}(Q_I^{\rm CRT})
 =\max_{\ell\mid N}d_\ell(m)
\]
is constant along the jet axis.
\end{corollary}

\begin{corollary}[No downstream readout growth along the jet axis]
\label{P2M-cor:v37-no-readout-growth}
For fixed exact support $I$ and squarefree $N$, increasing the finite jet
ceiling does not create additional Hall operation-tree factors, highest-face
Moore blocks, CRT detector factors, relative-face Karoubi summands or
Rees-homotopy blocks.  A new Confluent jet class may change the marked source
detector, primitive-line, carry, filler or pointing status, but every such
status is read on the already existing finite readout package.
\end{corollary}

\begin{warning}[Readout constancy is not independent higher-operation provenance]
\label{P2M-warn:v37-readout-not-new-provenance}
The theorem says that the constructed Hall, detector, CRT and Karoubi
readouts remain valid when the independent Confluent jet ledger is enlarged.
It does not assert that every newly detected higher arithmetic operation is
represented by a new geometric correspondence inside those readouts.  Such a
correspondence remains an independent provenance statement.  In particular
jetwise constancy must not be read as local fullness of the motivic operation
category.
\end{warning}

\begin{corollary}[End-to-end finite jet compatibility]
\label{P2M-cor:v37-end-to-end-jet-compatibility}
Within the odd-squarefree all-support scope, the complete constructed pipeline
\[
 \boxed{
 \begin{gathered}
 \text{Reedy reconstruction}
 \longrightarrow
 \text{reduced frontier history}\\
 \longrightarrow
 \text{secondary bar transgression}
 \longrightarrow
 \text{Hall/face/CRT/Karoubi readout}
 \end{gathered}}
\]
commutes with every finite jet successor.  The underlying motivic objects,
history apices, bar complexes, transgression pages and detector objects are
constant in the jet direction; only the compatible finite Confluent ledger
changes.
\end{corollary}


The following theorem is the closure statement for Layer II over the fixed Layer-I squarefree structural core.  It is not a coefficient-depth completion theorem.

\begin{theorem}[Finite motivic Moore--Reedy recursion closure]
\label{P2M-thm:v38-finite-motivic-recursion-closure}
Fix an odd squarefree coefficient order $N$, a finite ordered support $S$, and
a finite jet ceiling $M\ge3$.  The constructions of this paper together with
the finite Confluent ledger assemble into one finite marked
Moore--Reedy recursion package
\begin{equation}\label{P2M-eq:v38-finite-recursion-package}
 \boxed{
 \mathfrak T^{\Mot,\MR}_{N;S,\le M}
 :=\Bigl(
   \mathbb T^{\Mot,\MR}_{N;S,M},\
   \{\mathcal R^{\Mot,\MR,\mathrm{hist};\mathrm{jet}}_{N,n}\}_{3\le n\le M},\
   \{\mathbb{DBar}^{\Mot,\mathrm{jet}}_{\gamma,n}\}_{\gamma,n},\
   \{\mathfrak Q^{\Mot,\mathrm{jet}}_{I,i_0,n}\}_{I,i_0,n}
 \Bigr).}
\end{equation}
It has the following closure properties.
\begin{enumerate}[label=\textup{(\roman*)}]
\item \textbf{Support--jet closure.}
The object part is the strict finite functor of
\cref{P2M-thm:v34-support-jet-tower}; support deletion and jet successor commute,
and the underlying motivic object at $(I,n)$ is always
$\mathbb M_{I,N}^{\Mot,\MR}$.

\item \textbf{Reedy closure.}
At every jet level the proper-face latching object, Boolean saturation,
exact-support cofiber and structural antecedent are those of the original
Moore--Reedy reconstruction.  In particular, for $|I|\ge3$,
\[
 \mathbb M_{I,N}^{\Mot,\MR}
 \cong
 \widehat{\mathbb L}_I^{\Mot}\oplus\mathbb Q_I^{\Mot,\MR},
 \qquad
 d\gamma_I^{\Mot}=N\beta_I^{\Mot},
\]
independently of the jet level.

\item \textbf{History closure.}
For every reduced insertion history $\gamma$, the jet-level history
realizations have the same terminal structural apex and the same finite family
of closed history blocks.  The jet-successor maps form pseudonatural
transformations and commute with reduced horizontal composition, associators,
unitors, reorderings and mixed deletion--insertion $2$-cells.  No structural
Moore package is duplicated by either history composition or jet extension.

\item \textbf{Secondary-transgression closure.}
For every reduced history $\gamma$, the underlying normalized diagonal bar
complex, its bar-length filtration and all spectral pages are independent of
jet level.  If $|I|=m$, then at every $3\le n\le M$ the same first Moore
transgression satisfies
\begin{equation}\label{P2M-eq:v38-jet-constant-secondary-transgression}
 \boxed{
 d_{m-1}^{\Mot}
 \bigl[N^{m-3}\mathfrak J_I^{\rm bar}\bigr]
 =\pm[1_{\mathscr A}\otimes s\Omega_I],
 \qquad
 \ord[1_{\mathscr A}\otimes s\Omega_I]=N.}
\end{equation}
Thus support cardinality controls the bar page and Rees exponent, while jet
order only changes the external finite ledger.

\item \textbf{Readout closure.}
The Hall/root--Bass realization, highest-face detector, CRT-hybrid quotient and
cellular Karoubi splitting commute with every jet successor.  Their underlying
objects, differentials, projectors, resonance loci and minimal-rank profiles
are independent of jet level.

\item \textbf{No structural growth.}
Increasing the finite jet ceiling changes only the compatible markings
recording residual obstruction, effective carry, repeated-sector detection,
primitive-lattice, filler and pointing status.  It creates no new Boolean
carrier, exact-support channel, root Moore pair, frontier-history block, bar
complex, Hall factor, face detector, CRT factor or Karoubi summand.
\end{enumerate}
Consequently every mixed diagram obtained from support deletion, reduced
history transport, bar/transgression passage, downstream readout and finite
jet successor commutes in the appropriate sense: strictly on the ordinary
support--jet and readout diagrams and pseudonaturally on the bicategorical
history layer.  Thus \eqref{P2M-eq:v38-finite-recursion-package} is closed under all
finite transition operations constructed in this paper.
\end{theorem}

\begin{proof}
Parts \textup{(i)--(ii)} are
\cref{P2M-thm:v34-support-jet-tower,P2M-thm:v34-jetwise-reedy-reconstruction}.
Part \textup{(iii)} is
\cref{P2M-thm:v35-history-jet-pseudonatural-tower} together with the reduced
history realization of \cref{P2M-thm:v21-bicategorical-motivic-realization}.
Part \textup{(iv)} is \cref{P2M-thm:v36-bar-jet-interchange}; the displayed
transgression is the existing squarefree secondary-transgression theorem read
through the identity jet map.  Part \textup{(v)} is
\cref{P2M-thm:v37-readout-jet-constancy}.  The no-growth statement combines
\cref{P2M-cor:v34-no-jet-structural-growth,P2M-cor:v35-no-history-growth-along-jet-axis,P2M-cor:v37-no-readout-growth}
with the identity of the underlying bar complexes along the jet axis.  The
final interchange assertion follows by composing these already established
strict and pseudonatural squares.
\end{proof}

\begin{corollary}[Finite motivic recursion is structurally closed in the proved scope]
\label{P2M-cor:v38-finite-recursion-structurally-closed}
Within the odd-squarefree all-support scope and at every finite jet ceiling,
there is no remaining internal compatibility gate among the support Reedy
recursion, reduced frontier history, secondary bar transgression and the
Hall/face/CRT/Karoubi readouts.  Any further obstruction belongs to one of the
following external problems:
\begin{enumerate}[label=\textup{(\alph*)}]
\item supplying the compatible pointed primitive exact-depth source lines required by \cref{P2M-thm:prime-power-all-support-comparison} at a prescribed order $N_\nu$ when they are not already available;
\item supplying an independent geometric correspondence provenance theorem for
a newly detected higher arithmetic operation when more than carrier reuse is
required and no explicit central-extension/gerbe construction such as
\cref{W31-thm:terminal-unipotent-gerbe-provenance} is available;
\item deciding the arithmetic source, primitive-lattice, carry or pointing
statuses recorded by the imported Confluent ledger.
\end{enumerate}
None of these is a missing coherence law inside
$\mathfrak T^{\Mot,\MR}_{N;S,\le M}$.
\end{corollary}

\begin{warning}[Finite closure only]
\label{P2M-warn:v38-finite-not-motivic-adelic}
The closure theorem is entirely finite in support, jet ceiling and coefficient fibre.
It does not promote the coefficient-depth vector $\nu$ to obstruction height or to the
valuation-depth fibre of \(\tauone\), nor does it form an inverse limit in coefficient
depth.  It is a finite carrier-plus-ledger package, not the definition of the two
motivic levels.
\end{warning}


\part{Full-motive thickening and its internal lifting presentations}
\section{Completion gates inside the full-motive level}\label{sec:full-tower-gaps}
The two motivic levels are \(\FullMot\) and \(\RecOneMot\), related by
\(\tauone\).  The Moore--Reedy material below supplies finite presentations of the
realized full-motive object over a fixed support base.  Its support-recursive and
obstruction-height indices are technical construction coordinates inside that upper
level; they do not define additional motivic levels.  The preceding material gives the
support realization and a finite marked jet/confluence ledger.  We now close the
exact-support detector gate before recording the remaining lifting statements.

\begin{definition}[Diagonal full-support root class]\label{def:diagonal-full-support-root-class}
Let $I$ be a finite ordered support of cardinality $m\ge3$ and retain the standing odd
coefficient order $N$.  In the representation-valued exact-support detector
\[
 \mathcal D_I^{\Mot}=\mathcal M_N[I]\otimes_{\mathbf Z}\mathbf Z[I]
\]
of \cref{P2M-def:permutation-valued-moore-detector}, define the \emph{diagonal
full-support root class} by
\begin{equation}\label{eq:diagonal-full-support-root-class}
 \boxed{
 \alpha_{I}^{\mathrm{root},\Mot}
 :=[b^{\Mot}]\otimes\mathbf1_I
 \in H^{2-m}(\mathcal D_I^{\Mot}).}
\end{equation}
It is the full-support copy of the pure-weight root Moore line, with no scalar
compression of the face module.  In particular this definition remains valid at
$N$-resonant support cardinalities.
\end{definition}

\begin{definition}[Exact-support motivic latching class]\label{def:exact-support-motivic-latching-class}
Under the canonical $S_I$-equivariant packaging of the face-minimal exact-support
Moore detector from \cref{P2M-thm:equivariant-face-packaging}, let
\begin{equation}\label{eq:exact-support-motivic-latching-class}
 \boxed{
 \alpha_I^{\mathrm{lat},\Mot}:=\boldsymbol\beta_I
 \in H^{2-m}(\mathcal D_I^{\Mot})}
\end{equation}
be the class induced by the exact-support motivic latching action
$\Lambda_I^{\Mot}$ of \cref{P2M-def:v20-exact-support-latching-action}.
\end{definition}

\begin{theorem}[Exact-support root comparison]\label{thm:exact-support-root-comparison}
For every finite ordered support $I$ with $|I|\ge3$ in the standing odd-$N$ Moore--Reedy
sector, the exact-support motivic latching class and the diagonal full-support
pure-weight root class satisfy
\begin{equation}\label{eq:exact-support-root-comparison}
 \boxed{
 \alpha_I^{\mathrm{lat},\Mot}
 =2\,\alpha_I^{\mathrm{root},\Mot}.}
\end{equation}
Hence the two classes generate the same exact-order-$N$ cyclic root line.  The
comparison is canonical under support relabelling and strict support deletion, and the
factor $2\in(\mathbf Z/N)^\times$ is the normalization unit forced by the paired
highest-face channel convention.  No division by $|I|$ is used, so the comparison is
uniform at resonant and nonresonant support cardinalities.
\end{theorem}

\begin{proof}
By \cref{P2M-thm:equivariant-face-packaging}, the face-minimal exact-support detector is
canonically the permutation-valued complex $\mathcal D_I^{\Mot}$, and its latching class
is the normalized diagonal vector of
\cref{P2M-eq:diagonal-detector-vector}:
\[
 \alpha_I^{\mathrm{lat},\Mot}
 =\boldsymbol\beta_I
 =2[b^{\Mot}]\otimes\mathbf1_I.
\]
By \cref{eq:diagonal-full-support-root-class}, the last tensor is precisely
$2\alpha_I^{\mathrm{root},\Mot}$, proving
\eqref{eq:exact-support-root-comparison}.  The root Moore class $[b^{\Mot}]$ has exact
order $N$ by \cref{P2M-thm:v19-pure-weight-root-seed}; since $N$ is odd, multiplication
by $2$ is an automorphism of its cyclic order-$N$ line.  Relabelling equivariance follows
from \cref{P2M-thm:equivariant-face-packaging}.  Strict deletion removes the whole
exact-support-$I$ detector when $I$ is not retained, by
\cref{P2M-thm:v20-exact-support-motivic-moore-package}; when the support survives, the
root coefficient and the diagonal vector are unchanged.  Finally, the construction
never applies the augmentation $P_I\to\mathbf Z/N$, so the support-cardinality resonance
of \cref{P2M-thm:cardinality-resonance} creates no exceptional case.
\end{proof}

\begin{corollary}[Closure of the exact-support motivic gate]\label{cor:exact-support-gate-closed}
The exact-support Moore correction of
\cref{P2M-thm:v20-motivic-moore-reedy-reconstruction} has genuine pure-weight root
provenance at every finite support: its canonical latching class is the full-support
root Moore class up to a unit.  Thus the exact-support gate introduces no additional
hypothesis in the supportwise presentation of the full-motive realization.
\end{corollary}

\subsection{An obstruction-height lifting presentation}

The jet-successor maps constructed earlier are bookkeeping maps: they enlarge the finite
obstruction ledger without changing the underlying motivic object.  A chosen lifting
presentation for the next marked obstruction is obtained by a class-pulled root Moore
lift.  This is an obstruction-height presentation within \(\FullMot\), not an intrinsic
construction of another motive tower.  The lift must allow the next obstruction to have
order strictly smaller than the ambient coefficient order; forcing every stage into the
same order-$N$ root line would incorrectly identify obstruction height with coefficient
depth or with the valuation depth of \(\tauone\).

\begin{definition}[Cyclic decomposition of the next marked obstruction]
\label{def:recursive-obstruction-cyclic-decomposition}
Let $X^{(n)}$ be a finite marked motivic stage with a chosen obstruction-height
presentation through index $n$.  Let
\[
 \mathfrak o_{n+1}(X^{(n)})
\]
be the next marked obstruction supplied by the support-level obstruction ledger.  On
each connected marked branch write its finite marked coefficient module as the canonical
sum of the nonzero cyclic submodules generated by the retained marked components,
\[
 \mathscr O_{n+1}
 =\bigoplus_{a\in A_{n+1}}\mathscr L_a/d_a\mathscr L_a,
 \qquad d_a>1,
\]
where $\mathscr L_a$ is the corresponding invertible detector line and $d_a$ is the
exact additive order of its marked obstruction component.  Since the ambient coefficient
order is odd, every $d_a$ is odd.  A zero component is omitted from $A_{n+1}$ and creates
no new motivic cell.

Let $\sigma_a$ denote the cohomological/Tate shift carried by that marked Recursive
Secondary Cartan component.  This shift is part of the typed obstruction datum and is
independent of the coefficient order $d_a$.
\end{definition}

\begin{definition}[Obstruction-height motivic lifting step]
\label{def:genuine-recursive-motivic-successor}
For each $a\in A_{n+1}$ the marked obstruction component determines a section
\[
 \alpha_a:X^{(n)}\longrightarrow
 \mathfrak C_{d_a,X^{(n)}}(\mathscr L_a)[\sigma_a]
\]
of the relative generator Moore class stack of
\cref{P2R-def:v91-relative-moore-line}.  Define
\begin{equation}\label{eq:genuine-recursive-motivic-successor}
 \boxed{
 \mathfrak S_{\Mot}^{(n+1)}(X^{(n)})
 :=X^{(n)}
 \times^h_{\displaystyle\prod_{a\in A_{n+1}}
   \mathfrak C_{d_a,X^{(n)}}(\mathscr L_a)[\sigma_a]}
 \displaystyle\prod_{a\in A_{n+1}}
   \mathfrak R_{d_a,X^{(n)}}(\mathscr L_a)[\sigma_a].}
\end{equation}
The root legs are the line-valued cyclotomic Moore maps of
\cref{P2R-def:v92-relative-root-object}.  If $A_{n+1}=\varnothing$, the empty
homotopy pullback is defined to be $X^{(n)}$ itself.  The resulting stage is denoted
$X^{(n+1)}$.
\end{definition}

\begin{theorem}[Successor-stage lifting functor]
\label{thm:genuine-successor-stage-motivic-functor}
The construction \eqref{eq:genuine-recursive-motivic-successor} defines a canonical
successor functor on finite marked Recursive Secondary Cartan motivic stages,
\[
 \boxed{
 \mathfrak S_{\Mot}^{(n+1)}:
 \mathsf{MRMot}^{(n)}_{\rm Cart}
 \longrightarrow
 \mathsf{MRMot}^{(n+1)}_{\rm Cart}.}
\]
For every finite recursive height it has the following properties.
\begin{enumerate}[label=\textup{(\roman*)}]
\item \textbf{Motivic closure.}  $X^{(n+1)}$ is obtained by finite homotopy pullbacks
      of relative cyclotomic root Moore objects over relative Moore Picard stacks and is
      therefore again a Cartan--motivic stage.
\item \textbf{Moore exactness.}  On the $a$th fresh cyclic component the new root pair
      satisfies
      \[
       d\Gamma_a=d_a\Omega_a,
      \]
      and the derived path in the homotopy pullback identifies $\Omega_a$ with the
      corresponding next-obstruction class.  Thus the new motivic antecedent realizes
      exactly the universal Recursive Secondary Cartan cell relation.
\item \textbf{Latching compatibility.}  On every exact-support component the classifying
      section agrees with the Moore--Reedy latching class up to the canonical unit of
      \cref{thm:exact-support-root-comparison}; hence the new stage extends, rather than
      replaces, the already reconstructed proper-face motivic stage.
\item \textbf{Support functoriality.}  Ordered support deletion, relabelling and every
      direction-linear map carry the next obstruction to the corresponding target
      obstruction and commute with the successor homotopy pullback.  A fresh component
      whose exact support is deleted is removed; every surviving component is unchanged.
\item \textbf{Representative and frame independence.}  Replacing an obstruction cocycle
      by a cohomologous representative or changing a local frame of $\mathscr L_a$ yields
      canonically equivalent successor stages.  Thus the construction depends only on
      the marked recursive obstruction class and not on a chosen cocycle, cyclic
      generator, or trivialization.
\item \textbf{Coefficient-depth firewall.}  The order $d_a$ determines only which
      cyclotomic Moore object realizes the $a$th obstruction line.  The passage
      $n\mapsto n+1$ is governed by recursive obstruction height and does not identify
      $d_a$ (or any $q$-power coefficient depth) with recursive height.
\end{enumerate}
\end{theorem}

\begin{proof}
For each nonzero marked component, the Recursive Secondary Cartan higher-obstruction
recursion gives a canonical cohomology class, invariant under pointed gauge and natural
for controller maps and direction-linear maps; its universal recursive cell has
differential equal to minus a cocycle representative.  Passing to the cyclic submodule
actually generated by that marked component gives the exact order $d_a$ line
$\mathscr L_a/d_a\mathscr L_a$, hence a section of the relative generator Moore class
stack.  The general odd-order root-stack Moore construction supplies the corresponding
line-valued root leg.

The homotopy pullback in \eqref{eq:genuine-recursive-motivic-successor} is therefore the
line-valued Moore lift of \cref{P2R-def:v92-line-valued-lift}, applied
componentwise.  This proves (i).  The universal root relation
$d\Gamma_a=d_a\Omega_a$ and the path identifying the root top class with the marked
obstruction section give the required motivic image of the universal semi-free cell
relation, proving (ii).  On exact support, \cref{thm:exact-support-root-comparison}
identifies the latching class with the diagonal root class up to the unit $2$, so the
pullback attaches the fresh antecedent to the existing Moore--Reedy latching datum;
this is (iii).

Naturality of the next obstruction under deletion, permutation and direction-linear maps,
together with base-change functoriality of homotopy pullbacks and of the relative root
construction, gives (iv).  If a cocycle representative changes by $dY$, the universal
recursive-cell theorem identifies the two source attachments by the translation of the
fresh cell by $-Y$; the target-side statement is the homotopy invariance of the classifying
section.  Frame independence is exactly the $\mathbf G_m$-descent statement built into
\cref{P2R-def:v92-relative-root-object,P2R-thm:v92-relative-cylinder-presents-lift}.
This proves (v).  Finally the root object is indexed by the exact order of the cyclic
coefficient line, whereas the successor index is the PD/obstruction degree;
the multiplicity--depth and effective-weight separation in the source package
therefore gives (vi).
\end{proof}

\begin{corollary}[Finite obstruction-height lifting presentation]
\label{cor:finite-recursive-motivic-height-tower}
Starting from the low motivic seed, iterating
\cref{thm:genuine-successor-stage-motivic-functor} along any finite marked obstruction
branch produces a finite inverse system of motivic lifting spaces
\[
 \boxed{
 X^{(M)}\xrightarrow{p_M}X^{(M-1)}\longrightarrow\cdots
 \longrightarrow X^{(1)}\xrightarrow{p_1}X^{(0)}.}
\]
Here $p_{n+1}$ is the canonical projection from the homotopy pullback defining
$X^{(n+1)}$ to its previous stage $X^{(n)}$.  The presentation index therefore
increases in the opposite direction to the forgetful geometric maps.  At a stage at
which the next obstruction vanishes, $p_{n+1}$ is an equivalence and only the solved marking
advances; at a nonzero height its relative fibre is precisely the space of root Moore
antecedents of the nonzero cyclic obstruction summands.  No coefficient-depth inverse
limit is involved, and this inverse system is not the motivic
truncation \(\tauone:\FullMot\to\RecOneMot\).
\end{corollary}

\subsection{Typed multi-axis interchange}

The successor construction is functorial in the next obstruction class, but exact-order
compression forces one important distinction.  A support deletion either preserves a
fresh cyclic obstruction line or removes it.  By contrast, a confluence morphism may
multiply the line by a nonunit factorial scalar, and coefficient reduction may lower the
exact order of the same primitive line.  Thus the correct statement is not that all four
axes form a strict cartesian grid.  The presentation axis (obstruction height) is
pseudonatural over a typed base-change category; resonance is recorded by the canonical
cyclotomic reduction map rather than erased.

\begin{definition}[Typed Cartan base-change category]
\label{def:typed-cartan-base-change-category}
Fix an obstruction-height presentation index $n$.  Let
\[
 \mathsf B_n
\]
be the category generated by the following morphisms between marked finite Cartan
parameters $(I,\lambda,\nu)$:
\begin{enumerate}[label=\textup{(\roman*)}]
\item ordered support deletions and relabellings;
\item direction-linear confluence coarsenings $f:\lambda\to\mu$, carrying on each
      retained rank-one repeated-direction line the integral scalar $M(f)$;
\item finite coefficient reductions $\nu'\twoheadrightarrow\nu$, componentwise from
      deeper coefficient order to shallower coefficient order.
\end{enumerate}
No generator of $\mathsf B_n$ changes the presentation index or the motivic level.  In
particular there is no
canonical arrow in the opposite direction to a coefficient reduction unless a marked
coefficient lift is supplied as additional data.
The category \(\mathsf B_n\) is fibrewise over a fixed valuation-depth state of
\(\FullMot\).  It contains neither the truncation \(\tauone\) nor vertical
prime-depth motion; those belong to the overarching two-level architecture.
\end{definition}

\begin{lemma}[Cyclotomic exact-order reduction]
\label{lem:cyclotomic-exact-order-reduction}
Let a marked obstruction component of exact odd order $d$ map under a typed base change
to a component of exact order $d'$, with $d'\mid d$, or to zero.  If $d'>1$, there is
a canonical morphism of line-valued root Moore objects
\[
 \rho_{d,d'}:
 \mathfrak R_{d,X}(\mathscr L)
 \longrightarrow
 \mathfrak R_{d',X'}(\mathscr L')
\]
over the corresponding morphism of generator Moore class stacks.  It is induced
fiberwise by the standard root-stack map
\[
 [\mathbf A^1/\mu_d]\longrightarrow[\mathbf A^1/\mu_{d'}],
 \qquad
 u\longmapsto u^{d/d'},
 \qquad
 \zeta\longmapsto\zeta^{d/d'}.
\]
For $d'=d$ it is an equivalence.  If the image obstruction is zero, the target exact-order
summand is absent and the corresponding fresh root cell is deleted.
\end{lemma}

\begin{proof}
The displayed map sends a chosen $d$-th root to its $d/d'$-th power, hence to a
$d'$-th root of the same base section, and is equivariant for the quotient
$\mu_d\to\mu_{d'}$.  The ranked localization Moore class and its antecedent are natural
for this root-stack map, so the relation $d\Gamma=d\Omega$ descends to the exact-order
$d'$ Moore relation after passing to the cyclic image line.  The line-valued version
follows by the same $\mathbf G_m$-descent used in the relative root construction.
When $d'=d$ the map is the identity up to the canonical line descent.  A zero image
classifies the base point of the Moore class stack and therefore requires no fresh
exact-order root cell.
\end{proof}

\begin{theorem}[Typed multi-axis interchange]
\label{thm:typed-multi-axis-interchange}
Let $X^{(n)}_{I,\lambda,\nu}$ be a finite marked stage and let
\[
 \phi:(I,\lambda,\nu)\longrightarrow(J,\mu,\nu')
\]
be a morphism generated in $\mathsf B_n$.  Write
\[
 \phi_*X^{(n)}_{I,\lambda,\nu}=X^{(n)}_{J,\mu,\nu'}
\]
for the induced marked base change.  Then the genuine motivic successor of
\cref{thm:genuine-successor-stage-motivic-functor} admits a canonical comparison
\begin{equation}\label{eq:typed-successor-base-change}
 \boxed{
 \Phi_{\phi}^{(n+1)}:
 \phi_*\mathfrak S_{\Mot}^{(n+1)}(X^{(n)}_{I,\lambda,\nu})
 \longrightarrow
 \mathfrak S_{\Mot}^{(n+1)}(\phi_*X^{(n)}_{I,\lambda,\nu}).}
\end{equation}
These comparisons are compatible with composition in $\mathsf B_n$.  More precisely:
\begin{enumerate}[label=\textup{(\roman*)}]
\item \textbf{Support.}  Ordered support deletion commutes strictly with successor:
      a surviving exact-support component and its root antecedent are unchanged, while a
      deleted exact-support component and its fresh root cell disappear together.
\item \textbf{Confluence.}  For a confluence arrow $f$, a cyclic component of exact
      order $d$ is sent to the cyclic line generated by $M(f)\alpha$, of exact order
      \[
       d_f=\frac{d}{\gcd(d,M(f))}
      \]
      when nonzero.  The comparison $\Phi_f^{(n+1)}$ is the cyclotomic reduction
      $\rho_{d,d_f}$ on that component.  It is an equivalence exactly on the
      nonresonant/unit part $\gcd(d,M(f))=1$.  The lost subgroup is the already existing
      factorial/Moore confluence defect; it is not a new motivic level.
\item \textbf{Coefficient depth.}  A finite coefficient reduction carries each marked
      obstruction to its reduced cyclic image and induces the corresponding
      $\rho_{d,d'}$.  Hence coefficient reduction commutes with successor by a canonical
      comparison and by an equivalence whenever exact order is preserved.  There is no
      canonical reverse ``thickening'' interchange without a chosen marked lift.
\item \textbf{Obstruction height.}  The only operation that changes $n$ is
      $\mathfrak S_{\Mot}^{(n+1)}$.  Neither support size, confluence multiplicity nor
      coefficient exponent is reindexed as obstruction height, and \(n\) is not the
      valuation-depth coordinate of \(\FullMot\).
\item \textbf{Mixed squares.}  Support deletion commutes with confluence and coefficient
      reduction because all three act functorially on the marked obstruction class and
      exact-support line.  Confluence and coefficient reduction commute because scalar
      multiplication by $M(f)$ commutes with quotient of the coefficient line.  The
      resulting pairwise squares for \eqref{eq:typed-successor-base-change} agree
      canonically, and the same is true for every finite composite.
\end{enumerate}
\end{theorem}

\begin{proof}
The next Recursive Secondary Cartan obstruction is natural for support maps,
direction-linear maps and finite coefficient reduction.  Therefore every generator
$\phi$ of $\mathsf B_n$ sends each marked obstruction component $\alpha_a$ to a
well-defined cyclic image component $\phi_*\alpha_a$.  Its exact order $d'_a$ divides
the source order $d_a$.  Apply \cref{lem:cyclotomic-exact-order-reduction} to each
nonzero image component and take the product of the resulting root-stack maps.  Naturality
of the Moore classifying sections gives a morphism between the two diagrams defining the
homotopy pullbacks in \eqref{eq:genuine-recursive-motivic-successor}; the universal
property of the homotopy pullback yields \eqref{eq:typed-successor-base-change}.
Compatibility with composition follows from
\[
 \rho_{d',d''}\circ\rho_{d,d'}=\rho_{d,d''}
\]
and functoriality of the obstruction class.

For support deletion, exact-support functoriality says that a fresh component is either
retained unchanged or deleted, proving (i).  For confluence, the divided-power merge
calculus multiplies a normalized rank-one obstruction line by the path-independent scalar
$M(f)$.  In a cyclic group of order $d$, the element $M(f)\alpha$ has exact order
$d/\gcd(d,M(f))$, giving (ii).  Thus the positive resonance length records precisely
the order lost under confluence; it does not create another successor index.

Finite coefficient reduction is simply passage to the cyclic image in a quotient
coefficient line, so (iii) is the same exact-order reduction mechanism.  The absence of
a canonical reverse arrow is the coefficient-depth firewall: lifting a residue class is
additional marked data, not part of coefficient depth itself.  Statement (iv) is built
into the definition of $\mathsf B_n$ and the successor theorem.  Finally exact-support
deletion is linear, scalar confluence commutes with coefficient quotient, and all induced
root reductions satisfy the transitivity identity above.  Hence the mixed squares and
all finite composites agree canonically, proving (v).
\end{proof}

\begin{theorem}[Prime-power all-support comparison]
\label{P2M-thm:prime-power-all-support-comparison}
Let \(S\) be a finite ordered support and let \(\nu\) be a finite odd-prime
coefficient-depth datum, with
\[
 N_\nu=\prod_{p\in\operatorname{supp}(\nu)}p^{\nu_p},
 \qquad
 N_{\rm red}=\prod_{p\in\operatorname{supp}(\nu)}p.
\]
Assume that for every \(p\in\operatorname{supp}(\nu)\) the source Confluent
package supplies a pointed primitive filtered line
\[
 L^{\rm conf}_{p^{\nu_p}}\simeq\Z/p^{\nu_p}
\]
together with its universal Moore presentation, and that these pointed lines
are compatible under all coefficient reductions
\(p^{a}\twoheadrightarrow p^{b}\).  Assume also that their order-\(p\)
reductions are the repeated-\(p\) lines used by the Layer-I realization.

Then the CRT product
\[
 L^{\rm conf}_{\nu}
 :=\bigoplus_{p\in\operatorname{supp}(\nu)}
 L^{\rm conf}_{p^{\nu_p}}
 \simeq\Z/N_\nu
\]
has a canonical pointed universal Moore presentation.  For the general-odd
root-stack pair
\[
 dc_\nu^{\Mot}=N_\nu b_\nu^{\Mot},
 \qquad
 \operatorname{ord}[b_\nu^{\Mot}]=N_\nu,
\]
the pointed assignment \(B_\nu\mapsto b_\nu^{\Mot}\),
\(C_\nu\mapsto c_\nu^{\Mot}\) extends to a support-functorial dg realization
\begin{equation}\label{P2M-eq:prime-power-all-support-realization}
 \boxed{
 \rho_{S,\nu}^{\Mot,\MR}:
 \mathscr U^{\mathrm{rec}}_{S,N_\nu}
 \longrightarrow
 \mathscr M_{S,N_\nu}^{\Mot,\MR}.}
\end{equation}
It has the following properties.
\begin{enumerate}[label=\textup{(\roman*)}]
\item Every exact-support latching class has exact additive order \(N_\nu\)
      and has a genuine motivic antecedent obtained from \(c_\nu^{\Mot}\).
\item Ordered support deletion and relabelling are strict, and same-rank
      exact-support attachments commute.
\item If \(\mu\le\nu\), coefficient reduction gives a commutative square
      \[
      \begin{CD}
      \mathscr U^{\mathrm{rec}}_{S,N_\nu}
        @>{\rho_{S,\nu}^{\Mot,\MR}}>>
      \mathscr M_{S,N_\nu}^{\Mot,\MR}\\
      @V{\operatorname{red}_{\nu,\mu}}VV
        @VV{\operatorname{Red}_{\nu,\mu}^{\Mot}}V\\
      \mathscr U^{\mathrm{rec}}_{S,N_\mu}
        @>{\rho_{S,\mu}^{\Mot,\MR}}>>
      \mathscr M_{S,N_\mu}^{\Mot,\MR}.
      \end{CD}
      \]
      For \(\mu=\nu^{\rm sf}\), the lower row is the Layer-I squarefree
      realization after the fixed pointed identification.
\item A confluence coarsening \(f\) preserves its repeated-prime provenance
      and acts on a cyclic line by the divided-power scalar \(M(f)\).  Its
      motivic comparison is the cyclotomic reduction from order \(d\) to
      \[
       d_f=\frac{d}{\gcd(d,M(f))},
      \]
      so the prime-power depth lost at resonance is exactly the source Moore
      defect and not forgotten provenance.
\item The passage \(N_{\rm red}\leadsto N_\nu\) changes only the structural
      coefficient fibre.  It changes neither support, confluence multiplicity,
      obstruction height, nor valuation depth.
\end{enumerate}
Thus compatible primitive filtered confluence lines, rather than the bare
depth label \(\nu\), are precisely the extra input needed to remove the
squarefree coefficient restriction from the all-support realization.
\end{theorem}

\begin{proof}
The \(p\)-primary lines are pointed cyclic of pairwise coprime orders, so the
finite Chinese remainder theorem makes their direct sum pointed cyclic of
exact order \(N_\nu\).  The compatible universal Moore presentations therefore
assemble to
\[
 \Z C_\nu\xrightarrow{\,N_\nu\,}\Z B_\nu,
 \qquad dC_\nu=N_\nu B_\nu.
\]
The general-odd root-stack pair gives the displayed dg map on this initial
two-cell block.

After this initial block is fixed, the proof of
\cref{P2M-thm:channel-complete-realization} is coefficient-formal.  The Boolean
part is unchanged, every exact-support channel equation uses only
\(dc_\nu^{\Mot}=N_\nu b_\nu^{\Mot}\), and orthogonality of channel summands is
independent of the factorization of \(N_\nu\).  Projection to a highest-face
channel has transfer exponent one and sends the latching class to
\(\pm E_{13}b_\nu^{\Mot}\); hence its order, and therefore the order of the full
channel vector, is exactly \(N_\nu\).  The image of every support cell \(K_I\)
is assembled from placements of \(c_\nu^{\Mot}\), giving (i).  Exact-support
deletion, relabelling, and orthogonality are likewise coefficient-formal,
giving (ii).

For \(\mu\le\nu\), quotient of the pointed source lines commutes with their
universal Moore presentations.  On the target, the root-stack power map
\[
 [\mathbf A^1/\mu_{N_\nu}]
 \longrightarrow
 [\mathbf A^1/\mu_{N_\mu}],
 \qquad u\longmapsto u^{N_\nu/N_\mu},
\]
and the induced map on every channel give
\(\operatorname{Red}_{\nu,\mu}^{\Mot}\).  Naturality of the latching action and
antecedent makes the square in (iii) commute.  At squarefree depth this is the
fixed Layer-I root pair and hence recovers its realization.

Finally a confluence coarsening multiplies a normalized repeated-direction
line by \(M(f)\).  A cyclic class of order \(d\) then has image order
\(d/\gcd(d,M(f))\), and the same root-stack power map is exactly the comparison
of \cref{thm:typed-multi-axis-interchange}.  This proves (iv).  All maps are
fibrewise over the fixed support, obstruction-height, and valuation-depth
coordinates, proving (v).
\end{proof}

\begin{corollary}[Typed finite motivic lifting system]
\label{cor:typed-finite-recursive-motivic-tower}
For every finite obstruction-height ceiling $M$, the stages constructed by
\cref{cor:finite-recursive-motivic-height-tower} assemble over the typed base-change
categories into a finite motivic lifting system
\[
 \boxed{
 \mathbb M_{\rm Cart,\le M}^{\rm MR}
 =
 \{\mathcal M^{(n)}_{I,\lambda,\nu}\}_{0\le n\le M}.}
\]
The system is strict for support deletion and carries canonical comparison morphisms for
confluence and coefficient reduction.  On the nonresonant/order-preserving locus those
comparison morphisms are equivalences.  Obstruction height, confluence multiplicity and
coefficient depth are independent presentation coordinates of this system and remain
distinct from valuation depth.
\end{corollary}

\begin{proof}
Iterate the genuine successor and use
\cref{thm:typed-multi-axis-interchange} at every height.  Composition compatibility of
the comparison maps gives a well-defined typed finite system, while the stated
strictness/equivalence properties are exactly parts (i)--(iii) of that theorem.
\end{proof}

\subsection{Internal Cartan lifting tower and relative layers}

The homotopy-pullback definition of the lifting step supplies an internal filtration
presentation.  It is an inverse system of lifting spaces rather than an increasing
chain of subobjects.  It must not be confused with the structural truncation
\(\tauone:\FullMot\to\RecOneMot\).

\begin{definition}[Cartan lifting stages and relative layer]
\label{def:cartan-truncation-successor-layer}
For a finite lifting presentation of height $M$, write its stage-$n$ object as
\[
 \ell^{\rm Cart}_{\le n}\mathbb M:=X^{(n)},\qquad 0\le n\le M,
\]
and let
\[
 p_{n+1}:X^{(n+1)}\longrightarrow X^{(n)}
\]
be the canonical projection from the successor homotopy pullback.  Over a marked point
$x\in X^{(n)}$ with nonzero next-obstruction components
$\alpha_a(x)$ of exact orders $d_a$, define the relative Cartan successor layer by
\begin{equation}\label{eq:cartan-relative-successor-layer}
 \boxed{
 \operatorname{Gr}^{\rm Cart}_{n+1}(x)
 :=
 \prod_{a\in A_{n+1}}
 \operatorname{hofib}_{\alpha_a(x)}\!\left(
   \mathfrak R_{d_a,x}(\mathscr L_a)[\sigma_a]
   \longrightarrow
   \mathfrak C_{d_a,x}(\mathscr L_a)[\sigma_a]
 \right).}
\end{equation}
For a zero obstruction component the corresponding factor is omitted.  Thus
$\operatorname{Gr}^{\rm Cart}_{n+1}$ records root antecedent choices at obstruction
height $n+1$, not valuation-depth thickening, coefficient thickening, or confluence
multiplicity.
\end{definition}

\begin{theorem}[Cartan lifting presentation and relative layer]
\label{thm:intrinsic-cartan-truncation-layer}
The chosen lifting-step construction presents the finite motivic system as an inverse
system
\[
 X^{(M)}\to\cdots\to X^{(1)}\to X^{(0)}.
\]
For every $n<M$ and every marked base point $x\in X^{(n)}$, the homotopy fibre of
$p_{n+1}$ over $x$ is canonically the relative layer
\eqref{eq:cartan-relative-successor-layer}.  Consequently:
\begin{enumerate}[label=\textup{(\roman*)}]
\item the lifting-stage filtration is canonical after the obstruction presentation is
      fixed and requires no auxiliary choice of subobject filtration;
\item a vanishing next obstruction gives a contractible relative layer and an equivalence
      $X^{(n+1)}\simeq X^{(n)}$ on that marked branch;
\item a nonzero cyclic obstruction contributes exactly its line-valued root Moore lifting
      fibre and no extra primitive component;
\item the relative layers are functorial for the typed base-change comparisons of
      \cref{thm:typed-multi-axis-interchange}; on the order-preserving locus this
      functoriality is by equivalences;
\item obstruction height is the index of this lifting presentation, while support,
      confluence multiplicity and coefficient depth remain parameters of each stage;
      none is the two-level truncation index \(\tauone\).
\end{enumerate}
\end{theorem}

\begin{proof}
The defining square
\[
\begin{CD}
 X^{(n+1)} @>>> \prod_a\mathfrak R_{d_a,X^{(n)}}(\mathscr L_a)[\sigma_a]\\
 @V{p_{n+1}}VV @VVV\\
 X^{(n)} @>{(\alpha_a)_a}>>
 \prod_a\mathfrak C_{d_a,X^{(n)}}(\mathscr L_a)[\sigma_a]
\end{CD}
\]
is homotopy cartesian by \cref{def:genuine-recursive-motivic-successor}.  Taking the
homotopy fibre over $x$ gives exactly
\eqref{eq:cartan-relative-successor-layer}, proving the first assertion.  If the next
obstruction vanishes, the definition omits every zero component, so the successor is the
empty homotopy pullback $X^{(n)}$ itself; this proves (ii).  For a nonzero component, the
only new factor in the cartesian square is its relative root Moore object, which proves
(iii).  Part (iv) is base change of this homotopy-cartesian square together with
\cref{thm:typed-multi-axis-interchange}.  The final typing assertion is the
obstruction-height, valuation-depth, and coefficient-depth firewall already built into
the lifting construction.
\end{proof}

\begin{corollary}[Cartan-generated layer completeness]
\label{cor:cartan-generated-layer-completeness}
Within the chosen motivic lifting presentation, every new finite-height relative layer
is generated by a marked Secondary Cartan obstruction and its
line-valued root Moore antecedent.  There is no additional hidden finite-height layer in
that presentation.  This does not assert that the full motive itself is intrinsically
generated by recursion.  Any comparison with a larger classical motivic universe may
contain further primitive motivic directions, but those directions are not part of the
Cartan-generated sector treated here.
\end{corollary}

\section{Classical 0/1-motivic shadow and Yoneda comparison}
\label{sec:classical-one-motive-comparison}

The finite Moore line is torsion, so the correct classical target is not the heart of
free Deligne 1-motives by itself.  It is the bounded derived category, where finite
exact-order sectors occur canonically as cones of multiplication maps.  This also keeps
the classical comparison separate from the internal Cartan lifting presentation.

\begin{definition}[Classical 1-motivic target and standard Moore object]
\label{def:classical-one-motivic-target}
Let $k$ be a perfect field of exponential characteristic $p$, put
$\Lambda=\mathbf Z[1/p]$, and assume that every cyclic order used below is prime to $p$.
Let
\[
 \mathsf D_1(k;\Lambda):=D^b(\mathcal M_1(k)\otimes\Lambda)
\]
be the bounded derived category of Deligne 1-motives with the standard exact structure.
Write
\[
 \mathbf A:=[\Lambda\longrightarrow0],
 \qquad
 \mathbf T:=[0\longrightarrow\mathbf G_m]\otimes\Lambda
\]
for the rank-one Artin and toric 1-motives.  For an integer $d>1$ prime to $p$ define
\begin{equation}\label{eq:classical-torsion-moore-object}
 \boxed{
 \mathbf Q_d
 :=\operatorname{Cone}\!\left(
 d:\mathbf T\longrightarrow\mathbf T
 \right)[-1]\in\mathsf D_1(k;\Lambda).}
\end{equation}
We call $\mathbf Q_d$ the standard $d$-torsion Moore object in the classical
1-motivic derived category.  Exact order is carried by the marked primitive class in a
chosen cyclic summand, not by the unmarked cone object alone.  For objects $A,B$ of the
exact 1-motive model we use the notation
\[
 \operatorname{Ext}^m_{\mathsf D_1}(A,B)
 :=\operatorname{Hom}_{\mathsf D_1}(A,B[m]).
\]
\end{definition}

\begin{remark}[Why the derived target is classical]
\label{rem:bvk-classical-target}
Barbieri--Viale and Kahn construct, after inverting the exponential characteristic, a
fully faithful totalization functor from the derived category of Deligne 1-motives to the
\'etale version of Voevodsky's effective geometric motives
\cite{BarbieriVialeKahn2016}.  Thus $\mathsf D_1(k;\Lambda)$ is a genuine classical
motivic target.  Formula \eqref{eq:classical-torsion-moore-object} retains finite Moore
torsion without requiring that this torsion lie in the free 1-motive heart.
\end{remark}

\begin{lemma}[Exact-order reduction in the classical target]
\label{lem:classical-moore-order-reduction}
If $d'\mid d$, the commutative square
\[
\begin{CD}
 \mathbf T @>{d}>> \mathbf T\\
 @V{d/d'}VV @VV{1}V\\
 \mathbf T @>{d'}>> \mathbf T
\end{CD}
\]
induces a canonical morphism
\begin{equation}\label{eq:classical-moore-reduction}
 \boxed{q_{d,d'}:\mathbf Q_d\longrightarrow\mathbf Q_{d'}.}
\end{equation}
These maps are transitive:
$q_{d',d''}\circ q_{d,d'}=q_{d,d''}$ whenever $d''\mid d'\mid d$, and
$q_{d,d}=1$.
\end{lemma}

\begin{proof}
Apply the cone functor to the displayed square and shift by $[-1]$.  Functoriality of
cones gives \eqref{eq:classical-moore-reduction}; the transitivity identities follow from
the corresponding identities for multiplication on $\mathbf T$.
\end{proof}

\begin{definition}[Cartan--Moore classical shadow]
\label{def:cartan-moore-classical-shadow}
Let $\mathscr G^{\rm CM}_{\le1}$ denote the marked dg sector generated by:
\begin{enumerate}[label=\textup{(\roman*)}]
\item the rooted degree-zero support markings of
      \cref{P2L-thm:v48r19-rooted-zero-motive};
\item the exact-order root/Moore pairs $(b,c)$ satisfying
      $db=0$, $dc=d\,b$, and $\operatorname{ord}[b]=d$, where $d$ denotes the actual
      cyclic order of that marked component;
\item the exact-order reduction maps of
      \cref{lem:cyclotomic-exact-order-reduction};
\item the relative bar cells and their concatenations inside the Cartan--Moore dg
      comparison sector.
\end{enumerate}
For a marked rank-one detector/root line $\ell$, let $\mathbf A_\ell$ and
$\mathbf T_\ell$ denote the corresponding rank-one Artin and toric objects; changing a
frame of $\ell$ acts by a unit and hence by an isomorphism.  Define
\begin{equation}\label{eq:classical-shadow-functor}
 \boxed{
 \operatorname{cl}_1:\mathscr G^{\rm CM}_{\le1}
 \longrightarrow\mathsf D_1(k;\Lambda)}
\end{equation}
on generators by
\[
 \text{rooted degree-zero marking }\ell\longmapsto\mathbf A_\ell,
 \qquad
 (b,c;d)\longmapsto
 \mathbf Q_{d,\ell}:=
 \operatorname{Cone}(d:\mathbf T_\ell\to\mathbf T_\ell)[-1],
\]
and by sending an exact-order reduction $\rho_{d,d'}$ to $q_{d,d'}$.
On relative bar cells, $\operatorname{cl}_1$ is the induced morphism of the corresponding
cone presentations.
\end{definition}

\begin{theorem}[Classical low-sector comparison]
\label{thm:classical-low-sector-comparison}
The assignment \eqref{eq:classical-shadow-functor} is well defined on the framed
Cartan-generated root/Moore sector and descends, after forgetting the frame, up to
canonical isomorphism.  It has the following properties.
\begin{enumerate}[label=\textup{(\roman*)}]
\item The rooted arithmetic zero-motive of
      \cref{P2L-thm:v48r19-rooted-zero-motive} has classical shadow the corresponding
      Artin degree-zero object $\mathbf A_\ell$.  All Selmer, curvature, pointing and
      finite-transport data are forgotten by this shadow functor rather than being
      reinterpreted as part of a Deligne 1-motive.
\item The primitive motivic Moore seed of
      \cref{P2L-thm:v48r4-motivic-seed} maps to $\mathbf Q_{N,\ell}$, and its genuine
      root-stack antecedent maps to the standard cone antecedent.  Hence the relation
      $dc=N b$ becomes exactly the standard multiplication-by-$N$ Moore presentation.
\item Every successor component of exact order $d$ maps to the corresponding
      $\mathbf Q_{d,\ell}$.  If confluence or coefficient reduction lowers $d$ to $d'$, the
      typed comparison of \cref{thm:typed-multi-axis-interchange} maps to
      $q_{d,d'}$ of \cref{lem:classical-moore-order-reduction}.
\item Support deletion, unit change of detector frame, and all finite typed composites
      are respected.  In particular the classical shadow is compatible with the complete
      finite recursive Cartan--Moore tower constructed above.
\end{enumerate}
\end{theorem}

\begin{proof}
For (i), the rooted zero-stage contains a distinguished rank-one root marking; taking the
free Artin lattice on that marking is functorial for compatible pointed replacement.  The
remaining Selmer--Cartan decorations are intentionally outside the target of
$\operatorname{cl}_1$.

For (ii), the source Moore pair is an honest two-term relation with one closed generator
$b$ and one antecedent $c$ satisfying $dc=N b$.  The target
$\operatorname{Cone}(N:\mathbf T_\ell\to\mathbf T_\ell)[-1]$ is the standard two-term
Moore presentation with the same relation.  The root-stack construction fixes the
primitive cyclic generator up to a unit, while a unit change acts by an automorphism of
$\mathbf T_\ell$; hence the resulting object is independent of frame.  Part (iii) is
exactly \cref{lem:classical-moore-order-reduction} applied to the actual cyclic order of
each successor component.  Part (iv) follows from exact-support functoriality, the
transitivity of the $q_{d,d'}$, and the functoriality of cone presentations and relative
bar maps.
\end{proof}

\begin{theorem}[Bar composition equals Yoneda composition]
\label{thm:bar-yoneda-comparison}
Let
\[
 \eta_1\in
 \operatorname{Hom}_{\mathsf D_1}(A_0,A_1[1]),
 \qquad
 \eta_2\in
 \operatorname{Hom}_{\mathsf D_1}(A_1,A_2[1])
\]
be the classical shadows of two composable degree-one relative Cartan--Moore lifting/bar
classes.  Then the classical shadow of their length-two bar composition is
\begin{equation}\label{eq:bar-yoneda-product}
 \boxed{
 \operatorname{cl}_1(\eta_2\star\eta_1)
 =\eta_2[1]\circ\eta_1
 \in
 \operatorname{Hom}_{\mathsf D_1}(A_0,A_2[2])
 =\operatorname{Ext}^2_{\mathsf D_1}(A_0,A_2).}
\end{equation}
Under the standard identification of derived $\operatorname{Ext}$ with Yoneda extension
classes in the exact 1-motive model, the right-hand side is precisely the Yoneda product.
More generally a length-$m$ composable bar cell maps to the $m$-fold Yoneda product.
\end{theorem}

\begin{proof}
The two-sided bar differential and concatenation used in the Cartan--Moore comparison are
the standard dg model for derived composition.  By definition
$\operatorname{cl}_1$ sends the degree-one bar generators to the corresponding shifted
morphisms in $\mathsf D_1$.  Therefore a length-two bar word maps to composition of the
two shifted morphisms, with the same bar/Koszul sign convention already fixed in the
Morita comparison.  In a derived category,
$\operatorname{Hom}(A,B[1])$ represents $\operatorname{Ext}^1(A,B)$ and composition of
shifted morphisms represents splicing of the corresponding one-extensions.  This is the
Yoneda product, proving \eqref{eq:bar-yoneda-product}.  Iteration gives the last
statement.
\end{proof}

\begin{corollary}[Classical Cartan-generated extension sector]
\label{cor:classical-cartan-generated-extension-sector}
The bar-length filtration of the Cartan-generated root/Moore dg sector has a canonical
classical shadow in
\[
 \operatorname{Ext}^{\bullet}_{\mathsf D_1(k;\Lambda)}.
\]
In particular its secondary composable Moore/bar sector maps canonically to
$\operatorname{Ext}^2$ and its higher composable bar sector maps to the corresponding
iterated Yoneda products.  This statement concerns the generated root/Moore sector and
does not assert that every higher classical motivic extension is Cartan-generated.
\end{corollary}

\begin{firewall}[The whole Selmer--Cartan carrier is not a Deligne 1-motive]
\label{fw:cartan-layer-not-classical-motive}
The functor $\operatorname{cl}_1$ is a classical shadow of the marked root/Moore sector.
It deliberately forgets Selmer local conditions, Cartan curvature, support history and
higher coherence not represented by the Moore/bar generators.  Thus
\cref{thm:classical-low-sector-comparison,thm:bar-yoneda-comparison} close the
0/1-motivic and Yoneda comparison required for the Cartan-generated sector, but they do
not identify the full support-indexed Cartan stack with the category of Deligne 1-motives and
do not claim essential surjectivity onto Voevodsky motives.
\end{firewall}

\section{Finite lifting-system notation inside the full-motive tower}\label{sec:provisional-full-tower}
By \cref{cor:typed-finite-recursive-motivic-tower}, for every finite obstruction-height ceiling
$M$ we may write
\[
 \mathbb M_{\rm Cart,\le M}^{\rm MR}
 =\{\mathcal M^{(n)}_{I,\lambda,\nu}\}_{0\le n\le M}
\]
for the typed finite lifting system over a fixed full-motive state.  When no ceiling is displayed,
\[
 \mathbb M_{\rm Cart}^{\rm MR}
\]
denotes the compatible finite-height presentation, not the full motive tower itself
and not a coefficient-depth inverse limit.  The actual
two-level notation remains
\[
 \DepthFib\longrightarrow\FullMot
 \stackrel{\tauone}{\longrightarrow}\RecOneMot.
\]
Passage to an infinite obstruction-height colimit, when needed, is a separate
categorical operation and is not identified with valuation depth.

\part{Stack structure and marked Morita comparison}

\section{Moore-enhanced derived-stack globalization}
\label{sec:stackification}

A naive support-site stackification is not the correct construction.  The
exact-support Moore package is invisible on every proper face, so ordinary
Segal or \(2\)-Segal descent cannot reconstruct a completed support stage by
\cref{P2R-thm:v42-no-segal}.  We therefore fix one compatible cofibrant dg
presentation and supplement proper faces by the detector-oriented
exact-support Moore line.

\begin{theorem}[Moore-enhanced stack for a fixed presentation]
\label{thm:moore-enhanced-stack-globalization}
Fix a finite ordered support \(S\) and a compatible cofibrant dg presentation
of the finite Cartan--motivic system.  Then:
\begin{enumerate}[label=\textup{(\roman*)}]
\item at each exact support \(I\subseteq S\), proper-face data together with
      the detector-oriented Moore obstruction line form the enhanced matching
      datum;
\item the completed support-\(I\) stage is recovered by the root/Moore
      attachment of \cref{P2R-thm:v42-moore-matching-completion};
\item attachments of the same cardinality commute by
      \cref{P2R-prop:v42-same-rank-commute}, so cardinality-wise completion
      gives the single dg category
      \(\mathfrak G^{\rm rec}_{S,N}\) of the fixed presentation;
\item its moduli of pseudo-perfect modules
      \[
       \boxed{
       \mathfrak M^{\rm rec}_{S,N}
       :=\mathcal M_{\mathfrak G^{\rm rec}_{S,N}}}
      \]
      is the derived stack of \cref{P2R-thm:v42-stack-globalization};
\item exact-support truncation of the universal family recovers each finite
      completed Moore--Reedy stage of this presentation.
\end{enumerate}
The theorem itself is constructed in a chosen dg presentation; its transport
to another marked presentation is supplied by
\cref{P2M-thm:marked-morita-presentation-independence}.
\end{theorem}

\begin{proof}
The no-Segal proposition identifies the missing line.  The matching theorem
attaches its root/Moore antecedent, same-rank commutation permits
cardinality-wise iteration, and
\cref{P2R-thm:v42-single-dg-globalization} produces the fixed dg category.
Applying the moduli-of-pseudo-perfect-modules construction gives the displayed
stack.  Restriction deletes precisely the non-surviving exact-support
summands, proving the truncation assertion.
\end{proof}

\begin{corollary}[No additional fixed-presentation stack-effectivity gate]
\label{cor:no-additional-stack-effectivity-gate}
For the chosen dg presentation, stack globalization introduces no proof
obligation beyond the Moore-enhanced matching and single-dg globalization
statements.  Comparison with another presentation is covered by
\cref{P2M-thm:marked-morita-presentation-independence} when all its marking
hypotheses hold.  Comparison with an independently specified classical motivic
moduli stack remains a separate problem.
\end{corollary}

\section{Morita presentation comparison}
\label{sec:morita-invariance}

\begin{definition}[Pointed marked Morita equivalence]
\label{P2M-def:pointed-marked-morita-equivalence}
Let \(C\) and \(C'\) be compatible cofibrant/two-chart dg presentations.  A
\emph{pointed marked Morita equivalence} is an invertible perfect
\((C,C')\)-bimodule \(P\), together with coherent identifications under
\(-\otimes_C^{\mathbf L}P\) of:
\begin{enumerate}[label=\textup{(\roman*)}]
\item the distinguished root object, Boolean idempotent summands, the Moore
      pair \((b,c)\), and cyclotomic transfer;
\item every support deletion, permutation, exact-support channel, latching
      polynomial, and same-rank attachment map;
\item the role-separated objects \(d_{q_i}^{\rm sel},e_{q_j}^{\rm sel}\),
      their maps \(\lambda_{ij}^{\rm sel}\), and their marked higher cells;
\item transport, reduced-history, coefficient-reduction, jet, pointing, and
      coherence markings at every finite stage.
\end{enumerate}
The identifications are required to commute with cutoff restriction.  Thus the
word ``marked'' includes the structure needed in the conclusion; it is
strictly stronger than an unpointed equivalence of perfect categories.
\end{definition}

\begin{theorem}[Marked Morita independence of the complete tower]
\label{P2M-thm:marked-morita-presentation-independence}
\label{problem:morita-presentation-independence}
Let \(P:C\simeq_{\rm Mor}C'\) be a pointed marked Morita equivalence in the
sense of \cref{P2M-def:pointed-marked-morita-equivalence}.  Then, functorially
in every finite support \(S\), coefficient stage, and obstruction/jet ceiling,
there are equivalences
\begin{equation}
 \boxed{
 \mathfrak T^{\Mot,\MR}_{N;S,\le M}(C)
 \simeq
 \mathfrak T^{\Mot,\MR}_{N;S,\le M}(C')}
 \label{P2M-eq:marked-morita-finite-tower}
\end{equation}
which intertwine support deletion, same-rank Fubini, coefficient reduction,
jet successors, all standard readouts, and the object-level incidence maps
\[
 d_{q_i}^{\rm sel}\longrightarrow e_{q_j}^{\rm sel}.
\]
They induce a Morita equivalence of the completed dg categories
\[
 \mathfrak G^{\rm rec}_{S,N}(C)
 \simeq_{\rm Mor}
 \mathfrak G^{\rm rec}_{S,N}(C')
\]
and, relative to the chosen marked equivalence, a canonical equivalence of
pseudo-perfect-module stacks
\begin{equation}
 \boxed{
 \mathcal M^{\rm MR}(C)\simeq\mathcal M^{\rm MR}(C').}
 \label{P2M-eq:marked-morita-stack}
\end{equation}
Hence the complete marked tower is presentation independent inside the
groupoid of presentations and pointed marked Morita equivalences.
\end{theorem}

\begin{proof}
Derived Morita theory identifies the functor represented by \(P\) with an
exact equivalence
\[
 \Phi_P: \Perf(C)\xrightarrow{\sim}\Perf(C')
\]
whose inverse is represented by an inverse bimodule
\cite{ToenDerivedMorita}.  Exact equivalences of stable idempotent-complete dg
categories preserve finite sums, retracts, shifts, cones, and finite homotopy
limits and colimits.

At support cardinality zero the first marking clause identifies the root and
Moore seed.  Suppose all stages on proper supports have been identified.  The
enhanced matching object at support \(I\) is a finite homotopy limit of those
proper faces together with the marked exact-support Moore line.  The first two
marking clauses identify this diagram, so \(\Phi_P\) identifies its matching
object.  The completed \(I\)-stage is the corresponding finite root/Moore
attachment, hence a finite homotopy pushout (equivalently, in the stable
setting, a cone construction); exactness identifies it as well.  Induction on
\(|I|\), followed by induction on the finite obstruction/jet ceiling, proves
\eqref{P2M-eq:marked-morita-finite-tower}.  The cutoff-compatible coherences in
the definition make these equivalences natural for deletion and successor
maps.  The third clause and
\cref{P2M-thm:role-separated-objectification} give the asserted transport of
the independent role objects and control edges.  The fourth clause gives the
remaining history and readout compatibilities.

Same-rank attachment commutes before and after applying \(\Phi_P\), so the
finite equivalences assemble into an equivalence of the cardinality-wise
completed dg presentations.  Finally, a Morita equivalence identifies the
dg categories of modules and carries pseudo-perfect modules to
pseudo-perfect modules.  The moduli-of-objects construction is functorial for
this equivalence \cite{ToenVaquie2007}, yielding
\eqref{P2M-eq:marked-morita-stack}.
\end{proof}

\begin{remark}[Scope of presentation independence]
\label{P2M-rem:marked-morita-scope}
Bare Morita equivalence does not manufacture the Cartan, Moore, incidence, or
higher-operation markings.  The theorem says that once two presentations of
the same marked datum are connected by the stated coherent equivalence, every
finite construction and the stack are invariant.  It does not identify the
tower with an independently specified classical motivic moduli stack and does
not prove that an arbitrary external selected carrier admits the required
objectification/evaluation marking.
\end{remark}

\medskip
The structural two-level tower is now complete.  We close with an arithmetic
witness for its source-level repeated and terminal obstruction data.  The
special \(31\)-adic numbers certify non-vacuity; they are not additional
coordinates of the motivic tower.

\part{The \(31\)-adic arithmetic witness}

\section{Witness specification}
\label{sec:paper1-interface-contract}
The source tower above supplies the filtered repeated branch, the finite
\(\mathbf Z\)-linear Moore/PD carrier, and the universal higher-obstruction
calculus.  This part specializes those constructions; it does not repeat the
general confluence or realization theory.

Write \(\mathfrak R_{31}^{[31]}\) for the multiplicity-
\(31\) specialization of the universal source obstruction and
\(h_{31}^{\rm rec}\) for its next semi-free cell.  These are source-side
objects; no motivic or root-stack structure is inferred from the notation.

\begin{definition}[Arithmetic witness contract]\label{def:paper1-witness-contract}
A \(31\)-adic witness consists of:
\begin{enumerate}[label=\textup{(W\arabic*)}]
\item an explicit marked repeated cubic line satisfying the filtered source input;
\item an integral carry normalization at every finite coefficient depth used;
\item a coefficient-typed arithmetic target for the terminal obstruction;
\item a proof that the terminal class is locally trivial but globally nonzero.
\end{enumerate}
\end{definition}

\begin{warning}[Typed uses of \(31\)]\label{warn:typed-31-coordinates}
The support label \(31\), the repeated arithmetic prime \(q=31\), the finite
coefficient-depth prime in \(31^r\), and a local place above \(31\) are not
interchangeable coordinates.  Their numerical coincidence belongs to this
example, not to the general source tower.
\end{warning}

\section{Arithmetic inputs for the witness}\label{sec:external-inputs}

Only three nonstandard external results are used.
\begin{enumerate}[label=\textup{(X\arabic*)},leftmargin=3.2em]
\item\label{X:knospe-qi} The rank-one \(p\)-adic \(L\)-criteria of \cite{KnospeSpecialValues2026} and the split-prime Massey criterion of \cite{QiMasseyLambda}, used to certify exact \(\lambda_{31}=2\) for \(K_0=\mathbf Q(\sqrt{-331})\).
\item\label{X:llsww} The procyclic Bockstein--Massey theorem of \cite{LLSWW}, Theorem~4.3.1, used with \(R=\mathbf F_{31}\), \(C=C_{31}\), \(n=30\), and \(T=\mu_{31}\).
\item Standard class-field, Kummer, and Poitou--Tate duality facts are used in the form recorded in \cite{NSW}.
\end{enumerate}
Every numerical certificate needed below is displayed in the body.
\section{Explicit \texorpdfstring{\(q=31\) exact-\(\lambda=2\)}{q=31 exact-lambda=2} arithmetic branch}
\label{sec:v063-explicit-lambda2}

We now close the residual arithmetic existence gate of the source comparison by a
concrete imaginary quadratic calculation.  This section uses only the
rank-one \(p\)-adic \(L\)-function criteria of Knospe
\cite{KnospeSpecialValues2026} and the split-prime Massey criterion of Qi
\cite{QiMasseyLambda}.  No older arithmetic witness is used.

\subsection{The field, splitting, and class number}

Put
\begin{equation}
 \boxed{
 K:=\mathbf Q(\sqrt{-331}),
 \qquad D_K=-331,
 \qquad p=31.
 }
 \label{eq:v063-field}
\end{equation}
The discriminant \(-331\) is fundamental.  Since
\[
 -331\equiv 10\pmod{31},
 \qquad
 14^2\equiv10\pmod{31},
\]
we have
\begin{equation}
 \boxed{
 31\mathcal O_K=\mathfrak P_0\widetilde{\mathfrak P}_0,
 \qquad
 \mathfrak P_0=(31,\sqrt{-331}-14).
 }
 \label{eq:v063-splitting}
\end{equation}

The reduced primitive positive-definite binary quadratic forms of discriminant
\(-331\) are exactly
\[
 (1,1,83),\qquad (5,-3,17),\qquad (5,3,17).
\]
Hence
\begin{equation}
 \boxed{h_K=3,\qquad 31\nmid h_K.}
 \label{eq:v063-class-number}
\end{equation}
Thus the hypotheses of Qi's split-prime imaginary-quadratic criterion apply.

\subsection{First Bernoulli certificate: \texorpdfstring{\(\lambda_{31}>1\)}{lambda31 > 1}}

Let
\[
 \theta:=\chi_{-331}
\]
be the primitive odd quadratic Dirichlet character of conductor \(331\), and
put \(\chi=\theta\omega\), where \(\omega\) is the \(31\)-adic Teichmuller
character.  Because \(\theta(31)=1\), the rank-one \(31\)-adic \(L\)-function
\(L_{31}(s,\chi)\) has its trivial zero at \(s=0\).

Using the normalization
\[
 B_{31,\theta}
 =331^{30}\sum_{a=1}^{331}
   \theta(a)B_{31}\!\left(\frac a{331}\right),
\]
an exact integer calculation gives
\begin{equation}
 \boxed{
 v_{31}(B_{31,\theta})=3,
 \qquad
 \frac{B_{31,\theta}}{31^3}\equiv13\pmod{31}.
 }
 \label{eq:v063-B31-certificate}
\end{equation}
Equivalently,
\[
 B_{31,\theta}\equiv387283\pmod{31^4}.
\]
Knospe's rank-one criterion therefore gives
\begin{equation}
 \boxed{\lambda_{31}(\chi)>1.}
 \label{eq:v063-lambda-gt1}
\end{equation}

\subsection{Second-kind twist certificate: exact \texorpdfstring{\(\lambda_{31}=2\)}{lambda31 = 2}}

Let \(\psi\) be the second-kind character of conductor \(31^2\) and order
\(31\), normalized by
\[
 \psi(1+31)=\zeta_{31}
\]
and trivial on the Teichmuller subgroup.  Concretely, for a unit
\(a\bmod31^2\), put
\[
 \tau(a):=a^{31}\pmod{31^2},
 \qquad
 a\tau(a)^{-1}\equiv1+31j(a)\pmod{31^2},
\]
with \(j(a)\in\mathbf Z/31\), and set
\(\psi(a)=\zeta_{31}^{j(a)}\).

\begin{lemma}[Second-kind twist hypothesis audit]
\label{W31-lem:knospe-second-kind-hypotheses}
The characters used here satisfy the hypotheses of the first part of
Knospe's second-kind Bernoulli-value formula \cite{KnospeSpecialValues2026} with
\[
 p=31,\qquad i=k=n=1,\qquad e=1.
\]
More precisely:
\begin{enumerate}[label=\textup{(\roman*)}]
\item \(\psi\) has exact order \(31\) and primitive conductor \(31^2\);
\item \(\theta\) is primitive odd of conductor \(331\), \(\psi\) is even,
      and \(\theta\psi\) is primitive odd of conductor
      \(331\cdot31^2\);
\item \(\chi=\theta\omega\) is a nontrivial even character of the first kind
      and
      \(\mathcal O_\chi=\mathbf Z_{31}[\operatorname{Im}(\chi)]
      =\mathbf Z_{31}\), so the ramification index \(e\) in Knospe's theorem
      is \(1\).
\end{enumerate}
\end{lemma}

\begin{proof}
There is a direct product decomposition
\[
 (\mathbf Z/31^2)^\times
 \cong \mu_{30}\times
       (1+31\mathbf Z)/(1+31^2\mathbf Z).
\]
The second factor is cyclic of order \(31\), and the displayed definition of
\(j(a)\) is its additive coordinate.  Thus \(\psi\), trivial on
\(\mu_{30}\) and nontrivial on \(1+31\), has exact order \(31\) and cannot
factor modulo \(31\); its conductor is exactly \(31^2\).  In particular
\(\psi(-1)=1\).

The quadratic character \(\theta=\chi_{-331}\) is primitive odd because
\(-331\) is a fundamental discriminant.  Coprimality of \(331\) and \(31^2\)
then makes \(\theta\psi\) primitive of the product conductor, and its parity
is odd.  Finally \(\omega\) takes values in the Teichmuller group
\(\mu_{30}\subset\mathbf Z_{31}^{\times}\), while \(\theta\) takes values in
\(\{\pm1\}\subset\mu_{30}\).  Hence \(\mathcal O_\chi=\mathbf Z_{31}\) and
the theorem's coefficient-field ramification index is \(e=1\).  Since
\(\chi=\theta\omega^1\), its index is \(i=1\), so Knospe's convention gives
\(k=1\); the order \(31=31^1\) gives \(n=1\).
\end{proof}

Write
\[
 F:=331\cdot31^2=318091,
 \qquad
 A(\zeta_{31})
 :=\sum_{a=1}^{F}\theta(a)\psi(a)a.
\]
Since \(\theta\psi\) is nontrivial,
\[
 B_{1,\theta\psi}=\frac{A(\zeta_{31})}{F}.
\]
Put \(\pi=\zeta_{31}-1\).  Reduce \(A(1+\pi)\) modulo
\(\Phi_{31}(1+\pi)\) to degree at most \(29\):
\begin{equation}
 A(1+\pi)=\sum_{j=0}^{29}a_j\pi^j.
 \label{eq:v063-A-expansion}
\end{equation}
The exact integer coefficient calculation begins
\[
 a_0=-39443284,
 \qquad
 a_1=-512762692,
 \qquad
 a_2=-3006596132,
\]
and has valuation profile
\begin{equation}
 \boxed{
 v_{31}(a_0)=v_{31}(a_1)=3,
 \qquad
 v_{31}(a_j)=2\quad(2\le j\le29).
 }
 \label{eq:v063-coefficient-valuation-profile}
\end{equation}
For auditability, after division by \(31^2\), the coefficient vector modulo
\(31\) is
\begin{equation}
\begin{split}
 &(0,0,1,11,6,23,14,15,4,30,15,10,3,3,14,\\
 &\hspace{12mm}21,1,18,15,11,3,4,7,30,27,19,6,29,28,2).
\end{split}
\label{eq:v063-coefficient-vector}
\end{equation}
In particular the first nonzero residual coefficient occurs at \(\pi^2\).
Since \(v_{31}(\pi)=1/30\),
\[
 v_{31}(A)=2+\frac2{30}=\frac{31}{15},
 \qquad
 v_{31}(F)=2,
\]
and hence
\begin{equation}
 \boxed{
 v_{31}(B_{1,\theta\psi})=\frac1{15}<1.
 }
 \label{eq:v063-B1-valuation}
\end{equation}
By \cref{W31-lem:knospe-second-kind-hypotheses}, the first case of Knospe's second-kind twist formula
now gives
\begin{equation}
 \boxed{
 \lambda_{31}(K)
 =\lambda_{31}(\theta\omega)
 =30\,v_{31}(B_{1,\theta\psi})
 =2.
 }
 \label{eq:v063-lambda-exact2}
\end{equation}
Thus this is an exact-\(\lambda=2\), not merely a \(\lambda>1\), branch.

\subsection{Explicit principalizing Kummer element}

Write
\[
 \omega_K:=\frac{1+\sqrt{-331}}2,
 \qquad
 \mathcal O_K=\mathbf Z[\omega_K].
\]
Then
\[
 \operatorname N(a+b\omega_K)=a^2+ab+83b^2.
\]
The exact norm equation
\[
 164^2+164\cdot5+83\cdot5^2=29791=31^3
\]
gives
\begin{equation}
 \boxed{
 \alpha:=164+5\omega_K
 =\frac{333+5\sqrt{-331}}2,
 \qquad
 \operatorname N_{K/\mathbf Q}(\alpha)=31^3.
 }
 \label{eq:v063-alpha}
\end{equation}
Modulo \(\mathfrak P_0=(31,\sqrt{-331}-14)\), the numerator satisfies
\[
 333+5\cdot14=403=13\cdot31,
\]
whereas at the conjugate prime it is nonzero modulo~\(31\).  Therefore
\begin{equation}
 \boxed{(\alpha)=\mathfrak P_0^3.}
 \label{eq:v063-alpha-ideal}
\end{equation}
This is the principalizing Kummer element required in Qi's split-prime setup.

\subsection{Qi handoff: the literal repeated cubic is nonzero}

Let \(x\) denote the residual cyclotomic character and let the Kummer class of
\(\alpha\) be marked in an old direction slot.  Qi proves that for an
imaginary quadratic field with \(p\) split and \(p\nmid h_K\), assuming
\(\lambda\ge n-1\), one has
\[
 \lambda\ge n
 \quad\Longleftrightarrow\quad
 (\underbrace{x,\ldots,x}_{n-1},\alpha)_\rho=0
\]
for the proper defining system.  Since
\(\lambda_{31}(K)=2\), the binary gate vanishes while the next proper repeated
cubic does not:
\begin{equation}
 \boxed{
 (x,x,\alpha)_\rho\ne0.
 }
 \label{eq:v063-qi-cubic-nonzero}
\end{equation}

\begin{remark}[Proper value, not a zero-indeterminacy assertion]
\label{W31-rem:qi-proper-value-indeterminacy}
The subscript \(\rho\) denotes Qi's specified proper defining system.  The
criterion concerns the single value attached to that defining system.  We do
not replace it by the full unpointed Massey subset, and no vanishing of the
ordinary Massey indeterminacy is assumed or used.
\end{remark}
This is exactly a residual repeated-input cubic coefficient of shape
\(\gamma_2(e_{31})e_v\) after choosing an old marked slot \(v\in I_{10}\)
for the Kummer direction.

\begin{warning}[The old slot is a marking, not an arithmetic-prime identity]
\label{warn:v063-old-slot-marking}
The choice of \(v\in I_{10}\) labels the old controller direction receiving
the Kummer class of \(\alpha\).  It does not assert that \(\alpha\),
\(\mathfrak P_0\), or the field \(K\) is arithmetically equal to the rational
prime labelled by \(v\).  This is the same support-label versus realization
 firewall used throughout this witness.
\end{warning}

\subsection{Explicit arithmetic cubic incidence}

At residual depth one, the principal-unit torsor in the source construction disappears:
every pointed comparison between the arithmetic repeated line and the
 normalized source \(pqq\) line reduces to the same map modulo~\(31\).
Consequently the nonzero class
\eqref{eq:v063-qi-cubic-nonzero} has nonzero cubic incidence with the
repeated source line fixed above.  With the pointed normalization of
 the source comparison the residual incidence scalar is one.  The ownership firewall keeps this
statement at cubic order: no full high-jet arithmetic transport is inferred
from the line comparison alone.
\begin{theorem}[Explicit \((-331,31)\) arithmetic cubic seed]
\label{thm:v063-explicit-arithmetic-cubic}
For
\[
 K=\mathbf Q(\sqrt{-331}),\qquad p=31,
\]
with \(\alpha\) as in \eqref{eq:v063-alpha}, Qi's exact-\(\lambda=2\)
proper branch has a nonzero literal repeated cubic coefficient at \(q=31\).
Under the canonical pointed source comparison, this cubic generator
maps to the normalized \(pqq\) PD generator with scalar one.
\end{theorem}

\begin{proof}
The field-theoretic hypotheses are
\eqref{eq:v063-splitting}--\eqref{eq:v063-class-number}.  The exact Iwasawa
invariant is \eqref{eq:v063-lambda-exact2}, and
\eqref{eq:v063-alpha-ideal} supplies Qi's principalizing Kummer element.
Hence \eqref{eq:v063-qi-cubic-nonzero} gives the nonzero residual repeated
cubic seed.  Its pointed cyclic line is the arithmetic specialization of
\cref{P1C-thm:formal-filtered-alignment}.
\end{proof}

\section{Explicit integral carry normalization at \texorpdfstring{$(-331,31)$}{(-331,31)}}
\label{sec:v064-carry-normalization}

The residual calculation above determines the hit/no-hit question, but the
explicit branch allows the finer integral carry coordinate to be computed as
well.  We use the same prime
\[
 \mathfrak P_0=(31,\sqrt{-331}-14)
\]
and the same Kummer element
\[
 \alpha=\frac{333+5\sqrt{-331}}2,
 \qquad
 (\alpha)=\mathfrak P_0^3.
\]
Let
\[
 \iota_0:K\hookrightarrow\mathbf Q_{31}
\]
be the embedding determined by \(\mathfrak P_0\), and let
\(r_0=\iota_0(\sqrt{-331})\) be the Hensel root satisfying
\[
 r_0^2=-331,
 \qquad
 r_0\equiv14\pmod{31}.
\]
Since \(v_{31}(\iota_0(\alpha))=3\), write
\begin{equation}
 \boxed{
 \iota_0(\alpha)=31^3u,
 \qquad
 u\in\mathbf Z_{31}^{\times}.}
 \label{eq:v064-alpha-unit}
\end{equation}

\begin{lemma}[Finite Hensel certificate]
\label{lem:v064-hensel-certificate}
One may take
\begin{equation}
 \boxed{
 r_0\equiv795288553\pmod{31^6}.}
 \label{eq:v064-hensel-root}
\end{equation}
For the corresponding unit in \eqref{eq:v064-alpha-unit},
\begin{equation}
 \boxed{
 u\equiv7157\pmod{31^3},
 \qquad
 u^{30}-1\equiv26\cdot31^2\pmod{31^3}.}
 \label{eq:v064-unit-certificate}
\end{equation}
Consequently
\begin{equation}
 \boxed{v_{31}(u^{30}-1)=2.}
 \label{eq:v064-unit-fermat-depth}
\end{equation}
\end{lemma}

\begin{proof}
The displayed integer satisfies
\[
 795288553^2+331\equiv0\pmod{31^6}
\]
and is congruent to \(14\) modulo~\(31\), so Hensel uniqueness identifies it
with the required root to that precision.  Substituting into
\((333+5r_0)/2\), dividing by \(31^3\), and reducing modulo \(31^3\) gives
\(u\equiv7157\).  Direct modular exponentiation then gives
\[
 7157^{30}-1\equiv24986=26\cdot31^2\pmod{31^3},
\]
with \(26\ne0\) in \(\mathbf F_{31}\).
\end{proof}

\begin{proposition}[Exact carry depth]
\label{prop:v064-exact-carry-depth}
For the explicit \((-331,31)\) Qi branch, the distinguished integral binary
carry has exact depth
\begin{equation}
 \boxed{s_{\rm car}=1.}
 \label{eq:v064-scar-one}
\end{equation}
\end{proposition}

\begin{proof}
For an odd prime, the Iwasawa logarithm of a unit has the same valuation as
\(u^{p-1}-1\) whenever the latter has positive valuation.  Thus
\cref{lem:v064-hensel-certificate} gives
\[
 v_{31}(\log_{31}u)=2.
\]
the source construction identifies the root-projected integral binary carry
\[
 \gamma
 =\widetilde\chi\smile\kappa_{\mathbf Z_{31}}(\alpha)
 \in H^2(G_{K,S},\mathbf Z_{31}(1))
\]
with the additive cyclotomic character evaluated on the local unit part of
\(\alpha\), up to a fixed unit coming from the reciprocity orientation.  The
primitive \(\mathbf Z_{31}\)-valued cyclotomic character has valuation-one
period: equivalently it is a unit multiple of
\[
 x\longmapsto
 \frac{\log_{31}\langle x\rangle}{\log_{31}(1+31)}.
\]
Since
\(v_{31}(\log_{31}(1+31))=1\), the carry coordinate has valuation
\[
 2-1=1.
\]
This valuation is independent of the allowed unit normalization, hence is
precisely the intrinsic \(s_{\rm car}\).
\end{proof}

\begin{proposition}[Residual unit of the carry coordinate]
\label{prop:v064-carry-unit-five}
Fix the standard local normalization
\begin{equation}
 \ell_{31}(x)
 :=\frac{\log_{31}\langle x\rangle}{\log_{31}(32)}
 \label{eq:v064-standard-cyclotomic-coordinate}
\end{equation}
and choose the orientation of the global reciprocity line so that its
\(\mathfrak P_0\)-local invariant is the displayed coordinate.  If
\[
 \gamma E=31a_0\widetilde\kappa,
 \qquad
 a_0\in\mathbf Z_{31}^{\times},
\]
for a primitive lift \(\widetilde\kappa\) of the residual repeated cubic
point, then
\begin{equation}
 \boxed{a_0\equiv5\pmod{31}.}
 \label{eq:v064-a0-five}
\end{equation}
More precisely, with the normalization above,
\begin{equation}
 \ell_{31}(u)\equiv155=31\cdot5\pmod{31^2}.
 \label{eq:v064-log-coordinate-mod31sq}
\end{equation}
\end{proposition}

\begin{proof}
By \eqref{eq:v064-unit-certificate},
\[
 \log_{31}(u^{30})
 \equiv26\cdot31^2\pmod{31^3},
\]
because the next logarithmic term has valuation at least four.  Since
\(30\equiv-1\pmod{31}\), division by \(30\) gives
\[
 \frac{\log_{31}u}{31^2}\equiv5\pmod{31}.
\]
Also
\[
 \frac{\log_{31}(32)}{31}\equiv1\pmod{31}.
\]
Taking the quotient yields
\[
 \frac{\ell_{31}(u)}{31}\equiv5\pmod{31},
\]
which is the claimed unit coordinate.
\end{proof}

\begin{theorem}[Explicit carry-normalized \(31\)-adic pointing]
\label{thm:v064-explicit-carry-pointing}
Let \([5]\in\mathbf Z_{31}^{\times}\) be the Teichm\"uller lift of
\(5\in\mathbf F_{31}^{\times}\).  On the explicit \((-331,31)\) arithmetic
branch, the source carry-normalized primitive generator is uniquely determined
by
\begin{equation}
 \boxed{
 \gamma E
 =31[5]\,\kappa_{\rm Qi}^{\rm can}.}
 \label{eq:v064-canonical-carry}
\end{equation}
Equivalently, if \(a_0\) is the primitive-lift coefficient of
\cref{prop:v064-carry-unit-five}, then
\begin{equation}
 \boxed{
 \kappa_{\rm Qi}^{\rm can}
 =\frac{a_0}{[5]}\,\widetilde\kappa.}
 \label{eq:v064-canonical-kappa}
\end{equation}
At finite coefficient depth \(31^r\), the normalized bifiltered third-jet
relation therefore has binary coefficient
\begin{equation}
 \boxed{
 \alpha_{31^r}
 =
 \begin{cases}
  0,&r=1,\\[1mm]
  2^{-1}\,31[5]\pmod{31^r},&r\ge2,
 \end{cases}}
 \label{eq:v064-finite-carry-coefficient}
\end{equation}
and cubic coefficient one on
\(e_v\gamma_2(e_{31})\kappa_{\rm Qi}^{\rm can}\).
\end{theorem}

\begin{proof}
The exact depth is \cref{prop:v064-exact-carry-depth} and the residual unit is
\cref{prop:v064-carry-unit-five}.  The source carry-normalization removes the
principal-unit factor from a primitive lift and retains precisely the
Teichm\"uller unit.  This gives \eqref{eq:v064-canonical-carry} and
\eqref{eq:v064-canonical-kappa}.  The finite coefficient formula is the
normalized Qi bifiltered relation after the standard factor \(2^{-1}\); at
depth one the positive-filtration binary carry vanishes, while at every
\(r\ge2\) the depth-one carry survives exactly as displayed.
\end{proof}

\begin{corollary}[Carry-normalized third-jet datum is rigidified]
\label{cor:v064-thirdjet-roof-rigidified}
For the explicit \((-331,31)\) branch, the residual cubic generator, the
integral binary carry, its depth, its Teichm\"uller unit, and the canonical
cyclotomic-coefficient comparison pointing are all fixed.  Hence the arithmetic
\emph{third-jet} datum at every finite depth \(31^r\) is pointed without any
remaining principal-unit ambiguity.  This statement does not assert an
arithmetic lift of the subsequent source cells.
\end{corollary}

\section{The first coefficient-compatible classical point}
\label{sec:v073-root-typed}

The preliminary arithmetic search deliberately separated existence of a classical
one-dimensional Shafarevich line from the stronger problem of realizing the
\emph{same marked source obstruction detector}.  We now impose the missing
coefficient typing.  This changes the arithmetic point.

\subsection{Cyclotomic coefficient type of the source obstruction}

The source repeated cubic branch starts on the pointed cyclotomic
coefficient line
\[
 \mathbf F_{31}(1)E.
\]
The old and pure first tangents are diagonal/weight-zero directions, while
the distinguished repeated recursive cell remains in the same marked root
slot.  Thus the top stage-\(30\) factorization
\[
 \pi_{30}\mathfrak R_{31}^{[31]}
 =
 [\Lambda_{30}^{[31]},\bar h_{30}^{\rm rec}]_{\rm gr}
\]
contains exactly one root input and one diagonal input.

\begin{definition}[Cyclotomic-coefficient-typed marked realization]
\label{def:v073-root-typed-realization}
Let
\[
 T_{\rm AT}(F)
 =
 \operatorname{End}_{\mathbf F_{31}}
 \bigl(\mathbf F_{31}\oplus\mathbf F_{31}(1)\bigr).
\]
A marked dg realization into
\(C^\bullet(G_{F,S},T_{\rm AT}(F))\) is called
\emph{cyclotomic-coefficient-typed} if it preserves the four Artin--Tate matrix-weight
corners and sends the normalized source repeated line to the
cyclotomic coefficient corner
\[
 C^\bullet(G_{F,S},\mathbf F_{31}(1))E.
\]
\end{definition}

\begin{proposition}[Cyclotomic coefficient type of the height-\(31\) detector]
\label{prop:v073-one-root-typing}
Under every cyclotomic-coefficient-typed marked dg realization, the complete source
height-\(31\) detector has coefficient type
\begin{equation}
 \boxed{
 \rho(\mathfrak R_{31}^{[31]})
 \in
 H^2(G_{F,S},\mathbf F_{31}(1))E.}
 \label{eq:v073-root-valued-detector}
\end{equation}
In particular it cannot map to the opposite-root
\(\mathbf F_{31}(-1)F\)-corner.
\end{proposition}

\begin{proof}
At cubic order the assertion is the defining preservation of the marked
root line.  Every subsequent distinguished repeated cell is obtained by
relative semi-free attachment to the same one-old repeated PD sector.
Consequently each generator-specific repeated contribution has exactly one
root-type factor; all additional first-tangent factors are diagonal.

On the Artin--Tate matrix plane the coefficient identities are
\[
 D^2=D,\qquad DE=E,\qquad ED=E^2=0.
\]
Hence every nonzero product containing exactly one \(E\)-root input and
otherwise diagonal inputs again has coefficient type \(E\).
The complete height-\(31\) residue is a sum of such one-root words.  The
stage-\(30\) statement is visible directly from
\[
 [\Lambda_{30}^{[31]},\bar h_{30}^{\rm rec}]_{\rm gr};
\]
the lower cell stages obey the same typing by the repeated-sector
construction.  Therefore the complete detector lies in the cyclotomic coefficient corner.
\end{proof}



\subsection{The correct arithmetic search criterion}

For an imaginary quadratic field \(K\) in which \(31\) splits,
write
\[
 31\mathcal O_K=\mathfrak P_+\mathfrak P_-,
 \qquad
 S_{31}=\{\mathfrak P_+,\mathfrak P_-\}.
\]
Poitou--Tate duality gives
\[
 \Sha^2_{S_{31}}(K,\mu_{31})^\vee
 \simeq
 \Sha^1_{S_{31}}(K,\mathbf F_{31}).
\]
A strict class on the right is an everywhere-unramified
\(\mathbf F_{31}\)-character whose Frobenius is zero at the two primes
above \(31\).  Global class field theory therefore gives:

\begin{proposition}[Cyclotomic-coefficient class-group formula]
\label{prop:v073-root-classgroup-formula}
There is a canonical identification
\begin{equation}
 \boxed{
 \Sha^1_{S_{31}}(K,\mathbf F_{31})
 \simeq
 \operatorname{Hom}\!\left(
 \frac{\operatorname{Cl}(K)}
 {\langle[\mathfrak P_+],[\mathfrak P_-]\rangle},
 \mathbf F_{31}
 \right).}
 \label{eq:v073-root-classgroup-formula}
\end{equation}
Since
\(
[\mathfrak P_-]=[\mathfrak P_+]^{-1}
\),
the denominator is generated by one prime class.
\end{proposition}

Thus a cyclotomic-coefficient-compatible classical point requires
\begin{equation}
 \boxed{
 31\mid
 \frac{h_K}{\operatorname{ord}[\mathfrak P_+]}.}
 \label{eq:v073-root-point-criterion}
\end{equation}
This is different from the opposite-root \(\omega^2\)-criterion of
the preliminary search.


\subsection{Finite minimality search}

For a negative fundamental discriminant \(D\), the class number was computed
by the standard reduced-form inequalities
\[
 |b|\le a\le c,\qquad b^2-4ac=D,
\]
with the usual boundary sign convention.  For every split candidate with
\(31\mid h(D)\), the order of a prime over \(31\) was computed by Hensel
lifting a square root of \(D\) modulo \(31^n\), forming the corresponding
primitive form of leading coefficient \(31^n\), and reducing it.  The first
\(n\) for which the principal reduced form appears is
\(\operatorname{ord}[\mathfrak P_+]\).

\begin{proposition}[Minimal cyclotomic-coefficient-compatible discriminant]
\label{prop:v073-minimal-root-discriminant}
Among negative fundamental discriminants ordered by absolute value for which
\(31\) splits, the first one satisfying
\eqref{eq:v073-root-point-criterion} is
\begin{equation}
 \boxed{D_*=-15391.}
 \label{eq:v073-Dstar}
\end{equation}
For all earlier split discriminants with class number divisible by \(31\),
the quotient
\(
h(D)/\operatorname{ord}[\mathfrak P_+]
\)
is \(1\), \(2\), or \(4\), hence prime to \(31\).

More explicitly, before \(D_*\) the candidates with quotient \(2\) are
\[
 -9191,\ -10439,\ -11051,\ -11983,\ -13719,\ -14271,
\]
the candidates with quotient \(4\) are
\[
 -12815,\ -13655,
\]
and every other split candidate with \(31\mid h(D)\) below \(15391\) has
quotient \(1\).  At \(D_*\) the quotient jumps to \(31\).
\end{proposition}

\begin{proof}
This is a finite exact reduced-form/Hensel enumeration.  For auditability,
the complete list of the quotient-\(1\) candidates is
\[
\begin{gathered}
-719,-911,-2471,-2927,-3251,-4951,-6131,-6439,-6491,\\
-7639,-8647,-9203,-10427,-10484,-11863,-12799,-12923,\\
-13043,-13219,-14303,-14627,-14731,-15343.
\end{gathered}
\]
Together with the quotient-\(2\) and quotient-\(4\) lists above, these are
exactly the split negative fundamental discriminants of smaller absolute
value whose class number is divisible by \(31\).  Reducing the Hensel-lifted
\(31^n\)-forms gives the stated prime-class orders.  The explicit field
calculation below verifies the positive value \(31\) at \(D_*\).
\end{proof}


\subsection{The field \texorpdfstring{\(K_*=\mathbf Q(\sqrt{-15391})\)}{K*=Q(sqrt(-15391))}}

Put
\[
 \boxed{
 K_*:=\mathbf Q(\sqrt{-15391}),
 \qquad
 \omega_*:=\frac{1+\sqrt{-15391}}2.}
 \label{eq:v073-Kstar}
\]
The norm form is
\begin{equation}
 \boxed{
 N(a+b\omega_*)
 =a^2+ab+3848b^2.}
 \label{eq:v073-norm-form}
\end{equation}

\begin{theorem}[Class number, splitting, and the order-three \(31\)-prime]
\label{thm:v073-field-certificate}
One has
\begin{equation}
 \boxed{h_{K_*}=93=3\cdot31.}
 \label{eq:v073-class-number}
\end{equation}
Moreover
\begin{equation}
 \boxed{
 31\mathcal O_{K_*}
 =
 \mathfrak P_+\mathfrak P_-,
 \qquad
 \mathfrak P_+=(31,\omega_*-14),
 \quad
 \mathfrak P_-=(31,\omega_*-18),}
 \label{eq:v073-splitting}
\end{equation}
and
\begin{equation}
 \boxed{\operatorname{ord}[\mathfrak P_+]=3.}
 \label{eq:v073-P-order}
\end{equation}
\end{theorem}

\begin{proof}
The reduced positive-definite primitive forms of discriminant \(-15391\)
have the following leading-coefficient multiplicities:
\[
\begin{array}{c|rrrrrrrrrrr}
a&1&2&4&5&7&8&10&11&13&14&16\\ \hline
\#&1&2&2&2&2&2&4&2&2&4&2
\end{array}
\]
\[
\begin{array}{c|rrrrrrrrrrr}
a&20&22&25&26&28&31&32&35&37&40&41\\ \hline
\#&4&4&2&4&4&2&2&4&2&4&2
\end{array}
\]
\[
\begin{array}{c|rrrrrrrrrrr}
a&44&47&49&50&52&55&56&61&62&65&70\\ \hline
\#&4&2&2&4&4&4&4&2&4&2&2.
\end{array}
\]
The total is \(93\), proving \eqref{eq:v073-class-number}.

Modulo \(31\), the polynomial
\[
 x^2-x+3848
\]
has roots \(14\) and \(18\), giving \eqref{eq:v073-splitting}.
Now put
\begin{equation}
 \boxed{
 \alpha_*:=-121+2\omega_*
 =-120+\sqrt{-15391}.}
 \label{eq:v073-alpha}
\end{equation}
Then
\[
 N(\alpha_*)
 =120^2+15391
 =29791
 =31^3.
\]
Modulo \(\mathfrak P_+\),
\[
 -121+2\cdot14=-93\equiv0\pmod{31},
\]
whereas modulo \(\mathfrak P_-\),
\[
 -121+2\cdot18=-85\not\equiv0\pmod{31}.
\]
Thus
\begin{equation}
 \boxed{(\alpha_*)=\mathfrak P_+^3.}
 \label{eq:v073-P-cube}
\end{equation}

The prime \(\mathfrak P_+\) is not principal: if it were, the norm form
\eqref{eq:v073-norm-form} would represent \(31\).  For \(b=0\) this would
require a square equal to \(31\), while for \(b\ne0\)
\[
 a^2+ab+3848b^2
 =
 \left(a+\frac b2\right)^2+\frac{15391}{4}b^2
 >31.
\]
Hence \([\mathfrak P_+]\ne1\).  Together with
\eqref{eq:v073-P-cube}, this proves exact order \(3\).
\end{proof}

\begin{corollary}[The split class-group quotient has order \(31\)]
\label{cor:v073-class-quotient31}
One has
\begin{equation}
 \boxed{
 \left|
 \operatorname{Cl}(K_*)/
 \langle[\mathfrak P_+],[\mathfrak P_-]\rangle
 \right|
 =31.}
 \label{eq:v073-class-quotient31}
\end{equation}
\end{corollary}

\begin{proof}
The two prime classes are inverse, so they generate the order-\(3\) subgroup
\(\langle[\mathfrak P_+]\rangle\).  Divide
\(h_{K_*}=93\) by \(3\).
\end{proof}


\subsection{The first cyclotomic-coefficient-typed classical \texorpdfstring{\(\Sha^2\)}{Sha2}}

\begin{theorem}[Cyclotomic-coefficient-compatible classical Shafarevich line]
\label{thm:v073-root-sha-one}
For
\[
 S_{31}=\{\mathfrak P_+,\mathfrak P_-\}
\]
one has
\begin{equation}
 \boxed{
 \Sha^1_{S_{31}}(K_*,\mathbf F_{31})
 \simeq\mathbf F_{31},}
 \label{eq:v073-sha1-root}
\end{equation}
and therefore, by Poitou--Tate duality,
\begin{equation}
 \boxed{
 \Sha^2_{S_{31}}(K_*,\mu_{31})
 \simeq\mathbf F_{31}.}
 \label{eq:v073-sha2-root}
\end{equation}
This is the first classical Shafarevich line in the discriminant search that
has the same cyclotomic coefficient type as the source obstruction detector.
\end{theorem}

\begin{proof}
Apply \cref{prop:v073-root-classgroup-formula} and
\cref{cor:v073-class-quotient31}.  The quotient class group has order
\(31\), hence its nonzero homomorphisms to \(\mathbf F_{31}\) form a
one-dimensional space.  Poitou--Tate duality for
\(\mathbf F_{31}\) and its Tate dual \(\mu_{31}\) gives
\eqref{eq:v073-sha2-root}.
\end{proof}


\subsection{A small Frobenius pointing}

The quotient in \eqref{eq:v073-class-quotient31} can be pointed without a
large class-field calculation.  Since \(5\) splits in \(K_*\), put
\begin{equation}
 \boxed{
 \mathfrak Q_5:=(5,\omega_*-2).}
 \label{eq:v073-Q5}
\end{equation}
Its reduced binary quadratic form is
\[
 (5,3,770).
\]
The three classes in
\(\langle[\mathfrak P_+]\rangle\) have reduced representatives
\begin{equation}
 \boxed{
 (1,1,3848),\qquad
 (31,27,130),\qquad
 (31,-27,130).}
 \label{eq:v073-P-subgroup-forms}
\end{equation}
Therefore \([\mathfrak Q_5]\) does not belong to the order-\(3\) subgroup.

\begin{proposition}[The \(5\)-prime points the quotient]
\label{prop:v073-Q5-generator}
The image of \([\mathfrak Q_5]\) generates
\[
 \operatorname{Cl}(K_*)/
 \langle[\mathfrak P_+]\rangle
 \simeq C_{31}.
\]
Consequently there is a unique character
\begin{equation}
 \boxed{
 c_5:
 \operatorname{Cl}(K_*)/
 \langle[\mathfrak P_+]\rangle
 \xrightarrow{\sim}\mathbf F_{31},
 \qquad
 c_5([\mathfrak Q_5])=1.}
 \label{eq:v073-c5}
\end{equation}
\end{proposition}

\begin{proof}
The quotient has prime order \(31\).  The reduced representative
\((5,3,770)\) is distinct from all three representatives in
\eqref{eq:v073-P-subgroup-forms}, so the quotient class of
\(\mathfrak Q_5\) is nonzero and hence generates.
\end{proof}

By Artin reciprocity, \(c_5\) is the unique strict unramified character
\begin{equation}
 \boxed{
 z_5^{\rm root}
 \in
 \Sha^1_{S_{31}}(K_*,\mathbf F_{31}),
 \qquad
 z_5^{\rm root}(\operatorname{Frob}_{\mathfrak Q_5})=1.}
 \label{eq:v073-z5}
\end{equation}
The classes of \(\mathfrak P_\pm\) vanish by construction, so both primes in
\(S_{31}\) split completely in the corresponding cyclic degree-\(31\)
unramified extension.

\begin{definition}[Pointed cyclotomic-coefficient \(\Sha^2\) generator]
\label{def:v073-kappa-root}
Let
\[
 \kappa_5^{\rm root}
 \in
 \Sha^2_{S_{31}}(K_*,\mu_{31})
\]
be the unique class satisfying
\begin{equation}
 \boxed{
 \left\langle
 \kappa_5^{\rm root},z_5^{\rm root}
 \right\rangle_{\rm PT}=1.}
 \label{eq:v073-kappa-root}
\end{equation}
\end{definition}

\begin{corollary}[Pointed cyclotomic-coefficient-compatible detector-line comparison]
\label{cor:v073-root-pointed-line}
There is a unique coefficient-type-compatible pointed isomorphism of
one-dimensional cyclotomic-coefficient lines
\begin{equation}
 \boxed{
 \rho_{5}^{\det,\rm root}:
 \Sha^{2,\det}_{\rm SC,v,31}(31)
 \xrightarrow{\ \sim\ }
 \Sha^2_{S_{31}}(K_*,\mu_{31})}
 \label{eq:v073-root-pointed-line}
\end{equation}
such that
\begin{equation}
 \boxed{
 \rho_{5}^{\det,\rm root}
 \bigl(\mathfrak R_{31}^{[31]}\bigr)
 =
 \kappa_5^{\rm root}.}
 \label{eq:v073-root-pointed-generator}
\end{equation}
Source and target have the same cyclotomic coefficient type.
\end{corollary}


\subsection{The remaining strict dg provenance gate}

The first source obstruction detector uses the repeated-cell cascade only through
\[
 h_3^{\rm rec},\ldots,h_{30}^{\rm rec}.
\]
The height-\(31\) class is then read from that lower jet.  Adjoining an
additional cell \(h_{31}^{\rm rec}\) is a subsequent semi-free extension
problem and is not part of the data needed to define
\(\mathfrak R_{31}^{[31]}\).

\begin{definition}[Pre-resonance marked realization space]
\label{def:v073-pre31-mapspace}
Put
\[
 \mathscr C_{<31}^{\rm str}
 :=
 \overline{\mathscr C}^{\rm sf}_{\le30}.
\]
Let
\[
 \mathcal R_{<31}^{\rm root}(K_*)
\]
be the set of pointed cyclotomic-coefficient-typed strict dg realizations
\begin{equation}
 \rho_{<31}:
 \mathscr C_{<31}^{\rm str}
 \longrightarrow
 C^\bullet(G_{K_*,S_{31}},T_{\rm AT,*})
 \label{eq:v073-pre31-map}
\end{equation}
which preserve the two marked first tangents, the repeated cubic pointing,
and the distinguished cells through stage~30.
\end{definition}

Every such realization has a terminal detector evaluation
\begin{equation}
 \boxed{
 \operatorname{Ev}_{31}^{\det}(\rho_{<31})
 :=
 \rho_{<31,*}(\mathfrak R_{31}^{[31]})
 \in
 H^2(G_{K_*,S_{31}},\mu_{31})E.}
 \label{eq:v073-det-evaluation}
\end{equation}

\begin{theorem}[Exact strict provenance formulation]
\label{thm:v073-exact-provenance-gate}
A strict marked cyclotomic-coefficient-typed realization of the pointed detector exists if and
only if
\begin{equation}
 \boxed{
 \kappa_5^{\rm root}
 \in
 \operatorname{im}\!\left(
 \operatorname{Ev}_{31}^{\det}:
 \mathcal R_{<31}^{\rm root}(K_*)
 \to H^2(G_{K_*,S_{31}},\mu_{31})E
 \right).}
 \label{eq:v073-provenance-fiber}
\end{equation}
Equivalently, the remaining problem is the nonemptiness of the single fiber
\[
 \boxed{
 \bigl(\operatorname{Ev}_{31}^{\det}\bigr)^{-1}
 (\kappa_5^{\rm root}).}
\]
The later question of extending a chosen \(\rho_{<31}\) across an
\(h_{31}^{\rm rec}\)-attachment is a separate recursive-cell lifting
problem and is not required for the first classical \(\Sha^2\) realization.
\end{theorem}

\begin{proof}
The detector is a cohomology class of the controller through stage~30 and
is functorial under strict dg maps.  Thus a strict realization carrying the
normalized source obstruction detector to the pointed arithmetic class is exactly a
point of the displayed fiber.  Conversely any point of that fiber is, by
definition, a strict cyclotomic-coefficient-typed realization of all lower marked data whose
induced terminal detector is the desired class.

The source relative cell attachment has
\(dh_N=-\Omega_N\).  Hence adjoining \(h_{31}^{\rm rec}\) asks in addition
for a nullhomotopy of the image of its defining recursive relation.  That
condition concerns extension of the map after the detector has already been
formed, so it is logically distinct from
\eqref{eq:v073-provenance-fiber}.
\end{proof}

\begin{proposition}[Rank-one top-stage necessary gate]
\label{prop:v073-top-bracket-gate}
Let \(\rho_{<31}\) lie in the fiber
\eqref{eq:v073-provenance-fiber}.  Write
\[
 \Lambda_{30}^{\rm ar}
 :=\rho_{<31}(\Lambda_{30}^{[31]}),
 \qquad
 H_{30}^{\rm ar}:=\rho_{<31}(h_{30}^{\rm rec}).
\]
Then the stage-\(30\) contribution to the arithmetic terminal detector is
\begin{equation}
 \boxed{
 [\Lambda_{30}^{\rm ar},H_{30}^{\rm ar}]_{\rm gr}.}
 \label{eq:v073-arithmetic-top-bracket}
\end{equation}
After the exact triangular descent of lower cell stages has been fixed, the
final detector-faithfulness test is therefore a single cyclotomic-coefficient-valued
degree-two bracket whose total class must equal
\(\kappa_5^{\rm root}\).
\end{proposition}

\begin{proof}
Apply \(\rho_{<31}\) to the source identity established above
\[
 \pi_{30}\mathfrak R_{31}^{[31]}
 =
 [\Lambda_{30}^{[31]},\bar h_{30}^{\rm rec}]_{\rm gr}.
\]
Strict dg functoriality preserves multiplication, differential and the
marked root coefficient type.  The lower stages are handled by the exact
cell-stage descent already established by the source recursion above.
\end{proof}

\begin{remark}[What the coefficient-typing step changes]
\label{rem:v073-change}
The existence of a nonzero classical \(31\)-primary \(\Sha^2\) is not the
same as realization of the source obstruction detector.  The preliminary search found the
first small \emph{untyped} positive Artin--Tate point at \(D=-68\), and
The preliminary pointing step selected its opposite-root line.  The coefficient-typing step imposes the
actual root coefficient typing and finds the first compatible classical
point at \(D=-15391\).  The remaining provenance problem has now been
reduced to the explicit fiber
\eqref{eq:v073-provenance-fiber}, with the strongest associated-graded gate
given by the single bracket
\eqref{eq:v073-arithmetic-top-bracket}.
\end{remark}



\section{Kummer--Bockstein--Massey closure of the strict provenance fibre}
\label{sec:v074-kummer-bockstein-massey}

the coefficient-typing step leaves one precise problem:
\[
 \bigl(\operatorname{Ev}_{31}^{\det}\bigr)^{-1}
 (\kappa_5^{\rm root})
 \stackrel{?}{\ne}\varnothing.
\]
We now prove nonemptiness.  The proof has three ingredients which must be
kept distinct:
\begin{enumerate}[label=\textup{(\roman*)}]
\item an explicit Kummer representative of the pointed root
      Shafarevich class;
\item the top augmentation Bockstein for the cyclic unramified
      \(C_{31}\)-quotient;
\item the proper-Massey/Maurer--Cartan realization of that augmentation
      Bockstein.
\end{enumerate}

\subsection{An explicit Kummer representative of the cyclotomic coefficient class}

Retain
\[
 K_*=\mathbf Q(\sqrt{-15391}),
 \qquad
 \omega_*=\frac{1+\sqrt{-15391}}2,
\]
\[
 \mathfrak P_+=(31,\omega_*-14),
 \qquad
 \mathfrak P_-=(31,\omega_*-18),
 \qquad
 \mathfrak Q_5=(5,\omega_*-2).
\]
Put
\begin{equation}
 \boxed{
 \gamma_*
 :=
 71079143659-6025997842\,\omega_*.
 }
 \label{eq:v074-gamma}
\end{equation}

\begin{proposition}[Exact ideal relation for the pointing prime]
\label{prop:v074-explicit-ideal-relation}
One has
\begin{equation}
 \boxed{
 N_{K_*/\mathbf Q}(\gamma_*)
 =
 144354999065399169921875
 =
 31\cdot5^{31},}
 \label{eq:v074-gamma-norm}
\end{equation}
and
\begin{equation}
 \boxed{
 (\gamma_*)
 =
 \mathfrak Q_5^{31}\mathfrak P_-.
 }
 \label{eq:v074-gamma-ideal}
\end{equation}
Consequently, for
\begin{equation}
 \boxed{\beta_*:=\gamma_*/31,}
 \label{eq:v074-beta}
\end{equation}
one has
\begin{equation}
 \boxed{
 (\beta_*)
 =
 \mathfrak Q_5^{31}\mathfrak P_+^{-1}.}
 \label{eq:v074-beta-ideal}
\end{equation}
\end{proposition}

\begin{proof}
The norm form is
\[
 N(a+b\omega_*)=a^2+ab+3848b^2,
\]
and direct substitution gives \eqref{eq:v074-gamma-norm}.
Modulo \(5\), the polynomial of \(\omega_*\) has roots \(2\) and \(4\).
The element \(\gamma_*\) vanishes at \(\omega_*=2\) and is congruent to
\(1\) at \(\omega_*=4\).  Thus only \(\mathfrak Q_5\), not its conjugate,
divides \(\gamma_*\).

Modulo \(31\), the two roots are \(14\) and \(18\).  Substitution gives
\[
 \gamma_*(14)\equiv20\pmod{31},
 \qquad
 \gamma_*(18)\equiv0\pmod{31},
\]
so only \(\mathfrak P_-\) divides \(\gamma_*\).  The norm
\eqref{eq:v074-gamma-norm} leaves no remaining prime factors and forces the
exponents \(31\) and \(1\), proving \eqref{eq:v074-gamma-ideal}.
Finally
\[
 (31)=\mathfrak P_+\mathfrak P_-
\]
gives \eqref{eq:v074-beta-ideal}.
\end{proof}

Let
\[
 X_*:=\operatorname{Spec}\mathcal O_{K_*,S_{31}}.
\]
Both primes over \(31\) are invertible on \(X_*\), so
\eqref{eq:v074-beta-ideal} becomes
\begin{equation}
 \boxed{
 \operatorname{div}_{X_*}(\beta_*)
 =
 31\,\mathfrak Q_5.}
 \label{eq:v074-beta-divisor-X}
\end{equation}
Thus multiplication by \(\beta_*^{-1}\) gives a chosen trivialization
\[
 \mathfrak Q_5^{\otimes31}\xrightarrow{\sim}\mathcal O_{X_*}.
\]

\begin{definition}[Explicit Kummer lift]
\label{def:v074-kummer-lift}
Let
\begin{equation}
 \boxed{
 \lambda_5^{\rm Kum}
 \in H^1(X_*,\mu_{31})
 \simeq H^1(G_{K_*,S_{31}},\mu_{31})}
 \label{eq:v074-kummer-lift}
\end{equation}
be the Kummer torsor represented by
\[
 \left(
 \mathfrak Q_5,\;
 \mathfrak Q_5^{\otimes31}
 \xrightarrow{\ \beta_*^{-1}\ }
 \mathcal O_{X_*}
 \right).
\]
Equivalently it is represented by the Kummer element \(\beta_*\in K_*^\times\).
\end{definition}

By the coefficient-typing step,
\[
 \operatorname{Pic}(X_*)
 =
 \operatorname{Cl}(K_*)/
 \langle[\mathfrak P_+]\rangle
 \simeq C_{31},
\]
generated by \([\mathfrak Q_5]\).

\begin{definition}[Kummer Chern generator]
\label{def:v074-kummer-chern}
Write
\begin{equation}
 \boxed{
 \kappa_{\rm Kum}
 :=
 c_1^{(31)}(\mathfrak Q_5)
 \in H^2(X_*,\mu_{31})}
 \label{eq:v074-kummer-chern}
\end{equation}
for the image of \([\mathfrak Q_5]\) under
\[
 \operatorname{Pic}(X_*)/31
 \longrightarrow H^2(X_*,\mu_{31}).
\]
\end{definition}

\begin{proposition}[The Kummer Chern class is the pointed root
Shafarevich generator]
\label{prop:v074-kummer-equals-kappa}
One has
\begin{equation}
 \boxed{
 \kappa_{\rm Kum}
 =
 \kappa_5^{\rm root}.}
 \label{eq:v074-kummer-equals-kappa}
\end{equation}
In particular
\[
 \kappa_{\rm Kum}
 \in
 \Sha^2_{S_{31}}(K_*,\mu_{31})
\]
and is nonzero.
\end{proposition}

\begin{proof}
The restriction of a line bundle to \(\operatorname{Spec}K_{*,w}\) is
trivial, so its Kummer Chern class is zero at both \(w\in S_{31}\).
Since \([\mathfrak Q_5]\) generates
\(\operatorname{Pic}(X_*)\simeq C_{31}\), its Chern class is nonzero and
spans the one-dimensional root Shafarevich group of
\cref{thm:v073-root-sha-one}.

Poitou--Tate duality is compatible with global Artin reciprocity: the
pairing of \(c_1^{(31)}(\mathfrak a)\) with an unramified
\(\mathbf F_{31}\)-character \(c\) is \(c([\mathfrak a])\).  Therefore
\[
 \left\langle
 c_1^{(31)}(\mathfrak Q_5),
 z_5^{\rm root}
 \right\rangle_{\rm PT}
 =
 z_5^{\rm root}(\operatorname{Frob}_{\mathfrak Q_5})
 =1.
\]
This is the defining normalization
\eqref{eq:v073-kappa-root} of \(\kappa_5^{\rm root}\).
\end{proof}

\subsection{The coefficient Bockstein and the augmentation Bockstein are
different}

\begin{definition}[Coefficient Bockstein]
\label{def:v074-coeff-bockstein}
The short exact sequence
\[
 1\longrightarrow\mu_{31}
 \longrightarrow\mu_{31^2}
 \xrightarrow{(\cdot)^{31}}\mu_{31}
 \longrightarrow1
\]
defines
\[
 \delta_\mu:
 H^1(X_*,\mu_{31})
 \longrightarrow
 H^2(X_*,\mu_{31}).
\]
\end{definition}

\begin{proposition}[Coefficient Bockstein of the explicit Kummer torsor]
\label{prop:v074-coeff-bockstein}
With the preceding orientation,
\begin{equation}
 \boxed{
 \delta_\mu(\lambda_5^{\rm Kum})
 =
 \kappa_5^{\rm root}.}
 \label{eq:v074-coeff-bockstein}
\end{equation}
\end{proposition}

\begin{proof}
Choose an étale cover on which \(\mathfrak Q_5\) is trivial.
The chosen \(31\)-st-power trivialization defines a \(\mu_{31}\)-torsor.
Lifting its transition functions to \(\mu_{31^2}\), the failure of the
lifts to satisfy the cocycle identity on triple overlaps is the coefficient
Bockstein.

The same lifted transition functions compute the Kummer boundary
\(c_1^{(31)}(\mathfrak Q_5)\), since they are \(31\)-st roots of the
transition functions of the line bundle.  Thus
\[
 \delta_\mu(\lambda_5^{\rm Kum})
 =
 c_1^{(31)}(\mathfrak Q_5),
\]
and \cref{prop:v074-kummer-equals-kappa} finishes the proof.
\end{proof}

\begin{warning}[Bockstein firewall]
\label{warn:v074-two-bocksteins}
The operator \(\delta_\mu\) changes coefficient depth
\(\mu_{31}\leadsto\mu_{31^2}\).  It is \emph{not} the generalized
augmentation Bockstein used below.  The latter is attached to the
\(C_{31}\)-group-ring filtration and is the operator that generates the
proper \(31\)-fold defining system.
\end{warning}

\subsection{The cyclic \texorpdfstring{\(S\)}{S}-Hilbert quotient}

Let
\[
 x:=z_5^{\rm root}
 \in
 \Sha^1_{S_{31}}(K_*,\mathbf F_{31}),
 \qquad
 x(\operatorname{Frob}_{\mathfrak Q_5})=1.
\]
Let \(H_*/K_*\) be the cyclic degree-\(31\) extension cut out by \(x\).
It is the \(31\)-primary Hilbert \(S_{31}\)-class-field quotient.  Put
\[
 G:=G_{K_*,S_{31}},
 \qquad
 N:=G_{H_*,S_{H_*}},
 \qquad
 C:=G/N\simeq C_{31}.
\]

\begin{lemma}[\(S\)-capitulation of the pointing class]
\label{lem:v074-S-capitulation}
The pullback of \(\mathfrak Q_5\) to
\(\operatorname{Spec}\mathcal O_{H_*,S_{H_*}}\) is principal.  Consequently
\begin{equation}
 \boxed{
 \operatorname{res}_{N}
 (\kappa_5^{\rm root})
 =0
 \quad\text{in }H^2(N,\mu_{31}).}
 \label{eq:v074-kappa-restriction-zero}
\end{equation}
\end{lemma}

\begin{proof}
The explicit relation
\eqref{eq:v074-beta-ideal} gives in the ordinary class group
\[
 [\mathfrak Q_5]^{31}=[\mathfrak P_+].
\]
By \cref{thm:v073-field-certificate},
\(\operatorname{ord}[\mathfrak P_+]=3\), while
\cref{prop:v073-Q5-generator} says that the image of
\([\mathfrak Q_5]\) modulo \(\langle[\mathfrak P_+]\rangle\) has order
\(31\).  Since \(h_{K_*}=93\), it follows that
\[
 \boxed{
 \operatorname{Cl}(K_*)\simeq C_{93},
 \qquad
 \operatorname{ord}[\mathfrak Q_5]=93.}
\]
In particular \([\mathfrak Q_5]\) is an ordinary class-group generator.

The field \(H_*/K_*\) is the cyclic unramified degree-\(31\) extension
whose Artin kernel is
\(\langle[\mathfrak P_+]\rangle\).
Hilbert's theorem~94 says that the ordinary capitulation kernel in a cyclic
unramified degree-\(31\) extension has order divisible by \(31\).
The cyclic group \(C_{93}\) has a unique subgroup of order \(31\), namely
\[
 \langle[\mathfrak Q_5]^3\rangle,
\]
so this subgroup capitulates in \(H_*\).

Now
\[
 1=21\cdot3-2\cdot31.
\]
Hence
\[
 [\mathfrak Q_5]
 =
 [\mathfrak Q_5^3]^{21}
 [\mathfrak Q_5^{31}]^{-2}.
\]
The first factor becomes principal in \(H_*\) by Hilbert~94, while the
second factor is a power of \([\mathfrak P_+]\), which is trivial in the
\(S_{H_*}\)-class group because the primes over \(31\) are inverted.
Therefore the pullback of \([\mathfrak Q_5]\) is zero in
\(\operatorname{Pic}(\operatorname{Spec}\mathcal O_{H_*,S_{H_*}})\).

Kummer Chern classes commute with pullback, so
\[
 \operatorname{res}_{N}
 c_1^{(31)}(\mathfrak Q_5)=0.
\]
Use \cref{prop:v074-kummer-equals-kappa}.
\end{proof}

\subsection{The top augmentation Bockstein}

Choose
\[
 \sigma\in C,\qquad x(\sigma)=1,
\]
and put
\[
 \Omega:=\mathbf F_{31}[C],
 \qquad
 X:=\sigma-1,
 \qquad
 I:=(X).
\]
Then
\[
 \Omega\simeq\mathbf F_{31}[X]/(X^{31}),
 \qquad I^{31}=0,
\]
and
\begin{equation}
 \boxed{
 X^{30}
 =
 1+\sigma+\cdots+\sigma^{30}.}
 \label{eq:v074-top-norm-element}
\end{equation}

Put \(T:=\mu_{31}\).  The sequence
\[
 0\longrightarrow T\otimes I^{30}
 \longrightarrow T\otimes\Omega
 \longrightarrow T\otimes\Omega/I^{30}
 \longrightarrow0
\]
defines
\begin{equation}
 \boxed{
 \Psi_x^{(30)}:
 H^1(G,T\otimes\Omega/I^{30})
 \longrightarrow
 H^2(G,T)\otimes I^{30}.}
 \label{eq:v074-top-augmentation-bockstein}
\end{equation}
Orient \(I^{30}\) by \(X^{30}\).

\begin{theorem}[Top Bockstein image equals the restriction kernel]
\label{thm:v074-top-bockstein-image}
Under
\[
 T\otimes I^{30}\simeq T,
 \qquad
 t\otimes X^{30}\longleftrightarrow t,
\]
one has
\begin{equation}
 \boxed{
 \operatorname{im}\Psi_x^{(30)}
 =
 \ker\!\left(
 H^2(G,T)\xrightarrow{\operatorname{res}}
 H^2(N,T)
 \right)\otimes\mathbf F_{31}X^{30}.}
 \label{eq:v074-top-bockstein-kernel}
\end{equation}
\end{theorem}

\begin{proof}
The long exact sequence identifies the image with the kernel of
\[
 H^2(G,T\otimes I^{30})
 \longrightarrow
 H^2(G,T\otimes\Omega).
\]
By \eqref{eq:v074-top-norm-element}, the inclusion is
\[
 t\longmapsto
 t\otimes\sum_{c\in C}c.
\]
The \(G\)-module \(T\otimes\Omega\) is
\(\operatorname{Ind}_{N}^{G}(T|_N)\).
Under Shapiro's lemma, the induced cohomology map is ordinary restriction
\[
 H^2(G,T)\longrightarrow H^2(N,T).
\]
This proves the claim.
\end{proof}

\begin{corollary}[The pointed cyclotomic coefficient class is a top generalized Bockstein
value]
\label{cor:v074-kappa-is-top-bockstein}
There exists
\begin{equation}
 \boxed{
 [f_*]
 \in
 H^1(G,T\otimes\Omega/I^{30})}
 \label{eq:v074-fstar}
\end{equation}
such that
\begin{equation}
 \boxed{
 \Psi_x^{(30)}([f_*])
 =
 \kappa_5^{\rm root}\otimes X^{30}.}
 \label{eq:v074-fstar-bockstein}
\end{equation}
\end{corollary}

\begin{proof}
Combine
\cref{lem:v074-S-capitulation,thm:v074-top-bockstein-image}.
\end{proof}

\subsection{A nonzero marked cyclotomic-coefficient tangent can be retained}

Recall
\[
 \alpha_*=-120+\sqrt{-15391},
 \qquad
 (\alpha_*)=\mathfrak P_+^3.
\]
It gives a nonzero \(S_{31}\)-unit Kummer class
\begin{equation}
 \boxed{
 \lambda_\alpha:=[\alpha_*]
 \in H^1(G,\mu_{31}).}
 \label{eq:v074-lambda-alpha}
\end{equation}

\begin{lemma}[A nonzero corestriction direction]
\label{lem:v074-corestriction-direction}
The class \(\lambda_\alpha\) belongs to
\[
 \operatorname{cor}_{N}^{G}H^1(N,\mu_{31})
\]
and is nonzero.
\end{lemma}

\begin{proof}
The extension \(H_*/K_*\) is cyclic.  At the two places in \(S_{31}\) it
is split, so the local norm is surjective.  At every finite place outside
\(S_{31}\), the extension is unramified and \(\alpha_*\) is a unit; the norm
on the unit group of an unramified local extension is surjective.  There is
no archimedean obstruction.  The Hasse norm theorem therefore gives
\[
 \alpha_*=N_{H_*/K_*}(a)
\]
for some \(a\in H_*^\times\).  Kummer corestriction is induced by norm.

The class is nonzero because
\[
 v_{\mathfrak P_+}(\alpha_*)=3
\]
is not divisible by \(31\).
\end{proof}

Under Shapiro, choose
\[
 [g_\alpha]\in H^1(G,T\otimes\Omega)
\]
whose augmentation is \(\lambda_\alpha\).

\begin{proposition}[Pointed top-Bockstein lift]
\label{prop:v074-pointed-f}
One may choose \([f_*]\) in
\cref{cor:v074-kappa-is-top-bockstein} so that
\begin{equation}
 \boxed{
 \lambda_*:=
 \operatorname{aug}([f_*])
 \in H^1(G,\mu_{31})}
 \label{eq:v074-lambda-star}
\end{equation}
is nonzero.
\end{proposition}

\begin{proof}
Let \(f_0\) satisfy \eqref{eq:v074-fstar-bockstein}, and let
\(\bar g_\alpha\) be the image of \(g_\alpha\) modulo \(I^{30}\).
Exactness gives
\[
 \Psi_x^{(30)}(\bar g_\alpha)=0.
\]
Thus
\[
 f_c=f_0+c\bar g_\alpha
\]
has the same top Bockstein value for all \(c\in\mathbf F_{31}\), while
\[
 \operatorname{aug}(f_c)
 =
 \operatorname{aug}(f_0)+c\lambda_\alpha.
\]
At most one value of \(c\) makes this zero.
\end{proof}

\subsection{Proper \texorpdfstring{\(31\)}{31}-fold defining systems}

Choose a cocycle representative
\[
 f_*=\sum_{k=0}^{29}\lambda_kX^k
 \in Z^1(G,T\otimes\Omega/I^{30}),
 \qquad
 [\lambda_0]=\lambda_*\ne0.
\]
Let
\[
 b_j:=\binom{x}{j}
\]
denote the standard binomial cochains in the upper-left procyclic block.

\begin{theorem}[Proper-Massey value of the pointed top Bockstein]
\label{thm:v074-LLSWW-specialization}
The cocycle \(f_*\) determines a proper defining system
\(\rho_{f_*}\) for
\[
 \left(
 \underbrace{x,\ldots,x}_{30},
 \lambda_*
 \right),
\]
and
\begin{equation}
 \boxed{
 \left(
 \underbrace{x,\ldots,x}_{30},
 \lambda_*
 \right)_{\rho_{f_*}}
 =
 \kappa_5^{\rm root}.}
 \label{eq:v074-proper-31fold}
\end{equation}
\end{theorem}

\begin{proof}
This is Theorem~4.3.1 of Lam--Liu--Sharifi--Wake--Wang \cite{LLSWW}
(quoted as \ref{X:llsww} in \cref{sec:external-inputs}), specialized to
\[
 R=\mathbf F_{31},
 \quad C=C_{31},
 \quad n=30,
 \quad T=\mu_{31}.
\]
Their formula is
\[
 \Psi_x^{(30)}([f_*])
 =
 \left(
 x^{(30)},\lambda_*
 \right)_{\rho_{f_*}}
 \otimes X^{30}.
\]
Use \eqref{eq:v074-fstar-bockstein}.
\end{proof}

The last-column equations can be written, after the fixed controller sign
normalization, as
\begin{equation}
 \boxed{
 d\lambda_{s-1}
 =
 -\sum_{j=1}^{s-1}
 b_j\smile\lambda_{s-1-j},
 \qquad
 3\le s\le30.}
 \label{eq:v074-last-column-recursion}
\end{equation}
The terminal unfilled equation is
\begin{equation}
 \boxed{
 x\smile\lambda_{29}
 +b_2\smile\lambda_{28}
 +\cdots
 +b_{30}\smile\lambda_0
 \sim
 \kappa_5^{\rm root}.}
 \label{eq:v074-terminal-massey-cochain}
\end{equation}

\subsection{Strict realization through stage \texorpdfstring{\(30\)}{30}}

\begin{proposition}[Massey--Maurer--Cartan realization of the repeated spine]
\label{prop:v074-massey-mc-realization}
The proper defining system \(\rho_{f_*}\) determines a strict cyclotomic-coefficient-typed
realization of the selected repeated spine through stage~30.  Under the
fixed normalization,
\[
 \overline P_j\longmapsto b_jD,
 \qquad
 \overline R_s^{(v)}\longmapsto\lambda_{s-1}E,
\]
and the recursive cell \(h_s^{\rm rec}\) is sent to the corresponding
last-column filler for every \(3\le s\le30\).

The unfilled height-\(31\) repeated curvature is
\[
 \left(
 x^{(30)},\lambda_*
 \right)_{\rho_{f_*}}E.
\]
\end{proposition}

\begin{proof}
A pointed defining system is a strictly upper-triangular Maurer--Cartan
matrix with every nonterminal entry of
\[
 dU+U^2
\]
equal to zero.  Its upper-left block is the binomial block \(b_j\), and its
last column is \(\lambda_0,\ldots,\lambda_{29}\).
Equation \eqref{eq:v074-last-column-recursion} is exactly the recursive
differential needed to map the repeated semi-free cells.

The terminal \((1,32)\)-entry is the proper \(31\)-fold Massey cocycle.
The source Massey--Maurer--Cartan comparison identifies, after the normalized
PD word
\[
 \gamma_{30}(e_{31})e_v,
\]
this terminal curvature with
\(\mathfrak o^{\rm rec}_{30,1}\).
The semi-free universal property maps every \(h_s^{\rm rec}\) to the
corresponding filler.
\end{proof}

\begin{lemma}[Extension from the repeated spine to the selected source
controller]
\label{lem:v074-spine-to-full}
The preceding map extends to a strict cyclotomic-coefficient-typed realization
\begin{equation}
 \boxed{
 \rho_{<31}^{\rm BM}:
 \mathscr C_{<31}^{\rm str}
 \longrightarrow
 C^\bullet(G,T_{\rm AT,*}).}
 \label{eq:v074-rho-BM}
\end{equation}
\end{lemma}

\begin{proof}
Delete every old direction except the selected \(v\).  On the surviving
two-direction plane use
\[
 D^2=D,\qquad DE=E,\qquad ED=E^2=0.
\]
Every word with at least two root inputs therefore vanishes.  The pure words
are exactly the binomial block \(b_jD\), and the one-root words are exactly
the last-column system \(\lambda_kE\); their Maurer--Cartan equations hold
by construction.

The PD-summand semi-free refinement adjoins separate cells for complementary
direction coefficients.  After direction deletion and the square-zero root
specialization, the differentials of those complementary cells map to zero,
so those cells may be sent to zero.  The semi-free universal property gives
the claimed extension.
\end{proof}

\begin{theorem}[Strict provenance fibre is nonempty]
\label{thm:v074-provenance-fiber-closed}
The map \(\rho_{<31}^{\rm BM}\) satisfies
\begin{equation}
 \boxed{
 \operatorname{Ev}_{31}^{\det}
 \bigl(\rho_{<31}^{\rm BM}\bigr)
 =
 \kappa_5^{\rm root}.}
 \label{eq:v074-Ev-equals-kappa}
\end{equation}
Hence
\begin{equation}
 \boxed{
 \left(
 \operatorname{Ev}_{31}^{\det}
 \right)^{-1}
 \bigl(\kappa_5^{\rm root}\bigr)
 \ne\varnothing.}
 \label{eq:v074-fiber-nonempty}
\end{equation}
\end{theorem}

\begin{proof}
After the chosen direction specialization, only the one-old repeated
direction
\[
 \gamma_{30}(e_{31})e_v
\]
survives in the height-\(31\) detector.  The source PD recursion assigns
scalar \(1\) to this \(a=30\) summand.  Thus the complete detector evaluates to
the repeated terminal curvature.  By
\cref{prop:v074-massey-mc-realization,thm:v074-LLSWW-specialization}, this is
\[
 \left(
 x^{(30)},\lambda_*
 \right)_{\rho_{f_*}}E
 =
 \kappa_5^{\rm root}E.
\]
Suppress the fixed matrix basis vector \(E\).
\end{proof}

\subsection{The stage-\texorpdfstring{\(30\)}{30} bracket}

Let
\[
 H_{30}^{\rm ar}:=
 \rho_{<31}^{\rm BM}(h_{30}^{\rm rec}).
\]

\begin{corollary}[Realized top bracket]
\label{cor:v074-realized-top-bracket}
The stage-\(30\) associated-graded image is
\begin{equation}
 \boxed{
 \operatorname{gr}_{30}^{\rm cell}
 \rho_{<31}^{\rm BM}
 \bigl(\mathfrak R_{31}^{[31]}\bigr)
 =
 [xD,\lambda_{29}E]_{\rm gr}
 =
 (x\smile\lambda_{29})E.}
 \label{eq:v074-realized-top-bracket}
\end{equation}
The terms with \(b_j\), \(j\ge2\), in
\eqref{eq:v074-terminal-massey-cochain} are the lower cell-stage
corrections.
\end{corollary}

\begin{proof}
Apply \cref{P1C-thm:universal-higher-obstruction-recursion} together with
\cref{prop:v073-top-bracket-gate}.
In the \(D/E\)-plane, \(DE=E\) and \(ED=0\), so only the displayed ordered
cyclotomic-coefficient term survives in the highest cell stage.
\end{proof}

\subsection{The classical \texorpdfstring{\(\Sha^2\)}{Sha2} as the local--global cell-extension
obstruction}

Let
\[
 \mathscr C_{\le31}^{\rm rec}
 :=
 \mathscr C_{<31}^{\rm str}
 \langle h_{31}^{\rm rec}\rangle
\]
be the next formal repeated-cell attachment.

\begin{theorem}[Global nonextension at the first resonance]
\label{thm:v074-global-nonextension}
The map \(\rho_{<31}^{\rm BM}\) does not extend to a strict dg map
\[
 \mathscr C_{\le31}^{\rm rec}
 \longrightarrow
 C^\bullet(G,T_{\rm AT,*}).
\]
\end{theorem}

\begin{proof}
An extension would require a degree-one cochain
\(H_{31}^{\rm ar}\) with differential equal to minus the image of the
height-\(31\) repeated relation.  Its cohomology class would therefore
vanish.  But under the selected repeated direction this class is
\[
 \kappa_5^{\rm root}\ne0
\]
by \cref{thm:v074-provenance-fiber-closed}.
\end{proof}

\begin{theorem}[Local extension at both \(31\)-adic places]
\label{thm:v074-local-extension}
For each
\[
 w\in S_{31}=\{\mathfrak P_+,\mathfrak P_-\},
\]
the localized map
\[
 \rho_{<31,w}^{\rm BM}:
 \mathscr C_{<31}^{\rm str}
 \longrightarrow
 C^\bullet(G_{K_{*,w}},T_{\rm AT,*})
\]
does extend across \(h_{31}^{\rm rec}\).
\end{theorem}

\begin{proof}
The terminal localized obstruction is the restriction of
\(\kappa_5^{\rm root}\).  Since
\[
 \kappa_5^{\rm root}
 \in
 \Sha^2_{S_{31}}(K_*,\mu_{31}),
\]
this restriction vanishes at both places.  Choose a local degree-one
nullhomotopy and use the universal property of the semi-free attachment.
\end{proof}

\begin{corollary}[Shafarevich class as recursive local--global failure]
\label{cor:v074-sha-as-cell-failure}
The pointed class
\[
 \boxed{
 \kappa_5^{\rm root}
 \in
 \Sha^2_{S_{31}}(K_*,\mu_{31})}
\]
is precisely the obstruction to gluing the locally extendable
\(h_{31}^{\rm rec}\)-lifts into a global strict lift of
\(\rho_{<31}^{\rm BM}\).  Equivalently,
\begin{equation}
 \boxed{
 \Sha^2_{S_{31}}(K_*,\mu_{31})
 =
 \mathbf F_{31}\,
 [\,\text{global obstruction to the }h_{31}^{\rm rec}\text{ attachment}\,].}
 \label{eq:v074-sha-cell-obstruction}
\end{equation}
\end{corollary}

\begin{remark}[What the present closure argument closes]
\label{rem:v074-closed}
The first-resonance arithmetic realization now has the exact order
\[
 \boxed{
 \begin{gathered}
 \text{strict cyclotomic-coefficient-typed realization through stage }30\\
 \Downarrow\\
 \text{proper \(31\)-fold terminal curvature }
 =\kappa_5^{\rm root}\ne0\\
 \Downarrow\\
 \text{local extension across }h_{31}^{\rm rec}\text{ at both }31\text{-places}\\
 \text{but no global extension.}
 \end{gathered}}
\]
Thus the classical cyclotomic-coefficient \(\Sha^2\) is the actual local--global obstruction
class of the recursive Selmer--Cartan attachment, not merely an abstract
one-dimensional target.
\end{remark}



\section{Witness theorem and arithmetic conclusion}\label{sec:export-contract}

\begin{theorem}[The \(31\)-adic arithmetic witness]\label{thm:paper1-31adic-witness}
The branches \(K_0=\mathbf Q(\sqrt{-331})\) and \(K_*=\mathbf Q(\sqrt{-15391})\) satisfy the witness contract of \cref{def:paper1-witness-contract}.  More precisely:
\begin{enumerate}[label=\textup{(\roman*)}]
\item \(K_0\) supplies an exact-\(\lambda_{31}=2\) repeated cubic and a compatible family of pointed primitive reductions of exact orders \(31^r\), with canonical carry normalization at every finite \(r\);
\item \(K_*\) supplies a pointed nonzero class \(\kappa_5^{\rm root}\in\Sha^2_{S_{31}}(K_*,\mu_{31})\), where the superscript is retained as historical notation for the cyclotomic coefficient corner;
\item The displayed Kummer element represents the proper \(31\)-fold Massey value, equivalently the top augmentation Bockstein;
\item the associated recursive map exists through stage \(30\), extends at both local places above \(31\), and has no global extension across \(h_{31}^{\rm rec}\).
\end{enumerate}
Thus the terminal source obstruction has an explicit arithmetic local--global witness.
\end{theorem}

\section{Motivic closure of the \(31\)-adic witness}
\label{W31-sec:motivic-closure}

The arithmetic calculation has two coefficient profiles.  The repeated cubic
branch over \(K_0=\mathbf Q(\sqrt{-331})\) has a primitive carry-normalized
\(\mathbf Z_{31}\)-line and therefore compatible reductions of exact order
\(31^r\) for every finite \(r\).  The terminal proper-\(31\)-fold branch over
\(K_*=\mathbf Q(\sqrt{-15391})\) supplies only an exact-order-\(31\) arithmetic
line.  We realize these two profiles separately.

\subsection{Terminal unipotent-gerbe provenance}

Write \(\underline{\mathbf F}_{31}\) for the constant finite-etale group
scheme on \(X_*\) and regard \(\mu_{31}\) as an
\(\underline{\mathbf F}_{31}\)-module by exponentiation.  In additive module
notation define the cyclotomic generalized unipotent group scheme
\begin{equation}\label{W31-eq:twisted-unipotent-group}
 \boxed{
 \mathcal U_{32}^{\rm cyc}
 :=U_{31}(\underline{\mathbf F}_{31})
   \ltimes\mu_{31}^{31}.}
\end{equation}
It is the group of generalized upper-unitriangular \(32\)-by-\(32\) matrices:
the upper-left \(31\)-by-\(31\) block has entries in
\(\underline{\mathbf F}_{31}\), and the last column has entries in
\(\mu_{31}\).  The group law is
\[
 (A,v)(A',v')=(AA',v+Av').
\]
The top-right coordinate
\[
 \mathcal Z_{31}:=\mu_{31}E_{1,32}
\]
is central.  Put
\[
 \overline{\mathcal U}_{32}^{\rm cyc}
 :=\mathcal U_{32}^{\rm cyc}/\mathcal Z_{31}.
\]
Thus there is a central extension of finite-etale group schemes
\begin{equation}\label{W31-eq:twisted-unipotent-extension}
 \boxed{
 1\longrightarrow\mu_{31}
 \longrightarrow\mathcal U_{32}^{\rm cyc}
 \longrightarrow\overline{\mathcal U}_{32}^{\rm cyc}
 \longrightarrow1.}
\end{equation}
Let
\[
 \omega_{31}^{\rm univ}
 \in H^2_{\mathrm{\acute et}}
 \bigl(B\overline{\mathcal U}_{32}^{\rm cyc},\mu_{31}\bigr)
\]
be its extension class, oriented by the section with zero top-right coordinate
and the Maurer--Cartan sign convention of
\cref{prop:v074-massey-mc-realization}.

\begin{theorem}[Terminal proper-\(31\)-fold unipotent-gerbe provenance]
\label{W31-thm:terminal-unipotent-gerbe-provenance}
The proper defining system \(\rho_{f_*}\) of
\cref{thm:v074-LLSWW-specialization} assembles into a classifying morphism
\[
 \bar\rho_{f_*}:
 X_*\longrightarrow B\overline{\mathcal U}_{32}^{\rm cyc}
\]
whose first thirty adjacent coordinates are \(x\) and whose last adjacent
coordinate is \(\lambda_*\).  Its universal obstruction satisfies the pointed
identity
\begin{equation}\label{W31-eq:terminal-unipotent-pullback}
 \boxed{
 \bar\rho_{f_*}^{\,*}\omega_{31}^{\rm univ}
 =
 \bigl(x^{(30)},\lambda_*\bigr)_{\rho_{f_*}}
 =\kappa_5^{\rm root}
 =c_1^{(31)}(\mathfrak Q_5).}
\end{equation}
Consequently the homotopy pullback
\begin{equation}\label{W31-eq:terminal-obstruction-gerbe}
 \boxed{
 \mathcal G_{f_*}
 :=X_*\times_{B\overline{\mathcal U}_{32}^{\rm cyc}}
       B\mathcal U_{32}^{\rm cyc}}
\end{equation}
is an algebraic \(\mu_{31}\)-gerbe with class
\(c_1^{(31)}(\mathfrak Q_5)\).  The cartesian square
\[
\begin{CD}
 \mathcal G_{f_*}@>>>B\mathcal U_{32}^{\rm cyc}\\
 @VVV @VVV\\
 X_*@>{\bar\rho_{f_*}}>>
 B\overline{\mathcal U}_{32}^{\rm cyc}
\end{CD}
\]
is an independent algebraic-stack correspondence provenance for the terminal
proper-\(31\)-fold operation law.  Under the pointed motivic Moore--Reedy
specialization it carries the class in
\eqref{W31-eq:terminal-unipotent-pullback} to the distinguished terminal Moore
line of \cref{W31-thm:common-source-motivic-span}.
\end{theorem}

\begin{proof}
The upper-left binomial cochains \(b_j=\binom{x}{j}\) and the last-column
cochains \(\lambda_0,\ldots,\lambda_{29}\) satisfy the nonterminal
Maurer--Cartan equations by
\cref{thm:v074-LLSWW-specialization,prop:v074-massey-mc-realization}.  Matrix
multiplication in \(\mathcal U_{32}^{\rm cyc}\) is precisely those equations.
After deleting the \((1,32)\)-coordinate, the defining matrix is therefore a
continuous nonabelian cocycle with values in
\(\overline{\mathcal U}_{32}^{\rm cyc}\).  Since the group scheme is finite
etale, that cocycle is equivalently an algebraic torsor and hence the displayed
classifying morphism.

Dwyer's defining-system theorem identifies the terminal curvature cocycle with
the pullback of the class of the universal upper-unitriangular central
extension \cite{DwyerMassey}.  The generalized-matrix/module version used here
is exactly the proper defining-system formalism of Lam--Liu--Sharifi--Wake--Wang
\cite{LLSWW}.  It follows that
\[
 \bar\rho_{f_*}^{\,*}\omega_{31}^{\rm univ}
 =\bigl(x^{(30)},\lambda_*\bigr)_{\rho_{f_*}}.
\]
The second equality in
\eqref{W31-eq:terminal-unipotent-pullback} is
\cref{thm:v074-LLSWW-specialization}; the last is
\cref{prop:v074-kummer-equals-kappa}.  The normalization is literal, rather
than only up to a unit, because both sides pair to one with
\(z_5^{\rm root}\).

Pulling back the central extension gives
\eqref{W31-eq:terminal-obstruction-gerbe}; the standard classification of
banded gerbes identifies its class with the pullback of
\(\omega_{31}^{\rm univ}\).  All four stacks in the square are tame finite
quotient stacks over \(X_*\), so the square is an algebraic-stack
correspondence and lies in the motivic six-operation category used in
\cref{P2M-prop:appendix-root-moore-complex}, applied with \(U_d=X_*\) and
\(N=31\).  Finally
\(c_1^{(31)}(\mathfrak Q_5)\) is the mod-\(31\) cycle class of the explicit
divisor \(\mathfrak Q_5\), and the pointed carrier comparison sends that
generator to the terminal motivic Moore generator.  This proves the motivic
compatibility.
\end{proof}

\begin{remark}[Gerbe law versus its specialized cycle]
\label{W31-rem:gerbe-law-versus-cycle}
The universal operation law is the central extension class
\(\omega_{31}^{\rm univ}\) and is represented geometrically by the obstruction
gerbe, not asserted to be a single ordinary Chow correspondence.  Its
specialization on \(X_*\) is the ordinary algebraic divisor cycle
\([\mathfrak Q_5]\in\CH^1(X_*)/31\).  This distinction is preserved by the
motivic realization.
\end{remark}

\subsection{The prime-power repeated-confluence branch}

Let \(S_0\) be any finite ordered marked support carrying the repeated cubic
block of the \(K_0\)-branch, and let \(I_0\subseteq S_0\) be that exact-support
block.  For \(r\ge1\), put
\begin{equation}\label{W31-eq:prime-power-source-line}
 \boxed{
 L^{\rm conf}_{0,r}
 :=
 \frac{\mathbf Z_{31}\kappa_{\rm Qi}^{\rm can}}
      {31^r\mathbf Z_{31}\kappa_{\rm Qi}^{\rm can}}
 \simeq(\mathbf Z/31^r)
       \langle\kappa_{\rm Qi}^{\rm can}\bmod31^r\rangle .}
\end{equation}
The transition \(L^{\rm conf}_{0,r+1}\twoheadrightarrow
L^{\rm conf}_{0,r}\) is coefficient reduction.

\begin{theorem}[The \(31^r\) repeated-confluence motivic witness]
\label{W31-thm:prime-power-repeated-confluence}
For every finite \(r\ge1\), the line
\(L^{\rm conf}_{0,r}\) is pointed primitive of exact order \(31^r\), its
order-\(31\) reduction is the normalized repeated cubic source line, and its
bifiltered relation is
\[
 d\eta_r
 =
 \alpha_{31^r}e_ve_{31}
 +e_v\gamma_2(e_{31})\kappa_{\rm Qi}^{\rm can},
 \qquad
 \alpha_{31^r}
 =
 \begin{cases}
 0,&r=1,\\
 2^{-1}31[5]\pmod{31^r},&r\ge2.
 \end{cases}
\]
Here \(v\) is the retained old controller direction fixed in the explicit
branch.  Consequently \cref{P2M-thm:prime-power-all-support-comparison} gives a genuine
coefficient-faithful all-support realization
\begin{equation}\label{W31-eq:prime-power-motivic-realization}
 \boxed{
 \rho_{S_0,r}^{\Mot,\MR}:
 \mathscr U^{\mathrm{rec}}_{S_0,31^r}
 \longrightarrow
 \mathscr M_{S_0,31^r}^{\Mot,\MR}}
\end{equation}
whose exact-support \(I_0\)-line is generated by a pair
\[
 dc_{0,r}^{\Mot}=31^r b_{0,r}^{\Mot},
 \qquad
 \operatorname{ord}[b_{0,r}^{\Mot}]=31^r.
\]
For \(1\le s\le r\), the source reductions and cyclotomic root reductions form
commutative squares
\[
\begin{CD}
\mathscr U^{\mathrm{rec}}_{S_0,31^r}
 @>{\rho_{S_0,r}^{\Mot,\MR}}>>
\mathscr M_{S_0,31^r}^{\Mot,\MR}\\
@V{\operatorname{red}_{r,s}}VV
@VV{\operatorname{Red}_{r,s}^{\Mot}}V\\
\mathscr U^{\mathrm{rec}}_{S_0,31^s}
 @>{\rho_{S_0,s}^{\Mot,\MR}}>>
\mathscr M_{S_0,31^s}^{\Mot,\MR}.
\end{CD}
\]
Thus the \(K_0\)-branch gives a compatible finite-level
\(31^r\)-motivic Moore--Reedy family retaining the repeated-\(31\) confluence
provenance.
\end{theorem}

\begin{proof}
By \cref{thm:v064-explicit-carry-pointing},
\[
 \kappa_{\rm Qi}^{\rm can}
 =\frac{a_0}{[5]}\widetilde\kappa,
 \qquad
 \frac{a_0}{[5]}\in\mathbf Z_{31}^{\times},
\]
and \(\widetilde\kappa\) is primitive.  Hence
\(\kappa_{\rm Qi}^{\rm can}\) is a primitive generator of a free
\(\mathbf Z_{31}\)-line, proving that
\eqref{W31-eq:prime-power-source-line} has exact order \(31^r\) and that all
reductions are pointed and compatible.  Since \(a_0/[5]\equiv1\pmod{31}\),
the residual generator is the normalized repeated cubic line.  The displayed
carry coefficient is exactly
\eqref{eq:v064-finite-carry-coefficient}; in particular it changes the
filtered relation without replacing its repeated-direction provenance.

The hypotheses of
\cref{P2M-thm:prime-power-all-support-comparison} now hold for the one-prime
depth datum \(\nu_{31}=r\).  That theorem produces
\eqref{W31-eq:prime-power-motivic-realization}, the exact-order root/Moore pair,
and the reduction squares.  Strict support functoriality places the same
coefficient-faithful line in every surviving exact-support channel, including
the marked block \(I_0\).
\end{proof}

\begin{remark}[Finite-level family, not valuation-depth completion]
\label{W31-rem:prime-power-family-not-completion}
The compatible family in
\cref{W31-thm:prime-power-repeated-confluence} may be regarded as a pro-system
of finite coefficient fibres.  No identification is made with the
valuation-depth fibre of \(\FullMot\), with a new support axis, or with a
single \(31\)-adically completed motive.
\end{remark}

\subsection{The common-source arithmetic--motivic carrier span}

Fix a finite ordered marked support \(S_*\) for the terminal source controller,
and let \(I_*\subseteq S_*\) denote the underlying exact-support block on which
the terminal jet marking is recorded.  Neither \(S_*\) nor \(I_*\) is
identified with the arithmetic prime \(31\).

Put
\[
 \kappa_*^{\rm ar}:=\kappa_5^{\rm root}
 \in\Sha^2_{S_{31}}(K_*,\mu_{31}),
 \qquad
 L_*^{\rm ar}:=\mathbf F_{31}\kappa_*^{\rm ar},
\]
where the subscript \(5\) in the original symbol remains historical notation
for the selected cyclotomic coefficient corner.  Let
\[
 L_*^{\rm src}
 :=(\mathbf Z/31)
   \bigl\langle[\pi_{30}\mathfrak R_{31}^{[31]}]\bigr\rangle
\]
be the pointed terminal source line.  On the motivic side, specialize the
pure-weight Moore pair and the finite Moore--Reedy package to \(N=31\), and
write
\[
 db_*^{\Mot}=0,\qquad
 dc_*^{\Mot}=31b_*^{\Mot},\qquad
 \operatorname{ord}[b_*^{\Mot}]=31,
 \qquad
 L_*^{\Mot}:=(\mathbf Z/31)\langle[b_*^{\Mot}]\rangle .
\]

\begin{theorem}[Carrier-level motivic specialization of the witness]
\label{W31-thm:common-source-motivic-span}
Relative to the fixed source marking, the arithmetic realization constructed
above and the motivic realization of
\cref{P2L-thm:v48r4-motivic-seed,P2M-thm:channel-complete-realization}
give a pointed carrier span of cyclic order-\(31\) lines
\[
 \boxed{
 L_*^{\rm ar}
 \xleftarrow[\sim]{\ \rho_*^{\rm ar}\ }
 L_*^{\rm src}
 \xrightarrow[\sim]{\ \rho_*^{\Mot}\ }
 L_*^{\Mot}.}
\]
The right-hand line is the fixed Moore carrier in the \(I_*\)-block of the
finite marked Moore--Reedy package
\(\mathfrak T^{\Mot,\MR}_{31;S_*,\le31}\) of
\cref{P2M-thm:v38-finite-motivic-recursion-closure}.  Increasing the
obstruction/jet ceiling through the terminal stage changes the marked
Confluent ledger but not the underlying motivic Moore--Reedy object.
The proper \(31\)-fold arithmetic Massey value is therefore recorded as a
marking on this fixed motivic carrier.  The carrier identification alone does
not turn arithmetic cochains into motivic correspondences; the independent
operation-level provenance is instead the unipotent central-extension and
obstruction-gerbe correspondence of
\cref{W31-thm:terminal-unipotent-gerbe-provenance}.
\end{theorem}

\begin{proof}
The left arrow is the strict arithmetic realization through stage \(30\) and
the terminal identification proved in
\cref{thm:paper1-31adic-witness}.  The primitive motivic seed at \(N=31\)
has a genuine root-localization antecedent and exact order \(31\) by
\cref{P2L-thm:v48r4-motivic-seed}.  The marked cyclic source carrier therefore
has a unique generator-preserving identification with the motivic Moore line.
The channel-complete realization installs this carrier in the \(I_*\)-block,
while \cref{P2M-thm:v38-finite-motivic-recursion-closure} transports the finite
marked obstruction ledger without changing that underlying object along the
jet axis.  Both arrows preserve the distinguished generator and are pointed
isomorphisms of cyclic order-\(31\) carriers; the higher-operation marking is
additional ledger data.
\end{proof}

\subsection{The terminal full-motive lifting fibre}

Within the chosen marked full-motive presentation, let \(X_*^{(30)}\) be the
realized motivic stage carrying the lower marked ledger, let \(\mathscr L_*\)
be the invertible detector line of the terminal
component, and let \(\sigma_*\) be its typed cohomological/Tate shift.  The
motivic terminal marking determines a classifying section
\[
 \alpha_*:
 X_*^{(30)}
 \longrightarrow
 \mathfrak C_{31,X_*^{(30)}}(\mathscr L_*)[\sigma_*].
\]
Define
\begin{equation}
 \boxed{
 X_*^{(31)}
 :=
 \operatorname{Lift}_{31,X_*^{(30)}}
   (\mathscr L_*,\alpha_*)[\sigma_*]
 =
 X_*^{(30)}
 \times^h_{\mathfrak C_{31,X_*^{(30)}}(\mathscr L_*)[\sigma_*]}
 \mathfrak R_{31,X_*^{(30)}}(\mathscr L_*)[\sigma_*].}
 \label{W31-eq:terminal-motivic-lift}
\end{equation}

\begin{proposition}[Motivic terminal lifting stage]
\label{W31-prop:terminal-motivic-lift}
The object \(X_*^{(31)}\) is the obstruction-height-\(31\) Cartan--motivic
lifting stage attached to the terminal witness marking.  Its fresh relative
layer is
\[
 \operatorname{hofib}_{\alpha_*}\!\left(
 \mathfrak R_{31,X_*^{(30)}}(\mathscr L_*)[\sigma_*]
 \longrightarrow
 \mathfrak C_{31,X_*^{(30)}}(\mathscr L_*)[\sigma_*]
 \right),
\]
and its root/Moore pair satisfies
\[
 d\Gamma_*^{\Mot}=31\Omega_*^{\Mot}.
\]
This construction is compatible with support deletion and with the
Moore--Reedy latching presentation.  The arithmetic failure to extend
globally across \(h_{31}^{\rm rec}\) says that the marked arithmetic point
does not lift globally; it does not assert that the motivic lifting stack
\eqref{W31-eq:terminal-motivic-lift} is empty.
\end{proposition}

\begin{proof}
This is the specialization of
\cref{P2R-def:v92-line-valued-lift,thm:genuine-successor-stage-motivic-functor}
to the single nonzero cyclic component of order \(31\).  The description of
the relative layer follows from
\cref{thm:intrinsic-cartan-truncation-layer}.  The final distinction is the
difference between existence of the homotopy-pullback lifting object and
existence of a global arithmetic lift of its selected marked point.
\end{proof}

\subsection{Classical shadow and stack globalization}

\begin{corollary}[Motivic, classical, and stack readouts]
\label{W31-cor:motivic-classical-stack-readouts}
The exact-order-\(31\) motivic carrier specialization has the following
compatible readouts:
\begin{enumerate}[label=\textup{(\roman*)}]
\item its classical \(0/1\)-motivic shadow is the marked Moore object
      \[
       \boxed{
       \mathbf Q_{31,\mathscr L_*}
       =
       \operatorname{Cone}\!\left(
       31:\mathbf T_{\mathscr L_*}\to\mathbf T_{\mathscr L_*}
       \right)[-1];}
      \]
\item its Moore-enhanced exact-support completion belongs to the single dg
      category \(\mathfrak G^{\rm rec}_{S_*,31}\);
\item the associated fixed-presentation derived stack is
      \[
       \boxed{
       \mathfrak M^{\rm rec}_{S_*,31}
       =\mathcal M_{\mathfrak G^{\rm rec}_{S_*,31}}.}
      \]
\end{enumerate}
These are respectively the specializations of
\cref{thm:classical-low-sector-comparison,thm:moore-enhanced-stack-globalization}.
\end{corollary}

\begin{warning}[Prime-power and depth firewall]
\label{W31-warn:prime-power-motivic-firewall}
The prime-power family of
\cref{W31-thm:prime-power-repeated-confluence} belongs to the repeated cubic
\(K_0\)-branch.  It does not thicken the terminal proper-\(31\)-fold
\(K_*\)-branch: that branch supplies only the exact-order-\(31\) class
\(\kappa_5^{\rm root}\), and no compatible primitive terminal lines of orders
\(31^r\), \(r>1\), are asserted.  Likewise the finite-level family is not
valuation depth in \(\FullMot\), not a new support coordinate, and not a
single \(31^\infty\)-adic completed motive.  Finally neither the prime-power
comparison nor the bare carrier-level terminal span by itself constructs the
proper arithmetic \(31\)-fold operation.  Its independent geometric
provenance is exactly
\cref{W31-thm:terminal-unipotent-gerbe-provenance}; that gerbe theorem does not
produce compatible terminal classes of orders \(31^r\), \(r>1\).
\end{warning}

\begin{thebibliography}{99}

\bibitem{AMM}
F.~Amano, Y.~Mizusawa, and M.~Morishita,
\emph{On mod 3 triple Milnor invariants and triple cubic residue symbols in the Eisenstein number field},
Research in Number Theory \textbf{4} (2018), Paper No.~7, 29 pp.; arXiv:1412.6894v5.

\bibitem{MinacTan}
J.~Min\'a\v{c} and N.~D.~T\^an,
\emph{Triple Massey products vanish over all fields},
J. Lond. Math. Soc. (2) \textbf{94} (2016), no.~3, 909--932.

\bibitem{HarpazWittenberg}
Y.~Harpaz and O.~Wittenberg,
\emph{The Massey vanishing conjecture for number fields},
Duke Math. J. \textbf{172} (2023), no.~1, 1--41.

\bibitem{KimMorishita}
D.~Kim and M.~Morishita,
\emph{Triple symbols in arithmetic},
Research in Number Theory \textbf{11} (2025), Article 86; arXiv:2407.02063v2.

\bibitem{BrantnerMathew}
L.~Brantner and A.~Mathew,
\emph{Deformation theory and partition Lie algebras},
Acta Math. \textbf{235} (2025), no.~1, 1--148.

\bibitem{BCN}
L.~Brantner, R.~Campos, and J.~Nuiten,
\emph{PD Operads and Explicit Partition Lie Algebras},
Mem. Amer. Math. Soc. \textbf{315} (2025), no.~1597.

\bibitem{NeukirchANT}
J.~Neukirch,
\emph{Algebraic Number Theory},
Grundlehren der mathematischen Wissenschaften, vol.~322,
Springer, Berlin, 1999.

\bibitem{MilneCFT}
J.~S.~Milne,
\emph{Class Field Theory},
course notes, current online edition, Chapter III (local invariant map) and Chapter IV (Brauer groups).

\bibitem{Brown1982}
K.~S.~Brown,
\emph{Cohomology of Groups},
Graduate Texts in Mathematics, vol.~87,
Springer-Verlag, New York--Berlin, 1982.

\bibitem{ToenVaquie2007}
B.~To\"en and M.~Vaqui\'e,
\emph{Moduli of objects in dg-categories},
Ann. Sci. \'Ec. Norm. Sup\'er. (4) \textbf{40} (2007), no.~3, 387--444.

\bibitem{PTVV}
T.~Pantev, B.~To\"en, M.~Vaqui\'e, and G.~Vezzosi,
\emph{Shifted symplectic structures},
Publ. Math. Inst. Hautes \`Etudes Sci. \textbf{117} (2013), 271--328;
arXiv:1111.3209.

\bibitem{CalaqueCotangent}
D.~Calaque,
\emph{Shifted cotangent stacks are shifted symplectic},
Ann. Fac. Sci. Toulouse Math. (6) \textbf{28} (2019), no.~1, 67--90;
arXiv:1612.08101.

\bibitem{FuPositiveCharSymplectic}
J.~Fu,
\emph{$(-1)$-Shifted Darboux theorem of derived schemes in characteristic $p>2$},
arXiv:2511.12584, 2025.

\bibitem{BhattDerivedDR}
B.~Bhatt,
\emph{Completions and derived de Rham cohomology},
arXiv:1207.6193.

\bibitem{AntieauCrystallization}
B.~Antieau,
\emph{Filtrations and cohomology I: crystallization},
arXiv:2511.01567, 2025.

\bibitem{BergerMoerdijkReedy}
C.~Berger and I.~Moerdijk,
\emph{On an extension of the notion of Reedy category},
Math. Z. \textbf{269} (2011), 977--1004.
Preprint: arXiv:0809.3341.

\bibitem{GrasAmbiguous}
G.~Gras,
Sur les $\ell$-classes d'ideaux dans les extensions cycliques relatives de degre premier $\ell$,
\emph{Ann. Inst. Fourier (Grenoble)} 23 (1973), no.~4, 1--44.

\bibitem{QiMasseyLambda}
P.~Qi,
\emph{Iwasawa $\lambda$ invariant and Massey product},
arXiv:2402.06028 (2024).

\bibitem{LLSWW}
Y.~H.~J.~Lam, Y.~Liu, R.~Sharifi, P.~Wake, and J.~Wang,
\emph{Generalized Bockstein maps and Massey products},
Forum Math. Sigma \textbf{11} (2023), e5,
doi:10.1017/fms.2022.103.

\bibitem{SerreLF}
J.-P.~Serre,
\emph{Local Fields},
Graduate Texts in Mathematics 67,
Springer, New York, 1979.

\bibitem{DelbourgoKnospeLambda}
D.~Delbourgo and H.~Knospe,
\emph{On Iwasawa $\lambda$-invariants for abelian number fields and random matrix heuristics},
arXiv:2207.06287 (2022).

\bibitem{KnospeIwasawaData}
H.~Knospe,
\emph{iwasawa}: computational data accompanying the $\lambda$-invariant tables,
GitHub repository, file \texttt{lambda5-1000.txt}, 2026 snapshot.

\bibitem{HarpazWittenbergMassey}
Y.~Harpaz and O.~Wittenberg,
\emph{The Massey vanishing conjecture for number fields},
Duke Math. J. \textbf{172} (2023), no.~1, 1--41,
doi:10.1215/00127094-2022-0004.

\bibitem{MazurWiles}
B.~Mazur and A.~Wiles,
\emph{Class fields of abelian extensions of $\Q$},
Invent. Math. \textbf{76} (1984), no.~2, 179--330.

\bibitem{KnospeSpecialValues2026}
H.~Knospe,
\emph{Special values of \(p\)-adic \(L\)-functions and Iwasawa
\(\lambda\)-invariants of Dirichlet characters},
Research in Number Theory \textbf{12} (2026), Article 76.

\bibitem{NSW}
J.~Neukirch, A.~Schmidt and K.~Wingberg,
\emph{Cohomology of Number Fields},
2nd ed., Springer, 2008.

\bibitem{BarbieriVialeKahn2016}
L.~Barbieri-Viale and B.~Kahn,
\emph{On the derived category of 1-motives},
Ast\'erisque \textbf{381} (2016), xi+254 pp.; DOI 10.24033/ast.998.

\bibitem{Adams1956}
J.~F. Adams,
\emph{On the cobar construction},
Proc. Natl. Acad. Sci. USA \textbf{42} (1956), no.~7, 409--412.

\bibitem{HMS1974}
D.~Husemoller, J.~C. Moore, and J.~D. Stasheff,
\emph{Differential homological algebra and homogeneous spaces},
J. Pure Appl. Algebra \textbf{5} (1974), no.~2, 113--185.

\bibitem{LodayVallette2012}
J.-L. Loday and B.~Vallette,
\emph{Algebraic Operads},
Grundlehren der mathematischen Wissenschaften 346,
Springer, Berlin--Heidelberg, 2012.

\bibitem{ChungGraham1983}
F.~R.~K. Chung and R.~L. Graham,
\emph{Edge-colored complete graphs with precisely colored subgraphs},
Combinatorica \textbf{3} (1983), no.~3--4, 315--324.

\bibitem{GyarfasSimonyi2004}
A.~Gy\'arf\'as and G.~Simonyi,
\emph{Edge colorings of complete graphs without tricolored triangles},
J. Graph Theory \textbf{46} (2004), no.~3, 211--216.

\bibitem{DegliseGysin2011}
F.~D\'eglise,
\emph{Around the Gysin triangle I},
Doc. Math. \textbf{13} (2008), 613--675; revised arXiv version 2011.

\bibitem{BerestKhachatryanRamadoss}
Y.~Berest, G.~Khachatryan, and A.~Ramadoss,
\emph{Derived representation schemes and cyclic homology},
arXiv:1112.1449v2.

\bibitem{Letz2024}
J.~C.~Letz,
\emph{Transfer of $A_\infty$-structures to projective resolutions},
arXiv:2312.13696v2.

\bibitem{Burke2015}
J.~Burke,
\emph{Higher homotopies and Golod rings},
arXiv:1508.03782v2.

\bibitem{BoggiCorobCook}
M.~Boggi and G.~Corob Cook,
\emph{Continuous cohomology and homology of profinite groups},
Doc. Math. \textbf{21} (2016), 1269--1312; arXiv:math/0306381v4.

\bibitem{DwyerGreenleesIyengar}
W.~G.~Dwyer, J.~P.~C.~Greenlees and S.~Iyengar,
\emph{Duality in algebra and topology},
Adv. Math. \textbf{200} (2006), 357--402.

\bibitem{HoyoisEquivariantSix}
M.~Hoyois,
\emph{The six operations in equivariant motivic homotopy theory},
Adv. Math. \textbf{305} (2017), 197--279;
arXiv:1509.02145.

\bibitem{KhanRaviStacks}
A.~A. Khan and C.~Ravi,
\emph{Generalized cohomology theories for algebraic stacks},
arXiv:2106.15001.

\bibitem{ChoudhuryDeshmukhHogadi}
U.~Choudhury, N.~Deshmukh and A.~Hogadi,
\emph{The Nisnevich motive of an algebraic stack},
arXiv:2012.13304.

\bibitem{TotaroChowBG}
B.~Totaro,
\emph{The Chow ring of a classifying space},
in \emph{Algebraic $K$-theory (Seattle, WA, 1997)},
Proc. Sympos. Pure Math. \textbf{67}, AMS, 1999, 249--281;
arXiv:math/9802097.

\bibitem{DwyerMassey}
W.~G. Dwyer,
\emph{Homology, Massey products and maps between groups},
J. Pure Appl. Algebra \textbf{6} (1975), no.~2, 177--190;
doi:10.1016/0022-4049(75)90006-7.

\bibitem{ToenDerivedMorita}
B.~To\"en,
\emph{The homotopy theory of dg-categories and derived Morita theory},
Invent. Math. \textbf{167} (2007), no.~3, 615--667;
arXiv:math/0408337.

\end{thebibliography}

\end{document}
